% major-theorems.tex -- GENERATED by (write-major-theorems-tex! ...)
% from the table of major theorems.  Do not hand-edit: every edit is
% lost at the next load.  The statements, tactics, lemmas, oracles and
% bills below are read from the LIVE system, not from the table.
%
% Requires from the host preamble: amsmath, amssymb, amsthm, graphicx,
% longtable, and the theorem environments thm / prop / cor / lem.
%
% \vnbstmt -- a displayed statement that shrinks to the text width only
% when it would overflow it, so short statements are set at normal size and
% a wide one never runs into the margin.  Needs graphicx.
%
% The leading \leavevmode is not decoration.  amsthm defers the theorem head
% (`Theorem 1.1 (Title).') into \everypar, so it is injected at the start of
% the first paragraph of the body; without \leavevmode the opening \par is a
% no-op on an empty paragraph, the head lands in the SAME paragraph as the
% box below, and head-width + \linewidth overfills every statement by 300 pt.
% \leavevmode starts a paragraph for the head to attach to, and the \par then
% closes it before the box.
\providecommand{\vnbstmt}[1]{%
  \leavevmode\par\nobreak\smallskip\noindent
  \resizebox{\ifdim\width>\linewidth\linewidth\else\width\fi}{!}{$\displaystyle #1$}%
  \par\smallskip}
\providecommand{\vnbmissing}[1]{\textit{No theorem \texttt{#1} is installed in the library.}}
% \vnbprf -- one paragraph of the proof block.  Set RAGGED RIGHT: these are
% comma-separated lists of \texttt identifiers, which neither hyphenate nor
% break at their hyphens, so justified setting either overfills the measure or
% (under \sloppy) buys the fit with very loose interword glue.
\providecommand{\vnbprf}[1]{{\raggedright #1\par}}
% \vnbpart -- one PART of a statement that the library gives under several names:
% its own paragraph, the name as a run-in label (the user, 2026-09-21).
\providecommand{\vnbpart}[1]{\par\medskip{\raggedright\noindent\textup{\textbf{Part} #1.}\par}\nopagebreak\smallskip}
\providecommand{\vnbproofpart}[1]{\par\medskip\noindent\textbf{Proof of} #1\textbf{.}\par\nopagebreak}
%
\section{Calculus}
%
\renewcommand{\thethm}{1.1}
\begin{prop}[{Order compatible with addition and multiplication on Z}]\label{thm:rr-lt-add}
%
\leavevmode
\vnbpart{\texttt{rr-lt-add}}
\vnbstmt{\begin{array}{@{}l@{}} \forall x, y, u, v. \\ \quad (x \in \mathbb{R}) \wedge \\ \quad (y \in \mathbb{R}) \wedge \\ \quad (u \in \mathbb{R}) \wedge \\ \quad (v \in \mathbb{R}) \wedge \\ \quad (x < y) \wedge \\ \quad (u \leq v) \Rightarrow \\ \quad \quad ((x + u) < (y + v)) \end{array}}
\vnbpart{\texttt{rr-mul-le-right}}
\vnbstmt{\begin{array}{@{}l@{}} \forall a, b, c \in \mathbb{R}. \\ \quad (((a \leq b) \wedge (0 \leq c)) \Rightarrow ((a \cdot c) \leq (b \cdot c))) \end{array}}
\vnbpart{\texttt{rr-mul-pos}}
\vnbstmt{\begin{array}{@{}l@{}} \forall a, b \in \mathbb{R}. \\ \quad (((0 < a) \wedge (0 < b)) \Rightarrow (0 < (a \cdot b))) \end{array}}
%
\end{prop}
%
\begin{proof}
%
The proofs: the file that holds each, the lemmas it cites, its oracles and its bill.
%
\vnbprf{\vnbproofpart{\texttt{rr-lt-add}}}
%
\vnbprf{\textbf{Proof.} \texttt{theorem-\allowbreak{}library/rr-\allowbreak{}order-\allowbreak{}basics.scm} --- 8 recorded steps.}
%
\vnbprf{\textbf{Lemmas.} none cited by name.}
%
\vnbprf{\textbf{Oracles.} \texttt{ineq}.  \textbf{Bill:} \texttt{proven modulo 0}.}
%
\vnbprf{\vnbproofpart{\texttt{rr-mul-le-right}}}
%
\vnbprf{\textbf{Proof.} \texttt{theorem-\allowbreak{}library/rake-\allowbreak{}hb-\allowbreak{}leaves.scm} --- 29 recorded steps.}
%
\vnbprf{\textbf{Lemmas.} \texttt{rr-sub-in-rr}, \texttt{rr-mul-in-rr}, \texttt{rr-leq-mul-nonneg}.}
%
\vnbprf{\textbf{Oracles.} \texttt{ineq}, \texttt{crs}.  \textbf{Bill:} \texttt{proven modulo 0}.}
%
\vnbprf{\vnbproofpart{\texttt{rr-mul-pos}}}
%
\vnbprf{\textbf{Proof.} \texttt{theorem-\allowbreak{}library/rr-\allowbreak{}recip-\allowbreak{}order.scm} --- 24 recorded steps.}
%
\vnbprf{\textbf{Lemmas.} \texttt{rr-zero-in}, \texttt{rr-lt-scale-pos}, \texttt{rr-mul-closed}, \texttt{rr-mul-zero}.}
%
\vnbprf{\textbf{Oracles.} \texttt{ineq}, \texttt{crs}.  \textbf{Bill:} \texttt{proven modulo 0}.}
%
\end{proof}
%
\noindent\textbf{Deviation from the notes.} The additive half is rr-lt-add, stated over RR (x $<$ y and u $<$= v implies x + u $<$ y + v), not over ZZ; ZZ inherits it through zz-in-rr. The multiplicative half, a $<$ b and c in NN implies a * c $<$ b * c, is NOT stated anywhere: rr-mul-le-right is the NON-STRICT form (a $<$= b and 0 $<$= c implies a * c $<$= b * c) and rr-mul-pos gives only 0 $<$ a and 0 $<$ b implies 0 $<$ a * b. There is also no ZZ-specific order lemma of either shape; the notes derive both from the definition b - a in N, which the library does not use.
%
\renewcommand{\thethm}{1.2}
\begin{prop}[{Trichotomy of the integer ordering}]\label{thm:rr-lt-trichotomy}
%
\leavevmode
\vnbpart{\texttt{rr-lt-trichotomy}}
\vnbstmt{\forall u, v \in \mathbb{R}.\; ((u < v) \vee ((u = v) \vee (v < u)))}
\vnbpart{\texttt{zz-trichotomy}}
\vnbstmt{\forall a \in \mathbb{Z}.\; ((a < 0) \vee ((a = 0) \vee (0 < a)))}
%
\end{prop}
%
\begin{proof}
%
The proofs: the file that holds each, the lemmas it cites, its oracles and its bill.
%
\vnbprf{\vnbproofpart{\texttt{rr-lt-trichotomy}}}
%
\vnbprf{\textbf{Proof.} \texttt{theorem-\allowbreak{}library/rr-\allowbreak{}order-\allowbreak{}basics.scm} --- 55 recorded steps.}
%
\vnbprf{\textbf{Lemmas.} \texttt{rr-leq-total}, \texttt{neq-sym}.}
%
\vnbprf{\textbf{Oracles.} none.  \textbf{Bill:} \texttt{proven modulo 0}.}
%
\vnbprf{\vnbproofpart{\texttt{zz-trichotomy}}}
%
\vnbprf{\textbf{Proof.} \texttt{theorem-\allowbreak{}library/zz-\allowbreak{}order.scm} --- 5 recorded steps.}
%
\vnbprf{\textbf{Lemmas.} \texttt{zz-in-rr}, \texttt{rr-zero-in}, \texttt{rr-lt-trichotomy}.}
%
\vnbprf{\textbf{Oracles.} none.  \textbf{Bill:} \texttt{proven modulo 0}.}
%
\end{proof}
%
\noindent\textbf{Deviation from the notes.} rr-lt-trichotomy is stated over RR (u, v in RR implies u $<$ v or u = v or v $<$ u), not over ZZ; the ZZ statement in the tree, zz-trichotomy, is only the special case against 0 (a $<$ 0 or a = 0 or 0 $<$ a). The notes' pair form for integers has no library statement of its own.
%
\renewcommand{\thethm}{1.3}
\begin{prop}[{1 is the least natural number}]\label{thm:nn-one-le-succ}
%
\leavevmode
\vnbpart{\texttt{nn-one-le-succ}}
\vnbstmt{\forall n \in \mathbb{N}.\; (1 \leq (n + 1))}
\vnbpart{\texttt{zz-discrete}}
\vnbstmt{\forall a \in \mathbb{Z}.\; ((0 < a) \Rightarrow (1 \leq a))}
\vnbpart{\texttt{nn-pos-is-succ}}
\vnbstmt{\begin{array}{@{}l@{}} \forall n \in \mathbb{N}. \\ \quad ((1 \leq n) \Rightarrow \exists q \in \mathbb{N}.\; (n = (q + 1))) \end{array}}
%
\end{prop}
%
\begin{proof}
%
The proofs: the file that holds each, the lemmas it cites, its oracles and its bill.
%
\vnbprf{\vnbproofpart{\texttt{nn-one-le-succ}}}
%
\vnbprf{\textbf{Proof.} \texttt{theorem-\allowbreak{}library/nn-\allowbreak{}order-\allowbreak{}ord.scm} --- 10 recorded steps.}
%
\vnbprf{\textbf{Lemmas.} \texttt{nn-zero-in}, \texttt{nn-zero-le}, \texttt{nn-succ-mono}.}
%
\vnbprf{\textbf{Oracles.} \texttt{arith}.  \textbf{Bill:} \texttt{proven modulo 0}.}
%
\vnbprf{\vnbproofpart{\texttt{zz-discrete}}}
%
\vnbprf{\textbf{Proof.} \texttt{theorem-\allowbreak{}library/zz-\allowbreak{}order.scm} --- 27 recorded steps.}
%
\vnbprf{\textbf{Lemmas.} \texttt{zz-in-rr}, \texttt{rr-zero-in}, \texttt{rr-one-in}, \texttt{zz-generated-by-nn}, \texttt{nn-in-rr}, \texttt{nn-zero-le}, \texttt{rr-neg-closed}, \texttt{rr-pos-ne-zero}, \texttt{nn-nonzero-is-succ}, \texttt{nn-succ-closed}, \texttt{nn-one-le-succ}.}
%
\vnbprf{\textbf{Oracles.} \texttt{ineq}, \texttt{arith}.  \textbf{Bill:} \texttt{proven modulo 0}.}
%
\vnbprf{\vnbproofpart{\texttt{nn-pos-is-succ}}}
%
\vnbprf{\textbf{Proof.} \texttt{theorem-\allowbreak{}library/nn-\allowbreak{}pos-\allowbreak{}is-\allowbreak{}succ.scm} --- 17 recorded steps.}
%
\vnbprf{\textbf{Lemmas.} \texttt{nn-zero-in}, \texttt{nn-le-refl}, \texttt{nn-not-le-zero-pos}, \texttt{nn-nonzero-is-succ}.}
%
\vnbprf{\textbf{Oracles.} \texttt{ineq}, \texttt{arith}.  \textbf{Bill:} \texttt{proven modulo 0}.}
%
\end{proof}
%
\noindent\textbf{Deviation from the notes.} The notes' N starts at 1; VNB's NN CONTAINS 0 (axiom nn-zero-in), so forall c in NN. 1 $<$= c is FALSE in the library and no theorem states it. The nearest are nn-one-le-succ (forall n in NN. 1 $<$= succ(n)) and zz-discrete (a in ZZ and 0 $<$ a implies 1 $<$= a), i.e. the statement for the POSITIVE elements only. The notes derive this from the least-member axiom for subsets of N; the library's counterpart, nn-least-element, is proven.
%
\renewcommand{\thethm}{1.4}
\begin{cor}[{No integer strictly between m and m + 1}]\label{thm:zz-lt-succ-le}
%
\leavevmode
\vnbpart{\texttt{zz-lt-succ-le}}
\vnbstmt{\forall a, b \in \mathbb{Z}.\; ((a < b) \Rightarrow ((a + 1) \leq b))}
\vnbpart{\texttt{zz-discrete}}
\vnbstmt{\forall a \in \mathbb{Z}.\; ((0 < a) \Rightarrow (1 \leq a))}
%
\end{cor}
%
\begin{proof}
%
The proofs: the file that holds each, the lemmas it cites, its oracles and its bill.
%
\vnbprf{\vnbproofpart{\texttt{zz-lt-succ-le}}}
%
\vnbprf{\textbf{Proof.} \texttt{theorem-\allowbreak{}library/zz-\allowbreak{}order.scm} --- 12 recorded steps.}
%
\vnbprf{\textbf{Lemmas.} \texttt{zz-in-rr}, \texttt{zz-sub-in-zz}, \texttt{zz-discrete}.}
%
\vnbprf{\textbf{Oracles.} \texttt{ineq}, \texttt{arith}.  \textbf{Bill:} \texttt{proven modulo 0}.}
%
\vnbprf{\vnbproofpart{\texttt{zz-discrete}}}
%
\vnbprf{\textbf{Proof.} \texttt{theorem-\allowbreak{}library/zz-\allowbreak{}order.scm} --- 27 recorded steps.}
%
\vnbprf{\textbf{Lemmas.} \texttt{zz-in-rr}, \texttt{rr-zero-in}, \texttt{rr-one-in}, \texttt{zz-generated-by-nn}, \texttt{nn-in-rr}, \texttt{nn-zero-le}, \texttt{rr-neg-closed}, \texttt{rr-pos-ne-zero}, \texttt{nn-nonzero-is-succ}, \texttt{nn-succ-closed}, \texttt{nn-one-le-succ}.}
%
\vnbprf{\textbf{Oracles.} \texttt{ineq}, \texttt{arith}.  \textbf{Bill:} \texttt{proven modulo 0}.}
%
\end{proof}
%
\noindent\textbf{Deviation from the notes.} The library states the contrapositive form zz-lt-succ-le (a, b in ZZ and a $<$ b implies a + 1 $<$= b) rather than the negated-existence form of the notes; the two are interderivable in one step. No statement is phrased with the term m + 1 as an upper bound.
%
\renewcommand{\thethm}{1.9}
\begin{prop}[{Uniqueness of the least upper bound}]\label{thm:calculus-1.9}
%
\textit{For any set S, least upper bounds a, a' for S are identical.}
%
Not in the library.
%
\end{prop}
%
\begin{proof}
%
The proofs: the file that holds each, the lemmas it cites, its oracles and its bill.
%
\vnbprf{There is no library statement, and so no proof.}
%
\end{proof}
%
\noindent\textbf{Deviation from the notes.} Nothing states that two least upper bounds of a subset of RR coincide. The library takes sup as a total functoid governed by the three axioms rr-sup-in, rr-sup-upper, rr-sup-least (all three stamped asserted, warrant: reference), so uniqueness is built into the notation rather than proved. The nearest statements are esup-unique, which is the same claim for the EXTENDED nonnegative reals rr-pos-star with bounds quantified over rr-pos-star, and the axiom rr-leq-antisymmetric from which the notes' one-line proof would run.
%
\renewcommand{\thethm}{1.11}
\begin{prop}[{The nested interval property}]\label{thm:rr-nested-intervals}
%
\vnbstmt{\begin{array}{@{}l@{}} \forall o, i \in (\mathbb{N} \to \mathbb{R}). \\ \quad \forall k \in \mathbb{N}.\; (o(k) \leq o((k + 1))) \wedge \\ \quad \forall k \in \mathbb{N}.\; (i((k + 1)) \leq i(k)) \wedge \\ \quad \forall k \in \mathbb{N}.\; (o(k) \leq i(k)) \Rightarrow \\ \quad \quad \exists y \in \mathbb{R}. \forall n \in \mathbb{N}. \\ \quad \quad \quad ((o(n) \leq y) \wedge (y \leq i(n))) \end{array}}
%
\end{prop}
%
\begin{proof}
%
The proofs: the file that holds each, the lemmas it cites, its oracles and its bill.
%
\vnbprf{\textbf{Proof.} \texttt{theorem-\allowbreak{}library/zero-\allowbreak{}deriv-\allowbreak{}off-\allowbreak{}countable.scm} --- 135 recorded steps.}
%
\vnbprf{\textbf{Lemmas.} \texttt{nn-monotone-step-implies-le}, \texttt{nn-antitone-step-implies-le}, \texttt{nn-in-rr}, \texttt{fun-apply-type-c}, \texttt{rr-leq-total}, \texttt{subset-def}, \texttt{nn-zero-in}, \texttt{rr-sup-in}, \texttt{rr-sup-upper}, \texttt{rr-sup-least}.}
%
\vnbprf{\textbf{Oracles.} \texttt{ineq}.  \textbf{Bill:} \texttt{proven modulo 0}.}
%
\end{proof}
%
\noindent\textbf{Deviation from the notes.} UPDATED 2026-09-23 (batch 20-A, theorem-library/zero-deriv-off-countable.scm): rr-nested-intervals -- for sequences lo, hi in fun(nn, rr) with lo nondecreasing, hi nonincreasing and lo(k) $<$= hi(k) for every k, there is a real y with lo(n) $<$= y $<$= hi(n) for every n; the nesting of the notes is stated through the monotonicity of the endpoints rather than through set inclusion of ccint(lo(n), hi(n)), and the point is delivered by existence, the intersection is not formed. EARLIER NOTE: No statement about a nested sequence of closed real intervals ccint(a\_n, b\_n) with non-empty intersection. The nearest is nested-closed-balls-point, which is the corresponding fact in a COMPLETE metric space for closed BALLS and carries the extra hypothesis that the radii satisfy rds(succ(n)) $<$= recip(n + 1), so it is not the interval statement; also intersection-of-membership and heine-borel-ccint / compact-iff-fip, from which the interval statement would follow with work. Theorem 1.19 below, which the notes prove FROM 1.11, is itself proven in the library by a different route.
%
\renewcommand{\thethm}{1.13}
\begin{prop}[{Order isomorphism of (-1, 1) with the extended reals}]\label{thm:calculus-1.13}
%
\textit{f: x |-$>$ x / (1 - abs(x)) is an order preserving bijection from (-1, 1) onto R; its inverse is g: s |-$>$ s / (1 + abs(s)); f(-x) = -f(x); and f extends to an order preserving bijection from [-1, 1] onto the extended real line.}
%
Not in the library.
%
\end{prop}
%
\begin{proof}
%
The proofs: the file that holds each, the lemmas it cites, its oracles and its bill.
%
\vnbprf{There is no library statement, and so no proof.}
%
\end{proof}
%
\noindent\textbf{Deviation from the notes.} Nothing about the map x |-$>$ x / (1 - abs(x)), its inverse s |-$>$ s / (1 + abs(s)), its oddness, or any order isomorphism onto RR-STAR. The extended line exists only as the membership predicate rr-star-membership (x in RR-STAR iff x in RR or x = POS-INF or x = NEG-INF) with the two order axioms pos-inf-upper-bound and neg-inf-lower-bound; there is no order-isomorphism vocabulary for it and no extension of the map to the closed interval [-1, 1].
%
\renewcommand{\thethm}{1.17}
\begin{prop}[{Uniqueness of limits}]\label{thm:rr-limit-unique}
%
\leavevmode
\vnbpart{\texttt{rr-limit-unique}}
\vnbstmt{\begin{array}{@{}l@{}} \forall f \in (\mathbb{N} \to \mathbb{R}), \ell_{1}, \ell_{2} \in \mathbb{R}. \\ \quad \operatorname{converges\text{-}to}(\operatorname{rr\text{-}ms}, f, \ell_{1}) \wedge \\ \quad \operatorname{converges\text{-}to}(\operatorname{rr\text{-}ms}, f, \ell_{2}) \Rightarrow \\ \quad \quad (\ell_{1} = \ell_{2}) \end{array}}
\vnbpart{\texttt{metric-limit-unique}}
\vnbstmt{\begin{array}{@{}l@{}} \forall s, f, \ell_{1}, \ell_{2}. \\ \quad \operatorname{converges\text{-}to}(s, f, \ell_{1}) \wedge \\ \quad \operatorname{converges\text{-}to}(s, f, \ell_{2}) \Rightarrow \\ \quad \quad (\ell_{1} = \ell_{2}) \end{array}}
%
\end{prop}
%
\begin{proof}
%
The proofs: the file that holds each, the lemmas it cites, its oracles and its bill.
%
\vnbprf{\vnbproofpart{\texttt{rr-limit-unique}}}
%
\vnbprf{\textbf{Proof.} \texttt{theorem-\allowbreak{}library/metric-\allowbreak{}limit-\allowbreak{}unique.scm} --- 5 recorded steps.}
%
\vnbprf{\textbf{Lemmas.} \texttt{metric-limit-unique}.}
%
\vnbprf{\textbf{Oracles.} \texttt{ineq}, \texttt{crs}.  \textbf{Bill:} \texttt{proven modulo 0}.}
%
\vnbprf{\vnbproofpart{\texttt{metric-limit-unique}}}
%
\vnbprf{\textbf{Proof.} \texttt{theorem-\allowbreak{}library/metric-\allowbreak{}limit-\allowbreak{}unique.scm} --- 77 recorded steps.}
%
\vnbprf{\textbf{Lemmas.} \texttt{rr-pos-halvable}, \texttt{nn-max-closed}, \texttt{nn-in-rr}, \texttt{rr-le-max-left}, \texttt{rr-le-max-right}, \texttt{fun-apply-type-c}, \texttt{metric-sym}, \texttt{metric-triangle}, \texttt{metric-dist-real}, \texttt{rr-le-all-pos-nonpos}, \texttt{metric-pos}, \texttt{rr-zero-in}, \texttt{rr-leq-antisymmetric}, \texttt{metric-zero-eq}.}
%
\vnbprf{\textbf{Oracles.} \texttt{ineq}, \texttt{crs}.  \textbf{Bill:} \texttt{proven modulo 0}.}
%
\end{proof}
%
\renewcommand{\thethm}{1.19}
\begin{thm}[{Bolzano-Weierstrass on a closed interval}]\label{thm:rr-interval-seq-has-convergent-subseq}
%
\leavevmode
\vnbpart{\texttt{rr-interval-seq-has-convergent-subseq}}
\vnbstmt{\begin{array}{@{}l@{}} \forall v, \mathit{hiv} \in \mathbb{R}, \mathit{hqv} \in (\mathbb{N} \to \operatorname{ccint}(v, \mathit{hiv})). \exists \mathit{phv}. \\ \quad \operatorname{strictly\text{-}mono\text{-}nn}(\mathit{phv}) \wedge \\ \quad \exists \mathit{ptv} \in \mathbb{R}. \\ \quad \quad \operatorname{converges\text{-}to}(\operatorname{rr\text{-}ms}, \operatorname{subseq}(\mathit{hqv}, \mathit{phv}), \mathit{ptv}) \end{array}}
\vnbpart{\texttt{rr-bolzano-weierstrass}}
\vnbstmt{\begin{array}{@{}l@{}} \forall v \in (\mathbb{N} \to \mathbb{R}), \mathit{bdv} \in \mathbb{R}. \\ \quad \forall \mathit{iqv} \in \mathbb{N}.\; (\lvert v(\mathit{iqv}) \rvert \leq \mathit{bdv}) \Rightarrow  \\ \quad \quad \exists \mathit{phv}. \\ \quad \quad \quad \operatorname{strictly\text{-}mono\text{-}nn}(\mathit{phv}) \wedge \\ \quad \quad \quad \exists \mathit{ptv} \in \mathbb{R}. \\ \quad \quad \quad \quad \operatorname{converges\text{-}to}(\operatorname{rr\text{-}ms}, \operatorname{subseq}(v, \mathit{phv}), \mathit{ptv}) \end{array}}
%
\end{thm}
%
\begin{proof}
%
The proofs: the file that holds each, the lemmas it cites, its oracles and its bill.
%
\vnbprf{\vnbproofpart{\texttt{rr-interval-seq-has-convergent-subseq}}}
%
\vnbprf{\textbf{Proof.} \texttt{theorem-\allowbreak{}library/rake-\allowbreak{}bolzano-\allowbreak{}weierstrass-\allowbreak{}2.scm} --- 128 recorded steps.}
%
\vnbprf{\textbf{Lemmas.} \texttt{rr-is-set}, \texttt{ccint-subset-rr}, \texttt{subclass-of-set-is-set}, \texttt{fun-codomain-superset}, \texttt{null-rr-seq-exists}, \texttt{ccint-grid-cover-seq}, \texttt{block-family-combinatorial}, \texttt{diagonalization}, \texttt{nn-is-set}, \texttt{fun-apply-type-c}, \texttt{rr-cauchy-converges}, \texttt{rr-pos-rr-in-rr}, \texttt{nn-le-refl}, \texttt{rr-sub-in-rr}, \texttt{rr-abs-closed}, \texttt{rr-le-trans-c}, \texttt{subseq-apply}.}
%
\vnbprf{\textbf{Oracles.} \texttt{ineq}, \texttt{crs}, \texttt{arith}.  \textbf{Bill:} \texttt{proven modulo 0}.}
%
\vnbprf{\vnbproofpart{\texttt{rr-bolzano-weierstrass}}}
%
\vnbprf{\textbf{Proof.} \texttt{theorem-\allowbreak{}library/rake-\allowbreak{}bolzano-\allowbreak{}weierstrass-\allowbreak{}2.scm} --- 63 recorded steps.}
%
\vnbprf{\textbf{Lemmas.} \texttt{rr-abs-bound}, \texttt{fun-codomain-iff}, \texttt{fun-apply-type-c}, \texttt{ccint-membership}, \texttt{rr-neg-closed}, \texttt{rr-interval-seq-has-convergent-subseq}.}
%
\vnbprf{\textbf{Oracles.} \texttt{arith}, \texttt{ineq}, \texttt{crs}.  \textbf{Bill:} \texttt{proven modulo 0}.}
%
\end{proof}
%
\noindent\textbf{Deviation from the notes.} rr-interval-seq-has-convergent-subseq types the sequence as fun(nn, ccint(lov, hiv)) and asserts a strictly monotone phi and a limit ptv in RR; like the notes' statement it does NOT claim the limit lies in [a, b], although the notes' proof produces such a limit. It also carries no a $<$ b guard (vacuous when the interval is empty). rr-bolzano-weierstrass is the companion for a sequence merely bounded in absolute value.
%
\renewcommand{\thethm}{2.1}
\begin{prop}[{Derivative of the monomial x\textasciicircum{}n}]\label{thm:deriv-power}
%
\leavevmode
\vnbpart{\texttt{deriv-power}}
\vnbstmt{\begin{array}{@{}l@{}} \forall n \in \mathbb{N}, a \in \mathbb{R}. \\ \quad \operatorname{is\text{-}diff\text{-}at}((\lambda x \in \mathbb{R}.\; \operatorname{power}(x, (n + 1))), a, ((n + 1) \cdot \operatorname{power}(a, n))) \end{array}}
\vnbpart{\texttt{deriv-const}}
\vnbstmt{\begin{array}{@{}l@{}} \forall c, a. \\ \quad ((c \in \mathbb{R}) \wedge (a \in \mathbb{R})) \Rightarrow  \\ \quad \quad \operatorname{is\text{-}diff\text{-}at}((\lambda x \in \mathbb{R}.\; c), a, 0) \end{array}}
%
\end{prop}
%
\begin{proof}
%
The proofs: the file that holds each, the lemmas it cites, its oracles and its bill.
%
\vnbprf{\vnbproofpart{\texttt{deriv-power}}}
%
\vnbprf{\textbf{Proof.} \texttt{theorem-\allowbreak{}library/deriv-\allowbreak{}power.scm} --- 101 recorded steps.}
%
\vnbprf{\textbf{Lemmas.} \texttt{rr-subset-cc}, \texttt{nn-zero-in}, \texttt{nn-succ-closed}, \texttt{power-zero}, \texttt{deriv-identity}, \texttt{pow-lam-in-fun}, \texttt{power-succ}, \texttt{diff-transfer-ptwise-eq}, \texttt{nn-in-rr}, \texttt{power-closed-at}, \texttt{deriv-product}, \texttt{nn-succ-plus-one}, \texttt{pw-step-identity}.}
%
\vnbprf{\textbf{Oracles.} \texttt{arith}, \texttt{crs}, \texttt{ineq}.  \textbf{Bill:} \texttt{proven modulo 0}.}
%
\vnbprf{\vnbproofpart{\texttt{deriv-const}}}
%
\vnbprf{\textbf{Proof.} \texttt{theorem-\allowbreak{}library/differentiation.scm} --- 29 recorded steps.}
%
\vnbprf{\textbf{Lemmas.} \texttt{rr-zero-in}, \texttt{rr-is-set}, \texttt{const-continuous-at}.}
%
\vnbprf{\textbf{Oracles.} \texttt{crs}, \texttt{arith}, \texttt{ineq}.  \textbf{Bill:} \texttt{proven modulo 0}.}
%
\end{proof}
%
\noindent\textbf{Deviation from the notes.} deriv-power states the case n $>$= 1 in the shifted form is-diff-at(vnb-lambda(x, rr, x\textasciicircum{}succ(n)), a, succ(n) * a\textasciicircum{}n) for n in NN; the n = 0 case of the notes is a separate theorem, deriv-const. The notes also assert one-sided differentiability of x\textasciicircum{}n on [a, b) and (a, b] as part of the proof; the library has no one-sided notion (D2) and its is-diff-at requires the monomial as a TOTAL function on RR (D1), which it is.
%
\renewcommand{\thethm}{2.3}
\begin{prop}[{Higher derivatives of x\textasciicircum{}n}]\label{thm:nth-deriv-monomial}
%
\leavevmode
\vnbpart{\texttt{nth-deriv-monomial}}
\vnbstmt{\begin{array}{@{}l@{}} \forall n, k \in \mathbb{N}. \\ \quad (k \leq n) \Rightarrow  \\ \quad \quad (\operatorname{nth\text{-}deriv}((\lambda x \in \mathbb{R}.\; \operatorname{power}(x, n)), k) \simeq (\lambda x \in \mathbb{R}.\; ((\operatorname{factorial}(n) \cdot \operatorname{recip}(\operatorname{factorial}(\operatorname{nn\text{-}minus}(n, k)))) \cdot \operatorname{power}(x, \operatorname{nn\text{-}minus}(n, k))))) \end{array}}
\vnbpart{\texttt{nth-deriv-monomial-vanish}}
\vnbstmt{\begin{array}{@{}l@{}} \forall n, k \in \mathbb{N}. \\ \quad (n < k) \Rightarrow  \\ \quad \quad (\operatorname{nth\text{-}deriv}((\lambda x \in \mathbb{R}.\; \operatorname{power}(x, n)), k) \simeq (\lambda x \in \mathbb{R}.\; 0)) \end{array}}
%
\end{prop}
%
\begin{proof}
%
The proofs: the file that holds each, the lemmas it cites, its oracles and its bill.
%
\vnbprf{\vnbproofpart{\texttt{nth-deriv-monomial}}}
%
\vnbprf{\textbf{Proof.} \texttt{theorem-\allowbreak{}library/taylor-\allowbreak{}cluster.scm} --- 38 recorded steps.}
%
\vnbprf{\textbf{Lemmas.} \texttt{nn-le-gap}, \texttt{nn-in-rr}, \texttt{rr-add-in-rr}, \texttt{rr-sub-in-rr}, \texttt{nn-add-comm}, \texttt{nth-deriv-monomial-gap}.}
%
\vnbprf{\textbf{Oracles.} \texttt{crs}, \texttt{arith}, \texttt{ineq}.  \textbf{Bill:} \texttt{proven modulo 0}.}
%
\vnbprf{\vnbproofpart{\texttt{nth-deriv-monomial-vanish}}}
%
\vnbprf{\textbf{Proof.} \texttt{theorem-\allowbreak{}library/taylor-\allowbreak{}cluster.scm} --- 12 recorded steps.}
%
\vnbprf{\textbf{Lemmas.} \texttt{nn-lt-succ-le}, \texttt{nn-succ-closed}, \texttt{nn-le-gap}, \texttt{tab-succ-add}, \texttt{nth-deriv-monomial-beyond}.}
%
\vnbprf{\textbf{Oracles.} \texttt{arith}, \texttt{ineq}, \texttt{crs}.  \textbf{Bill:} \texttt{proven modulo 0}.}
%
\end{proof}
%
\noindent\textbf{Deviation from the notes.} UPDATED 2026-09-24 (batch 25): PROVEN, theorem-library/taylor-cluster.scm. THE NOTES' FORMULA (8) IS MISPRINTED: the coefficient of x\textasciicircum{}(n-k) is n!/(n-k)!, not k! (the printed one gives x for the derivative of x\textasciicircum{}2). EARLIER NOTE: No statement computing nth-deriv of a monomial, and no notion of infinite differentiability. The tree has nth-deriv only through its two defining equations nth-deriv-zero and nth-deriv-succ and through the Taylor development; the nearest computational fact is deriv-power for the single derivative and deriv-polynomial for the derivative of a partial-sum polynomial.
%
\renewcommand{\thethm}{2.4}
\begin{prop}[{Differentiability implies continuity}]\label{thm:diff-implies-continuous}
%
\leavevmode
\vnbpart{\texttt{diff-implies-continuous}}
\vnbstmt{\begin{array}{@{}l@{}} \forall f, a, l. \\ \quad \operatorname{is\text{-}diff\text{-}at}(f, a, l) \Rightarrow  \\ \quad \quad \operatorname{is\text{-}continuous\text{-}at}(\operatorname{rr\text{-}ms}, \operatorname{rr\text{-}ms}, f, a) \end{array}}
\vnbpart{\texttt{diff-on-implies-continuous}}
\vnbstmt{\begin{array}{@{}l@{}} \forall k, u, f, a, l. \\ \quad \operatorname{is\text{-}diff\text{-}on}(k, u, f, a, l) \Rightarrow  \\ \quad \quad \operatorname{is\text{-}continuous\text{-}at}(\operatorname{subspace\text{-}ms}(\operatorname{nf\text{-}metric\text{-}space}(k), u), \operatorname{nf\text{-}metric\text{-}space}(k), f, a) \end{array}}
%
\end{prop}
%
\begin{proof}
%
The proofs: the file that holds each, the lemmas it cites, its oracles and its bill.
%
\vnbprf{\vnbproofpart{\texttt{diff-implies-continuous}}}
%
\vnbprf{\textbf{Proof.} \texttt{theorem-\allowbreak{}library/differentiation.scm} --- 28 recorded steps.}
%
\vnbprf{\textbf{Lemmas.} \texttt{fun-apply-type-c}, \texttt{const-continuous-at}, \texttt{identity-continuous-at}, \texttt{sub-continuous-at}, \texttt{product-continuous-at}, \texttt{sum-continuous-at}, \texttt{eq-sym}, \texttt{cont-transfer-ptwise-eq}.}
%
\vnbprf{\textbf{Oracles.} \texttt{crs}, \texttt{arith}, \texttt{ineq}.  \textbf{Bill:} \texttt{proven modulo 0}.}
%
\vnbprf{\vnbproofpart{\texttt{diff-on-implies-continuous}}}
%
\vnbprf{\textbf{Proof.} \texttt{theorem-\allowbreak{}library/diff-\allowbreak{}on-\allowbreak{}laws.scm} --- 201 recorded steps.}
%
\vnbprf{\textbf{Lemmas.} \texttt{diff-on-normed-field}, \texttt{diff-on-open}, \texttt{diff-on-in-fun}, \texttt{diff-on-pt-in}, \texttt{diff-on-deriv-in-carr}, \texttt{nf-open-subset-carr}, \texttt{nf-metric-space-is-metric-space}, \texttt{subspace-is-metric-space}, \texttt{subspace-pts}, \texttt{rr-pos-rr-in-rr}, \texttt{rr-one-in}, \texttt{nf-norm-in-rr}, \texttt{nf-norm-nonneg}, \texttt{rr-add-in-rr}, \texttt{rr-pos-rr-of-lt}, \texttt{rr-scale-eps}, \texttt{rr-lt-of-pos-rr}, \texttt{rr-min-pos}, \texttt{subset-mem-fwd}, \texttt{fun-apply-type-c}, \texttt{nf-neg-in-carr}, \texttt{nf-add-in-carr}, \texttt{metric-dist-real}, \texttt{subspace-dist}, \texttt{metric-sym}, \texttt{nf-norm-le-add}, \texttt{nf-norm-mult}, \texttt{rr-mul-le-right}, \texttt{rr-mul-in-rr}.}
%
\vnbprf{\textbf{Oracles.} \texttt{ineq}, \texttt{crs}, \texttt{arith}.  \textbf{Bill:} \texttt{proven modulo 0}.}
%
\end{proof}
%
\noindent\textbf{Deviation from the notes.} The notes state the ONE-SIDED form (right differentiable at x implies right continuous at x, and likewise on the left); the library has no one-sided derivative and no one-sided continuity (D2), so diff-implies-continuous is the two-sided statement only, and it requires f total on RR (D1). diff-on-implies-continuous is the same fact for is-diff-on over a normed field, where f in fun(u, carr(k)) with u OPEN, so it still does not reach a half-neighbourhood.
%
\renewcommand{\thethm}{2.5}
\begin{prop}[{Derivative of a sum}]\label{thm:deriv-sum}
%
\leavevmode
\vnbpart{\texttt{deriv-sum}}
\vnbstmt{\begin{array}{@{}l@{}} \forall f, g, a, l, m. \\ \quad \operatorname{is\text{-}diff\text{-}at}(f, a, l) \wedge \\ \quad \operatorname{is\text{-}diff\text{-}at}(g, a, m) \Rightarrow \\ \quad \quad \operatorname{is\text{-}diff\text{-}at}((\lambda x \in \mathbb{R}.\; (f(x) + g(x))), a, (l + m)) \end{array}}
\vnbpart{\texttt{diff-on-sum}}
\vnbstmt{\begin{array}{@{}l@{}} \forall k, u, f, g, a, l, m. \\ \quad \operatorname{is\text{-}normed\text{-}field}(k) \wedge \\ \quad \operatorname{is\text{-}diff\text{-}on}(k, u, f, a, l) \wedge \\ \quad \operatorname{is\text{-}diff\text{-}on}(k, u, g, a, m) \Rightarrow \\ \quad \quad \operatorname{is\text{-}diff\text{-}on}(k, u, (\lambda x \in u.\; \operatorname{add}(k)(f(x), g(x))), a, \operatorname{add}(k)(l, m)) \end{array}}
%
\end{prop}
%
\begin{proof}
%
The proofs: the file that holds each, the lemmas it cites, its oracles and its bill.
%
\vnbprf{\vnbproofpart{\texttt{deriv-sum}}}
%
\vnbprf{\textbf{Proof.} \texttt{theorem-\allowbreak{}library/deriv-\allowbreak{}sum-\allowbreak{}product.scm} --- 69 recorded steps.}
%
\vnbprf{\textbf{Lemmas.} \texttt{rr-add-closed}, \texttt{sum-lam-in-fun}, \texttt{sum-continuous-at}, \texttt{fun-apply-type-c}, \texttt{rr-sub-in-rr}, \texttt{eq-sym}.}
%
\vnbprf{\textbf{Oracles.} \texttt{crs}, \texttt{ineq}, \texttt{arith}.  \textbf{Bill:} \texttt{proven modulo 0}.}
%
\vnbprf{\vnbproofpart{\texttt{diff-on-sum}}}
%
\vnbprf{\textbf{Proof.} \texttt{theorem-\allowbreak{}library/diff-\allowbreak{}on-\allowbreak{}laws-\allowbreak{}2.scm} --- 173 recorded steps.}
%
\vnbprf{\textbf{Lemmas.} \texttt{diff-on-open}, \texttt{diff-on-pt-in}, \texttt{diff-on-in-fun}, \texttt{diff-on-deriv-in-carr}, \texttt{nf-metric-space-is-metric-space}, \texttt{nf-open-subset-carr}, \texttt{subspace-is-metric-space}, \texttt{subspace-pts}, \texttt{nf-open-is-set}, \texttt{subset-mem-fwd}, \texttt{nf-add-in-carr}, \texttt{nf-lam-sum-in-fun}, \texttt{nf-sum-continuous-on}, \texttt{fun-apply-type-c}, \texttt{commutative-ring-is-ring}, \texttt{normed-field-ring-view-carr}, \texttt{normed-field-ring-view-add}, \texttt{normed-field-ring-view-mul}, \texttt{normed-field-ring-view-neg}, \texttt{normed-field-ring-view-zero}, \texttt{normed-field-ring-view-one}.}
%
\vnbprf{\textbf{Oracles.} \texttt{crs}, \texttt{ineq}, \texttt{arith}.  \textbf{Bill:} \texttt{proven modulo 0}.}
%
\end{proof}
%
\noindent\textbf{Deviation from the notes.} The notes state the ONE-SIDED form; the library has only the two-sided derivative (D2). deriv-sum requires f and g total on RR (D1) and concludes about the term vnb-lambda(x, rr, f(x) + g(x)). diff-on-sum is the same law for is-diff-on over a normed field k with U open, which covers the notes' domain hypothesis for the two-sided case but not the one-sided one.
%
\renewcommand{\thethm}{2.6}
\begin{prop}[{Derivative of a product}]\label{thm:deriv-product}
%
\leavevmode
\vnbpart{\texttt{deriv-product}}
\vnbstmt{\begin{array}{@{}l@{}} \forall f, g, a, l, m. \\ \quad \operatorname{is\text{-}diff\text{-}at}(f, a, l) \wedge \\ \quad \operatorname{is\text{-}diff\text{-}at}(g, a, m) \Rightarrow \\ \quad \quad \operatorname{is\text{-}diff\text{-}at}((\lambda x \in \mathbb{R}.\; (f(x) \cdot g(x))), a, ((l \cdot g(a)) + (f(a) \cdot m))) \end{array}}
\vnbpart{\texttt{diff-on-product}}
\vnbstmt{\begin{array}{@{}l@{}} \forall k, u, f, g, a, l, m. \\ \quad \operatorname{is\text{-}normed\text{-}field}(k) \wedge \\ \quad \operatorname{is\text{-}diff\text{-}on}(k, u, f, a, l) \wedge \\ \quad \operatorname{is\text{-}diff\text{-}on}(k, u, g, a, m) \Rightarrow \\ \quad \quad \operatorname{is\text{-}diff\text{-}on}(k, u, (\lambda x \in u.\; \operatorname{mul}(k)(f(x), g(x))), a, \\ \quad \quad \quad \operatorname{add}(k)(\operatorname{mul}(k)(l, g(a)), \operatorname{mul}(k)(m, f(a)))) \end{array}}
%
\end{prop}
%
\begin{proof}
%
The proofs: the file that holds each, the lemmas it cites, its oracles and its bill.
%
\vnbprf{\vnbproofpart{\texttt{deriv-product}}}
%
\vnbprf{\textbf{Proof.} \texttt{theorem-\allowbreak{}library/deriv-\allowbreak{}sum-\allowbreak{}product.scm} --- 94 recorded steps.}
%
\vnbprf{\textbf{Lemmas.} \texttt{diff-implies-continuous}, \texttt{fun-apply-type-c}, \texttt{rr-mul-closed}, \texttt{rr-add-closed}, \texttt{const-lam-in-fun}, \texttt{prod-lam-in-fun}, \texttt{sum-lam-in-fun}, \texttt{const-continuous-at}, \texttt{product-continuous-at}, \texttt{sum-continuous-at}, \texttt{rr-sub-in-rr}, \texttt{eq-sym}.}
%
\vnbprf{\textbf{Oracles.} \texttt{crs}, \texttt{arith}, \texttt{ineq}.  \textbf{Bill:} \texttt{proven modulo 0}.}
%
\vnbprf{\vnbproofpart{\texttt{diff-on-product}}}
%
\vnbprf{\textbf{Proof.} \texttt{theorem-\allowbreak{}library/diff-\allowbreak{}on-\allowbreak{}laws-\allowbreak{}2.scm} --- 198 recorded steps.}
%
\vnbprf{\textbf{Lemmas.} \texttt{diff-on-open}, \texttt{diff-on-pt-in}, \texttt{diff-on-in-fun}, \texttt{diff-on-deriv-in-carr}, \texttt{diff-on-implies-continuous}, \texttt{nf-metric-space-is-metric-space}, \texttt{nf-open-subset-carr}, \texttt{subspace-is-metric-space}, \texttt{subspace-pts}, \texttt{nf-open-is-set}, \texttt{subset-mem-fwd}, \texttt{fun-apply-type-c}, \texttt{nf-mul-in-carr}, \texttt{nf-add-in-carr}, \texttt{nf-lam-prod-in-fun}, \texttt{nf-lam-const-z-in-fun}, \texttt{ms-const-continuous-on}, \texttt{nf-lam-prod-v-in-fun}, \texttt{nf-product-continuous-v}, \texttt{nf-lam-sum-in-fun}, \texttt{nf-sum-continuous-on}, \texttt{commutative-ring-is-ring}, \texttt{normed-field-ring-view-carr}, \texttt{normed-field-ring-view-add}, \texttt{normed-field-ring-view-mul}, \texttt{normed-field-ring-view-neg}, \texttt{normed-field-ring-view-zero}, \texttt{normed-field-ring-view-one}.}
%
\vnbprf{\textbf{Oracles.} \texttt{crs}, \texttt{ineq}, \texttt{arith}.  \textbf{Bill:} \texttt{proven modulo 0}.}
%
\end{proof}
%
\noindent\textbf{Deviation from the notes.} The notes state the ONE-SIDED form; the library has only the two-sided derivative (D2). deriv-product requires f and g total on RR (D1). diff-on-product is the same law for is-diff-on over a normed field k with U open. Note the notes give no hypothesis on g beyond right differentiability at x, and neither does the library.
%
\renewcommand{\thethm}{2.8}
\begin{prop}[{Chain Rule}]\label{thm:has-deriv-at-chain}
%
\leavevmode
\vnbpart{\texttt{has-deriv-at-chain}}
\vnbstmt{\begin{array}{@{}l@{}} \forall r, f, g, h, x, l, m. \\ \quad \operatorname{pos\text{-}rr}(r) \wedge \\ \quad \operatorname{restrict}(h, \operatorname{ooint}((x - r), (x + r))) \\ \quad \in\; (\operatorname{ooint}((x - r), (x + r)) \to \mathbb{R}) \wedge \\ \quad \forall y \in \operatorname{ooint}((x - r), (x + r)).\; (h(y) \simeq g(f(y))) \wedge \\ \quad \operatorname{has\text{-}deriv\text{-}at}(f, x, l) \wedge \\ \quad \operatorname{has\text{-}deriv\text{-}at}(g, f(x), m) \Rightarrow \\ \quad \quad \operatorname{has\text{-}deriv\text{-}at}(h, x, (m \cdot l)) \end{array}}
\vnbpart{\texttt{deriv-chain}}
\vnbstmt{\begin{array}{@{}l@{}} \forall f, g, a, l, m. \\ \quad \operatorname{is\text{-}diff\text{-}at}(f, a, l) \wedge \\ \quad \operatorname{is\text{-}diff\text{-}at}(g, f(a), m) \Rightarrow \\ \quad \quad \operatorname{is\text{-}diff\text{-}at}(\operatorname{compose}(g, f), a, (m \cdot l)) \end{array}}
\vnbpart{\texttt{diff-on-chain}}
\vnbstmt{\begin{array}{@{}l@{}} \forall k, u, s, f, g, a, l, m. \\ \quad \operatorname{is\text{-}normed\text{-}field}(k) \wedge \\ \quad \operatorname{is\text{-}diff\text{-}on}(k, u, f, a, l) \wedge \\ \quad \operatorname{is\text{-}diff\text{-}on}(k, s, g, f(a), m) \wedge \\ \quad \forall y \in u.\; (f(y) \in s) \Rightarrow \\ \quad \quad \operatorname{is\text{-}diff\text{-}on}(k, u, (\lambda x \in u.\; g(f(x))), a, \operatorname{mul}(k)(m, l)) \end{array}}
%
\end{prop}
%
\begin{proof}
%
The proofs: the file that holds each, the lemmas it cites, its oracles and its bill.
%
\vnbprf{\vnbproofpart{\texttt{has-deriv-at-chain}}}
%
\vnbprf{\textbf{Proof.} \texttt{theorem-\allowbreak{}library/has-\allowbreak{}deriv-\allowbreak{}at-\allowbreak{}chain.scm} --- 126 recorded steps.}
%
\vnbprf{\textbf{Lemmas.} \texttt{rr-is-normed-field}, \texttt{rr-nf-carr}, \texttt{has-deriv-at-value-in-rr}, \texttt{has-deriv-at-pt-in-rr}, \texttt{has-deriv-at-in-rr}, \texttt{rr-pos-rr-in-rr}, \texttt{rr-lt-of-pos-rr}, \texttt{rr-one-in}, \texttt{rr-pos-rr-of-lt}, \texttt{has-deriv-at-shrink}, \texttt{rr-sub-in-rr}, \texttt{rr-add-in-rr}, \texttt{has-deriv-at-maps-into}, \texttt{diff-on-ooint-shrink}, \texttt{ooint-center}, \texttt{restrict-apply}, \texttt{rr-nf-mul-apply}, \texttt{diff-on-chain}, \texttt{subset-def}, \texttt{ooint-membership}, \texttt{restrict-in-fun-of-restrict}, \texttt{subset-mem-fwd}, \texttt{diff-on-transfer-ptwise-eq}, \texttt{has-deriv-at-intro}.}
%
\vnbprf{\textbf{Oracles.} \texttt{ineq}, \texttt{arith}, \texttt{crs}.  \textbf{Bill:} \texttt{proven modulo 0}.}
%
\vnbprf{\vnbproofpart{\texttt{deriv-chain}}}
%
\vnbprf{\textbf{Proof.} \texttt{theorem-\allowbreak{}library/chain-\allowbreak{}rule.scm} --- 100 recorded steps.}
%
\vnbprf{\textbf{Lemmas.} \texttt{diff-implies-continuous}, \texttt{rr-is-set}, \texttt{compose-type}, \texttt{compose-apply}, \texttt{compose-continuous-at}, \texttt{prod-lam-in-fun}, \texttt{product-continuous-at}, \texttt{rr-mul-closed}, \texttt{fun-apply-type-c}.}
%
\vnbprf{\textbf{Oracles.} \texttt{crs}, \texttt{arith}, \texttt{ineq}.  \textbf{Bill:} \texttt{proven modulo 0}.}
%
\vnbprf{\vnbproofpart{\texttt{diff-on-chain}}}
%
\vnbprf{\textbf{Proof.} \texttt{theorem-\allowbreak{}library/diff-\allowbreak{}on-\allowbreak{}laws-\allowbreak{}2.scm} --- 164 recorded steps.}
%
\vnbprf{\textbf{Lemmas.} \texttt{diff-on-open}, \texttt{diff-on-pt-in}, \texttt{diff-on-in-fun}, \texttt{diff-on-deriv-in-carr}, \texttt{diff-on-implies-continuous}, \texttt{nf-metric-space-is-metric-space}, \texttt{nf-open-subset-carr}, \texttt{subspace-is-metric-space}, \texttt{subspace-pts}, \texttt{nf-open-is-set}, \texttt{subset-mem-fwd}, \texttt{nf-mul-in-carr}, \texttt{nf-lam-comp-z-in-fun}, \texttt{ms-corestrict-continuous}, \texttt{ms-compose-continuous-on}, \texttt{nf-lam-prod-v-in-fun}, \texttt{nf-product-continuous-v}, \texttt{fun-apply-type-c}, \texttt{commutative-ring-is-ring}, \texttt{normed-field-ring-view-carr}, \texttt{normed-field-ring-view-add}, \texttt{normed-field-ring-view-mul}, \texttt{normed-field-ring-view-neg}, \texttt{normed-field-ring-view-zero}, \texttt{normed-field-ring-view-one}.}
%
\vnbprf{\textbf{Oracles.} \texttt{crs}, \texttt{ineq}, \texttt{arith}.  \textbf{Bill:} \texttt{proven modulo 0}.}
%
\end{proof}
%
\noindent\textbf{Deviation from the notes.} Restated 2026-09-22 for the LOCAL derivative (the first name): f differentiable at x, g differentiable at f(x), as in the notes; that f maps a neighbourhood of x into the domain of g is PROVEN from the two hypotheses, not assumed. The composite g o f is written as any function h that agrees with t |-$>$ g(f(t)) on an open interval about x. The second name is the older form for functions TOTAL on RR; the third is the rule over a normed field on open sets.
%
\renewcommand{\thethm}{2.10}
\begin{prop}[{Vanishing of the derivative at an interior extremum}]\label{thm:interval-local-max-right-deriv-nonpos}
%
\leavevmode
\vnbpart{\texttt{interval-local-max-right-deriv-nonpos}}
\vnbstmt{\begin{array}{@{}l@{}} \forall a, b \in \mathbb{R}, f, x, l. \\ \quad \operatorname{is\text{-}local\text{-}max\text{-}at}(f, \operatorname{ccint}(a, b), x) \wedge \\ \quad \operatorname{has\text{-}right\text{-}deriv\text{-}at}(f, x, l) \Rightarrow \\ \quad \quad (l \leq 0) \end{array}}
\vnbpart{\texttt{interval-local-max-left-deriv-nonneg}}
\vnbstmt{\begin{array}{@{}l@{}} \forall a, b \in \mathbb{R}, f, x, l. \\ \quad \operatorname{is\text{-}local\text{-}max\text{-}at}(f, \operatorname{ccint}(a, b), x) \wedge \\ \quad \operatorname{has\text{-}left\text{-}deriv\text{-}at}(f, x, l) \Rightarrow \\ \quad \quad (0 \leq l) \end{array}}
\vnbpart{\texttt{interval-local-max-deriv-zero}}
\vnbstmt{\begin{array}{@{}l@{}} \forall a, b \in \mathbb{R}, f.; \forall c \in \operatorname{ooint}(a, b), l. \\ \quad \operatorname{is\text{-}local\text{-}max\text{-}at}(f, \operatorname{ccint}(a, b), c) \wedge \\ \quad \operatorname{has\text{-}deriv\text{-}at}(f, c, l) \Rightarrow \\ \quad \quad (l = 0) \end{array}}
\vnbpart{\texttt{interior-max-deriv-zero}}
\vnbstmt{\begin{array}{@{}l@{}} \forall f, a, b, c, l. \\ \quad (f \in (\mathbb{R} \to \mathbb{R})) \wedge \\ \quad (a \in \mathbb{R}) \wedge \\ \quad (b \in \mathbb{R}) \wedge \\ \quad (c \in \mathbb{R}) \wedge \\ \quad (a < c) \wedge \\ \quad (c < b) \wedge \\ \quad \forall x \in \operatorname{ccint}(a, b).\; (f(x) \leq f(c)) \wedge \\ \quad \operatorname{is\text{-}diff\text{-}at}(f, c, l) \Rightarrow \\ \quad \quad (l = 0) \end{array}}
%
\end{prop}
%
\begin{proof}
%
The proofs: the file that holds each, the lemmas it cites, its oracles and its bill.
%
\vnbprf{\vnbproofpart{\texttt{interval-local-max-right-deriv-nonpos}}}
%
\vnbprf{\textbf{Proof.} \texttt{theorem-\allowbreak{}library/one-\allowbreak{}sided-\allowbreak{}derivative-\allowbreak{}laws.scm} --- 5 recorded steps.}
%
\vnbprf{\textbf{Lemmas.} \texttt{local-max-right-deriv-nonpos}.}
%
\vnbprf{\textbf{Oracles.} \texttt{ineq}, \texttt{arith}, \texttt{crs}.  \textbf{Bill:} \texttt{proven modulo 0}.}
%
\vnbprf{\vnbproofpart{\texttt{interval-local-max-left-deriv-nonneg}}}
%
\vnbprf{\textbf{Proof.} \texttt{theorem-\allowbreak{}library/one-\allowbreak{}sided-\allowbreak{}derivative-\allowbreak{}laws.scm} --- 5 recorded steps.}
%
\vnbprf{\textbf{Lemmas.} \texttt{local-max-left-deriv-nonneg}.}
%
\vnbprf{\textbf{Oracles.} \texttt{ineq}, \texttt{arith}, \texttt{crs}.  \textbf{Bill:} \texttt{proven modulo 0}.}
%
\vnbprf{\vnbproofpart{\texttt{interval-local-max-deriv-zero}}}
%
\vnbprf{\textbf{Proof.} \texttt{theorem-\allowbreak{}library/interval-\allowbreak{}extremum.scm} --- 141 recorded steps.}
%
\vnbprf{\textbf{Lemmas.} \texttt{ooint-membership}, \texttt{rr-pos-rr-in-rr}, \texttt{rr-lt-of-pos-rr}, \texttt{rr-sub-in-rr}, \texttt{rr-min-pos}, \texttt{rr-pos-rr-of-lt}, \texttt{rr-pos-halvable}, \texttt{rr-add-in-rr}, \texttt{extend-const-in-fun}, \texttt{extend-const-deriv-fwd}, \texttt{extend-const-ccint-fixes}, \texttt{ccint-membership}, \texttt{rr-abs-bound}, \texttt{interior-max-deriv-zero}.}
%
\vnbprf{\textbf{Oracles.} \texttt{ineq}, \texttt{crs}, \texttt{arith}.  \textbf{Bill:} \texttt{proven modulo 0}.}
%
\vnbprf{\vnbproofpart{\texttt{interior-max-deriv-zero}}}
%
\vnbprf{\textbf{Proof.} \texttt{theorem-\allowbreak{}library/interior-\allowbreak{}extremum-\allowbreak{}proof.scm} --- 221 recorded steps.}
%
\vnbprf{\textbf{Lemmas.} \texttt{fun-apply-type-c}, \texttt{rr-sub-in-rr}, \texttt{rr-lt-trans}, \texttt{rr-lt-implies-le}, \texttt{rr-lt-diff-pos}, \texttt{ccint-membership}, \texttt{rr-le-diff-nonpos}, \texttt{rr-prod-nonpos-pos}, \texttt{rr-lt-diff-neg}, \texttt{rr-prod-nonpos-neg}, \texttt{rr-zero-in}, \texttt{continuous-nonpos-right}, \texttt{continuous-nonneg-left}, \texttt{rr-leq-antisymmetric}.}
%
\vnbprf{\textbf{Oracles.} \texttt{ineq}, \texttt{crs}, \texttt{arith}.  \textbf{Bill:} \texttt{proven modulo 0}.}
%
\end{proof}
%
\noindent\textbf{Deviation from the notes.} Proven 2026-09-22 as the notes state it, sentence by sentence (the first three names): at a LOCAL maximum of a function ON [a, b], a right derivative is $<$= 0, a left derivative is $>$= 0, and at an interior point a derivative is 0. The one-sided derivatives are has-right-deriv-at and has-left-deriv-at (the notes' (5) and (6), in Caratheodory form on a half-window; the limit form of (5)/(6) is not stated: the library has no limit of a function given only on a set). The notes' Definition 2.9 as printed has the two inequalities exchanged; is-local-max-at uses f(x) $<$= f(t), the sense the notes' proof uses. The last name is the older form: f TOTAL on RR and a GLOBAL maximum on [a, b].
%
\renewcommand{\thethm}{2.11}
\begin{thm}[{Generalized Mean Value Theorem}]\label{thm:generalized-mvt-on-interval}
%
\leavevmode
\vnbpart{\texttt{generalized-mvt-on-interval}}
\vnbstmt{\begin{array}{@{}l@{}} \forall a, b \in \mathbb{R}, f, g. \\ \quad (a < b) \wedge \\ \quad (f \in (\operatorname{ccint}(a, b) \to \mathbb{R})) \wedge \\ \quad (g \in (\operatorname{ccint}(a, b) \to \mathbb{R})) \wedge \\ \quad \operatorname{is\text{-}continuous\text{-}on}(f, \operatorname{ccint}(a, b)) \wedge \\ \quad \operatorname{is\text{-}continuous\text{-}on}(g, \operatorname{ccint}(a, b)) \wedge \\ \quad \forall x \in \operatorname{ooint}(a, b). \exists l. \\ \quad \quad \operatorname{has\text{-}deriv\text{-}at}(f, x, l) \wedge \\ \quad \forall x \in \operatorname{ooint}(a, b). \exists m. \\ \quad \quad \operatorname{has\text{-}deriv\text{-}at}(g, x, m) \Rightarrow \\ \quad \quad \exists c \in \operatorname{ooint}(a, b), l, m. \\ \quad \quad \quad \operatorname{has\text{-}deriv\text{-}at}(f, c, l) \wedge \\ \quad \quad \quad \operatorname{has\text{-}deriv\text{-}at}(g, c, m) \wedge \\ \quad \quad \quad ((l \cdot (g(b) - g(a))) = (m \cdot (f(b) - f(a)))) \end{array}}
\vnbpart{\texttt{generalized-mvt}}
\vnbstmt{\begin{array}{@{}l@{}} \forall f, g, a, b. \\ \quad (f \in (\mathbb{R} \to \mathbb{R})) \wedge \\ \quad (g \in (\mathbb{R} \to \mathbb{R})) \wedge \\ \quad (a \in \mathbb{R}) \wedge \\ \quad (b \in \mathbb{R}) \wedge \\ \quad (a < b) \wedge \\ \quad \forall x \in \operatorname{ccint}(a, b). \\ \quad \quad \operatorname{is\text{-}continuous\text{-}at}(\operatorname{rr\text{-}ms}, \operatorname{rr\text{-}ms}, f, x) \wedge \\ \quad \quad \operatorname{is\text{-}continuous\text{-}at}(\operatorname{rr\text{-}ms}, \operatorname{rr\text{-}ms}, g, x) \wedge \\ \quad \forall x \in \mathbb{R}. \\ \quad \quad ((a < x) \wedge (x < b)) \Rightarrow  \\ \quad \quad \quad \exists l.\; \operatorname{is\text{-}diff\text{-}at}(f, x, l) \wedge \\ \quad \quad \quad \exists m.\; \operatorname{is\text{-}diff\text{-}at}(g, x, m) \Rightarrow \\ \quad \quad \exists c \in \mathbb{R}. \\ \quad \quad \quad (a < c) \wedge \\ \quad \quad \quad (c < b) \wedge \\ \quad \quad \quad \exists l, m. \\ \quad \quad \quad \quad \operatorname{is\text{-}diff\text{-}at}(f, c, l) \wedge \\ \quad \quad \quad \quad \operatorname{is\text{-}diff\text{-}at}(g, c, m) \wedge \\ \quad \quad \quad \quad ((l \cdot (g(b) - g(a))) = (m \cdot (f(b) - f(a)))) \end{array}}
%
\end{thm}
%
\begin{proof}
%
The proofs: the file that holds each, the lemmas it cites, its oracles and its bill.
%
\vnbprf{\vnbproofpart{\texttt{generalized-mvt-on-interval}}}
%
\vnbprf{\textbf{Proof.} \texttt{theorem-\allowbreak{}library/interval-\allowbreak{}mvt.scm} --- 144 recorded steps.}
%
\vnbprf{\textbf{Lemmas.} \texttt{extend-const-in-fun}, \texttt{ccint-elt-in-rr}, \texttt{extend-const-continuous-at}, \texttt{ooint-membership}, \texttt{extend-const-deriv-fwd}, \texttt{extend-const-fixes}, \texttt{generalized-mvt}, \texttt{extend-const-deriv-bwd}.}
%
\vnbprf{\textbf{Oracles.} \texttt{ineq}, \texttt{arith}, \texttt{crs}.  \textbf{Bill:} \texttt{proven modulo 0}.}
%
\vnbprf{\vnbproofpart{\texttt{generalized-mvt}}}
%
\vnbprf{\textbf{Proof.} \texttt{theorem-\allowbreak{}library/generalized-\allowbreak{}mvt-\allowbreak{}proof.scm} --- 305 recorded steps.}
%
\vnbprf{\textbf{Lemmas.} \texttt{rr-is-set}, \texttt{fun-apply-type-c}, \texttt{rr-sub-in-rr}, \texttt{rr-mul-closed}, \texttt{gmvt-aux-cont}, \texttt{gmvt-aux-diff}, \texttt{rolle}, \texttt{diff-value-real}, \texttt{derivative-unique}, \texttt{rr-diff-zero-eq}.}
%
\vnbprf{\textbf{Oracles.} \texttt{crs}, \texttt{arith}, \texttt{ineq}.  \textbf{Bill:} \texttt{proven modulo 0}.}
%
\end{proof}
%
\noindent\textbf{Deviation from the notes.} Restated 2026-09-21 for functions ON [a, b] (the first name). The derivative is the local relation has-deriv-at(f, x, l), so f'(theta) appears as a value l. The second name is the older form for functions TOTAL on RR with two-sided continuity at the endpoints; it remains as the lemma behind the first.
%
\renewcommand{\thethm}{2.12}
\begin{lem}[{Rolle's Lemma}]\label{thm:rolle-on-interval}
%
\leavevmode
\vnbpart{\texttt{rolle-on-interval}}
\vnbstmt{\begin{array}{@{}l@{}} \forall a, b \in \mathbb{R}, h \in (\operatorname{ccint}(a, b) \to \mathbb{R}). \\ \quad (a < b) \wedge \\ \quad \operatorname{is\text{-}continuous\text{-}on}(h, \operatorname{ccint}(a, b)) \wedge \\ \quad \forall x \in \operatorname{ooint}(a, b). \exists l. \\ \quad \quad \operatorname{has\text{-}deriv\text{-}at}(h, x, l) \wedge \\ \quad (h(a) = h(b)) \Rightarrow \\ \quad \quad \exists c \in \operatorname{ooint}(a, b). \\ \quad \quad \quad \operatorname{has\text{-}deriv\text{-}at}(h, c, 0) \end{array}}
\vnbpart{\texttt{rolle}}
\vnbstmt{\begin{array}{@{}l@{}} \forall h, a, b. \\ \quad (h \in (\mathbb{R} \to \mathbb{R})) \wedge \\ \quad (a \in \mathbb{R}) \wedge \\ \quad (b \in \mathbb{R}) \wedge \\ \quad (a < b) \wedge \\ \quad \forall x \in \operatorname{ccint}(a, b). \\ \quad \quad \operatorname{is\text{-}continuous\text{-}at}(\operatorname{rr\text{-}ms}, \operatorname{rr\text{-}ms}, h, x) \wedge \\ \quad \forall x \in \mathbb{R}. \\ \quad \quad ((a < x) \wedge (x < b)) \Rightarrow  \\ \quad \quad \quad \exists l.\; \operatorname{is\text{-}diff\text{-}at}(h, x, l) \wedge \\ \quad (h(a) = h(b)) \Rightarrow \\ \quad \quad \exists c \in \mathbb{R}. \\ \quad \quad \quad ((a < c) \wedge ((c < b) \wedge \operatorname{is\text{-}diff\text{-}at}(h, c, 0))) \end{array}}
%
\end{lem}
%
\begin{proof}
%
The proofs: the file that holds each, the lemmas it cites, its oracles and its bill.
%
\vnbprf{\vnbproofpart{\texttt{rolle-on-interval}}}
%
\vnbprf{\textbf{Proof.} \texttt{theorem-\allowbreak{}library/interval-\allowbreak{}mvt.scm} --- 106 recorded steps.}
%
\vnbprf{\textbf{Lemmas.} \texttt{extend-const-in-fun}, \texttt{ccint-elt-in-rr}, \texttt{extend-const-continuous-at}, \texttt{ooint-membership}, \texttt{extend-const-deriv-fwd}, \texttt{extend-const-fixes}, \texttt{rolle}, \texttt{extend-const-deriv-bwd}.}
%
\vnbprf{\textbf{Oracles.} \texttt{ineq}, \texttt{arith}, \texttt{crs}.  \textbf{Bill:} \texttt{proven modulo 0}.}
%
\vnbprf{\vnbproofpart{\texttt{rolle}}}
%
\vnbprf{\textbf{Proof.} \texttt{theorem-\allowbreak{}library/rolle-\allowbreak{}proof.scm} --- 354 recorded steps.}
%
\vnbprf{\textbf{Lemmas.} \texttt{rr-lt-implies-le}, \texttt{extreme-value-max}, \texttt{extreme-value-min}, \texttt{ccint-membership}, \texttt{rr-leq-reflexive}, \texttt{fun-apply-type-c}, \texttt{rr-le-cases}, \texttt{max-val-strict-interior}, \texttt{interior-max-deriv-zero}, \texttt{min-val-strict-interior}, \texttt{interior-min-deriv-zero}, \texttt{rr-midpoint-between}, \texttt{rr-leq-antisymmetric}.}
%
\vnbprf{\textbf{Oracles.} \texttt{ineq}, \texttt{arith}, \texttt{crs}.  \textbf{Bill:} \texttt{proven modulo 0}.}
%
\end{proof}
%
\noindent\textbf{Deviation from the notes.} Restated 2026-09-21 for functions ON [a, b] (the first name). The derivative is the local relation has-deriv-at(f, x, l), so f'(theta) appears as a value l. The second name is the older form for functions TOTAL on RR with two-sided continuity at the endpoints; it remains as the lemma behind the first.
%
\renewcommand{\thethm}{2.13}
\begin{thm}[{Mean Value Theorem}]\label{thm:mvt-on-interval}
%
\leavevmode
\vnbpart{\texttt{mvt-on-interval}}
\vnbstmt{\begin{array}{@{}l@{}} \forall a, b \in \mathbb{R}, f \in (\operatorname{ccint}(a, b) \to \mathbb{R}). \\ \quad (a < b) \wedge \\ \quad \operatorname{is\text{-}continuous\text{-}on}(f, \operatorname{ccint}(a, b)) \wedge \\ \quad \forall x \in \operatorname{ooint}(a, b). \exists l. \\ \quad \quad \operatorname{has\text{-}deriv\text{-}at}(f, x, l) \Rightarrow \\ \quad \quad \exists c \in \operatorname{ooint}(a, b), l. \\ \quad \quad \quad \operatorname{has\text{-}deriv\text{-}at}(f, c, l) \wedge \\ \quad \quad \quad (l = ((f(b) - f(a)) \cdot \operatorname{recip}((b - a)))) \end{array}}
\vnbpart{\texttt{mvt}}
\vnbstmt{\begin{array}{@{}l@{}} \forall f, a, b. \\ \quad (f \in (\mathbb{R} \to \mathbb{R})) \wedge \\ \quad (a \in \mathbb{R}) \wedge \\ \quad (b \in \mathbb{R}) \wedge \\ \quad (a < b) \wedge \\ \quad \forall x \in \operatorname{ccint}(a, b). \\ \quad \quad \operatorname{is\text{-}continuous\text{-}at}(\operatorname{rr\text{-}ms}, \operatorname{rr\text{-}ms}, f, x) \wedge \\ \quad \forall x \in \mathbb{R}. \\ \quad \quad ((a < x) \wedge (x < b)) \Rightarrow  \\ \quad \quad \quad \exists l.\; \operatorname{is\text{-}diff\text{-}at}(f, x, l) \Rightarrow \\ \quad \quad \exists c \in \mathbb{R}. \\ \quad \quad \quad (a < c) \wedge \\ \quad \quad \quad (c < b) \wedge \\ \quad \quad \quad \exists l. \\ \quad \quad \quad \quad (\operatorname{is\text{-}diff\text{-}at}(f, c, l) \wedge ((l \cdot (b - a)) = (f(b) - f(a)))) \end{array}}
%
\end{thm}
%
\begin{proof}
%
The proofs: the file that holds each, the lemmas it cites, its oracles and its bill.
%
\vnbprf{\vnbproofpart{\texttt{mvt-on-interval}}}
%
\vnbprf{\textbf{Proof.} \texttt{theorem-\allowbreak{}library/interval-\allowbreak{}mvt.scm} --- 171 recorded steps.}
%
\vnbprf{\textbf{Lemmas.} \texttt{extend-const-in-fun}, \texttt{ccint-elt-in-rr}, \texttt{extend-const-continuous-at}, \texttt{ooint-membership}, \texttt{extend-const-deriv-fwd}, \texttt{extend-const-fixes}, \texttt{ccint-membership}, \texttt{fun-apply-type-c}, \texttt{rr-sub-in-rr}, \texttt{rr-pos-ne-zero}, \texttt{rr-recip-closed}, \texttt{rr-mul-comm}, \texttt{mvt}, \texttt{extend-const-deriv-bwd}, \texttt{has-deriv-at-in-rr}, \texttt{rr-recip-solve}.}
%
\vnbprf{\textbf{Oracles.} \texttt{ineq}, \texttt{crs}, \texttt{arith}.  \textbf{Bill:} \texttt{proven modulo 0}.}
%
\vnbprf{\vnbproofpart{\texttt{mvt}}}
%
\vnbprf{\textbf{Proof.} \texttt{theorem-\allowbreak{}library/mvt-\allowbreak{}proof.scm} --- 258 recorded steps.}
%
\vnbprf{\textbf{Lemmas.} \texttt{rr-is-set}, \texttt{fun-apply-type-c}, \texttt{rr-sub-in-rr}, \texttt{rr-mul-closed}, \texttt{mvt-aux-cont}, \texttt{mvt-aux-diff}, \texttt{rolle}, \texttt{diff-value-real}, \texttt{derivative-unique}, \texttt{rr-diff-zero-eq}.}
%
\vnbprf{\textbf{Oracles.} \texttt{crs}, \texttt{arith}, \texttt{ineq}.  \textbf{Bill:} \texttt{proven modulo 0}.}
%
\end{proof}
%
\noindent\textbf{Deviation from the notes.} Restated 2026-09-21 for functions ON [a, b] (the first name). The derivative is the local relation has-deriv-at(f, x, l), so f'(theta) appears as a value l. The second name is the older form for functions TOTAL on RR with two-sided continuity at the endpoints; it remains as the lemma behind the first. The quotient of the notes is stored as (f(b) - f(a)) * recip(b - a): the reader desugars `/' when it parses.
%
\renewcommand{\thethm}{2.14}
\begin{cor}[{Increment bounds from derivative bounds}]\label{thm:mvt-upper-bound-on-interval}
%
\leavevmode
\vnbpart{\texttt{mvt-upper-bound-on-interval}}
\vnbstmt{\begin{array}{@{}l@{}} \forall a, b, m \in \mathbb{R}, f \in (\operatorname{ccint}(a, b) \to \mathbb{R}). \\ \quad (a < b) \wedge \\ \quad \operatorname{is\text{-}continuous\text{-}on}(f, \operatorname{ccint}(a, b)) \wedge \\ \quad \forall x \in \operatorname{ooint}(a, b). \exists l. \\ \quad \quad (\operatorname{has\text{-}deriv\text{-}at}(f, x, l) \wedge (l \leq m)) \Rightarrow \\ \quad \quad ((f(b) - f(a)) \leq (m \cdot (b - a))) \end{array}}
\vnbpart{\texttt{mvt-lower-bound-on-interval}}
\vnbstmt{\begin{array}{@{}l@{}} \forall a, b, m \in \mathbb{R}, f \in (\operatorname{ccint}(a, b) \to \mathbb{R}). \\ \quad (a < b) \wedge \\ \quad \operatorname{is\text{-}continuous\text{-}on}(f, \operatorname{ccint}(a, b)) \wedge \\ \quad \forall x \in \operatorname{ooint}(a, b). \exists l. \\ \quad \quad (\operatorname{has\text{-}deriv\text{-}at}(f, x, l) \wedge (m \leq l)) \Rightarrow \\ \quad \quad ((m \cdot (b - a)) \leq (f(b) - f(a))) \end{array}}
\vnbpart{\texttt{mvt-abs-bound-on-interval}}
\vnbstmt{\begin{array}{@{}l@{}} \forall a, b, m \in \mathbb{R}, f \in (\operatorname{ccint}(a, b) \to \mathbb{R}). \\ \quad (a < b) \wedge \\ \quad \operatorname{is\text{-}continuous\text{-}on}(f, \operatorname{ccint}(a, b)) \wedge \\ \quad \forall x \in \operatorname{ooint}(a, b). \exists l. \\ \quad \quad (\operatorname{has\text{-}deriv\text{-}at}(f, x, l) \wedge (\lvert l \rvert \leq m)) \Rightarrow \\ \quad \quad (\lvert (f(b) - f(a)) \rvert \leq (m \cdot (b - a))) \end{array}}
\vnbpart{\texttt{mvt-upper-bound}}
\vnbstmt{\begin{array}{@{}l@{}} \forall f, a, b, m. \\ \quad (f \in (\mathbb{R} \to \mathbb{R})) \wedge \\ \quad (a \in \mathbb{R}) \wedge \\ \quad (b \in \mathbb{R}) \wedge \\ \quad (m \in \mathbb{R}) \wedge \\ \quad (a < b) \wedge \\ \quad \forall x \in \operatorname{ccint}(a, b). \\ \quad \quad \operatorname{is\text{-}continuous\text{-}at}(\operatorname{rr\text{-}ms}, \operatorname{rr\text{-}ms}, f, x) \wedge \\ \quad \forall x \in \mathbb{R}. \\ \quad \quad ((a < x) \wedge (x < b)) \Rightarrow  \\ \quad \quad \quad \exists l. \\ \quad \quad \quad \quad (\operatorname{is\text{-}diff\text{-}at}(f, x, l) \wedge (l \leq m)) \Rightarrow \\ \quad \quad ((f(b) - f(a)) \leq (m \cdot (b - a))) \end{array}}
\vnbpart{\texttt{mvt-lower-bound}}
\vnbstmt{\begin{array}{@{}l@{}} \forall f, a, b, m. \\ \quad (f \in (\mathbb{R} \to \mathbb{R})) \wedge \\ \quad (a \in \mathbb{R}) \wedge \\ \quad (b \in \mathbb{R}) \wedge \\ \quad (m \in \mathbb{R}) \wedge \\ \quad (a < b) \wedge \\ \quad \forall x \in \operatorname{ccint}(a, b). \\ \quad \quad \operatorname{is\text{-}continuous\text{-}at}(\operatorname{rr\text{-}ms}, \operatorname{rr\text{-}ms}, f, x) \wedge \\ \quad \forall x \in \mathbb{R}. \\ \quad \quad ((a < x) \wedge (x < b)) \Rightarrow  \\ \quad \quad \quad \exists l. \\ \quad \quad \quad \quad (\operatorname{is\text{-}diff\text{-}at}(f, x, l) \wedge (m \leq l)) \Rightarrow \\ \quad \quad ((m \cdot (b - a)) \leq (f(b) - f(a))) \end{array}}
\vnbpart{\texttt{mvt-abs-bound}}
\vnbstmt{\begin{array}{@{}l@{}} \forall f, a, b, m. \\ \quad (f \in (\mathbb{R} \to \mathbb{R})) \wedge \\ \quad (a \in \mathbb{R}) \wedge \\ \quad (b \in \mathbb{R}) \wedge \\ \quad (m \in \mathbb{R}) \wedge \\ \quad (a < b) \wedge \\ \quad \forall x \in \operatorname{ccint}(a, b). \\ \quad \quad \operatorname{is\text{-}continuous\text{-}at}(\operatorname{rr\text{-}ms}, \operatorname{rr\text{-}ms}, f, x) \wedge \\ \quad \forall x \in \mathbb{R}. \\ \quad \quad ((a < x) \wedge (x < b)) \Rightarrow  \\ \quad \quad \quad \exists l \in \mathbb{R}. \\ \quad \quad \quad \quad (\operatorname{is\text{-}diff\text{-}at}(f, x, l) \wedge (\lvert l \rvert \leq m)) \Rightarrow \\ \quad \quad (\lvert (f(b) - f(a)) \rvert \leq (m \cdot (b - a))) \end{array}}
%
\end{cor}
%
\begin{proof}
%
The proofs: the file that holds each, the lemmas it cites, its oracles and its bill.
%
\vnbprf{\vnbproofpart{\texttt{mvt-upper-bound-on-interval}}}
%
\vnbprf{\textbf{Proof.} \texttt{theorem-\allowbreak{}library/interval-\allowbreak{}mvt.scm} --- 86 recorded steps.}
%
\vnbprf{\textbf{Lemmas.} \texttt{extend-const-in-fun}, \texttt{ccint-elt-in-rr}, \texttt{extend-const-continuous-at}, \texttt{ooint-membership}, \texttt{has-deriv-at-in-rr}, \texttt{extend-const-deriv-fwd}, \texttt{extend-const-fixes}, \texttt{mvt-upper-bound}.}
%
\vnbprf{\textbf{Oracles.} \texttt{ineq}, \texttt{arith}, \texttt{crs}.  \textbf{Bill:} \texttt{proven modulo 0}.}
%
\vnbprf{\vnbproofpart{\texttt{mvt-lower-bound-on-interval}}}
%
\vnbprf{\textbf{Proof.} \texttt{theorem-\allowbreak{}library/interval-\allowbreak{}mvt.scm} --- 86 recorded steps.}
%
\vnbprf{\textbf{Lemmas.} \texttt{extend-const-in-fun}, \texttt{ccint-elt-in-rr}, \texttt{extend-const-continuous-at}, \texttt{ooint-membership}, \texttt{has-deriv-at-in-rr}, \texttt{extend-const-deriv-fwd}, \texttt{extend-const-fixes}, \texttt{mvt-lower-bound}.}
%
\vnbprf{\textbf{Oracles.} \texttt{ineq}, \texttt{arith}, \texttt{crs}.  \textbf{Bill:} \texttt{proven modulo 0}.}
%
\vnbprf{\vnbproofpart{\texttt{mvt-abs-bound-on-interval}}}
%
\vnbprf{\textbf{Proof.} \texttt{theorem-\allowbreak{}library/interval-\allowbreak{}mvt.scm} --- 66 recorded steps.}
%
\vnbprf{\textbf{Lemmas.} \texttt{extend-const-in-fun}, \texttt{ccint-elt-in-rr}, \texttt{extend-const-continuous-at}, \texttt{ooint-membership}, \texttt{has-deriv-at-in-rr}, \texttt{extend-const-deriv-fwd}, \texttt{extend-const-fixes}, \texttt{mvt-abs-bound}.}
%
\vnbprf{\textbf{Oracles.} \texttt{ineq}, \texttt{arith}, \texttt{crs}.  \textbf{Bill:} \texttt{proven modulo 0}.}
%
\vnbprf{\vnbproofpart{\texttt{mvt-upper-bound}}}
%
\vnbprf{\textbf{Proof.} \texttt{theorem-\allowbreak{}library/mvt-\allowbreak{}bounds-\allowbreak{}proof.scm} --- 103 recorded steps.}
%
\vnbprf{\textbf{Lemmas.} \texttt{mvt}, \texttt{derivative-unique}, \texttt{rr-sub-in-rr}, \texttt{rr-zero-in}, \texttt{rr-lt-diff-pos}, \texttt{rr-lt-implies-le}, \texttt{diff-value-real}, \texttt{rr-le-scale-nonneg-right}.}
%
\vnbprf{\textbf{Oracles.} \texttt{crs}, \texttt{arith}, \texttt{ineq}.  \textbf{Bill:} \texttt{proven modulo 0}.}
%
\vnbprf{\vnbproofpart{\texttt{mvt-lower-bound}}}
%
\vnbprf{\textbf{Proof.} \texttt{theorem-\allowbreak{}library/mvt-\allowbreak{}bounds-\allowbreak{}proof.scm} --- 103 recorded steps.}
%
\vnbprf{\textbf{Lemmas.} \texttt{mvt}, \texttt{derivative-unique}, \texttt{rr-sub-in-rr}, \texttt{rr-zero-in}, \texttt{rr-lt-diff-pos}, \texttt{rr-lt-implies-le}, \texttt{diff-value-real}, \texttt{rr-le-scale-nonneg-right}.}
%
\vnbprf{\textbf{Oracles.} \texttt{crs}, \texttt{arith}, \texttt{ineq}.  \textbf{Bill:} \texttt{proven modulo 0}.}
%
\vnbprf{\vnbproofpart{\texttt{mvt-abs-bound}}}
%
\vnbprf{\textbf{Proof.} \texttt{theorem-\allowbreak{}library/mvt-\allowbreak{}abs-\allowbreak{}bound.scm} --- 111 recorded steps.}
%
\vnbprf{\textbf{Lemmas.} \texttt{rr-neg-closed}, \texttt{rr-sub-in-rr}, \texttt{fun-apply-type-c}, \texttt{rr-mul-closed}, \texttt{rr-le-abs}, \texttt{rr-neg-abs-le}, \texttt{rr-abs-closed}, \texttt{mvt-upper-bound}, \texttt{mvt-lower-bound}, \texttt{rr-abs-bound}.}
%
\vnbprf{\textbf{Oracles.} \texttt{ineq}, \texttt{crs}, \texttt{arith}.  \textbf{Bill:} \texttt{proven modulo 0}.}
%
\end{proof}
%
\noindent\textbf{Deviation from the notes.} Restated 2026-09-21 for functions ON [a, b] (the first three names), with the local derivative has-deriv-at. The last three names are the older forms for functions TOTAL on RR with two-sided continuity at the endpoints.
%
\renewcommand{\thethm}{2.15}
\begin{cor}[{Zero derivative implies constant}]\label{thm:deriv-zero-constant-on-interval}
%
\leavevmode
\vnbpart{\texttt{deriv-zero-constant-on-interval}}
\vnbstmt{\begin{array}{@{}l@{}} \forall a, b \in \mathbb{R}, h \in (\operatorname{ccint}(a, b) \to \mathbb{R}), u, v \in \operatorname{ccint}(a, b). \\ \quad (a < b) \wedge \\ \quad \operatorname{is\text{-}continuous\text{-}on}(h, \operatorname{ccint}(a, b)) \wedge \\ \quad \forall t \in \operatorname{ooint}(a, b). \\ \quad \quad \operatorname{has\text{-}deriv\text{-}at}(h, t, 0) \Rightarrow \\ \quad \quad (h(u) = h(v)) \end{array}}
\vnbpart{\texttt{deriv-zero-implies-constant}}
\vnbstmt{\begin{array}{@{}l@{}} \forall f, a, b, u, v. \\ \quad (f \in (\mathbb{R} \to \mathbb{R})) \wedge \\ \quad (a \in \mathbb{R}) \wedge \\ \quad (b \in \mathbb{R}) \wedge \\ \quad (a < b) \wedge \\ \quad \forall x \in \operatorname{ccint}(a, b). \\ \quad \quad \operatorname{is\text{-}continuous\text{-}at}(\operatorname{rr\text{-}ms}, \operatorname{rr\text{-}ms}, f, x) \wedge \\ \quad \forall x \in \mathbb{R}. \\ \quad \quad ((a < x) \wedge (x < b)) \Rightarrow  \\ \quad \quad \quad \operatorname{is\text{-}diff\text{-}at}(f, x, 0) \wedge \\ \quad (u \in \operatorname{ccint}(a, b)) \wedge \\ \quad (v \in \operatorname{ccint}(a, b)) \Rightarrow \\ \quad \quad (f(u) = f(v)) \end{array}}
\vnbpart{\texttt{deriv-zero-const-up}}
\vnbstmt{\begin{array}{@{}l@{}} \forall f, a, b, u, v. \\ \quad (f \in (\mathbb{R} \to \mathbb{R})) \wedge \\ \quad (a \in \mathbb{R}) \wedge \\ \quad (b \in \mathbb{R}) \wedge \\ \quad (a < b) \wedge \\ \quad \forall x \in \operatorname{ccint}(a, b). \\ \quad \quad \operatorname{is\text{-}continuous\text{-}at}(\operatorname{rr\text{-}ms}, \operatorname{rr\text{-}ms}, f, x) \wedge \\ \quad \forall x \in \mathbb{R}. \\ \quad \quad ((a < x) \wedge (x < b)) \Rightarrow  \\ \quad \quad \quad \operatorname{is\text{-}diff\text{-}at}(f, x, 0) \wedge \\ \quad (u \in \operatorname{ccint}(a, b)) \wedge \\ \quad (v \in \operatorname{ccint}(a, b)) \wedge \\ \quad (u < v) \Rightarrow \\ \quad \quad (f(u) = f(v)) \end{array}}
%
\end{cor}
%
\begin{proof}
%
The proofs: the file that holds each, the lemmas it cites, its oracles and its bill.
%
\vnbprf{\vnbproofpart{\texttt{deriv-zero-constant-on-interval}}}
%
\vnbprf{\textbf{Proof.} \texttt{theorem-\allowbreak{}library/pw-\allowbreak{}antiderivative-\allowbreak{}laws.scm} --- 84 recorded steps.}
%
\vnbprf{\textbf{Lemmas.} \texttt{extend-const-in-fun}, \texttt{ccint-elt-in-rr}, \texttt{extend-const-continuous-at}, \texttt{ooint-membership}, \texttt{extend-const-deriv-fwd}, \texttt{deriv-zero-implies-constant}, \texttt{ccint-membership}, \texttt{extend-const-fixes}.}
%
\vnbprf{\textbf{Oracles.} \texttt{ineq}, \texttt{arith}, \texttt{crs}.  \textbf{Bill:} \texttt{proven modulo 0}.}
%
\vnbprf{\vnbproofpart{\texttt{deriv-zero-implies-constant}}}
%
\vnbprf{\textbf{Proof.} \texttt{theorem-\allowbreak{}library/deriv-\allowbreak{}constant-\allowbreak{}proof.scm} --- 117 recorded steps.}
%
\vnbprf{\textbf{Lemmas.} \texttt{ccint-membership}, \texttt{rr-lt-trichotomy}, \texttt{deriv-zero-const-up}, \texttt{fun-apply-type-c}.}
%
\vnbprf{\textbf{Oracles.} \texttt{crs}, \texttt{ineq}, \texttt{arith}.  \textbf{Bill:} \texttt{proven modulo 0}.}
%
\vnbprf{\vnbproofpart{\texttt{deriv-zero-const-up}}}
%
\vnbprf{\textbf{Proof.} \texttt{theorem-\allowbreak{}library/deriv-\allowbreak{}constant-\allowbreak{}proof.scm} --- 179 recorded steps.}
%
\vnbprf{\textbf{Lemmas.} \texttt{ccint-membership}, \texttt{rr-le-trans-c}, \texttt{rr-le-lt-trans}, \texttt{rr-lt-le-trans}, \texttt{mvt}, \texttt{derivative-unique}, \texttt{fun-apply-type-c}, \texttt{rr-diff-zero-eq}.}
%
\vnbprf{\textbf{Oracles.} \texttt{crs}, \texttt{ineq}, \texttt{arith}.  \textbf{Bill:} \texttt{proven modulo 0}.}
%
\end{proof}
%
\noindent\textbf{Deviation from the notes.} Restated 2026-09-22 for a function ON [a, b] (the first name). `f is a constant' is written as `f takes the same value at any two points of [a, b]'. The remaining names are the older forms for functions TOTAL on RR with two-sided continuity at the endpoints; they remain as lemmas.
%
\renewcommand{\thethm}{2.16}
\begin{prop}[{Uniqueness of the interpolating Taylor polynomial at 0}]\label{thm:taylor-interpolant-at-zero}
%
\vnbstmt{\begin{array}{@{}l@{}} \forall n \in \mathbb{N}, f. \\ \quad \forall \mathit{tck} \in \mathbb{N}. \\ \quad \quad (\mathit{tck} \leq n) \Rightarrow  \\ \quad \quad \quad (\operatorname{nth\text{-}deriv}(f, \mathit{tck})(0) \in \mathbb{R}) \Rightarrow \\ \quad \quad \exists a \in (\mathbb{N} \to \mathbb{R}). \\ \quad \quad \quad \forall \mathit{tck} \in \mathbb{N}. \\ \quad \quad \quad \quad (\mathit{tck} \leq n) \Rightarrow  \\ \quad \quad \quad \quad \quad a(\mathit{tck}) \\ \quad \quad \quad \quad \quad =\; (\operatorname{nth\text{-}deriv}(f, \mathit{tck})(0) \cdot \operatorname{recip}(\operatorname{factorial}(\mathit{tck}))) \wedge \\ \quad \quad \quad \forall \mathit{tck} \in \mathbb{N}. \\ \quad \quad \quad \quad (\mathit{tck} \leq n) \Rightarrow  \\ \quad \quad \quad \quad \quad \operatorname{nth\text{-}deriv}((\lambda x \in \mathbb{R}.\; \operatorname{series\text{-}partial\text{-}sum}((\lambda k \in \mathbb{N}.\; (a(k) \cdot \operatorname{power}(x, k))), (n + 1))), \mathit{tck})(0) \\ \quad \quad \quad \quad \quad =\; \operatorname{nth\text{-}deriv}(f, \mathit{tck})(0) \wedge \\ \quad \quad \quad \forall g \in (\mathbb{N} \to \mathbb{R}), \mathit{tck} \in \mathbb{N}. \\ \quad \quad \quad \quad \forall \mathit{tck} \in \mathbb{N}. \\ \quad \quad \quad \quad \quad (\mathit{tck} \leq n) \Rightarrow  \\ \quad \quad \quad \quad \quad \quad \operatorname{nth\text{-}deriv}((\lambda x \in \mathbb{R}.\; \operatorname{series\text{-}partial\text{-}sum}((\lambda k \in \mathbb{N}.\; (g(k) \cdot \operatorname{power}(x, k))), (n + 1))), \mathit{tck})(0) \\ \quad \quad \quad \quad \quad \quad =\; \operatorname{nth\text{-}deriv}(f, \mathit{tck})(0) \wedge \\ \quad \quad \quad \quad (\mathit{tck} \leq n) \Rightarrow \\ \quad \quad \quad \quad \quad (g(\mathit{tck}) = a(\mathit{tck})) \end{array}}
%
\end{prop}
%
\begin{proof}
%
The proofs: the file that holds each, the lemmas it cites, its oracles and its bill.
%
\vnbprf{\textbf{Proof.} \texttt{theorem-\allowbreak{}library/taylor-\allowbreak{}cluster.scm} --- 150 recorded steps.}
%
\vnbprf{\textbf{Lemmas.} \texttt{nn-is-set}, \texttt{recip-factorial-in-rr}, \texttt{rr-mul-in-rr}, \texttt{rr-zero-in}, \texttt{taylor-derivs-from-coef}, \texttt{taylor-coef-from-derivs}.}
%
\vnbprf{\textbf{Oracles.} \texttt{arith}, \texttt{ineq}, \texttt{crs}.  \textbf{Bill:} \texttt{proven modulo 0}.}
%
\end{proof}
%
\noindent\textbf{Deviation from the notes.} UPDATED 2026-09-24 (batch 25): PROVEN: existence, and uniqueness stated on the COEFFICIENT sequences (two partial-sum polynomials of degree below n whose first n - 1 derivatives at 0 agree have equal coefficients), theorem-library/taylor-cluster.scm. EARLIER NOTE: Nothing states existence and uniqueness of a polynomial of degree n - 1 whose first n - 1 derivatives at 0 agree with those of f. The tree has the Taylor polynomial as a functoid, taylor-poly, with the unfold r8e-tp-unfold and the two agreement facts taylor-poly-at-center and taylor-center-partial-sum (the polynomial's VALUE at the centre is f of the centre); it has no statement that its derivatives at the centre match f's, and no uniqueness statement, nor any notion of the DEGREE of a polynomial function.
%
\renewcommand{\thethm}{2.18}
\begin{thm}[{Taylor's formula with a general auxiliary function G}]\label{thm:taylor-formula-on-interval}
%
\leavevmode
\vnbpart{\texttt{taylor-formula-on-interval}}
\vnbstmt{\begin{array}{@{}l@{}} \forall a, b \in \mathbb{R}, d.; \forall m \in \mathbb{N}, g \in (\operatorname{ccint}(a, b) \to \mathbb{R}). \\ \quad \operatorname{is\text{-}taylor\text{-}family}(d, (m + 1), a, b) \wedge \\ \quad \operatorname{is\text{-}continuous\text{-}on}(g, \operatorname{ccint}(a, b)) \wedge \\ \quad \forall x \in \operatorname{ooint}(a, b). \exists l. \\ \quad \quad \operatorname{has\text{-}deriv\text{-}at}(g, x, l) \Rightarrow \\ \quad \quad \exists i \in \operatorname{ooint}(a, b), u. \\ \quad \quad \quad \operatorname{has\text{-}deriv\text{-}at}(g, i, u) \wedge \\ \quad \quad \quad ((u \cdot \operatorname{taylor\text{-}rem}(d, (m + 1), a, b)) \cdot \operatorname{factorial}(m)) \\ \quad \quad \quad =\; ((d((m + 1))(i) \cdot \operatorname{power}((b - i), m)) \cdot (g(b) - g(a))) \wedge \\ \quad \quad \quad (u \cdot \operatorname{taylor\text{-}rem}(d, (m + 1), a, b)) \\ \quad \quad \quad =\; (((d((m + 1))(i) \cdot \operatorname{recip}(\operatorname{factorial}(m))) \cdot \operatorname{power}((b - i), m)) \cdot (g(b) - g(a))) \end{array}}
\vnbpart{\texttt{taylor-lagrange}}
\vnbstmt{\begin{array}{@{}l@{}} \forall f, a, x, n. \\ \quad (f \in (\mathbb{R} \to \mathbb{R})) \wedge \\ \quad (a \in \mathbb{R}) \wedge \\ \quad (x \in \mathbb{R}) \wedge \\ \quad (n \in \mathbb{N}) \wedge \\ \quad (a < x) \wedge \\ \quad \operatorname{taylor\text{-}differentiable}(f, a, x, n) \Rightarrow \\ \quad \quad \exists b \in \mathbb{R}. \\ \quad \quad \quad (a < b) \wedge \\ \quad \quad \quad (b < x) \wedge \\ \quad \quad \quad (\operatorname{factorial}((n + 1)) \cdot (f(x) - \operatorname{taylor\text{-}poly}(f, a, n, x))) \\ \quad \quad \quad =\; (\operatorname{nth\text{-}deriv}(f, (n + 1))(b) \cdot \operatorname{power}((x - a), (n + 1))) \end{array}}
%
\end{thm}
%
\begin{proof}
%
The proofs: the file that holds each, the lemmas it cites, its oracles and its bill.
%
\vnbprf{\vnbproofpart{\texttt{taylor-formula-on-interval}}}
%
\vnbprf{\textbf{Proof.} \texttt{theorem-\allowbreak{}library/interval-\allowbreak{}taylor.scm} --- 12 recorded steps.}
%
\vnbprf{\textbf{Lemmas.} \texttt{nn-succ-closed}, \texttt{nn-le-refl}, \texttt{taylor-aux-continuous-on}, \texttt{taylor-aux-has-deriv}, \texttt{taylor-family-sum-at-endpoint}, \texttt{taylor-formula-on-interval-from-aux}.}
%
\vnbprf{\textbf{Oracles.} \texttt{ineq}, \texttt{crs}, \texttt{arith}.  \textbf{Bill:} \texttt{proven modulo 0}.}
%
\vnbprf{\vnbproofpart{\texttt{taylor-lagrange}}}
%
\vnbprf{\textbf{Proof.} \texttt{theorem-\allowbreak{}library/taylor-\allowbreak{}proof.scm} --- 225 recorded steps.}
%
\vnbprf{\textbf{Lemmas.} \texttt{taylor-derivs-in-fun}, \texttt{taylor-g-in-fun}, \texttt{taylor-h-in-fun}, \texttt{taylor-gmvt-cont}, \texttt{taylor-gmvt-diff}, \texttt{generalized-mvt}, \texttt{rr-sub-in-rr}, \texttt{rr-zero-in}, \texttt{nn-succ-closed}, \texttt{nn-in-rr}, \texttt{recip-factorial-in-rr}, \texttt{taylor-deriv-real}, \texttt{power-real-closed}, \texttt{rr-mul-closed}, \texttt{taylor-g-diff}, \texttt{taylor-h-diff}, \texttt{derivative-unique}, \texttt{taylor-g-at-x}, \texttt{taylor-g-at-a}, \texttt{taylor-h-at-x}, \texttt{taylor-h-at-a}, \texttt{rr-lt-diff-pos}, \texttt{rr-power-pos}, \texttt{rr-pos-ne-zero}, \texttt{taylor-poly-in-rr}, \texttt{fun-apply-type-c}, \texttt{taylor-clear}.}
%
\vnbprf{\textbf{Oracles.} \texttt{ineq}, \texttt{arith}, \texttt{crs}.  \textbf{Bill:} \texttt{proven modulo 0}.}
%
\end{proof}
%
\noindent\textbf{Deviation from the notes.} Proven 2026-09-22 for functions ON [a, b] with a general G, as in the notes (the first name); its conclusion gives (20) twice, multiplied through by (n-1)! and literally with the division. n is written succ(m). The derivatives f, f', ..., f\textasciicircum{}(n) enter as a finite family d(0), ..., d(n) (is-taylor-family): each d(k), k $<$ n, continuous on [a, b] with derivative d(k+1) on (a, b); no derivative is taken at an endpoint, where the notes speak of f\textasciicircum{}(n-1) `defined on [a, b]' (a one-sided derivative at a and b). The second name is the older Lagrange form for functions TOTAL on RR.
%
\renewcommand{\thethm}{2.19}
\begin{cor}[{Taylor's formula with Lagrange remainder}]\label{thm:taylor-lagrange-on-interval}
%
\leavevmode
\vnbpart{\texttt{taylor-lagrange-on-interval}}
\vnbstmt{\begin{array}{@{}l@{}} \forall a, b \in \mathbb{R}, d.; \forall m \in \mathbb{N}. \\ \quad \operatorname{is\text{-}taylor\text{-}family}(d, (m + 1), a, b) \Rightarrow  \\ \quad \quad \exists i \in \operatorname{ooint}(a, b). \\ \quad \quad \quad \operatorname{taylor\text{-}rem}(d, (m + 1), a, b) \\ \quad \quad \quad =\; ((d((m + 1))(i) \cdot \operatorname{recip}(\operatorname{factorial}((m + 1)))) \cdot \operatorname{power}((b - a), (m + 1))) \end{array}}
\vnbpart{\texttt{taylor-lagrange}}
\vnbstmt{\begin{array}{@{}l@{}} \forall f, a, x, n. \\ \quad (f \in (\mathbb{R} \to \mathbb{R})) \wedge \\ \quad (a \in \mathbb{R}) \wedge \\ \quad (x \in \mathbb{R}) \wedge \\ \quad (n \in \mathbb{N}) \wedge \\ \quad (a < x) \wedge \\ \quad \operatorname{taylor\text{-}differentiable}(f, a, x, n) \Rightarrow \\ \quad \quad \exists b \in \mathbb{R}. \\ \quad \quad \quad (a < b) \wedge \\ \quad \quad \quad (b < x) \wedge \\ \quad \quad \quad (\operatorname{factorial}((n + 1)) \cdot (f(x) - \operatorname{taylor\text{-}poly}(f, a, n, x))) \\ \quad \quad \quad =\; (\operatorname{nth\text{-}deriv}(f, (n + 1))(b) \cdot \operatorname{power}((x - a), (n + 1))) \end{array}}
%
\end{cor}
%
\begin{proof}
%
The proofs: the file that holds each, the lemmas it cites, its oracles and its bill.
%
\vnbprf{\vnbproofpart{\texttt{taylor-lagrange-on-interval}}}
%
\vnbprf{\textbf{Proof.} \texttt{theorem-\allowbreak{}library/interval-\allowbreak{}taylor.scm} --- 257 recorded steps.}
%
\vnbprf{\textbf{Lemmas.} \texttt{taylor-family-order-pos}, \texttt{taylor-family-endpoints-lt}, \texttt{nn-succ-closed}, \texttt{nn-in-rr}, \texttt{nn-le-refl}, \texttt{nn-zero-in}, \texttt{rr-zero-in}, \texttt{rr-one-in}, \texttt{bt-lt-succ}, \texttt{rr-neg-closed}, \texttt{taylor-family-in-fun}, \texttt{fun-domain-in-set}, \texttt{ccint-subset-rr}, \texttt{ooint-subset-ccint}, \texttt{rr-is-set}, \texttt{recip-factorial-in-rr}, \texttt{rr-sub-in-rr}, \texttt{power-closed-at}, \texttt{ccint-membership}, \texttt{fun-apply-type-c}, \texttt{taylor-sum-in-rr}, \texttt{ccint-elt-in-rr}, \texttt{power-sub-succ-continuous}, \texttt{cont-on-of-total}, \texttt{ooint-elt-in-rr}, \texttt{power-sub-succ-has-deriv}, \texttt{ooint-interior-radius}, \texttt{rr-pos-rr-in-rr}, \texttt{rr-add-in-rr}, \texttt{subset-trans}, \texttt{restrict-in-fun}, \texttt{subset-mem-fwd}, \texttt{has-deriv-at-local}, \texttt{taylor-formula-on-interval}, \texttt{taylor-family-top-in-fun}, \texttt{rr-mul-closed}, \texttt{has-deriv-at-unique}, \texttt{power-zero-base}, \texttt{tgd-recip-factorial-succ}, \texttt{ooint-membership}, \texttt{rr-power-pos}, \texttt{rr-mul-pos}, \texttt{rr-pos-ne-zero}, \texttt{rr-cancel-mul-left}.}
%
\vnbprf{\textbf{Oracles.} \texttt{ineq}, \texttt{crs}, \texttt{arith}.  \textbf{Bill:} \texttt{proven modulo 0}.}
%
\vnbprf{\vnbproofpart{\texttt{taylor-lagrange}}}
%
\vnbprf{\textbf{Proof.} \texttt{theorem-\allowbreak{}library/taylor-\allowbreak{}proof.scm} --- 225 recorded steps.}
%
\vnbprf{\textbf{Lemmas.} \texttt{taylor-derivs-in-fun}, \texttt{taylor-g-in-fun}, \texttt{taylor-h-in-fun}, \texttt{taylor-gmvt-cont}, \texttt{taylor-gmvt-diff}, \texttt{generalized-mvt}, \texttt{rr-sub-in-rr}, \texttt{rr-zero-in}, \texttt{nn-succ-closed}, \texttt{nn-in-rr}, \texttt{recip-factorial-in-rr}, \texttt{taylor-deriv-real}, \texttt{power-real-closed}, \texttt{rr-mul-closed}, \texttt{taylor-g-diff}, \texttt{taylor-h-diff}, \texttt{derivative-unique}, \texttt{taylor-g-at-x}, \texttt{taylor-g-at-a}, \texttt{taylor-h-at-x}, \texttt{taylor-h-at-a}, \texttt{rr-lt-diff-pos}, \texttt{rr-power-pos}, \texttt{rr-pos-ne-zero}, \texttt{taylor-poly-in-rr}, \texttt{fun-apply-type-c}, \texttt{taylor-clear}.}
%
\vnbprf{\textbf{Oracles.} \texttt{ineq}, \texttt{arith}, \texttt{crs}.  \textbf{Bill:} \texttt{proven modulo 0}.}
%
\end{proof}
%
\noindent\textbf{Deviation from the notes.} Proven 2026-09-22 for functions ON [a, b] (the first name), from 2.18 with G(x) = (b - x)\textasciicircum{}n; derivatives as the finite family of 2.18. The second name is the older form for functions TOTAL on RR, with the remainder cross-multiplied by the factorial.
%
\renewcommand{\thethm}{2.20}
\begin{cor}[{Supremum bound for the Taylor remainder}]\label{thm:taylor-remainder-sup-bound}
%
\vnbstmt{\begin{array}{@{}l@{}} \forall a, b \in \mathbb{R}, f.; \forall m \in \mathbb{N}. \\ \quad \operatorname{is\text{-}taylor\text{-}family}(f, (m + 1), a, b) \wedge \\ \quad \operatorname{rr\text{-}bounded\text{-}above}(\operatorname{image}((\lambda v \in \operatorname{ooint}(a, b).\; \lvert f((m + 1))(v) \rvert), \operatorname{ooint}(a, b))) \Rightarrow \\ \quad \quad \lvert \operatorname{taylor\text{-}rem}(f, (m + 1), a, b) \rvert \\ \quad \quad \leq\; ((\operatorname{power}((b - a), (m + 1)) \cdot \operatorname{recip}(\operatorname{factorial}((m + 1)))) \cdot \operatorname{sup}(\operatorname{image}((\lambda v \in \operatorname{ooint}(a, b).\; \lvert f((m + 1))(v) \rvert), \operatorname{ooint}(a, b)))) \end{array}}
%
\end{cor}
%
\begin{proof}
%
The proofs: the file that holds each, the lemmas it cites, its oracles and its bill.
%
\vnbprf{\textbf{Proof.} \texttt{theorem-\allowbreak{}library/taylor-\allowbreak{}cluster.scm} --- 147 recorded steps.}
%
\vnbprf{\textbf{Lemmas.} \texttt{taylor-lagrange-on-interval}, \texttt{taylor-family-top-in-fun}, \texttt{fun-apply-type-c}, \texttt{taylor-family-endpoints-lt}, \texttt{rr-sub-in-rr}, \texttt{nn-succ-closed}, \texttt{power-real-closed}, \texttt{recip-factorial-in-rr}, \texttt{rr-mul-in-rr}, \texttt{rr-abs-closed}, \texttt{rr-power-nonneg}, \texttt{factorial-real-pos}, \texttt{rr-recip-pos}, \texttt{rr-leq-mul-nonneg}, \texttt{rr-abs-mult}, \texttt{rr-abs-of-nonneg}, \texttt{fun-domain-in-set}, \texttt{subset-def}, \texttt{image-subset-codomain}, \texttt{membership-implies-sethood}, \texttt{rr-sup-upper}, \texttt{rr-mul-le-right}.}
%
\vnbprf{\textbf{Oracles.} \texttt{ineq}, \texttt{crs}, \texttt{arith}.  \textbf{Bill:} \texttt{proven modulo 0}.}
%
\end{proof}
%
\noindent\textbf{Deviation from the notes.} UPDATED 2026-09-24 (batch 25): PROVEN with ONE explicit hypothesis the notes leave implicit: the set of |f\textasciicircum{}(n)(xi)|, xi in (a, b), is bounded above (it does not follow from 2.18's hypotheses). EARLIER NOTE: No statement bounding the Taylor remainder by (b - a)\textasciicircum{}n / n! times the SUPREMUM of abs(f\textasciicircum{}(n)) over the open interval. The library has only the pointwise-witness form taylor-lagrange, and the normed-space analogues vector-taylor-remainder-bound and nvs-taylor-remainder-bound, which bound by the norm of the derivative AT THE WITNESS theta, not by a supremum; no sup over an open interval of a derivative family is formed anywhere.
%
\renewcommand{\thethm}{2.21}
\begin{cor}[{Continuous remainder factor S\_n}]\label{thm:taylor-remainder-factor}
%
\leavevmode
\vnbpart{\texttt{taylor-remainder-factor}}
\vnbstmt{\begin{array}{@{}l@{}} \forall a, b \in \mathbb{R}, f.; \forall m \in \mathbb{N}. \\ \quad (a < b) \wedge \\ \quad \forall k \in \mathbb{N}. \\ \quad \quad (k \leq (m + 1)) \Rightarrow  \\ \quad \quad \quad (f(k) \in (\operatorname{ccint}(a, b) \to \mathbb{R})) \wedge \\ \quad \forall k \in \mathbb{N}. \\ \quad \quad (k < (m + 1)) \Rightarrow  \\ \quad \quad \quad \operatorname{is\text{-}continuous\text{-}on}(f(k), \operatorname{ccint}(a, b)) \wedge \\ \quad \forall k \in \mathbb{N}, x \in \operatorname{ooint}(a, b). \\ \quad \quad (k < (m + 1)) \Rightarrow  \\ \quad \quad \quad \operatorname{has\text{-}deriv\text{-}at}(f(k), x, f((k + 1))(x)) \wedge \\ \quad \operatorname{is\text{-}continuous\text{-}on}(f((m + 1)), \operatorname{ccint}(a, b)) \Rightarrow \\ \quad \quad \exists s \in (\operatorname{ccint}(a, b) \to \mathbb{R}). \\ \quad \quad \quad \operatorname{is\text{-}continuous\text{-}on}(s, \operatorname{ccint}(a, b)) \wedge \\ \quad \quad \quad (s(a) = f((m + 1))(a)) \wedge \\ \quad \quad \quad \forall t \in \operatorname{ccint}(a, b). \\ \quad \quad \quad \quad \operatorname{taylor\text{-}rem}(f, (m + 1), a, t) \\ \quad \quad \quad \quad =\; ((s(t) \cdot \operatorname{power}((t - a), (m + 1))) \cdot \operatorname{recip}(\operatorname{factorial}((m + 1)))) \wedge \\ \quad \quad \quad \forall t \in \operatorname{ccint}(a, b). \\ \quad \quad \quad \quad \lvert s(t) \rvert \\ \quad \quad \quad \quad \leq\; \operatorname{sup}(\operatorname{image}((\lambda v \in \operatorname{ccint}(a, b).\; \lvert f((m + 1))(v) \rvert), \operatorname{ccint}(a, b))) \end{array}}
\vnbpart{\texttt{taylor-lagrange-subinterval}}
\vnbstmt{\begin{array}{@{}l@{}} \forall a, b \in \mathbb{R}, f.; \forall m \in \mathbb{N}, t \in \mathbb{R}. \\ \quad \forall k \in \mathbb{N}. \\ \quad \quad (k \leq (m + 1)) \Rightarrow  \\ \quad \quad \quad (f(k) \in (\operatorname{ccint}(a, b) \to \mathbb{R})) \wedge \\ \quad \forall k \in \mathbb{N}. \\ \quad \quad (k < (m + 1)) \Rightarrow  \\ \quad \quad \quad \operatorname{is\text{-}continuous\text{-}on}(f(k), \operatorname{ccint}(a, b)) \wedge \\ \quad \forall k \in \mathbb{N}, x \in \operatorname{ooint}(a, b). \\ \quad \quad (k < (m + 1)) \Rightarrow  \\ \quad \quad \quad \operatorname{has\text{-}deriv\text{-}at}(f(k), x, f((k + 1))(x)) \wedge \\ \quad (a < t) \wedge \\ \quad (t \leq b) \Rightarrow \\ \quad \quad \exists i \in \operatorname{ooint}(a, t). \\ \quad \quad \quad \operatorname{taylor\text{-}rem}(f, (m + 1), a, t) \\ \quad \quad \quad =\; ((f((m + 1))(i) \cdot \operatorname{recip}(\operatorname{factorial}((m + 1)))) \cdot \operatorname{power}((t - a), (m + 1))) \end{array}}
%
\end{cor}
%
\begin{proof}
%
The proofs: the file that holds each, the lemmas it cites, its oracles and its bill.
%
\vnbprf{\vnbproofpart{\texttt{taylor-remainder-factor}}}
%
\vnbprf{\textbf{Proof.} \texttt{theorem-\allowbreak{}library/taylor-\allowbreak{}cluster.scm} --- 524 recorded steps.}
%
\vnbprf{\textbf{Lemmas.} \texttt{nn-succ-closed}, \texttt{nn-in-rr}, \texttt{factorial-real-pos}, \texttt{recip-factorial-in-rr}, \texttt{ccint-membership}, \texttt{nn-le-refl}, \texttt{fun-apply-type-c}, \texttt{rr-recip-factorial}, \texttt{nn-zero-in}, \texttt{nn-one-le-succ}, \texttt{taylor-factor-fn-continuous}, \texttt{taylor-factor-glue-continuous}, \texttt{taylor-factor-fn-value}, \texttt{ccint-subset-rr}, \texttt{restrict-in-fun}, \texttt{ccint-elt-in-rr}, \texttt{restrict-apply}, \texttt{cont-on-of-total}, \texttt{nn-le-succ}, \texttt{nn-le-trans-guarded}, \texttt{rr-le-cases}, \texttt{rr-lt-not-le}, \texttt{nn-zero-le}, \texttt{extend-const-fixes}, \texttt{rr-sub-in-rr}, \texttt{rr-power-pos}, \texttt{power-real-closed}, \texttt{rr-pos-ne-zero}, \texttt{recip-star-value}, \texttt{rr-recip-closed}, \texttt{rr-recip-inverse}, \texttt{taylor-sum-continuous}, \texttt{rr-mul-in-rr}, \texttt{rr-one-in}, \texttt{taylor-sum-at-base}, \texttt{rr-zero-in}, \texttt{power-zero-base}, \texttt{fun-domain-in-set}, \texttt{rr-abs-closed}, \texttt{subset-def}, \texttt{image-subset-codomain}, \texttt{membership-implies-sethood}, \texttt{extend-const-in-fun}, \texttt{extend-const-continuous-at}, \texttt{continuous-abs-bounded-on-ccint}, \texttt{rr-sup-upper}, \texttt{ooint-membership}.}
%
\vnbprf{\textbf{Oracles.} \texttt{ineq}, \texttt{crs}, \texttt{arith}.  \textbf{Bill:} \texttt{proven modulo 0}.}
%
\vnbprf{\vnbproofpart{\texttt{taylor-lagrange-subinterval}}}
%
\vnbprf{\textbf{Proof.} \texttt{theorem-\allowbreak{}library/taylor-\allowbreak{}cluster.scm} --- 335 recorded steps.}
%
\vnbprf{\textbf{Lemmas.} \texttt{nn-succ-closed}, \texttt{nn-in-rr}, \texttt{nn-one-le-succ}, \texttt{ccint-subset-ccint}, \texttt{ooint-subset-ccint}, \texttt{subset-def}, \texttt{subset-mem-fwd}, \texttt{rr-is-metric-space}, \texttt{ccint-subset-rr}, \texttt{rr-lt-irrefl}, \texttt{nn-le-refl}, \texttt{restrict-in-fun}, \texttt{subspace-restrict-continuous-at}, \texttt{ooint-membership}, \texttt{ooint-inner-radius}, \texttt{rr-pos-rr-in-rr}, \texttt{rr-sub-in-rr}, \texttt{rr-add-in-rr}, \texttt{restrict-apply}, \texttt{has-deriv-at-local}, \texttt{taylor-lagrange-on-interval}, \texttt{ccint-membership}, \texttt{fun-apply-type-c}, \texttt{nn-zero-in}, \texttt{nn-zero-le}, \texttt{power-real-closed}, \texttt{rr-mul-in-rr}, \texttt{recip-factorial-in-rr}, \texttt{series-partial-sum-agree-below}, \texttt{series-partial-sum-in-rr-below}.}
%
\vnbprf{\textbf{Oracles.} \texttt{ineq}, \texttt{arith}, \texttt{crs}.  \textbf{Bill:} \texttt{proven modulo 0}.}
%
\end{proof}
%
\noindent\textbf{Deviation from the notes.} UPDATED 2026-09-24 (batch 25): PROVEN in full (a continuous S on [a, b] with R\_n(t) = S(t)(t - a)\textasciicircum{}n / n!, S(a) = f\textasciicircum{}(n)(a), |S(t)| $<$= sup of |f\textasciicircum{}(n)| over [a, b], the boundedness proven); the family hypothesis is written out as four conjuncts because IS-TAYLOR-FAMILY puts f\textasciicircum{}(n) on the open interval. THE NOTES' "R\_n(a) = f\textasciicircum{}(n)(a)" IS A MISPRINT for S\_n(a) = f\textasciicircum{}(n)(a). EARLIER NOTE: Nothing constructs a continuous function S\_n on [a, b] with R\_n(t) = S\_n(t)(t - a)\textasciicircum{}n / n!, nothing states S\_n(a) = f\textasciicircum{}(n)(a) or the bound abs(S\_n(t)) $<$= sup over [a, b] of abs(f\textasciicircum{}(n)). The nearest machinery is the Caratheodory-style witness used for is-diff-at itself (the phi of the definition, and diff-at-witness-family), which produces a continuous factor for the FIRST-order increment only.
%
\renewcommand{\thethm}{2.22}
\begin{prop}[{Taylor's formula is exact for a polynomial}]\label{thm:taylor-exact-polynomial}
%
\vnbstmt{\begin{array}{@{}l@{}} \forall m \in \mathbb{N}, f \in (\mathbb{N} \to \mathbb{R}), a, h \in \mathbb{R}. \\ \quad (\lambda x \in \mathbb{R}.\; \operatorname{series\text{-}partial\text{-}sum}((\lambda k \in \mathbb{N}.\; (f(k) \cdot \operatorname{power}(x, k))), (m + 1)))((a + h)) \\ \quad =\; \operatorname{series\text{-}partial\text{-}sum}((\lambda i \in \mathbb{N}.\; \\ \quad \quad \quad ((\operatorname{nth\text{-}deriv}((\lambda x \in \mathbb{R}.\; \operatorname{series\text{-}partial\text{-}sum}((\lambda k \in \mathbb{N}.\; (f(k) \cdot \operatorname{power}(x, k))), (m + 1))), i)(a) \cdot \operatorname{power}(h, i)) \cdot \operatorname{recip}(\operatorname{factorial}(i)))), \\ \quad \quad (m + 1)) \end{array}}
%
\end{prop}
%
\begin{proof}
%
The proofs: the file that holds each, the lemmas it cites, its oracles and its bill.
%
\vnbprf{\textbf{Proof.} \texttt{theorem-\allowbreak{}library/taylor-\allowbreak{}cluster.scm} --- 779 recorded steps.}
%
\vnbprf{\textbf{Lemmas.} \texttt{nn-zero-in}, \texttt{nn-succ-closed}, \texttt{rr-add-in-rr}, \texttt{poly-degree-zero-value}, \texttt{nn-is-set}, \texttt{poly-nth-deriv-in-fun}, \texttt{fun-apply-type-c}, \texttt{power-real-closed}, \texttt{rr-mul-in-rr}, \texttt{recip-factorial-in-rr}, \texttt{series-partial-sum-in-rr}, \texttt{series-partial-sum-succ}, \texttt{series-partial-sum-zero}, \texttt{rr-subset-cc}, \texttt{power-zero}, \texttt{nn-one-is-succ-zero}, \texttt{rr-recip-one}, \texttt{rr-zero-in}, \texttt{rr-one-in}, \texttt{poly-deriv-coef-exists}, \texttt{poly-deriv-fn}, \texttt{nn-in-rr}, \texttt{poly-lam-in-fun}, \texttt{rr-is-set}, \texttt{anti-lam-in-fun}, \texttt{compose-type}, \texttt{compose-apply}, \texttt{rr-sub-in-rr}, \texttt{deriv-polynomial}, \texttt{deriv-affine}, \texttt{deriv-chain}, \texttt{poly-antiderivative}, \texttt{deriv-difference}, \texttt{fun-codomain-iff}, \texttt{fun-domain-extensionality}, \texttt{nth-deriv-succ-inner}, \texttt{tgd-recip-factorial-succ}, \texttt{diff-implies-continuous}, \texttt{anti-term-lam-in-fun}, \texttt{rr-lt-trichotomy}, \texttt{pw-zero-deriv-total}, \texttt{ooint-elt-in-rr}, \texttt{power-zero-base}, \texttt{nn-succ-nonzero}, \texttt{rr-recip-closed}, \texttt{series-partial-sum-zero-seq}, \texttt{series-partial-sum-shift}, \texttt{recip-succ-cancel}.}
%
\vnbprf{\textbf{Oracles.} \texttt{crs}, \texttt{arith}, \texttt{ineq}.  \textbf{Bill:} \texttt{proven modulo 0}.}
%
\end{proof}
%
\noindent\textbf{Deviation from the notes.} UPDATED 2026-09-24 (batch 25): PROVEN for h of EITHER sign, by induction on the degree (the notes' route through the Lagrange form covers only a $<$ b, i.e. h $>$ 0). EARLIER NOTE: No statement that a polynomial function of degree n equals its own order-n Taylor expansion about an arbitrary centre. There is no predicate for being a polynomial function nor any notion of degree; polynomial functions appear only as the explicit term vnb-lambda(x, rr, series-partial-sum(vnb-lambda(k, nn, cf(k) * x\textasciicircum{}k), succ(n))), for which deriv-polynomial and poly-antiderivative are proven. The carrier carr(poly(a)) is the ALGEBRAIC polynomial ring (finitely supported sequences), which is not connected to that function term by any theorem.
%
\renewcommand{\thethm}{2.24}
\begin{prop}[{Mean Value Theorem for a directional derivative}]\label{thm:dir-deriv-mvt}
%
\leavevmode
\vnbpart{\texttt{dir-deriv-mvt}}
\vnbstmt{\begin{array}{@{}l@{}} \forall m, f, a, b. \\ \quad \operatorname{is\text{-}normed\text{-}vector\text{-}space}(m) \wedge \\ \quad (f \in (\operatorname{vec}(m) \to \mathbb{R})) \wedge \\ \quad (a \in \operatorname{vec}(m)) \wedge \\ \quad (b \in \operatorname{vec}(m)) \wedge \\ \quad \forall x \in \operatorname{ccint}(0, 1). \\ \quad \quad \operatorname{is\text{-}continuous\text{-}at}(\operatorname{rr\text{-}ms}, \operatorname{rr\text{-}ms}, \operatorname{seg\text{-}curve}(m, f, a, b), x) \wedge \\ \quad \forall x \in \mathbb{R}. \\ \quad \quad ((0 < x) \wedge (x < 1)) \Rightarrow  \\ \quad \quad \quad \exists l. \\ \quad \quad \quad \quad \operatorname{is\text{-}diff\text{-}at}(\operatorname{seg\text{-}curve}(m, f, a, b), x, l) \Rightarrow \\ \quad \quad \exists h \in \mathbb{R}. \\ \quad \quad \quad (0 < h) \wedge \\ \quad \quad \quad (h < 1) \wedge \\ \quad \quad \quad \exists v. \\ \quad \quad \quad \quad \operatorname{is\text{-}dir\text{-}diff\text{-}at}(m, f, \operatorname{vadd}(m)(a, \operatorname{act}(m)(h, b)), b, v) \wedge \\ \quad \quad \quad \quad (v = (f(\operatorname{vadd}(m)(a, b)) - f(a))) \end{array}}
\vnbpart{\texttt{dir-deriv-increment}}
\vnbstmt{\begin{array}{@{}l@{}} \forall m.; \forall f \in (\operatorname{vec}(m) \to \mathbb{R}), b, a \in \operatorname{vec}(m). \\ \quad \operatorname{is\text{-}normed\text{-}vector\text{-}space}(m) \wedge \\ \quad \forall x \in \operatorname{vec}(m). \exists k. \\ \quad \quad \operatorname{is\text{-}dir\text{-}diff\text{-}at}(m, f, x, a, k) \Rightarrow \\ \quad \quad \exists h \in \mathbb{R}. \\ \quad \quad \quad (0 < h) \wedge \\ \quad \quad \quad (h < 1) \wedge \\ \quad \quad \quad (f(\operatorname{vadd}(m)(b, a)) - f(b)) \\ \quad \quad \quad =\; \operatorname{dir\text{-}deriv}(m, f, \operatorname{vadd}(m)(b, \operatorname{act}(m)(h, a)), a) \end{array}}
%
\end{prop}
%
\begin{proof}
%
The proofs: the file that holds each, the lemmas it cites, its oracles and its bill.
%
\vnbprf{\vnbproofpart{\texttt{dir-deriv-mvt}}}
%
\vnbprf{\textbf{Proof.} \texttt{theorem-\allowbreak{}library/directional-\allowbreak{}derivative.scm} --- 106 recorded steps.}
%
\vnbprf{\textbf{Lemmas.} \texttt{rr-zero-in}, \texttt{rr-one-in}, \texttt{seg-curve-in-fun}, \texttt{mvt}, \texttt{nvs-act-in-vec}, \texttt{nvs-vadd-in-vec}, \texttt{rr-add-closed}, \texttt{seg-curve-apply}, \texttt{nvs-act-collect}, \texttt{nvs-act-one}, \texttt{nvs-act-zero}, \texttt{diff-value-real}, \texttt{dir-diff-shift}, \texttt{eq-sym}.}
%
\vnbprf{\textbf{Oracles.} \texttt{arith}, \texttt{crs}, \texttt{ineq}.  \textbf{Bill:} \texttt{proven modulo 0}.}
%
\vnbprf{\vnbproofpart{\texttt{dir-deriv-increment}}}
%
\vnbprf{\textbf{Proof.} \texttt{theorem-\allowbreak{}library/directional-\allowbreak{}derivative.scm} --- 64 recorded steps.}
%
\vnbprf{\textbf{Lemmas.} \texttt{seg-curve-in-fun}, \texttt{nvs-act-in-vec}, \texttt{nvs-vadd-in-vec}, \texttt{seg-curve-diff-at}, \texttt{ccint-membership}, \texttt{diff-implies-continuous}, \texttt{dir-deriv-mvt}, \texttt{dir-deriv-value}, \texttt{eq-sym}.}
%
\vnbprf{\textbf{Oracles.} \texttt{crs}, \texttt{arith}, \texttt{ineq}.  \textbf{Bill:} \texttt{proven modulo 0}.}
%
\end{proof}
%
\noindent\textbf{Deviation from the notes.} There is no open set U in the library statements. dir-deriv-mvt requires f in fun(vec(m), rr), i.e. f TOTAL on the normed vector space, and replaces the hypothesis that the partial derivative is defined on U by hypotheses on the SEGMENT CURVE seg-curve(m, f, a, eta): two-sided continuity at every point of ccint(0, 1), including the endpoints 0 and 1, and differentiability at every interior point. dir-deriv-increment is the closer match in shape -- it concludes f(b + zeta) - f(b) = dir-deriv(m, f, b + th * zeta, zeta) for some th in (0, 1) -- but it asks that the directional derivative in direction zeta exist at EVERY point of vec(m), where the notes ask only that it be defined on U with the segment [a, a + eta] inside U. The ambient space is an arbitrary normed vector space, not the specific E = R\textasciicircum{}n of the notes.
%
\renewcommand{\thethm}{2.25}
\begin{lem}[{Homogeneity of the directional derivative}]\label{thm:dir-deriv-homogeneous}
%
\vnbstmt{\begin{array}{@{}l@{}} \forall m, f, a, b, c, l. \\ \quad \operatorname{is\text{-}normed\text{-}vector\text{-}space}(m) \wedge \\ \quad (f \in (\operatorname{vec}(m) \to \mathbb{R})) \wedge \\ \quad (a \in \operatorname{vec}(m)) \wedge \\ \quad (b \in \operatorname{vec}(m)) \wedge \\ \quad (c \in \mathbb{R}) \wedge \\ \quad \operatorname{is\text{-}dir\text{-}diff\text{-}at}(m, f, a, b, l) \Rightarrow \\ \quad \quad \operatorname{is\text{-}dir\text{-}diff\text{-}at}(m, f, a, \operatorname{act}(m)(c, b), (c \cdot l)) \end{array}}
%
\end{lem}
%
\begin{proof}
%
The proofs: the file that holds each, the lemmas it cites, its oracles and its bill.
%
\vnbprf{\textbf{Proof.} \texttt{theorem-\allowbreak{}library/directional-\allowbreak{}derivative.scm} --- 38 recorded steps.}
%
\vnbprf{\textbf{Lemmas.} \texttt{nvs-act-in-vec}, \texttt{seg-curve-in-fun}, \texttt{rr-mul-closed}, \texttt{seg-curve-apply}, \texttt{nvs-act-scale-assoc}, \texttt{dir-diff-scale}.}
%
\vnbprf{\textbf{Oracles.} \texttt{crs}, \texttt{arith}, \texttt{ineq}.  \textbf{Bill:} \texttt{proven modulo 0}.}
%
\end{proof}
%
\noindent\textbf{Deviation from the notes.} f is required to be in fun(vec(m), rr), i.e. TOTAL on the normed vector space, where the notes ask that f and d\_eta f be defined on a neighbourhood of a. In the other direction the library is MORE general: it needs only is-dir-diff-at(m, f, a, eta, dl) at the single point a, not that d\_eta f be defined on a whole neighbourhood. The ambient space is an arbitrary normed vector space, not E = R\textasciicircum{}n. Neither statement guards lambda against 0, although the notes' proof divides by lambda.
%
\renewcommand{\thethm}{2.26}
\begin{lem}[{Additivity of the directional derivative under continuity}]\label{thm:dir-deriv-additive}
%
\vnbstmt{\begin{array}{@{}l@{}} \forall m, f, a, b, i, l, c. \\ \quad \operatorname{is\text{-}normed\text{-}vector\text{-}space}(m) \wedge \\ \quad (f \in (\operatorname{vec}(m) \to \mathbb{R})) \wedge \\ \quad (a \in \operatorname{vec}(m)) \wedge \\ \quad (b \in \operatorname{vec}(m)) \wedge \\ \quad (i \in \operatorname{vec}(m)) \wedge \\ \quad \operatorname{is\text{-}dir\text{-}diff\text{-}at}(m, f, a, b, l) \wedge \\ \quad \operatorname{is\text{-}dir\text{-}diff\text{-}at}(m, f, a, i, c) \wedge \\ \quad \forall x \in \operatorname{vec}(m). \exists k. \\ \quad \quad \operatorname{is\text{-}dir\text{-}diff\text{-}at}(m, f, x, b, k) \wedge \\ \quad \forall x \in \operatorname{vec}(m). \exists k. \\ \quad \quad \operatorname{is\text{-}dir\text{-}diff\text{-}at}(m, f, x, i, k) \wedge \\ \quad \operatorname{is\text{-}continuous\text{-}at}(\operatorname{nvs\text{-}metric\text{-}space}(m), \operatorname{rr\text{-}ms}, \\ \quad \quad (\lambda x \in \operatorname{vec}(m).\; \operatorname{dir\text{-}deriv}(m, f, x, b)), a) \wedge \\ \quad \operatorname{is\text{-}continuous\text{-}at}(\operatorname{nvs\text{-}metric\text{-}space}(m), \operatorname{rr\text{-}ms}, \\ \quad \quad (\lambda x \in \operatorname{vec}(m).\; \operatorname{dir\text{-}deriv}(m, f, x, i)), a) \Rightarrow \\ \quad \quad \operatorname{is\text{-}dir\text{-}diff\text{-}at}(m, f, a, \operatorname{vadd}(m)(b, i), (l + c)) \end{array}}
%
\end{lem}
%
\begin{proof}
%
The proofs: the file that holds each, the lemmas it cites, its oracles and its bill.
%
\vnbprf{\textbf{Proof.} \texttt{theorem-\allowbreak{}library/directional-\allowbreak{}derivative.scm} --- 51 recorded steps.}
%
\vnbprf{\textbf{Lemmas.} \texttt{rr-zero-in}, \texttt{rr-one-in}, \texttt{dir-diff-value-in-rr}, \texttt{nvs-vadd-in-vec}, \texttt{seg-curve-in-fun}, \texttt{rr-add-closed}, \texttt{dir-quotient-in-fun}, \texttt{dir-quotient-at-zero}, \texttt{dir-quotient-continuous}, \texttt{dir-quotient-identity}, \texttt{diff-at-local}.}
%
\vnbprf{\textbf{Oracles.} \texttt{arith}, \texttt{ineq}, \texttt{crs}.  \textbf{Bill:} \texttt{proven modulo 0}.}
%
\end{proof}
%
\noindent\textbf{Deviation from the notes.} f is required to be in fun(vec(m), rr), i.e. TOTAL on the normed vector space, and the existence hypothesis is global: forall x in vec(m). forsome k. is-dir-diff-at(m, f, x, eta, k), and likewise for xi, where the notes ask only that d\_eta f and d\_xi f exist in a NEIGHBOURHOOD of a. The continuity hypothesis is stated at the single point a -- is-continuous-at of the maps x |-$>$ dir-deriv(m, f, x, eta) and x |-$>$ dir-deriv(m, f, x, xi) at a -- where the notes ask for continuity throughout the neighbourhood. The ambient space is an arbitrary normed vector space, not E = R\textasciicircum{}n. The conclusion, that d\_(eta + xi) f(a) is DEFINED and equals the sum, is carried by the relational is-dir-diff-at and so is faithful.
%
\renewcommand{\thethm}{2.27}
\begin{prop}[{Linearity of the directional derivative on the span}]\label{thm:dir-deriv-linear}
%
\vnbstmt{\begin{array}{@{}l@{}} \forall m, f, a, b, i, r, s, l, c. \\ \quad \operatorname{is\text{-}normed\text{-}vector\text{-}space}(m) \wedge \\ \quad (f \in (\operatorname{vec}(m) \to \mathbb{R})) \wedge \\ \quad (a \in \operatorname{vec}(m)) \wedge \\ \quad (b \in \operatorname{vec}(m)) \wedge \\ \quad (i \in \operatorname{vec}(m)) \wedge \\ \quad (r \in \mathbb{R}) \wedge \\ \quad (s \in \mathbb{R}) \wedge \\ \quad \operatorname{is\text{-}dir\text{-}diff\text{-}at}(m, f, a, b, l) \wedge \\ \quad \operatorname{is\text{-}dir\text{-}diff\text{-}at}(m, f, a, i, c) \wedge \\ \quad \forall x \in \operatorname{vec}(m). \exists k. \\ \quad \quad \operatorname{is\text{-}dir\text{-}diff\text{-}at}(m, f, x, \operatorname{act}(m)(r, b), k) \wedge \\ \quad \forall x \in \operatorname{vec}(m). \exists k. \\ \quad \quad \operatorname{is\text{-}dir\text{-}diff\text{-}at}(m, f, x, \operatorname{act}(m)(s, i), k) \wedge \\ \quad \operatorname{is\text{-}continuous\text{-}at}(\operatorname{nvs\text{-}metric\text{-}space}(m), \operatorname{rr\text{-}ms}, \\ \quad \quad (\lambda x \in \operatorname{vec}(m).\; \\ \quad \quad \quad \operatorname{dir\text{-}deriv}(m, f, x, \operatorname{act}(m)(r, b))), \\ \quad \quad a) \wedge \\ \quad \operatorname{is\text{-}continuous\text{-}at}(\operatorname{nvs\text{-}metric\text{-}space}(m), \operatorname{rr\text{-}ms}, \\ \quad \quad (\lambda x \in \operatorname{vec}(m).\; \\ \quad \quad \quad \operatorname{dir\text{-}deriv}(m, f, x, \operatorname{act}(m)(s, i))), \\ \quad \quad a) \Rightarrow \\ \quad \quad \operatorname{is\text{-}dir\text{-}diff\text{-}at}(m, f, a, \operatorname{vadd}(m)(\operatorname{act}(m)(r, b), \operatorname{act}(m)(s, i)), ((r \cdot l) + (s \cdot c))) \end{array}}
%
\end{prop}
%
\begin{proof}
%
The proofs: the file that holds each, the lemmas it cites, its oracles and its bill.
%
\vnbprf{\textbf{Proof.} \texttt{theorem-\allowbreak{}library/directional-\allowbreak{}derivative.scm} --- 118 recorded steps.}
%
\vnbprf{\textbf{Lemmas.} \texttt{rr-mul-closed}, \texttt{nvs-act-in-vec}, \texttt{dir-deriv-homogeneous}, \texttt{dir-deriv-additive}.}
%
\vnbprf{\textbf{Oracles.} \texttt{crs}, \texttt{arith}, \texttt{ineq}.  \textbf{Bill:} \texttt{proven modulo 0}.}
%
\end{proof}
%
\noindent\textbf{Deviation from the notes.} The notes quantify over ALL x in U and ALL phi in the vector space V generated by eta and xi, and assert that phi |-$>$ d\_phi f(x) is a linear MAP. dir-deriv-linear states only the single instance at the single point a: from is-dir-diff-at at a in directions eta and xi it concludes is-dir-diff-at(m, f, a, r*eta + s*xi, r*dl + s*dm), and it does so under extra hypotheses the notes do not carry at that level -- that the directional derivatives in directions r*eta and s*xi exist at EVERY point of vec(m) and that the two derivative maps are continuous at a. f is required total on vec(m). No span or generated-subspace vocabulary is used, and no statement says the assignment is linear as a map.
%
\renewcommand{\thethm}{2.28}
\begin{prop}[{Continuous partials give a directional derivative in every direction}]\label{thm:calculus-2.28}
%
\textit{U contained in E, f: U --$>$ R, and d\_(e\_i) f(x) for i = 1, ..., n defined and continuous as functions of x in U: then for all eta in E and x in U, d\_eta f(x) is defined, is linear in eta and continuous in x.}
%
Not in the library.
%
\end{prop}
%
\begin{proof}
%
The proofs: the file that holds each, the lemmas it cites, its oracles and its bill.
%
\vnbprf{There is no library statement, and so no proof.}
%
\end{proof}
%
\noindent\textbf{Deviation from the notes.} The space E = R\textasciicircum{}n with the sup norm max(abs(eta\_1), ..., abs(eta\_n)) does not exist in the library, and neither do the unit basis vectors e\_i nor the partial derivatives d\_(e\_i) f. The nearest facts are the two-direction lemmas dir-deriv-additive and dir-deriv-homogeneous over an abstract normed vector space, which are what an induction on k would use; the induction over a basis, the conclusion that d\_eta f is linear in eta, and the conclusion that it is continuous in x are all absent.
%
\renewcommand{\thethm}{2.29}
\begin{prop}[{Continuous partials imply differentiability}]\label{thm:calculus-2.29}
%
\textit{f such that the partial derivatives d\_(e\_i) f(x) are defined and continuous in x in a neighborhood of a: then for any eps $>$ 0 there is delta such that for norm(eta) $<$= delta, abs(f(a + eta) - f(a) - sum over i of eta\_i d\_(e\_i) f(a)) $<$= norm(eta) eps.}
%
Not in the library.
%
\end{prop}
%
\begin{proof}
%
The proofs: the file that holds each, the lemmas it cites, its oracles and its bill.
%
\vnbprf{There is no library statement, and so no proof.}
%
\end{proof}
%
\noindent\textbf{Deviation from the notes.} Nothing states total differentiability of a function on R\textasciicircum{}n from the continuity of its partial derivatives; there is no R\textasciicircum{}n with the sup norm, no basis vectors e\_i, no Df(a) = (d\_(e\_1) f(a), ..., d\_(e\_n) f(a)), and no estimate of the form abs(f(a + eta) - f(a) - sum eta\_i d\_(e\_i) f(a)) $<$= norm(eta) eps. The whole first-order differentiability theory in the tree is scalar: is-diff-at over RR and is-diff-on over a normed field. CLAUDE.md records Props 2.28 and 2.29 as the next terrain, after the metric subspace structure.
%
\renewcommand{\thethm}{3.3}
\begin{prop}[{Unions and finite intersections of open sets}]\label{thm:union-of-opens-open}
%
\leavevmode
\vnbpart{\texttt{union-of-opens-open}}
\vnbstmt{\begin{array}{@{}l@{}} \forall s, a, g. \\ \quad \operatorname{is\text{-}metric\text{-}space}(s) \wedge \\ \quad \forall i \in a.\; \operatorname{is\text{-}open}(s, g(i)) \Rightarrow \\ \quad \quad \operatorname{is\text{-}open}(s, \operatorname{big\text{-}union}(i, a, g(i))) \end{array}}
\vnbpart{\texttt{inter-of-opens-open}}
\vnbstmt{\begin{array}{@{}l@{}} \forall s, u, w. \\ \quad \operatorname{is\text{-}metric\text{-}space}(s) \wedge \\ \quad \operatorname{is\text{-}open}(s, u) \wedge \\ \quad \operatorname{is\text{-}open}(s, w) \Rightarrow \\ \quad \quad \operatorname{is\text{-}open}(s, (u \cap w)) \end{array}}
%
\end{prop}
%
\begin{proof}
%
The proofs: the file that holds each, the lemmas it cites, its oracles and its bill.
%
\vnbprf{\vnbproofpart{\texttt{union-of-opens-open}}}
%
\vnbprf{\textbf{Proof.} \texttt{theorem-\allowbreak{}library/rake-\allowbreak{}open-\allowbreak{}sets.scm} --- 43 recorded steps.}
%
\vnbprf{\textbf{Lemmas.} \texttt{subset-def}, \texttt{subset-mem-fwd}.}
%
\vnbprf{\textbf{Oracles.} none.  \textbf{Bill:} \texttt{proven modulo 0}.}
%
\vnbprf{\vnbproofpart{\texttt{inter-of-opens-open}}}
%
\vnbprf{\textbf{Proof.} \texttt{theorem-\allowbreak{}library/rake-\allowbreak{}open-\allowbreak{}sets.scm} --- 114 recorded steps.}
%
\vnbprf{\textbf{Lemmas.} \texttt{subset-def}, \texttt{intersection-membership}, \texttt{subset-mem-fwd}, \texttt{rr-pos-rr-in-rr}, \texttt{rr-lt-of-pos-rr}, \texttt{rr-min-pos}, \texttt{rr-pos-rr-of-lt}, \texttt{ball-sep-unfold}, \texttt{metric-dist-real}.}
%
\vnbprf{\textbf{Oracles.} \texttt{ineq}.  \textbf{Bill:} \texttt{proven modulo 0}.}
%
\end{proof}
%
\noindent\textbf{Deviation from the notes.} The intersection half is proven only for TWO open sets (inter-of-opens-open: is-open(s,u) and is-open(s,w) implies is-open(s, intersection(u,w))); there is no statement for the intersection of a finite family, which is what the notes assert. The union half is stated for a family presented as an index class a with an index map g, forall i in a, is-open(s, g(i)) implies is-open(s, big-union(i, a, g(i))), rather than for a family of sets; both carry is-metric-space(s) as an explicit guard.
%
\renewcommand{\thethm}{3.7}
\begin{prop}[{A monotone bounded sequence converges to its sup (resp. inf)}]\label{thm:monotone-convergence-rr}
%
\vnbstmt{\begin{array}{@{}l@{}} \forall f \in (\mathbb{N} \to \mathbb{R}). \\ \quad \forall k \in \mathbb{N}.\; (f(k) \leq f((k + 1))) \wedge \\ \quad \exists d \in \mathbb{R}.\; \forall k \in \mathbb{N}.\; (f(k) \leq d) \Rightarrow \\ \quad \quad \operatorname{converges}(\operatorname{rr\text{-}ms}, f) \end{array}}
%
\end{prop}
%
\begin{proof}
%
The proofs: the file that holds each, the lemmas it cites, its oracles and its bill.
%
\vnbprf{\textbf{Proof.} \texttt{theorem-\allowbreak{}library/monotone-\allowbreak{}convergence-\allowbreak{}proof.scm} --- 123 recorded steps.}
%
\vnbprf{\textbf{Lemmas.} \texttt{nn-zero-in}, \texttt{fun-apply-type-c}, \texttt{subset-def}, \texttt{rr-sup-in}, \texttt{rr-sup-upper}, \texttt{rr-is-metric-space}, \texttt{rr-sup-approx}, \texttt{nn-monotone-step-implies-le}, \texttt{rr-sub-in-rr}, \texttt{rr-ms-dist}, \texttt{rr-abs-bound}.}
%
\vnbprf{\textbf{Oracles.} \texttt{ineq}, \texttt{arith}, \texttt{crs}.  \textbf{Bill:} \texttt{proven modulo 0}.}
%
\end{proof}
%
\noindent\textbf{Deviation from the notes.} The library proves only that a non-decreasing sequence bounded above CONVERGES; it does not identify the limit with sup\_n x\_n, which is the content of the notes' statement. The non-increasing / bounded-below / inf half is absent entirely. The sequence must be typed f in fun(nn, rr), the monotonicity is given in step form forall k in nn, f(k) $<$= f(succ(k)), and the bound as forsome bnd in rr, forall k in nn, f(k) $<$= bnd.
%
\renewcommand{\thethm}{3.9}
\begin{prop}[{Closed sets are exactly those stable under limits of sequences}]\label{thm:closed-limit-in}
%
\leavevmode
\vnbpart{\texttt{closed-limit-in}}
\vnbstmt{\begin{array}{@{}l@{}} \forall s, a.; \forall f \in (\mathbb{N} \to a), l. \\ \quad \operatorname{is\text{-}metric\text{-}space}(s) \wedge \\ \quad \operatorname{is\text{-}closed}(s, a) \wedge \\ \quad \operatorname{converges\text{-}to}(s, f, l) \Rightarrow \\ \quad \quad (l \in a) \end{array}}
\vnbpart{\texttt{closed-set-limit-in}}
\vnbstmt{\begin{array}{@{}l@{}} \forall s, a, f, l. \\ \quad \operatorname{is\text{-}closed}(s, a) \wedge \\ \quad \operatorname{converges\text{-}to}(s, f, l) \wedge \\ \quad \forall n \in \mathbb{N}.\; (f(n) \in a) \Rightarrow \\ \quad \quad (l \in a) \end{array}}
\vnbpart{\texttt{closed-set-limit-in-tail}}
\vnbstmt{\begin{array}{@{}l@{}} \forall s, v, \mathit{fv}, \mathit{lv}, b. \\ \quad \operatorname{is\text{-}closed}(s, v) \wedge \\ \quad \operatorname{converges\text{-}to}(s, \mathit{fv}, \mathit{lv}) \wedge \\ \quad (b \in \mathbb{N}) \wedge \\ \quad \forall n \in \mathbb{N}. \\ \quad \quad ((b \leq n) \Rightarrow (\operatorname{fv}(n) \in v)) \Rightarrow \\ \quad \quad (\mathit{lv} \in v) \end{array}}
%
\end{prop}
%
\begin{proof}
%
The proofs: the file that holds each, the lemmas it cites, its oracles and its bill.
%
\vnbprf{\vnbproofpart{\texttt{closed-limit-in}}}
%
\vnbprf{\textbf{Proof.} \texttt{theorem-\allowbreak{}library/metric-\allowbreak{}subspace-\allowbreak{}laws.scm} --- 105 recorded steps.}
%
\vnbprf{\textbf{Lemmas.} \texttt{complement-in-membership}, \texttt{rr-pos-halvable}, \texttt{rr-lt-of-pos-rr}, \texttt{rr-pos-rr-in-rr}, \texttt{nn-in-rr}, \texttt{rr-leq-reflexive}, \texttt{fun-apply-type-c}, \texttt{subset-mem-fwd}, \texttt{metric-dist-real}, \texttt{metric-sym}, \texttt{ball-mem-from-le}.}
%
\vnbprf{\textbf{Oracles.} \texttt{ineq}, \texttt{crs}.  \textbf{Bill:} \texttt{proven modulo 0}.}
%
\vnbprf{\vnbproofpart{\texttt{closed-set-limit-in}}}
%
\vnbprf{\textbf{Proof.} \texttt{theorem-\allowbreak{}library/metric-\allowbreak{}closure-\allowbreak{}laws.scm} --- 93 recorded steps.}
%
\vnbprf{\textbf{Lemmas.} \texttt{complement-in-membership}, \texttt{rr-pos-rr-in-rr}, \texttt{rr-lt-of-pos-rr}, \texttt{ball-is-set}, \texttt{rr-pos-halvable}, \texttt{nn-le-refl}, \texttt{subset-mem-fwd}, \texttt{metric-dist-real}, \texttt{metric-sym}, \texttt{ball-mem-from-le}.}
%
\vnbprf{\textbf{Oracles.} \texttt{ineq}, \texttt{crs}.  \textbf{Bill:} \texttt{proven modulo 0}.}
%
\vnbprf{\vnbproofpart{\texttt{closed-set-limit-in-tail}}}
%
\vnbprf{\textbf{Proof.} \texttt{theorem-\allowbreak{}library/rake-\allowbreak{}baire-\allowbreak{}2.scm} --- 98 recorded steps.}
%
\vnbprf{\textbf{Lemmas.} \texttt{complement-in-membership}, \texttt{rr-pos-rr-in-rr}, \texttt{rr-lt-of-pos-rr}, \texttt{ball-is-set}, \texttt{rr-pos-halvable}, \texttt{nn-max-closed}, \texttt{nn-in-rr}, \texttt{rr-le-max-left}, \texttt{rr-le-max-right}, \texttt{subset-mem-fwd}, \texttt{metric-dist-real}, \texttt{metric-sym}, \texttt{ball-mem-from-le}.}
%
\vnbprf{\textbf{Oracles.} \texttt{ineq}, \texttt{crs}.  \textbf{Bill:} \texttt{proven modulo 0}.}
%
\end{proof}
%
\noindent\textbf{Deviation from the notes.} Only the FORWARD direction is proven: if A is closed and a sequence of elements of A converges, its limit lies in A. The CONVERSE -- if every convergent sequence of elements of A converges to an element of A then A is closed -- is absent, and it is the half whose proof in the notes constructs the sequence. closed-limit-in types the sequence as f in fun(nn, a); closed-set-limit-in states membership pointwise as forall n in nn, f(n) in a; closed-set-limit-in-tail only from some index onwards.
%
\renewcommand{\thethm}{3.10}
\begin{prop}[{A set is closed if and only if it contains its accumulation points}]\label{thm:calculus-3.10}
%
\textit{A necessary and sufficient condition that A subset X be closed is that it contain all its accumulation points, x being an accumulation point of A when every ball B(x,r) contains an element of A distinct from x.}
%
Not in the library.
%
\end{prop}
%
\begin{proof}
%
The proofs: the file that holds each, the lemmas it cites, its oracles and its bill.
%
\vnbprf{There is no library statement, and so no proof.}
%
\end{proof}
%
\noindent\textbf{Deviation from the notes.} The library has no notion of accumulation point (cluster-point is for SEQUENCES, not for sets), and no statement of this equivalence. The nearest facts are about the closure: closure-membership (x in closure(s,a) iff x in pts(s) and every open u containing x meets a), closure-is-closed, closure-contains and closure-in-closed.
%
\renewcommand{\thethm}{3.12}
\begin{prop}[{Four equivalent characterizations of a compact metric space}]\label{thm:compact-iff-fip}
%
\leavevmode
\vnbpart{\texttt{compact-iff-fip}}
\vnbstmt{\begin{array}{@{}l@{}} \forall s. \\ \quad \operatorname{is\text{-}metric\text{-}space}(s) \Rightarrow  \\ \quad \quad (\operatorname{is\text{-}compact}(s) \Leftrightarrow \forall c.\; (\operatorname{has\text{-}fip}(s, c) \Rightarrow \exists p.\; (p \in \operatorname{intersection\text{-}of}(c)))) \end{array}}
\vnbpart{\texttt{compact-iff-cluster-point}}
\vnbstmt{\begin{array}{@{}l@{}} \forall s. \\ \quad \operatorname{is\text{-}metric\text{-}space}(s) \Rightarrow  \\ \quad \quad (\operatorname{is\text{-}compact}(s) \Leftrightarrow \forall f \in (\mathbb{N} \to \operatorname{pts}(s)).\; \exists x.\; \operatorname{cluster\text{-}point}(s, f, x)) \end{array}}
\vnbpart{\texttt{compact-iff-seq-compact}}
\vnbstmt{\begin{array}{@{}l@{}} \forall s. \\ \quad \operatorname{is\text{-}metric\text{-}space}(s) \Rightarrow  \\ \quad \quad (\operatorname{is\text{-}compact}(s) \Leftrightarrow \operatorname{seq\text{-}compact}(s)) \end{array}}
\vnbpart{\texttt{compact-iff-tb-complete}}
\vnbstmt{\begin{array}{@{}l@{}} \forall s. \\ \quad \operatorname{is\text{-}metric\text{-}space}(s) \Rightarrow  \\ \quad \quad (\operatorname{is\text{-}compact}(s) \Leftrightarrow (\operatorname{totally\text{-}bounded}(s) \wedge \operatorname{is\text{-}complete}(s))) \end{array}}
%
\end{prop}
%
\begin{proof}
%
The proofs: the file that holds each, the lemmas it cites, its oracles and its bill.
%
\vnbprf{\vnbproofpart{\texttt{compact-iff-fip}}}
%
\vnbprf{\textbf{Proof.} \texttt{theorem-\allowbreak{}library/rake-\allowbreak{}compact-\allowbreak{}fip.scm} --- 6 recorded steps.}
%
\vnbprf{\textbf{Lemmas.} \texttt{compact-implies-fip}, \texttt{fip-implies-compact}.}
%
\vnbprf{\textbf{Oracles.} \texttt{ineq}, \texttt{arith}.  \textbf{Bill:} \texttt{proven modulo 0}.}
%
\vnbprf{\vnbproofpart{\texttt{compact-iff-cluster-point}}}
%
\vnbprf{\textbf{Proof.} \texttt{theorem-\allowbreak{}library/rake-\allowbreak{}compact-\allowbreak{}equivalences.scm} --- 58 recorded steps.}
%
\vnbprf{\textbf{Lemmas.} \texttt{compact-seq-has-cluster}, \texttt{cluster-point-has-convergent-subseq}, \texttt{seq-compact-implies-compact}.}
%
\vnbprf{\textbf{Oracles.} \texttt{ineq}, \texttt{arith}, \texttt{crs}.  \textbf{Bill:} \texttt{proven modulo 0}.}
%
\vnbprf{\vnbproofpart{\texttt{compact-iff-seq-compact}}}
%
\vnbprf{\textbf{Proof.} \texttt{theorem-\allowbreak{}library/rake-\allowbreak{}lebesgue-\allowbreak{}number.scm} --- 6 recorded steps.}
%
\vnbprf{\textbf{Lemmas.} \texttt{compact-implies-seq-compact}, \texttt{seq-compact-implies-compact}.}
%
\vnbprf{\textbf{Oracles.} \texttt{ineq}, \texttt{arith}, \texttt{crs}.  \textbf{Bill:} \texttt{proven modulo 0}.}
%
\vnbprf{\vnbproofpart{\texttt{compact-iff-tb-complete}}}
%
\vnbprf{\textbf{Proof.} \texttt{theorem-\allowbreak{}library/rake-\allowbreak{}compact-\allowbreak{}equivalences.scm} --- 68 recorded steps.}
%
\vnbprf{\textbf{Lemmas.} \texttt{compact-implies-totally-bounded}, \texttt{compact-implies-complete}, \texttt{totally-bounded-has-cauchy-subsequence}, \texttt{subseq-is-fun}, \texttt{seq-compact-implies-compact}.}
%
\vnbprf{\textbf{Oracles.} \texttt{ineq}, \texttt{arith}, \texttt{crs}.  \textbf{Bill:} \texttt{proven modulo 0}.}
%
\end{proof}
%
\noindent\textbf{Deviation from the notes.} All four equivalences exist, but as four separate two-way statements each guarded by is-metric-space(s), not as one chain. The finite intersection property is relativised to the library predicate has-fip(s,c), which adds two hypotheses the notes do not make: the family c must be INHABITED (forsome a0, a0 in c), and only INHABITED finite subfamilies f are required to have a common point, so the notes' case F = empty (which forces X non-empty) is not covered. The library also proves the sequentially-compact form compact-iff-seq-compact (every sequence has a CONVERGENT SUBSEQUENCE), which the notes do not state; their clause (3) is compact-iff-cluster-point.
%
\renewcommand{\thethm}{3.14}
\begin{prop}[{Continuity at a point is sequential continuity}]\label{thm:continuous-at-iff-sequential}
%
\leavevmode
\vnbpart{\texttt{continuous-at-iff-sequential}}
\vnbstmt{\begin{array}{@{}l@{}} \forall s, t, f, a. \\ \quad \operatorname{is\text{-}metric\text{-}space}(s) \wedge \\ \quad \operatorname{is\text{-}metric\text{-}space}(t) \wedge \\ \quad (f \in (\operatorname{pts}(s) \to \operatorname{pts}(t))) \wedge \\ \quad (a \in \operatorname{pts}(s)) \Rightarrow \\ \quad \quad (\operatorname{is\text{-}continuous\text{-}at}(s, t, f, a) \Leftrightarrow \forall q \in \operatorname{pts}(s)^{\mathbb{N}}.\; (\operatorname{converges\text{-}to}(s, q, a) \Rightarrow \operatorname{converges\text{-}to}(t, \operatorname{compose}(f, q), f(a)))) \end{array}}
\vnbpart{\texttt{continuous-at-implies-sequential}}
\vnbstmt{\begin{array}{@{}l@{}} \forall s, t, f, a.; \forall q \in \operatorname{pts}(s)^{\mathbb{N}}. \\ \quad \operatorname{is\text{-}metric\text{-}space}(s) \wedge \\ \quad \operatorname{is\text{-}metric\text{-}space}(t) \wedge \\ \quad (f \in (\operatorname{pts}(s) \to \operatorname{pts}(t))) \wedge \\ \quad (a \in \operatorname{pts}(s)) \wedge \\ \quad \operatorname{is\text{-}continuous\text{-}at}(s, t, f, a) \wedge \\ \quad \operatorname{converges\text{-}to}(s, q, a) \Rightarrow \\ \quad \quad \operatorname{converges\text{-}to}(t, \operatorname{compose}(f, q), f(a)) \end{array}}
%
\end{prop}
%
\begin{proof}
%
The proofs: the file that holds each, the lemmas it cites, its oracles and its bill.
%
\vnbprf{\vnbproofpart{\texttt{continuous-at-iff-sequential}}}
%
\vnbprf{\textbf{Proof.} \texttt{theorem-\allowbreak{}library/sequential-\allowbreak{}continuity.scm} --- 16 recorded steps.}
%
\vnbprf{\textbf{Lemmas.} \texttt{continuous-at-implies-sequential}, \texttt{sequential-implies-continuous-at}.}
%
\vnbprf{\textbf{Oracles.} \texttt{ineq}, \texttt{crs}, \texttt{arith}.  \textbf{Bill:} \texttt{proven modulo 0}.}
%
\vnbprf{\vnbproofpart{\texttt{continuous-at-implies-sequential}}}
%
\vnbprf{\textbf{Proof.} \texttt{theorem-\allowbreak{}library/sequential-\allowbreak{}continuity.scm} --- 72 recorded steps.}
%
\vnbprf{\textbf{Lemmas.} \texttt{fun-apply-type-c}, \texttt{nn-is-set}, \texttt{compose-type}, \texttt{metric-sym}, \texttt{compose-apply}.}
%
\vnbprf{\textbf{Oracles.} none.  \textbf{Bill:} \texttt{proven modulo 0}.}
%
\end{proof}
%
\noindent\textbf{Deviation from the notes.} The library carries the typings as explicit antecedents -- is-metric-space(s), is-metric-space(t), f in fun(pts(s), pts(t)) and a in pts(s) -- where the notes get them from the ambient sentence f : X --$>$ Y. The sequence is quantified over sqn(pts(s)) and the image sequence is written compose(f, sq).
%
\renewcommand{\thethm}{3.15}
\begin{prop}[{Continuity equals closedness (openness) of preimages}]\label{thm:continuous-implies-closed-preimage}
%
\leavevmode
\vnbpart{\texttt{continuous-implies-closed-preimage}}
\vnbstmt{\begin{array}{@{}l@{}} \forall s, t, f, a. \\ \quad (\operatorname{is\text{-}continuous}(s, t, f) \wedge \operatorname{is\text{-}closed}(t, a)) \Rightarrow  \\ \quad \quad \operatorname{is\text{-}closed}(s, \operatorname{preimage}(s, f, a)) \end{array}}
\vnbpart{\texttt{closed-preimage-implies-continuous}}
\vnbstmt{\begin{array}{@{}l@{}} \forall s, t, f. \\ \quad \operatorname{is\text{-}metric\text{-}space}(s) \wedge \\ \quad \operatorname{is\text{-}metric\text{-}space}(t) \wedge \\ \quad (f \in (\operatorname{pts}(s) \to \operatorname{pts}(t))) \wedge \\ \quad \forall a. \\ \quad \quad \operatorname{is\text{-}closed}(t, a) \Rightarrow  \\ \quad \quad \quad \operatorname{is\text{-}closed}(s, \operatorname{preimage}(s, f, a)) \Rightarrow \\ \quad \quad \operatorname{is\text{-}continuous}(s, t, f) \end{array}}
\vnbpart{\texttt{continuous-implies-open-preimage}}
\vnbstmt{\begin{array}{@{}l@{}} \forall s, t, f, v. \\ \quad (\operatorname{is\text{-}continuous}(s, t, f) \wedge \operatorname{is\text{-}open}(t, v)) \Rightarrow  \\ \quad \quad \operatorname{is\text{-}open}(s, \operatorname{preimage}(s, f, v)) \end{array}}
\vnbpart{\texttt{open-preimage-implies-continuous}}
\vnbstmt{\begin{array}{@{}l@{}} \forall s, t, f. \\ \quad \operatorname{is\text{-}metric\text{-}space}(s) \wedge \\ \quad \operatorname{is\text{-}metric\text{-}space}(t) \wedge \\ \quad (f \in (\operatorname{pts}(s) \to \operatorname{pts}(t))) \wedge \\ \quad \forall v. \\ \quad \quad \operatorname{is\text{-}open}(t, v) \Rightarrow  \\ \quad \quad \quad \operatorname{is\text{-}open}(s, \operatorname{preimage}(s, f, v)) \Rightarrow \\ \quad \quad \operatorname{is\text{-}continuous}(s, t, f) \end{array}}
%
\end{prop}
%
\begin{proof}
%
The proofs: the file that holds each, the lemmas it cites, its oracles and its bill.
%
\vnbprf{\vnbproofpart{\texttt{continuous-implies-closed-preimage}}}
%
\vnbprf{\textbf{Proof.} \texttt{theorem-\allowbreak{}library/rake-\allowbreak{}open-\allowbreak{}sets.scm} --- 46 recorded steps.}
%
\vnbprf{\textbf{Lemmas.} \texttt{subset-def}, \texttt{preimage-complement}, \texttt{eq-sym}, \texttt{continuous-implies-open-preimage}.}
%
\vnbprf{\textbf{Oracles.} \texttt{ineq}, \texttt{crs}.  \textbf{Bill:} \texttt{proven modulo 0}.}
%
\vnbprf{\vnbproofpart{\texttt{closed-preimage-implies-continuous}}}
%
\vnbprf{\textbf{Proof.} \texttt{theorem-\allowbreak{}library/rake-\allowbreak{}open-\allowbreak{}sets.scm} --- 110 recorded steps.}
%
\vnbprf{\textbf{Lemmas.} \texttt{complement-in-membership}, \texttt{subset-mem-fwd}, \texttt{class-extensionality}, \texttt{complement-in-subset}, \texttt{preimage-complement}, \texttt{subset-def}, \texttt{eq-sym}, \texttt{open-preimage-implies-continuous}.}
%
\vnbprf{\textbf{Oracles.} \texttt{ineq}, \texttt{crs}.  \textbf{Bill:} \texttt{proven modulo 0}.}
%
\vnbprf{\vnbproofpart{\texttt{continuous-implies-open-preimage}}}
%
\vnbprf{\textbf{Proof.} \texttt{theorem-\allowbreak{}library/rake-\allowbreak{}cont-\allowbreak{}preimage.scm} --- 112 recorded steps.}
%
\vnbprf{\textbf{Lemmas.} \texttt{subset-def}, \texttt{rr-pos-rr-in-rr}, \texttt{rr-pos-halvable}, \texttt{rr-lt-of-pos-rr}, \texttt{ball-membership}, \texttt{fun-apply-type-c}, \texttt{ball-mem-from-le}, \texttt{subset-mem-fwd}.}
%
\vnbprf{\textbf{Oracles.} \texttt{ineq}, \texttt{crs}.  \textbf{Bill:} \texttt{proven modulo 0}.}
%
\vnbprf{\vnbproofpart{\texttt{open-preimage-implies-continuous}}}
%
\vnbprf{\textbf{Proof.} \texttt{theorem-\allowbreak{}library/rake-\allowbreak{}open-\allowbreak{}sets.scm} --- 127 recorded steps.}
%
\vnbprf{\textbf{Lemmas.} \texttt{fun-apply-type-c}, \texttt{rr-pos-rr-in-rr}, \texttt{rr-lt-of-pos-rr}, \texttt{metric-self-zero}, \texttt{ball-is-open}, \texttt{rr-pos-halvable}, \texttt{metric-dist-real}, \texttt{subset-mem-fwd}, \texttt{ball-sep-unfold}.}
%
\vnbprf{\textbf{Oracles.} \texttt{ineq}, \texttt{crs}.  \textbf{Bill:} \texttt{proven modulo 0}.}
%
\end{proof}
%
\noindent\textbf{Deviation from the notes.} Stated as four separate implications rather than as one three-way equivalence, so the notes' (2) iff (3) is not a library statement. The two converse directions carry the extra typing hypotheses is-metric-space(s), is-metric-space(t) and f in fun(pts(s), pts(t)); the forward directions read them off is-continuous(s,t,f), whose defining iff contains them.
%
\renewcommand{\thethm}{3.16}
\begin{lem}[{Subsequence principle}]\label{thm:subsequence-principle}
%
\vnbstmt{\begin{array}{@{}l@{}} \forall s, f, l. \\ \quad \operatorname{is\text{-}metric\text{-}space}(s) \wedge \\ \quad (f \in (\mathbb{N} \to \operatorname{pts}(s))) \wedge \\ \quad (l \in \operatorname{pts}(s)) \wedge \\ \quad \forall i. \\ \quad \quad \operatorname{strictly\text{-}mono\text{-}nn}(i) \Rightarrow  \\ \quad \quad \quad \exists a. \\ \quad \quad \quad \quad \operatorname{strictly\text{-}mono\text{-}nn}(a) \wedge \\ \quad \quad \quad \quad \operatorname{converges\text{-}to}(s, \operatorname{subseq}(\operatorname{subseq}(f, i), a), l) \Rightarrow \\ \quad \quad \operatorname{converges\text{-}to}(s, f, l) \end{array}}
%
\end{lem}
%
\begin{proof}
%
The proofs: the file that holds each, the lemmas it cites, its oracles and its bill.
%
\vnbprf{\textbf{Proof.} \texttt{theorem-\allowbreak{}library/subsequence-\allowbreak{}principle.scm} --- 128 recorded steps.}
%
\vnbprf{\textbf{Lemmas.} \texttt{subset-def}, \texttt{nn-finite-subset-bounded}, \texttt{subsequence-capture}, \texttt{fun-codomain-superset}, \texttt{nn-le-refl}, \texttt{fun-apply-type-c}.}
%
\vnbprf{\textbf{Oracles.} \texttt{ineq}, \texttt{arith}.  \textbf{Bill:} \texttt{proven modulo 0}.}
%
\end{proof}
%
\noindent\textbf{Deviation from the notes.} Explicit typings is-metric-space(s), f in fun(nn, pts(s)) and l in pts(s) are added as antecedents, and subsequences are presented as strictly-mono-nn reindexings composed by subseq(subseq(f, phi), psi) rather than as sequences in their own right.
%
\renewcommand{\thethm}{3.19}
\begin{prop}[{The continuous image of a compact set is compact}]\label{thm:continuous-image-of-compact}
%
\leavevmode
\vnbpart{\texttt{continuous-image-of-compact}}
\vnbstmt{\begin{array}{@{}l@{}} \forall s, t, f, a. \\ \quad \operatorname{is\text{-}continuous}(s, t, f) \wedge \\ \quad (a \subseteq \operatorname{pts}(s)) \wedge \\ \quad \operatorname{is\text{-}compact}(\operatorname{subspace\text{-}ms}(s, a)) \Rightarrow \\ \quad \quad \operatorname{is\text{-}compact}(\operatorname{subspace\text{-}ms}(t, \operatorname{image}(f, a))) \end{array}}
\vnbpart{\texttt{compact-image}}
\vnbstmt{\begin{array}{@{}l@{}} \forall s, t, f. \\ \quad (\operatorname{is\text{-}continuous}(s, t, f) \wedge \operatorname{is\text{-}compact}(s)) \Rightarrow  \\ \quad \quad \operatorname{is\text{-}compact}(\operatorname{subspace\text{-}ms}(t, \operatorname{image}(f, \operatorname{pts}(s)))) \end{array}}
%
\end{prop}
%
\begin{proof}
%
The proofs: the file that holds each, the lemmas it cites, its oracles and its bill.
%
\vnbprf{\vnbproofpart{\texttt{continuous-image-of-compact}}}
%
\vnbprf{\textbf{Proof.} \texttt{theorem-\allowbreak{}library/compact-\allowbreak{}image.scm} --- 60 recorded steps.}
%
\vnbprf{\textbf{Lemmas.} \texttt{subclass-of-set-is-set}, \texttt{subspace-pts}, \texttt{restrict-in-fun}, \texttt{subset-mem-fwd}, \texttt{restrict-continuous-at}, \texttt{compact-image}, \texttt{image-restrict}.}
%
\vnbprf{\textbf{Oracles.} \texttt{arith}, \texttt{ineq}, \texttt{crs}.  \textbf{Bill:} \texttt{proven modulo 0}.}
%
\vnbprf{\vnbproofpart{\texttt{compact-image}}}
%
\vnbprf{\textbf{Proof.} \texttt{theorem-\allowbreak{}library/compact-\allowbreak{}image.scm} --- 300 recorded steps.}
%
\vnbprf{\textbf{Lemmas.} \texttt{fun-apply-type-c}, \texttt{membership-implies-sethood}, \texttt{subset-def}, \texttt{subspace-is-metric-space}, \texttt{image-set}, \texttt{power-set}, \texttt{subspace-pts}, \texttt{preimage-subspace-open}, \texttt{preimage-in-power}, \texttt{power-mem-in}, \texttt{class-extensionality}, \texttt{subclass-of-set-is-set}, \texttt{card-image-finite}, \texttt{card-subset-nn}, \texttt{subset-mem-fwd}, \texttt{image-of-preimage-onto}.}
%
\vnbprf{\textbf{Oracles.} \texttt{arith}, \texttt{ineq}, \texttt{crs}.  \textbf{Bill:} \texttt{proven modulo 0}.}
%
\end{proof}
%
\noindent\textbf{Deviation from the notes.} Proven 2026-09-22. is-compact is a predicate of a metric SPACE, so `A is compact' reads is-compact(subspace-ms(s, A)) and `f(A) is compact' reads is-compact(subspace-ms(t, image(f, A))); f is continuous on all of s (is-continuous(s, t, f)), as in the notes. The second name is the case A = the whole space.
%
\renewcommand{\thethm}{3.20}
\begin{cor}[{A continuous bijection from a compact space is a homeomorphism}]\label{thm:calculus-3.20}
%
\textit{If X is compact and f : X --$>$ Y is a continuous bijective map, then f is a homeomorphism.}
%
Not in the library.
%
\end{cor}
%
\begin{proof}
%
The proofs: the file that holds each, the lemmas it cites, its oracles and its bill.
%
\vnbprf{There is no library statement, and so no proof.}
%
\end{proof}
%
\noindent\textbf{Deviation from the notes.} Not in the library. Since 2026-09-22 calculus 3.19 is proven, and the statement needs no new notion (`the inverse bijection is continuous'); the missing step is that a compact subspace of a metric space is CLOSED, the converse of subspace-complete.
%
\renewcommand{\thethm}{3.21}
\begin{prop}[{A continuous function on a compact space is uniformly continuous}]\label{thm:continuous-uniformly-continuous-on-ccint}
%
\vnbstmt{\begin{array}{@{}l@{}} \forall f, a, b, s. \\ \quad (f \in (\mathbb{R} \to \mathbb{R})) \wedge \\ \quad (a \in \mathbb{R}) \wedge \\ \quad (b \in \mathbb{R}) \wedge \\ \quad (a \leq b) \wedge \\ \quad \operatorname{pos\text{-}rr}(s) \wedge \\ \quad \forall x \in \operatorname{ccint}(a, b). \\ \quad \quad \operatorname{is\text{-}continuous\text{-}at}(\operatorname{rr\text{-}ms}, \operatorname{rr\text{-}ms}, f, x) \Rightarrow \\ \quad \quad \exists d \in \mathbb{R}. \\ \quad \quad \quad (0 < d) \wedge \\ \quad \quad \quad \forall p, q \in \operatorname{ccint}(a, b). \\ \quad \quad \quad \quad (\lvert (p - q) \rvert \leq d) \Rightarrow  \\ \quad \quad \quad \quad \quad (\lvert (f(p) - f(q)) \rvert \leq s) \end{array}}
%
\end{prop}
%
\begin{proof}
%
The proofs: the file that holds each, the lemmas it cites, its oracles and its bill.
%
\vnbprf{\textbf{Proof.} \texttt{theorem-\allowbreak{}library/uniform-\allowbreak{}continuity-\allowbreak{}ccint.scm} --- 421 recorded steps.}
%
\vnbprf{\textbf{Lemmas.} \texttt{subset-def}, \texttt{rr-leq-reflexive}, \texttt{fun-apply-type-c}, \texttt{ccint-membership}, \texttt{rr-leq-antisymmetric}, \texttt{rr-sub-in-rr}, \texttt{rr-abs-bound}, \texttt{rr-pos-halvable}, \texttt{rr-le-ne-lt}, \texttt{rr-min-pos}, \texttt{rr-leq-total}, \texttt{rr-ms-dist}, \texttt{ccint-creep}.}
%
\vnbprf{\textbf{Oracles.} \texttt{arith}, \texttt{ineq}, \texttt{crs}.  \textbf{Bill:} \texttt{proven modulo 0}.}
%
\end{proof}
%
\noindent\textbf{Deviation from the notes.} The general metric-space statement is absent. What exists is the single special case X = a closed bounded interval of RR, Y = RR, and it deviates in three ways: f is required to be in fun(rr, rr), i.e. TOTAL on RR, with continuity assumed at every x in ccint(a, b) and hence TWO-SIDED at the endpoints a and b; the conclusion is spelled out epsilon-delta inline (for a given eps there is d\_ $>$ 0 with abs(f(p\_) - f(q\_)) $<$= eps for p\_, q\_ in ccint(a, b) at distance $<$= d\_) and does NOT use the library's own is-uniformly-continuous predicate; and the epsilon is a PARAMETER of the theorem rather than universally quantified inside it. The guard is a $<$= b, so the degenerate interval a = b is allowed.
%
\renewcommand{\thethm}{3.22}
\begin{prop}[{Continuity plus uniform continuity off a compact set gives uniform continuity}]\label{thm:calculus-3.22}
%
\textit{Suppose X, Y are metric spaces and f : X --$>$ Y is continuous.  If K subset X is compact and f is uniformly continuous on X \textbackslash{} K, then f is uniformly continuous.}
%
Not in the library.
%
\end{prop}
%
\begin{proof}
%
The proofs: the file that holds each, the lemmas it cites, its oracles and its bill.
%
\vnbprf{There is no library statement, and so no proof.}
%
\end{proof}
%
\noindent\textbf{Deviation from the notes.} Nothing near it. The library has is-uniformly-continuous and is-compact but no theorem combining them beyond the two trunc-metric identity maps (trunc-metric-id-uniformly-continuous and its converse).
%
\renewcommand{\thethm}{3.23}
\begin{cor}[{A continuous function constant outside a compact set is uniformly continuous}]\label{thm:calculus-3.23}
%
\textit{If f : X --$>$ Y is continuous and constant outside a compact set K, then f is uniformly continuous.}
%
Not in the library.
%
\end{cor}
%
\begin{proof}
%
The proofs: the file that holds each, the lemmas it cites, its oracles and its bill.
%
\vnbprf{There is no library statement, and so no proof.}
%
\end{proof}
%
\noindent\textbf{Deviation from the notes.} Absent, as is Proposition 3.22 from which the notes derive it.
%
\renewcommand{\thethm}{3.26}
\begin{prop}[{Unique uniformly continuous extension from a dense subspace}]\label{thm:calculus-3.26}
%
\textit{Suppose X, Y are metric spaces, A subset X is a dense subspace and Y is complete.  Then every uniformly continuous function f : A --$>$ Y has a unique extension to a uniformly continuous function fbar : X --$>$ Y.}
%
Not in the library.
%
\end{prop}
%
\begin{proof}
%
The proofs: the file that holds each, the lemmas it cites, its oracles and its bill.
%
\vnbprf{There is no library statement, and so no proof.}
%
\end{proof}
%
\noindent\textbf{Deviation from the notes.} No extension theorem for uniformly continuous maps, and no EXISTENCE or UNIQUENESS statement for an extension along a dense subspace. The library does have the metric completion (completion-is-metric-space, completion-is-complete, embed-isometry, embed-in-fun) and a density notion for sequences (is-dense-seq), but no theorem connects them to an extension; the only extends-on facts in the whole library are the Hahn-Banach ones (hahn-banach, hahn-banach-extend-one, hb-extend-construct).
%
\renewcommand{\thethm}{3.27}
\begin{prop}[{min(d,1) is a metric, with the same convergent sequences}]\label{thm:trunc-metric-is-metric-space}
%
\leavevmode
\vnbpart{\texttt{trunc-metric-is-metric-space}}
\vnbstmt{\begin{array}{@{}l@{}} \forall s. \\ \quad \operatorname{is\text{-}metric\text{-}space}(s) \Rightarrow  \\ \quad \quad \operatorname{is\text{-}metric\text{-}space}(\operatorname{trunc\text{-}metric}(s)) \end{array}}
\vnbpart{\texttt{trunc-metric-dist}}
\vnbstmt{\begin{array}{@{}l@{}} \forall s, u, v. \\ \quad ((u \in \operatorname{pts}(s)) \wedge (v \in \operatorname{pts}(s))) \Rightarrow  \\ \quad \quad (\operatorname{dist}(\operatorname{trunc\text{-}metric}(s))(u, v) \simeq \operatorname{min}(1, \operatorname{dist}(s)(u, v))) \end{array}}
\vnbpart{\texttt{trunc-metric-id-uniformly-continuous}}
\vnbstmt{\begin{array}{@{}l@{}} \forall s. \\ \quad \operatorname{is\text{-}metric\text{-}space}(s) \Rightarrow  \\ \quad \quad \operatorname{is\text{-}uniformly\text{-}continuous}(s, \operatorname{trunc\text{-}metric}(s), (\lambda a \in \operatorname{pts}(s).\; a)) \end{array}}
\vnbpart{\texttt{trunc-metric-id-uniformly-continuous-back}}
\vnbstmt{\begin{array}{@{}l@{}} \forall s. \\ \quad \operatorname{is\text{-}metric\text{-}space}(s) \Rightarrow  \\ \quad \quad \operatorname{is\text{-}uniformly\text{-}continuous}(\operatorname{trunc\text{-}metric}(s), s, (\lambda a \in \operatorname{pts}(s).\; a)) \end{array}}
\vnbpart{\texttt{bdd-metric-converges-iff}}
\vnbstmt{\begin{array}{@{}l@{}} \forall s.; \forall f \in (\mathbb{N} \to \operatorname{pts}(s)), \ell \in \operatorname{pts}(s). \\ \quad \operatorname{is\text{-}metric\text{-}space}(s) \Rightarrow  \\ \quad \quad (\operatorname{converges\text{-}to}(\operatorname{bdd\text{-}metric}(s), f, \ell) \Leftrightarrow \operatorname{converges\text{-}to}(s, f, \ell)) \end{array}}
%
\end{prop}
%
\begin{proof}
%
The proofs: the file that holds each, the lemmas it cites, its oracles and its bill.
%
\vnbprf{\vnbproofpart{\texttt{trunc-metric-is-metric-space}}}
%
\vnbprf{\textbf{Proof.} \texttt{theorem-\allowbreak{}library/trunc-\allowbreak{}metric-\allowbreak{}proof.scm} --- 154 recorded steps.}
%
\vnbprf{\textbf{Lemmas.} \texttt{rr-one-in}, \texttt{metric-dist-real}, \texttt{rr-min-closed}, \texttt{cartesian-set-iff}, \texttt{rr-zero-in}, \texttt{rr-zero-lt-one}, \texttt{rr-min-le-left}, \texttt{rr-min-le-right}, \texttt{metric-pos}, \texttt{rr-le-min}, \texttt{metric-self-zero}, \texttt{rr-min-cases}, \texttt{metric-zero-eq}, \texttt{metric-sym}, \texttt{metric-triangle}, \texttt{rr-add-closed}, \texttt{rr-min-one-mono}, \texttt{rr-min-one-subadditive}.}
%
\vnbprf{\textbf{Oracles.} \texttt{arith}, \texttt{ineq}.  \textbf{Bill:} \texttt{proven modulo 0}.}
%
\vnbprf{\vnbproofpart{\texttt{trunc-metric-dist}}}
%
\vnbprf{\textbf{Proof.} \texttt{theorem-\allowbreak{}library/trunc-\allowbreak{}metric-\allowbreak{}proof.scm} --- 6 recorded steps.}
%
\vnbprf{\textbf{Lemmas.} none cited by name.}
%
\vnbprf{\textbf{Oracles.} none.  \textbf{Bill:} \texttt{proven modulo 0}.}
%
\vnbprf{\vnbproofpart{\texttt{trunc-metric-id-uniformly-continuous}}}
%
\vnbprf{\textbf{Proof.} \texttt{theorem-\allowbreak{}library/trunc-\allowbreak{}metric-\allowbreak{}proof.scm} --- 46 recorded steps.}
%
\vnbprf{\textbf{Lemmas.} \texttt{trunc-metric-is-metric-space}, \texttt{rr-one-in}, \texttt{metric-dist-real}, \texttt{rr-min-closed}, \texttt{rr-min-le-right}.}
%
\vnbprf{\textbf{Oracles.} \texttt{ineq}, \texttt{arith}.  \textbf{Bill:} \texttt{proven modulo 0}.}
%
\vnbprf{\vnbproofpart{\texttt{trunc-metric-id-uniformly-continuous-back}}}
%
\vnbprf{\textbf{Proof.} \texttt{theorem-\allowbreak{}library/trunc-\allowbreak{}metric-\allowbreak{}proof.scm} --- 109 recorded steps.}
%
\vnbprf{\textbf{Lemmas.} \texttt{trunc-metric-is-metric-space}, \texttt{rr-zero-in}, \texttt{rr-one-in}, \texttt{rr-pos-halvable}, \texttt{rr-min-pos}, \texttt{metric-dist-real}, \texttt{rr-min-closed}, \texttt{rr-min-cases}.}
%
\vnbprf{\textbf{Oracles.} \texttt{arith}, \texttt{ineq}, \texttt{crs}.  \textbf{Bill:} \texttt{proven modulo 0}.}
%
\vnbprf{\vnbproofpart{\texttt{bdd-metric-converges-iff}}}
%
\vnbprf{\textbf{Proof.} \texttt{theorem-\allowbreak{}library/bdd-\allowbreak{}metric-\allowbreak{}convergence.scm} --- 6 recorded steps.}
%
\vnbprf{\textbf{Lemmas.} \texttt{bdd-metric-converges-fwd}, \texttt{bdd-metric-converges-bwd}.}
%
\vnbprf{\textbf{Oracles.} \texttt{ineq}, \texttt{arith}, \texttt{crs}.  \textbf{Bill:} \texttt{proven modulo 0}.}
%
\end{proof}
%
\noindent\textbf{Deviation from the notes.} The two halves are proven for DIFFERENT truncations. trunc-metric(s) is literally the notes' min(1, d) (trunc-metric-dist) and is proven to be a metric space, but for it the library offers only the uniform continuity of the identity in both directions, never the convergence equivalence. The convergence equivalence that IS proven, bdd-metric-converges-iff, is for bdd-metric(s) = d/(1 + d), which is not the notes' d'. So no single library object carries both conclusions of Proposition 3.27.
%
\renewcommand{\thethm}{3.28}
\begin{prop}[{Tychonov Theorem for Metric Spaces}]\label{thm:compact-countable-product}
%
\leavevmode
\vnbpart{\texttt{compact-countable-product}}
\vnbstmt{\begin{array}{@{}l@{}} \forall s. \\ \quad \operatorname{is\text{-}ms\text{-}sequence}(s) \wedge \\ \quad \forall n \in \mathbb{N}.\; \operatorname{is\text{-}compact}(s(n)) \Rightarrow \\ \quad \quad \operatorname{is\text{-}compact}(\operatorname{product\text{-}metric}(s)) \end{array}}
\vnbpart{\texttt{seq-compact-countable-product}}
\vnbstmt{\begin{array}{@{}l@{}} \forall s. \\ \quad \operatorname{is\text{-}ms\text{-}sequence}(s) \wedge \\ \quad \forall n \in \mathbb{N}.\; \operatorname{seq\text{-}compact}(s(n)) \Rightarrow \\ \quad \quad \operatorname{seq\text{-}compact}(\operatorname{product\text{-}metric}(s)) \end{array}}
%
\end{prop}
%
\begin{proof}
%
The proofs: the file that holds each, the lemmas it cites, its oracles and its bill.
%
\vnbprf{\vnbproofpart{\texttt{compact-countable-product}}}
%
\vnbprf{\textbf{Proof.} \texttt{theorem-\allowbreak{}library/tychonoff-\allowbreak{}proof.scm} --- 22 recorded steps.}
%
\vnbprf{\textbf{Lemmas.} \texttt{compact-iff-seq-compact-rev}, \texttt{product-metric-default-summable}, \texttt{product-is-metric-space}, \texttt{seq-compact-countable-product}, \texttt{compact-iff-seq-compact}.}
%
\vnbprf{\textbf{Oracles.} \texttt{ineq}, \texttt{arith}, \texttt{crs}.  \textbf{Bill:} \texttt{proven modulo 0}.}
%
\vnbprf{\vnbproofpart{\texttt{seq-compact-countable-product}}}
%
\vnbprf{\textbf{Proof.} \texttt{theorem-\allowbreak{}library/rake-\allowbreak{}diagonal-\allowbreak{}subseq.scm} --- 50 recorded steps.}
%
\vnbprf{\textbf{Lemmas.} \texttt{product-metric-default-summable}, \texttt{product-metric-carrier}, \texttt{product-is-metric-space}, \texttt{coordinatewise-diagonal-subseq}, \texttt{subseq-is-fun}, \texttt{product-convergence-coordinatewise}.}
%
\vnbprf{\textbf{Oracles.} \texttt{arith}, \texttt{ineq}, \texttt{crs}.  \textbf{Bill:} \texttt{proven modulo 0}.}
%
\end{proof}
%
\noindent\textbf{Deviation from the notes.} The product metric is FIXED by the library rather than chosen in the proof: product-metric(ms) is product-metric-w with the weight n |--$>$ 1/2\textasciicircum{}(n+1) and, in each factor, the truncation bdd-metric = d/(1+d), NOT the notes' min(d, 1) with weights 2\textasciicircum{}-k. The countable family is presented as is-ms-sequence(ms), a function on nn, rather than as a countable family \{X\_k\}. Compactness of every factor is the hypothesis forall n in nn, is-compact(ms(n)).
%
\renewcommand{\thethm}{3.33}
\begin{prop}[{On a compact space pointwise and uniform convergence agree on an equicontinuous set}]\label{thm:ptwise-cauchy-compact-equicont-unif}
%
\leavevmode
\vnbpart{\texttt{ptwise-cauchy-compact-equicont-unif}}
\vnbstmt{\begin{array}{@{}l@{}} \forall s, m. \\ \quad \operatorname{is\text{-}compact}(s) \wedge \\ \quad (m \in (\mathbb{N} \to (\operatorname{pts}(s) \to \mathbb{R}))) \wedge \\ \quad \operatorname{is\text{-}equicontinuous}(s, \operatorname{rr\text{-}ms}, m) \wedge \\ \quad \operatorname{is\text{-}ptwise\text{-}cauchy}(s, m) \Rightarrow \\ \quad \quad \operatorname{is\text{-}unif\text{-}cauchy}(s, m) \end{array}}
\vnbpart{\texttt{unif-cauchy-cont-implies-uniform-limit}}
\vnbstmt{\begin{array}{@{}l@{}} \forall s.; \forall m \in (\mathbb{N} \to (\operatorname{pts}(s) \to \mathbb{R})). \\ \quad \forall k \in \mathbb{N}. \\ \quad \quad \operatorname{is\text{-}continuous}(s, \operatorname{rr\text{-}ms}, m(k)) \wedge \\ \quad \operatorname{is\text{-}unif\text{-}cauchy}(s, m) \Rightarrow \\ \quad \quad \exists g \in (\operatorname{pts}(s) \to \mathbb{R}). \\ \quad \quad \quad \operatorname{is\text{-}continuous}(s, \operatorname{rr\text{-}ms}, g) \wedge \\ \quad \quad \quad \operatorname{converges\text{-}uniformly}(s, m, g) \end{array}}
\vnbpart{\texttt{uniform-limit-continuous}}
\vnbstmt{\begin{array}{@{}l@{}} \forall s, m, g. \\ \quad (m \in (\mathbb{N} \to (\operatorname{pts}(s) \to \mathbb{R}))) \wedge \\ \quad (g \in (\operatorname{pts}(s) \to \mathbb{R})) \wedge \\ \quad \forall k \in \mathbb{N}. \\ \quad \quad \operatorname{is\text{-}continuous}(s, \operatorname{rr\text{-}ms}, m(k)) \wedge \\ \quad \operatorname{converges\text{-}uniformly}(s, m, g) \Rightarrow \\ \quad \quad \operatorname{is\text{-}continuous}(s, \operatorname{rr\text{-}ms}, g) \end{array}}
%
\end{prop}
%
\begin{proof}
%
The proofs: the file that holds each, the lemmas it cites, its oracles and its bill.
%
\vnbprf{\vnbproofpart{\texttt{ptwise-cauchy-compact-equicont-unif}}}
%
\vnbprf{\textbf{Proof.} \texttt{theorem-\allowbreak{}library/ptwise-\allowbreak{}cauchy-\allowbreak{}unif.scm} --- 333 recorded steps.}
%
\vnbprf{\textbf{Lemmas.} \texttt{rr-pos-halvable}, \texttt{rr-lt-of-pos-rr}, \texttt{rr-pos-rr-in-rr}, \texttt{class-extensionality}, \texttt{power-set-membership}, \texttt{metric-self-zero}, \texttt{ball-sep-unfold}, \texttt{ball-is-open}, \texttt{fun-apply-type-c}, \texttt{rr-ms-dist}, \texttt{rr-sub-in-rr}, \texttt{rr-abs-triangle-c}, \texttt{rr-abs-sub-sym}, \texttt{power-set}, \texttt{subclass-of-set-is-set}, \texttt{card-image-finite}, \texttt{subset-def}, \texttt{subset-mem-fwd}, \texttt{choice-axiom}, \texttt{nn-finite-subset-bounded}, \texttt{open-cover-covers-point}, \texttt{membership-implies-sethood}, \texttt{nn-in-rr}, \texttt{rr-leq-total}, \texttt{nn-le-trans-guarded}.}
%
\vnbprf{\textbf{Oracles.} \texttt{crs}, \texttt{ineq}, \texttt{arith}.  \textbf{Bill:} \texttt{proven modulo 0}.}
%
\vnbprf{\vnbproofpart{\texttt{unif-cauchy-cont-implies-uniform-limit}}}
%
\vnbprf{\textbf{Proof.} \texttt{theorem-\allowbreak{}library/ascoli-\allowbreak{}bridge.scm} --- 9 recorded steps.}
%
\vnbprf{\textbf{Lemmas.} \texttt{unif-cauchy-has-uniform-limit}, \texttt{uniform-limit-continuous}.}
%
\vnbprf{\textbf{Oracles.} \texttt{ineq}, \texttt{crs}, \texttt{arith}.  \textbf{Bill:} \texttt{proven modulo 0}.}
%
\vnbprf{\vnbproofpart{\texttt{uniform-limit-continuous}}}
%
\vnbprf{\textbf{Proof.} \texttt{theorem-\allowbreak{}library/ascoli-\allowbreak{}analytic-\allowbreak{}cores.scm} --- 166 recorded steps.}
%
\vnbprf{\textbf{Lemmas.} \texttt{rr-is-metric-space}, \texttt{rr-pos-halvable}, \texttt{fun-apply-type-c}, \texttt{nn-le-refl}, \texttt{rr-ms-dist}, \texttt{rr-sub-in-rr}, \texttt{rr-abs-triangle-c}, \texttt{rr-abs-sub-sym}.}
%
\vnbprf{\textbf{Oracles.} \texttt{crs}, \texttt{ineq}, \texttt{arith}.  \textbf{Bill:} \texttt{proven modulo 0}.}
%
\end{proof}
%
\noindent\textbf{Deviation from the notes.} The equivalence itself is not stated. The library proves only the CAUCHY form -- an equicontinuous family on a compact space that is pointwise Cauchy is uniformly Cauchy -- plus, separately, that a uniformly Cauchy family of continuous functions has a continuous uniform limit. Two further restrictions: the target is fixed to rr-ms, the family being typed fam in fun(nn, fun(pts(s), rr)), where the notes allow an arbitrary metric space Y; and the object is a SEQUENCE indexed by nn, where the notes speak of an arbitrary SET E of functions and its topology. The easy converse (uniform implies pointwise) is not stated either.
%
\renewcommand{\thethm}{3.34}
\begin{prop}[{Ascoli-Arzela}]\label{thm:ascoli-arzela-sequential}
%
\leavevmode
\vnbpart{\texttt{ascoli-arzela-sequential}}
\vnbstmt{\begin{array}{@{}l@{}} \forall s, m. \\ \quad \operatorname{is\text{-}compact}(s) \wedge \\ \quad \exists x.\; (x \in \operatorname{pts}(s)) \wedge \\ \quad (m \in (\mathbb{N} \to (\operatorname{pts}(s) \to \mathbb{R}))) \wedge \\ \quad \forall k \in \mathbb{N}. \\ \quad \quad \operatorname{is\text{-}continuous}(s, \operatorname{rr\text{-}ms}, m(k)) \wedge \\ \quad \operatorname{is\text{-}equicontinuous}(s, \operatorname{rr\text{-}ms}, m) \wedge \\ \quad \operatorname{pointwise\text{-}bounded}(s, m) \Rightarrow \\ \quad \quad \exists l. \\ \quad \quad \quad \operatorname{strictly\text{-}mono\text{-}nn}(l) \wedge \\ \quad \quad \quad \exists g \in (\operatorname{pts}(s) \to \mathbb{R}). \\ \quad \quad \quad \quad \operatorname{is\text{-}continuous}(s, \operatorname{rr\text{-}ms}, g) \wedge \\ \quad \quad \quad \quad \operatorname{converges\text{-}uniformly}(s, \operatorname{subseq}(m, l), g) \end{array}}
\vnbpart{\texttt{ascoli-sequential-from-diagonal}}
\vnbstmt{\begin{array}{@{}l@{}} \forall s, m. \\ \quad \operatorname{is\text{-}compact}(s) \wedge \\ \quad \exists x.\; (x \in \operatorname{pts}(s)) \wedge \\ \quad (m \in (\mathbb{N} \to (\operatorname{pts}(s) \to \mathbb{R}))) \wedge \\ \quad \forall k \in \mathbb{N}. \\ \quad \quad \operatorname{is\text{-}continuous}(s, \operatorname{rr\text{-}ms}, m(k)) \wedge \\ \quad \operatorname{is\text{-}equicontinuous}(s, \operatorname{rr\text{-}ms}, m) \wedge \\ \quad \operatorname{pointwise\text{-}bounded}(s, m) \wedge \\ \quad \forall q. \\ \quad \quad \operatorname{is\text{-}dense\text{-}seq}(s, q) \Rightarrow  \\ \quad \quad \quad \exists l. \\ \quad \quad \quad \quad \operatorname{strictly\text{-}mono\text{-}nn}(l) \wedge \\ \quad \quad \quad \quad \operatorname{converges\text{-}on}(s, \operatorname{subseq}(m, l), q) \Rightarrow \\ \quad \quad \exists l. \\ \quad \quad \quad \operatorname{strictly\text{-}mono\text{-}nn}(l) \wedge \\ \quad \quad \quad \exists g \in (\operatorname{pts}(s) \to \mathbb{R}). \\ \quad \quad \quad \quad \operatorname{is\text{-}continuous}(s, \operatorname{rr\text{-}ms}, g) \wedge \\ \quad \quad \quad \quad \operatorname{converges\text{-}uniformly}(s, \operatorname{subseq}(m, l), g) \end{array}}
\vnbpart{\texttt{ascoli-pointwise-diagonal}}
\vnbstmt{\begin{array}{@{}l@{}} \forall s, m, q. \\ \quad \operatorname{pointwise\text{-}bounded}(s, m) \wedge \\ \quad \operatorname{is\text{-}dense\text{-}seq}(s, q) \Rightarrow \\ \quad \quad \exists l. \\ \quad \quad \quad \operatorname{strictly\text{-}mono\text{-}nn}(l) \wedge \\ \quad \quad \quad \operatorname{converges\text{-}on}(s, \operatorname{subseq}(m, l), q) \end{array}}
%
\end{prop}
%
\begin{proof}
%
The proofs: the file that holds each, the lemmas it cites, its oracles and its bill.
%
\vnbprf{\vnbproofpart{\texttt{ascoli-arzela-sequential}}}
%
\vnbprf{\textbf{Proof.} \texttt{theorem-\allowbreak{}library/rake-\allowbreak{}ascoli-\allowbreak{}diagonal.scm} --- 10 recorded steps.}
%
\vnbprf{\textbf{Lemmas.} \texttt{ascoli-pointwise-diagonal}, \texttt{ascoli-sequential-from-diagonal}.}
%
\vnbprf{\textbf{Oracles.} \texttt{arith}, \texttt{ineq}, \texttt{crs}.  \textbf{Bill:} \texttt{proven modulo 0}.}
%
\vnbprf{\vnbproofpart{\texttt{ascoli-sequential-from-diagonal}}}
%
\vnbprf{\textbf{Proof.} \texttt{theorem-\allowbreak{}library/rake-\allowbreak{}ascoli.scm} --- 91 recorded steps.}
%
\vnbprf{\textbf{Lemmas.} \texttt{compact-metric-is-separable}, \texttt{nn-is-set}, \texttt{fun-apply-type-c}, \texttt{subseq-apply}, \texttt{equicont-subseq}, \texttt{ascoli-dense-bridge}.}
%
\vnbprf{\textbf{Oracles.} \texttt{ineq}, \texttt{arith}, \texttt{crs}.  \textbf{Bill:} \texttt{proven modulo 0}.}
%
\vnbprf{\vnbproofpart{\texttt{ascoli-pointwise-diagonal}}}
%
\vnbprf{\textbf{Proof.} \texttt{theorem-\allowbreak{}library/rake-\allowbreak{}ascoli-\allowbreak{}diagonal.scm} --- 224 recorded steps.}
%
\vnbprf{\textbf{Lemmas.} \texttt{rr-is-metric-space}, \texttt{fun-apply-type-c}, \texttt{nn-is-set}, \texttt{bounded-block-converges}, \texttt{inf-subsets-is-set}, \texttt{nn-in-inf-subsets}, \texttt{dc-on-nn-pred}, \texttt{nn-succ-closed}, \texttt{diagonalization}, \texttt{coord-block-estimate}, \texttt{subseq-is-fun}, \texttt{subseq-apply}, \texttt{converges-to-transfer}.}
%
\vnbprf{\textbf{Oracles.} \texttt{arith}, \texttt{ineq}, \texttt{crs}.  \textbf{Bill:} \texttt{proven modulo 0}.}
%
\end{proof}
%
\noindent\textbf{Deviation from the notes.} A SEQUENTIAL form, not the notes' relative compactness in the topology of uniform convergence: the library states that an equicontinuous, pointwise-bounded sequence of continuous functions on a compact X has a subsequence converging UNIFORMLY to a continuous limit. Four differences. (a) Y is fixed to rr-ms; the notes allow any compact metric Y. (b) pointwise-bounded(s, fam) is an added hypothesis, free in the notes because Y is compact. (c) X must be INHABITED -- forsome x, x in pts(s) -- an added guard against the empty metric space. (d) Only SUFFICIENCY is proven; the notes' necessity half, that a relatively compact set is equicontinuous, is absent, and with no relative-compactness notion in the library it cannot be stated.
%
\renewcommand{\thethm}{3.36}
\begin{lem}[{Equicontinuous sequence converging on a dense set converges everywhere, with continuous limit}]\label{thm:ascoli-dense-bridge}
%
\leavevmode
\vnbpart{\texttt{ascoli-dense-bridge}}
\vnbstmt{\begin{array}{@{}l@{}} \forall s, m. \\ \quad \operatorname{is\text{-}compact}(s) \wedge \\ \quad (m \in (\mathbb{N} \to (\operatorname{pts}(s) \to \mathbb{R}))) \wedge \\ \quad \forall k \in \mathbb{N}. \\ \quad \quad \operatorname{is\text{-}continuous}(s, \operatorname{rr\text{-}ms}, m(k)) \wedge \\ \quad \operatorname{is\text{-}equicontinuous}(s, \operatorname{rr\text{-}ms}, m) \wedge \\ \quad \exists q. \\ \quad \quad (\operatorname{is\text{-}dense\text{-}seq}(s, q) \wedge \operatorname{converges\text{-}on}(s, m, q)) \Rightarrow \\ \quad \quad \exists g \in (\operatorname{pts}(s) \to \mathbb{R}). \\ \quad \quad \quad \operatorname{is\text{-}continuous}(s, \operatorname{rr\text{-}ms}, g) \wedge \\ \quad \quad \quad \operatorname{converges\text{-}uniformly}(s, m, g) \end{array}}
\vnbpart{\texttt{equicont-dense-conv-implies-unif-cauchy}}
\vnbstmt{\begin{array}{@{}l@{}} \forall s, m. \\ \quad \operatorname{is\text{-}compact}(s) \wedge \\ \quad (m \in (\mathbb{N} \to (\operatorname{pts}(s) \to \mathbb{R}))) \wedge \\ \quad \operatorname{is\text{-}equicontinuous}(s, \operatorname{rr\text{-}ms}, m) \wedge \\ \quad \exists q. \\ \quad \quad (\operatorname{is\text{-}dense\text{-}seq}(s, q) \wedge \operatorname{converges\text{-}on}(s, m, q)) \Rightarrow \\ \quad \quad \operatorname{is\text{-}unif\text{-}cauchy}(s, m) \end{array}}
\vnbpart{\texttt{equicont-dense-conv-ptwise-cauchy}}
\vnbstmt{\begin{array}{@{}l@{}} \forall s.; \forall m \in (\mathbb{N} \to (\operatorname{pts}(s) \to \mathbb{R})). \\ \quad \operatorname{is\text{-}equicontinuous}(s, \operatorname{rr\text{-}ms}, m) \wedge \\ \quad \exists q. \\ \quad \quad (\operatorname{is\text{-}dense\text{-}seq}(s, q) \wedge \operatorname{converges\text{-}on}(s, m, q)) \Rightarrow \\ \quad \quad \operatorname{is\text{-}ptwise\text{-}cauchy}(s, m) \end{array}}
%
\end{lem}
%
\begin{proof}
%
The proofs: the file that holds each, the lemmas it cites, its oracles and its bill.
%
\vnbprf{\vnbproofpart{\texttt{ascoli-dense-bridge}}}
%
\vnbprf{\textbf{Proof.} \texttt{theorem-\allowbreak{}library/ascoli-\allowbreak{}bridge.scm} --- 9 recorded steps.}
%
\vnbprf{\textbf{Lemmas.} \texttt{equicont-dense-conv-implies-unif-cauchy}, \texttt{unif-cauchy-cont-implies-uniform-limit}.}
%
\vnbprf{\textbf{Oracles.} \texttt{crs}, \texttt{ineq}, \texttt{arith}.  \textbf{Bill:} \texttt{proven modulo 0}.}
%
\vnbprf{\vnbproofpart{\texttt{equicont-dense-conv-implies-unif-cauchy}}}
%
\vnbprf{\textbf{Proof.} \texttt{theorem-\allowbreak{}library/ascoli-\allowbreak{}assembly.scm} --- 8 recorded steps.}
%
\vnbprf{\textbf{Lemmas.} \texttt{equicont-dense-conv-ptwise-cauchy}, \texttt{ptwise-cauchy-compact-equicont-unif}.}
%
\vnbprf{\textbf{Oracles.} \texttt{crs}, \texttt{ineq}, \texttt{arith}.  \textbf{Bill:} \texttt{proven modulo 0}.}
%
\vnbprf{\vnbproofpart{\texttt{equicont-dense-conv-ptwise-cauchy}}}
%
\vnbprf{\textbf{Proof.} \texttt{theorem-\allowbreak{}library/ascoli-\allowbreak{}analytic-\allowbreak{}cores.scm} --- 145 recorded steps.}
%
\vnbprf{\textbf{Lemmas.} \texttt{rr-pos-halvable}, \texttt{fun-apply-type-c}, \texttt{rr-ms-dist}, \texttt{rr-sub-in-rr}, \texttt{rr-abs-triangle-c}, \texttt{rr-abs-sub-sym}.}
%
\vnbprf{\textbf{Oracles.} \texttt{crs}, \texttt{ineq}, \texttt{arith}.  \textbf{Bill:} \texttt{proven modulo 0}.}
%
\end{proof}
%
\noindent\textbf{Deviation from the notes.} Y is fixed to rr-ms, so the notes' hypothesis that Y be a complete metric space is discharged by the completeness of RR rather than assumed; a general complete Y is not available. The dense set D is presented as a dense SEQUENCE is-dense-seq(s, dseq) rather than as a dense subset, and convergence on it as converges-on(s, fam, dseq). The library's conclusion is UNIFORM convergence to a continuous g, which is stronger than the notes' pointwise conclusion plus continuity of the limit. fam must be typed fam in fun(nn, fun(pts(s), rr)) and each fam(k) assumed is-continuous, where the notes assume only equicontinuity of the sequence.
%
\renewcommand{\thethm}{3.37}
\begin{prop}[{The isometry group of a compact metric space is compact}]\label{thm:calculus-3.37}
%
\textit{If X is a compact metric space and Gamma is the group of isometric isomorphisms phi : X --$>$ X with d\_Gamma(phi, psi) = sup\_\{x in X\} d\_X(phi(x), psi(x)), then Gamma with d\_Gamma is a compact metric space.}
%
Not in the library.
%
\end{prop}
%
\begin{proof}
%
The proofs: the file that holds each, the lemmas it cites, its oracles and its bill.
%
\vnbprf{There is no library statement, and so no proof.}
%
\end{proof}
%
\noindent\textbf{Deviation from the notes.} is-isometry(s, t, f) is defined, but there is no group of isometric isomorphisms of a space, no sup metric d\_Gamma(phi, psi) = sup\_x d(phi(x), psi(x)) on such a set, and no compactness statement about either.
%
\renewcommand{\thethm}{3.38}
\begin{prop}[{In the isometry group uniform and pointwise convergence agree}]\label{thm:calculus-3.38}
%
\textit{Let Gamma be as in 3.37 and \{g\_i : i in N\} a sequence of elements of Gamma.  Then g\_i --$>$ g uniformly if and only if g\_i(x) --$>$ g(x) for all x in X.}
%
Not in the library.
%
\end{prop}
%
\begin{proof}
%
The proofs: the file that holds each, the lemmas it cites, its oracles and its bill.
%
\vnbprf{There is no library statement, and so no proof.}
%
\end{proof}
%
\noindent\textbf{Deviation from the notes.} Absent, together with the group Gamma of Proposition 3.37 on which it is stated.
%
\renewcommand{\thethm}{4.5}
\begin{prop}[{Regulated functions are exactly the uniform limits of step functions}]\label{thm:regulated-on-is-step-limit}
%
\leavevmode
\vnbpart{\texttt{regulated-on-is-step-limit}}
\vnbstmt{\begin{array}{@{}l@{}} \forall a, b \in \mathbb{R}, f. \\ \quad \operatorname{is\text{-}regulated\text{-}on}(f, a, b) \Rightarrow  \\ \quad \quad \exists m \in (\mathbb{N} \to (\mathbb{R} \to \mathbb{R})). \\ \quad \quad \quad \forall k \in \mathbb{N}.\; \operatorname{is\text{-}step\text{-}fn}(m(k), a, b) \wedge \\ \quad \quad \quad \operatorname{is\text{-}unif\text{-}limit\text{-}on}(m, \operatorname{extend\text{-}const}(f, a, b), a, b) \end{array}}
\vnbpart{\texttt{regulated-on-step-approx}}
\vnbstmt{\begin{array}{@{}l@{}} \forall a, b, f, e. \\ \quad (\operatorname{is\text{-}regulated\text{-}on}(f, a, b) \wedge \operatorname{pos\text{-}rr}(e)) \Rightarrow  \\ \quad \quad \exists h \in (\mathbb{R} \to \mathbb{R}). \\ \quad \quad \quad \operatorname{is\text{-}step\text{-}fn}(h, a, b) \wedge \\ \quad \quad \quad \forall z \in \operatorname{ccint}(a, b). \\ \quad \quad \quad \quad (\lvert (\operatorname{extend\text{-}const}(f, a, b)(z) - h(z)) \rvert \leq e) \end{array}}
\vnbpart{\texttt{uniform-limit-of-regulated}}
\vnbstmt{\begin{array}{@{}l@{}} \forall m, g, a, b. \\ \quad \operatorname{is\text{-}unif\text{-}limit\text{-}on}(m, g, a, b) \wedge \\ \quad \forall k \in \mathbb{N}.\; \operatorname{is\text{-}regulated}(m(k), a, b) \Rightarrow \\ \quad \quad \operatorname{is\text{-}regulated}(g, a, b) \end{array}}
\vnbpart{\texttt{continuous-is-regulated}}
\vnbstmt{\begin{array}{@{}l@{}} \forall f, a, b. \\ \quad (f \in (\mathbb{R} \to \mathbb{R})) \wedge \\ \quad (a \in \mathbb{R}) \wedge \\ \quad (b \in \mathbb{R}) \wedge \\ \quad (a < b) \wedge \\ \quad \forall x \in \operatorname{ccint}(a, b). \\ \quad \quad \operatorname{is\text{-}continuous\text{-}at}(\operatorname{rr\text{-}ms}, \operatorname{rr\text{-}ms}, f, x) \Rightarrow \\ \quad \quad \operatorname{is\text{-}regulated}(f, a, b) \end{array}}
%
\end{prop}
%
\begin{proof}
%
The proofs: the file that holds each, the lemmas it cites, its oracles and its bill.
%
\vnbprf{\vnbproofpart{\texttt{regulated-on-is-step-limit}}}
%
\vnbprf{\textbf{Proof.} \texttt{theorem-\allowbreak{}library/regulated-\allowbreak{}step-\allowbreak{}approx.scm} --- 135 recorded steps.}
%
\vnbprf{\textbf{Lemmas.} \texttt{rr-lt-implies-le}, \texttt{extend-const-in-fun}, \texttt{regulated-on-step-approx}, \texttt{nn-recip-succ-pos}, \texttt{choice-axiom}, \texttt{nn-is-set}, \texttt{nn-recip-succ-small}, \texttt{rr-pos-rr-in-rr}, \texttt{nn-recip-succ-antitone}, \texttt{ccint-elt-in-rr}, \texttt{fun-apply-type-c}, \texttt{rr-sub-in-rr}, \texttt{rr-abs-closed}.}
%
\vnbprf{\textbf{Oracles.} \texttt{ineq}, \texttt{arith}, \texttt{crs}.  \textbf{Bill:} \texttt{proven modulo 0}.}
%
\vnbprf{\vnbproofpart{\texttt{regulated-on-step-approx}}}
%
\vnbprf{\textbf{Proof.} \texttt{theorem-\allowbreak{}library/regulated-\allowbreak{}step-\allowbreak{}approx.scm} --- 58 recorded steps.}
%
\vnbprf{\textbf{Lemmas.} \texttt{rr-lt-implies-le}, \texttt{extend-const-in-fun}, \texttt{ccint-elt-in-rr}, \texttt{ccint-membership}, \texttt{extend-const-fixes}, \texttt{fun-apply-type-c}, \texttt{regulated-on-step-approx-creep}, \texttt{step-approx-q-step-fn}.}
%
\vnbprf{\textbf{Oracles.} \texttt{ineq}, \texttt{arith}.  \textbf{Bill:} \texttt{proven modulo 0}.}
%
\vnbprf{\vnbproofpart{\texttt{uniform-limit-of-regulated}}}
%
\vnbprf{\textbf{Proof.} \texttt{theorem-\allowbreak{}library/uniform-\allowbreak{}limit-\allowbreak{}regulated.scm} --- 645 recorded steps.}
%
\vnbprf{\textbf{Lemmas.} \texttt{ccint-elt-in-rr}, \texttt{cauchy-criterion-right}, \texttt{rr-pos-halvable}, \texttt{nn-le-refl}, \texttt{rr-sub-in-rr}, \texttt{rr-lt-diff-pos}, \texttt{rr-min-pos}, \texttt{rr-pos-rr-of-lt}, \texttt{ccint-membership}, \texttt{fun-apply-type-c}, \texttt{rr-abs-bound}, \texttt{cauchy-criterion-left}.}
%
\vnbprf{\textbf{Oracles.} \texttt{ineq}, \texttt{crs}, \texttt{arith}.  \textbf{Bill:} \texttt{proven modulo 0}.}
%
\vnbprf{\vnbproofpart{\texttt{continuous-is-regulated}}}
%
\vnbprf{\textbf{Proof.} \texttt{theorem-\allowbreak{}library/continuous-\allowbreak{}is-\allowbreak{}regulated.scm} --- 41 recorded steps.}
%
\vnbprf{\textbf{Lemmas.} \texttt{ccint-membership}, \texttt{fun-apply-type-c}, \texttt{continuous-at-right-limit}, \texttt{continuous-at-left-limit}.}
%
\vnbprf{\textbf{Oracles.} \texttt{ineq}, \texttt{arith}, \texttt{crs}.  \textbf{Bill:} \texttt{proven modulo 0}.}
%
\end{proof}
%
\noindent\textbf{Deviation from the notes.} UPDATED 2026-09-23 (batch 22-A): NECESSITY is proven ON [a, b] -- regulated-on-is-step-limit: a function regulated on [a, b] (is-regulated-on, structure-library/regulated-primitive.scm, for f in fun(ccint(a, b), rr)) is the uniform limit on [a, b] of a sequence of step functions (is-step-fn, total), the limit stated through extend-const(f, a, b) because is-unif-limit-on is total; by the creeping argument, no finite subcover. SUFFICIENCY remains only in the total form uniform-limit-of-regulated. EARLIER NOTE: NEITHER direction of the notes' equivalence is stated.  is-step-fn is defined in the library but NO theorem anywhere mentions it, so necessity (a regulated function is a uniform limit of step functions) is absent, and sufficiency is present only in the weakened form uniform-limit-of-regulated: a uniform limit of REGULATED functions is regulated. Separately, is-regulated(f, a, b) types f in fun(rr, rr), i.e. TOTAL on RR, with a and b as parameters, where the notes have f : [a,b] --$>$ R; it demands a $<$ b strictly; and it demands a right limit only where x\_ $<$ b and a left limit only where a $<$ x\_, where the notes (Remark 4.2) instead declare f(a-) = f(a) and f(b+) = f(b) so that all four one-sided limits exist.
%
\renewcommand{\thethm}{4.8}
\begin{prop}[{The antiderivable functions on [a,b] form a vector space}]\label{thm:antiderivable-add}
%
\leavevmode
\vnbpart{\texttt{antiderivable-add}}
\vnbstmt{\begin{array}{@{}l@{}} \forall i, c, a, b. \\ \quad \operatorname{is\text{-}antiderivable}(i, a, b) \wedge \\ \quad \operatorname{is\text{-}antiderivable}(c, a, b) \Rightarrow \\ \quad \quad \operatorname{is\text{-}antiderivable}((\lambda x \in \mathbb{R}.\; (i(x) + c(x))), a, b) \end{array}}
\vnbpart{\texttt{antiderivable-scale}}
\vnbstmt{\begin{array}{@{}l@{}} \forall c, i, a, b. \\ \quad ((c \in \mathbb{R}) \wedge \operatorname{is\text{-}antiderivable}(i, a, b)) \Rightarrow  \\ \quad \quad \operatorname{is\text{-}antiderivable}((\lambda x \in \mathbb{R}.\; (c \cdot i(x))), a, b) \end{array}}
\vnbpart{\texttt{antiderivable-sub}}
\vnbstmt{\begin{array}{@{}l@{}} \forall i, c, a, b. \\ \quad \operatorname{is\text{-}antiderivable}(i, a, b) \wedge \\ \quad \operatorname{is\text{-}antiderivable}(c, a, b) \Rightarrow \\ \quad \quad \operatorname{is\text{-}antiderivable}((\lambda x \in \mathbb{R}.\; (i(x) - c(x))), a, b) \end{array}}
%
\end{prop}
%
\begin{proof}
%
The proofs: the file that holds each, the lemmas it cites, its oracles and its bill.
%
\vnbprf{\vnbproofpart{\texttt{antiderivable-add}}}
%
\vnbprf{\textbf{Proof.} \texttt{theorem-\allowbreak{}library/antiderivative.scm} --- 11 recorded steps.}
%
\vnbprf{\textbf{Lemmas.} \texttt{antiderivative-add}.}
%
\vnbprf{\textbf{Oracles.} \texttt{ineq}, \texttt{crs}, \texttt{arith}.  \textbf{Bill:} \texttt{proven modulo 0}.}
%
\vnbprf{\vnbproofpart{\texttt{antiderivable-scale}}}
%
\vnbprf{\textbf{Proof.} \texttt{theorem-\allowbreak{}library/antiderivative.scm} --- 9 recorded steps.}
%
\vnbprf{\textbf{Lemmas.} \texttt{antiderivative-scale}.}
%
\vnbprf{\textbf{Oracles.} \texttt{crs}, \texttt{arith}, \texttt{ineq}.  \textbf{Bill:} \texttt{proven modulo 0}.}
%
\vnbprf{\vnbproofpart{\texttt{antiderivable-sub}}}
%
\vnbprf{\textbf{Proof.} \texttt{theorem-\allowbreak{}library/antiderivative-\allowbreak{}transfer.scm} --- 11 recorded steps.}
%
\vnbprf{\textbf{Lemmas.} \texttt{antiderivative-sub}.}
%
\vnbprf{\textbf{Oracles.} \texttt{crs}, \texttt{ineq}, \texttt{arith}.  \textbf{Bill:} \texttt{proven modulo 0}.}
%
\end{proof}
%
\noindent\textbf{Deviation from the notes.} Stated as three closure properties -- sum, scalar multiple and difference -- and not as a vector space: the library has no vector-space structure on a class of functions to name, so the conclusion of the notes has no counterpart. The functions are in fun(rr, rr), TOTAL on RR, with a and b as parameters, and is-antiderivable(phi, a, b) demands a $<$ b strictly, so the degenerate interval is excluded. The sum is written as the explicit term vnb-lambda(x, rr, phi(x) + psi(x)).
%
\renewcommand{\thethm}{4.10}
\begin{prop}[{Two antiderivatives of the same function differ by a constant}]\label{thm:pw-antiderivative-differ-by-some-constant}
%
\leavevmode
\vnbpart{\texttt{pw-antiderivative-differ-by-some-constant}}
\vnbstmt{\begin{array}{@{}l@{}} \forall f, w, i, a, b. \\ \quad \operatorname{is\text{-}primitive}(f, i, a, b) \wedge \\ \quad \operatorname{is\text{-}primitive}(w, i, a, b) \Rightarrow \\ \quad \quad \exists c \in \mathbb{R}. \forall s \in \operatorname{ccint}(a, b). \\ \quad \quad \quad (f(s) = (w(s) + c)) \end{array}}
\vnbpart{\texttt{pw-antiderivative-differ-by-constant}}
\vnbstmt{\begin{array}{@{}l@{}} \forall f, w, i, a, b.; \forall s \in \operatorname{ccint}(a, b). \\ \quad \operatorname{is\text{-}primitive}(f, i, a, b) \wedge \\ \quad \operatorname{is\text{-}primitive}(w, i, a, b) \Rightarrow \\ \quad \quad (f(s) = (w(s) + (f(a) - w(a)))) \end{array}}
\vnbpart{\texttt{antiderivative-differ-by-constant}}
\vnbstmt{\begin{array}{@{}l@{}} \forall f, g, i, a, b. \\ \quad \operatorname{is\text{-}antiderivative}(f, i, a, b) \wedge \\ \quad \operatorname{is\text{-}antiderivative}(g, i, a, b) \Rightarrow \\ \quad \quad \exists t \in \mathbb{R}. \forall s \in \operatorname{ccint}(a, b). \\ \quad \quad \quad (f(s) = (g(s) + t)) \end{array}}
%
\end{prop}
%
\begin{proof}
%
The proofs: the file that holds each, the lemmas it cites, its oracles and its bill.
%
\vnbprf{\vnbproofpart{\texttt{pw-antiderivative-differ-by-some-constant}}}
%
\vnbprf{\textbf{Proof.} \texttt{theorem-\allowbreak{}library/pw-\allowbreak{}int-\allowbreak{}laws-\allowbreak{}2.scm} --- 52 recorded steps.}
%
\vnbprf{\textbf{Lemmas.} \texttt{primitive-endpoints}, \texttt{rr-leq-reflexive}, \texttt{ccint-membership}, \texttt{primitive-in-fun}, \texttt{fun-apply-type-c}, \texttt{rr-sub-in-rr}, \texttt{pw-antiderivative-differ-by-constant}.}
%
\vnbprf{\textbf{Oracles.} \texttt{ineq}, \texttt{crs}, \texttt{arith}.  \textbf{Bill:} \texttt{proven modulo 0}.}
%
\vnbprf{\vnbproofpart{\texttt{pw-antiderivative-differ-by-constant}}}
%
\vnbprf{\textbf{Proof.} \texttt{theorem-\allowbreak{}library/pw-\allowbreak{}int-\allowbreak{}laws-\allowbreak{}2.scm} --- 212 recorded steps.}
%
\vnbprf{\textbf{Lemmas.} \texttt{primitive-endpoints}, \texttt{rr-leq-reflexive}, \texttt{ccint-membership}, \texttt{primitive-in-fun}, \texttt{primitive-integrand-in-fun}, \texttt{primitive-continuous}, \texttt{extend-const-in-fun}, \texttt{fun-apply-type-c}, \texttt{ooint-subset-ccint}, \texttt{rr-sub-in-rr}, \texttt{rr-is-set}, \texttt{extend-const-continuous-at}, \texttt{sub-continuous-at}, \texttt{primitive-exceptional-set}, \texttt{countable-union-2}, \texttt{subset-def}, \texttt{union-membership}, \texttt{subset-mem-fwd}, \texttt{ccint-elt-in-rr}, \texttt{extend-const-deriv-fwd}, \texttt{deriv-difference}, \texttt{ccint-parts}, \texttt{rr-le-cases}, \texttt{ooint-membership}, \texttt{zero-deriv-off-countable}, \texttt{extend-const-fixes}.}
%
\vnbprf{\textbf{Oracles.} \texttt{ineq}, \texttt{crs}, \texttt{arith}.  \textbf{Bill:} \texttt{proven modulo 0}.}
%
\vnbprf{\vnbproofpart{\texttt{antiderivative-differ-by-constant}}}
%
\vnbprf{\textbf{Proof.} \texttt{theorem-\allowbreak{}library/antiderivative.scm} --- 120 recorded steps.}
%
\vnbprf{\textbf{Lemmas.} \texttt{sub-lam-in-fun}, \texttt{sub-continuous-at}, \texttt{fun-apply-type-c}, \texttt{deriv-sub}, \texttt{deriv-zero-implies-constant}, \texttt{rr-lt-implies-le}, \texttt{rr-leq-reflexive}, \texttt{ccint-membership}, \texttt{rr-sub-in-rr}, \texttt{eq-sym}.}
%
\vnbprf{\textbf{Oracles.} \texttt{crs}, \texttt{ineq}, \texttt{arith}.  \textbf{Bill:} \texttt{proven modulo 0}.}
%
\end{proof}
%
\noindent\textbf{Deviation from the notes.} Restated 2026-09-22 for functions ON [a, b] (the first name(s)), for the PIECEWISE primitive is-pw-antiderivative: continuous on [a, b], with derivative phi(t) at every t of (a, b) outside a finite set (Dieudonne 8.7 with a finite exceptional set; the notes' Definition 4.6 is the case of an empty exceptional set, which the library does not single out yet). The remaining names are the older forms for functions TOTAL on RR (is-antiderivative, c-int); they remain in the library.
%
\renewcommand{\thethm}{4.11}
\begin{cor}[{The endpoint difference is the same for all antiderivatives}]\label{thm:pw-antiderivative-endpoint-difference}
%
\leavevmode
\vnbpart{\texttt{pw-antiderivative-endpoint-difference}}
\vnbstmt{\begin{array}{@{}l@{}} \forall f, w, i, a, b. \\ \quad \operatorname{is\text{-}primitive}(f, i, a, b) \wedge \\ \quad \operatorname{is\text{-}primitive}(w, i, a, b) \Rightarrow \\ \quad \quad ((f(b) - f(a)) = (w(b) - w(a))) \end{array}}
\vnbpart{\texttt{antiderivative-endpoint-difference}}
\vnbstmt{\begin{array}{@{}l@{}} \forall f, g, i, a, b. \\ \quad \operatorname{is\text{-}antiderivative}(f, i, a, b) \wedge \\ \quad \operatorname{is\text{-}antiderivative}(g, i, a, b) \Rightarrow \\ \quad \quad ((f(b) - f(a)) = (g(b) - g(a))) \end{array}}
%
\end{cor}
%
\begin{proof}
%
The proofs: the file that holds each, the lemmas it cites, its oracles and its bill.
%
\vnbprf{\vnbproofpart{\texttt{pw-antiderivative-endpoint-difference}}}
%
\vnbprf{\textbf{Proof.} \texttt{theorem-\allowbreak{}library/pw-\allowbreak{}antiderivative-\allowbreak{}laws.scm} --- 157 recorded steps.}
%
\vnbprf{\textbf{Lemmas.} \texttt{primitive-endpoints}, \texttt{rr-leq-reflexive}, \texttt{primitive-in-fun}, \texttt{primitive-integrand-in-fun}, \texttt{primitive-continuous}, \texttt{extend-const-in-fun}, \texttt{extend-const-fixes}, \texttt{ccint-membership}, \texttt{fun-apply-type-c}, \texttt{rr-sub-in-rr}, \texttt{rr-is-set}, \texttt{extend-const-continuous-at}, \texttt{sub-continuous-at}, \texttt{primitive-exceptional-set}, \texttt{countable-union-2}, \texttt{subset-def}, \texttt{union-membership}, \texttt{subset-mem-fwd}, \texttt{ccint-subset-rr}, \texttt{subset-trans}, \texttt{ooint-subset-ccint}, \texttt{extend-const-deriv-fwd}, \texttt{deriv-difference}, \texttt{zero-deriv-off-countable}.}
%
\vnbprf{\textbf{Oracles.} \texttt{ineq}, \texttt{crs}, \texttt{arith}.  \textbf{Bill:} \texttt{proven modulo 0}.}
%
\vnbprf{\vnbproofpart{\texttt{antiderivative-endpoint-difference}}}
%
\vnbprf{\textbf{Proof.} \texttt{theorem-\allowbreak{}library/antiderivative.scm} --- 33 recorded steps.}
%
\vnbprf{\textbf{Lemmas.} \texttt{antiderivative-endpoints}, \texttt{antiderivative-map-in-fun}, \texttt{rr-lt-implies-le}, \texttt{rr-leq-reflexive}, \texttt{ccint-membership}, \texttt{fun-apply-type-c}, \texttt{antiderivative-differ-by-constant}.}
%
\vnbprf{\textbf{Oracles.} \texttt{crs}, \texttt{ineq}, \texttt{arith}.  \textbf{Bill:} \texttt{proven modulo 0}.}
%
\end{proof}
%
\noindent\textbf{Deviation from the notes.} Restated 2026-09-22 for functions ON [a, b] (the first name(s)), for the PIECEWISE primitive is-pw-antiderivative: continuous on [a, b], with derivative phi(t) at every t of (a, b) outside a finite set (Dieudonne 8.7 with a finite exceptional set; the notes' Definition 4.6 is the case of an empty exceptional set, which the library does not single out yet). The remaining names are the older forms for functions TOTAL on RR (is-antiderivative, c-int); they remain in the library.
%
\renewcommand{\thethm}{4.12}
\begin{prop}[{The definite integral is a linear functional}]\label{thm:pw-antiderivative-linear}
%
\leavevmode
\vnbpart{\texttt{pw-antiderivative-linear}}
\vnbstmt{\begin{array}{@{}l@{}} \forall a, b, c, d \in \mathbb{R}, f, w, i, \mathit{ppsi}.; \forall h, \mathit{pchi} \in (\operatorname{ccint}(a, b) \to \mathbb{R}). \\ \quad \operatorname{is\text{-}primitive}(f, i, a, b) \wedge \\ \quad \operatorname{is\text{-}primitive}(w, \mathit{ppsi}, a, b) \wedge \\ \quad \forall y \in \operatorname{ccint}(a, b). \\ \quad \quad (h(y) \simeq ((c \cdot f(y)) + (d \cdot w(y)))) \wedge \\ \quad \forall y \in \operatorname{ccint}(a, b). \\ \quad \quad (\operatorname{pchi}(y) \simeq ((c \cdot i(y)) + (d \cdot \operatorname{ppsi}(y)))) \Rightarrow \\ \quad \quad \operatorname{is\text{-}primitive}(h, \mathit{pchi}, a, b) \wedge \\ \quad \quad (\int_{a}^{b} \mathit{pchi} = ((c \cdot \int_{a}^{b} i) + (d \cdot \int_{a}^{b} \mathit{ppsi}))) \end{array}}
\vnbpart{\texttt{pw-antiderivative-sum}}
\vnbstmt{\begin{array}{@{}l@{}} \forall a, b \in \mathbb{R}, f, w, i, \mathit{ppsi}.; \forall h, \mathit{pchi} \in (\operatorname{ccint}(a, b) \to \mathbb{R}). \\ \quad \operatorname{is\text{-}primitive}(f, i, a, b) \wedge \\ \quad \operatorname{is\text{-}primitive}(w, \mathit{ppsi}, a, b) \wedge \\ \quad \forall y \in \operatorname{ccint}(a, b).\; (h(y) \simeq (f(y) + w(y))) \wedge \\ \quad \forall y \in \operatorname{ccint}(a, b). \\ \quad \quad (\operatorname{pchi}(y) \simeq (i(y) + \operatorname{ppsi}(y))) \Rightarrow \\ \quad \quad \operatorname{is\text{-}primitive}(h, \mathit{pchi}, a, b) \wedge \\ \quad \quad (\int_{a}^{b} \mathit{pchi} = (\int_{a}^{b} i + \int_{a}^{b} \mathit{ppsi})) \end{array}}
\vnbpart{\texttt{pw-antiderivative-real-mul}}
\vnbstmt{\begin{array}{@{}l@{}} \forall a, b \in \mathbb{R}, f, i.; \forall c \in \mathbb{R}, h, \mathit{pchi} \in (\operatorname{ccint}(a, b) \to \mathbb{R}). \\ \quad \operatorname{is\text{-}primitive}(f, i, a, b) \wedge \\ \quad \forall y \in \operatorname{ccint}(a, b).\; (h(y) \simeq (c \cdot f(y))) \wedge \\ \quad \forall y \in \operatorname{ccint}(a, b). \\ \quad \quad (\operatorname{pchi}(y) \simeq (c \cdot i(y))) \Rightarrow \\ \quad \quad \operatorname{is\text{-}primitive}(h, \mathit{pchi}, a, b) \wedge \\ \quad \quad (\int_{a}^{b} \mathit{pchi} = (c \cdot \int_{a}^{b} i)) \end{array}}
\vnbpart{\texttt{c-int-add}}
\vnbstmt{\begin{array}{@{}l@{}} \forall i, c, a, b. \\ \quad \operatorname{is\text{-}antiderivable}(i, a, b) \wedge \\ \quad \operatorname{is\text{-}antiderivable}(c, a, b) \Rightarrow \\ \quad \quad (\int_{a}^{b} (\lambda x \in \mathbb{R}.\; (i(x) + c(x))) = (\int_{a}^{b} i + \int_{a}^{b} c)) \end{array}}
\vnbpart{\texttt{c-int-scale}}
\vnbstmt{\begin{array}{@{}l@{}} \forall c, i, a, b. \\ \quad ((c \in \mathbb{R}) \wedge \operatorname{is\text{-}antiderivable}(i, a, b)) \Rightarrow  \\ \quad \quad (\int_{a}^{b} (\lambda x \in \mathbb{R}.\; (c \cdot i(x))) = (c \cdot \int_{a}^{b} i)) \end{array}}
%
\end{prop}
%
\begin{proof}
%
The proofs: the file that holds each, the lemmas it cites, its oracles and its bill.
%
\vnbprf{\vnbproofpart{\texttt{pw-antiderivative-linear}}}
%
\vnbprf{\textbf{Proof.} \texttt{theorem-\allowbreak{}library/pw-\allowbreak{}int-\allowbreak{}laws-\allowbreak{}3.scm} --- 222 recorded steps.}
%
\vnbprf{\textbf{Lemmas.} \texttt{primitive-endpoints}, \texttt{rr-leq-reflexive}, \texttt{ccint-membership}, \texttt{rr-is-set}, \texttt{ccint-subset-rr}, \texttt{subclass-of-set-is-set}, \texttt{primitive-in-fun}, \texttt{primitive-integrand-in-fun}, \texttt{fun-apply-type-c}, \texttt{rr-mul-closed}, \texttt{ccint-elt-in-rr}, \texttt{pw-antiderivative-real-mul}, \texttt{pw-antiderivative-sum}, \texttt{pw-int-value}.}
%
\vnbprf{\textbf{Oracles.} \texttt{ineq}, \texttt{crs}, \texttt{arith}.  \textbf{Bill:} \texttt{proven modulo 0}.}
%
\vnbprf{\vnbproofpart{\texttt{pw-antiderivative-sum}}}
%
\vnbprf{\textbf{Proof.} \texttt{theorem-\allowbreak{}library/pw-\allowbreak{}antiderivative-\allowbreak{}laws.scm} --- 216 recorded steps.}
%
\vnbprf{\textbf{Lemmas.} \texttt{primitive-endpoints}, \texttt{rr-is-set}, \texttt{ccint-subset-rr}, \texttt{subclass-of-set-is-set}, \texttt{primitive-in-fun}, \texttt{primitive-integrand-in-fun}, \texttt{primitive-continuous}, \texttt{extend-const-in-fun}, \texttt{ooint-subset-ccint}, \texttt{fun-apply-type-c}, \texttt{rr-add-closed}, \texttt{extend-const-continuous-at}, \texttt{sum-continuous-at}, \texttt{ccint-elt-in-rr}, \texttt{ccint-membership}, \texttt{extend-const-fixes}, \texttt{pw-restrict-continuous-on}, \texttt{primitive-exceptional-set}, \texttt{countable-union-2}, \texttt{subset-def}, \texttt{union-membership}, \texttt{subset-mem-fwd}, \texttt{ooint-inner-radius}, \texttt{rr-pos-rr-in-rr}, \texttt{rr-sub-in-rr}, \texttt{rr-add-in-rr}, \texttt{restrict-in-fun}, \texttt{has-deriv-at-sum}, \texttt{pw-int-value}, \texttt{rr-leq-reflexive}.}
%
\vnbprf{\textbf{Oracles.} \texttt{ineq}, \texttt{crs}, \texttt{arith}.  \textbf{Bill:} \texttt{proven modulo 0}.}
%
\vnbprf{\vnbproofpart{\texttt{pw-antiderivative-real-mul}}}
%
\vnbprf{\textbf{Proof.} \texttt{theorem-\allowbreak{}library/pw-\allowbreak{}int-\allowbreak{}laws-\allowbreak{}2.scm} --- 175 recorded steps.}
%
\vnbprf{\textbf{Lemmas.} \texttt{primitive-endpoints}, \texttt{rr-leq-reflexive}, \texttt{ccint-membership}, \texttt{rr-is-set}, \texttt{ccint-subset-rr}, \texttt{subclass-of-set-is-set}, \texttt{primitive-in-fun}, \texttt{primitive-integrand-in-fun}, \texttt{primitive-continuous}, \texttt{extend-const-in-fun}, \texttt{ooint-subset-ccint}, \texttt{fun-apply-type-c}, \texttt{rr-mul-closed}, \texttt{extend-const-continuous-at}, \texttt{scale-continuous-at}, \texttt{ccint-elt-in-rr}, \texttt{ccint-parts}, \texttt{extend-const-fixes}, \texttt{pw-restrict-continuous-on}, \texttt{primitive-exceptional-set}, \texttt{subset-mem-fwd}, \texttt{ooint-inner-radius}, \texttt{rr-pos-rr-in-rr}, \texttt{rr-sub-in-rr}, \texttt{rr-add-in-rr}, \texttt{restrict-in-fun}, \texttt{has-deriv-at-real-mul}, \texttt{pw-int-value}.}
%
\vnbprf{\textbf{Oracles.} \texttt{ineq}, \texttt{crs}, \texttt{arith}.  \textbf{Bill:} \texttt{proven modulo 0}.}
%
\vnbprf{\vnbproofpart{\texttt{c-int-add}}}
%
\vnbprf{\textbf{Proof.} \texttt{theorem-\allowbreak{}library/c-\allowbreak{}int.scm} --- 27 recorded steps.}
%
\vnbprf{\textbf{Lemmas.} \texttt{antiderivative-endpoints}, \texttt{antiderivative-map-in-fun}, \texttt{fun-apply-type-c}, \texttt{rr-sub-in-rr}, \texttt{antiderivative-add}, \texttt{c-int-value}.}
%
\vnbprf{\textbf{Oracles.} \texttt{crs}, \texttt{ineq}, \texttt{arith}.  \textbf{Bill:} \texttt{proven modulo 0}.}
%
\vnbprf{\vnbproofpart{\texttt{c-int-scale}}}
%
\vnbprf{\textbf{Proof.} \texttt{theorem-\allowbreak{}library/c-\allowbreak{}int.scm} --- 19 recorded steps.}
%
\vnbprf{\textbf{Lemmas.} \texttt{antiderivative-endpoints}, \texttt{antiderivative-map-in-fun}, \texttt{fun-apply-type-c}, \texttt{rr-sub-in-rr}, \texttt{antiderivative-scale}, \texttt{c-int-value}.}
%
\vnbprf{\textbf{Oracles.} \texttt{crs}, \texttt{arith}, \texttt{ineq}.  \textbf{Bill:} \texttt{proven modulo 0}.}
%
\end{proof}
%
\noindent\textbf{Deviation from the notes.} UPDATED 2026-09-23: pw-antiderivative-linear (batch 18-A) states Prop 4.12 in ONE formula; since batch 20 (2026-09-23) the primitive is is-primitive with a COUNTABLE exceptional set (Dieudonne VIII.7), of which the notes' Definition 4.6 (no exceptional set) is a special case. EARLIER NOTE: Restated 2026-09-22 for functions ON [a, b] (the first name(s)), for the PIECEWISE primitive is-pw-antiderivative: continuous on [a, b], with derivative phi(t) at every t of (a, b) outside a finite set (Dieudonne 8.7 with a finite exceptional set; the notes' Definition 4.6 is the case of an empty exceptional set, which the library does not single out yet). `Linear functional' is given as two theorems, the sum and the real multiple: the space of antiderivable functions is not a structure in the library. The remaining names are the older forms for functions TOTAL on RR (is-antiderivative, c-int); they remain in the library.
%
\renewcommand{\thethm}{4.13}
\begin{prop}[{The integral of a sum is the sum of the integrals}]\label{thm:pw-antiderivative-sum}
%
\leavevmode
\vnbpart{\texttt{pw-antiderivative-sum}}
\vnbstmt{\begin{array}{@{}l@{}} \forall a, b \in \mathbb{R}, f, w, i, \mathit{ppsi}.; \forall h, \mathit{pchi} \in (\operatorname{ccint}(a, b) \to \mathbb{R}). \\ \quad \operatorname{is\text{-}primitive}(f, i, a, b) \wedge \\ \quad \operatorname{is\text{-}primitive}(w, \mathit{ppsi}, a, b) \wedge \\ \quad \forall y \in \operatorname{ccint}(a, b).\; (h(y) \simeq (f(y) + w(y))) \wedge \\ \quad \forall y \in \operatorname{ccint}(a, b). \\ \quad \quad (\operatorname{pchi}(y) \simeq (i(y) + \operatorname{ppsi}(y))) \Rightarrow \\ \quad \quad \operatorname{is\text{-}primitive}(h, \mathit{pchi}, a, b) \wedge \\ \quad \quad (\int_{a}^{b} \mathit{pchi} = (\int_{a}^{b} i + \int_{a}^{b} \mathit{ppsi})) \end{array}}
\vnbpart{\texttt{c-int-add}}
\vnbstmt{\begin{array}{@{}l@{}} \forall i, c, a, b. \\ \quad \operatorname{is\text{-}antiderivable}(i, a, b) \wedge \\ \quad \operatorname{is\text{-}antiderivable}(c, a, b) \Rightarrow \\ \quad \quad (\int_{a}^{b} (\lambda x \in \mathbb{R}.\; (i(x) + c(x))) = (\int_{a}^{b} i + \int_{a}^{b} c)) \end{array}}
\vnbpart{\texttt{antiderivative-add}}
\vnbstmt{\begin{array}{@{}l@{}} \forall f, g, i, c, a, b. \\ \quad \operatorname{is\text{-}antiderivative}(f, i, a, b) \wedge \\ \quad \operatorname{is\text{-}antiderivative}(g, c, a, b) \Rightarrow \\ \quad \quad \operatorname{is\text{-}antiderivative}((\lambda x \in \mathbb{R}.\; (f(x) + g(x))), (\lambda x \in \mathbb{R}.\; (i(x) + c(x))), \\ \quad \quad \quad a, b) \end{array}}
%
\end{prop}
%
\begin{proof}
%
The proofs: the file that holds each, the lemmas it cites, its oracles and its bill.
%
\vnbprf{\vnbproofpart{\texttt{pw-antiderivative-sum}}}
%
\vnbprf{\textbf{Proof.} \texttt{theorem-\allowbreak{}library/pw-\allowbreak{}antiderivative-\allowbreak{}laws.scm} --- 216 recorded steps.}
%
\vnbprf{\textbf{Lemmas.} \texttt{primitive-endpoints}, \texttt{rr-is-set}, \texttt{ccint-subset-rr}, \texttt{subclass-of-set-is-set}, \texttt{primitive-in-fun}, \texttt{primitive-integrand-in-fun}, \texttt{primitive-continuous}, \texttt{extend-const-in-fun}, \texttt{ooint-subset-ccint}, \texttt{fun-apply-type-c}, \texttt{rr-add-closed}, \texttt{extend-const-continuous-at}, \texttt{sum-continuous-at}, \texttt{ccint-elt-in-rr}, \texttt{ccint-membership}, \texttt{extend-const-fixes}, \texttt{pw-restrict-continuous-on}, \texttt{primitive-exceptional-set}, \texttt{countable-union-2}, \texttt{subset-def}, \texttt{union-membership}, \texttt{subset-mem-fwd}, \texttt{ooint-inner-radius}, \texttt{rr-pos-rr-in-rr}, \texttt{rr-sub-in-rr}, \texttt{rr-add-in-rr}, \texttt{restrict-in-fun}, \texttt{has-deriv-at-sum}, \texttt{pw-int-value}, \texttt{rr-leq-reflexive}.}
%
\vnbprf{\textbf{Oracles.} \texttt{ineq}, \texttt{crs}, \texttt{arith}.  \textbf{Bill:} \texttt{proven modulo 0}.}
%
\vnbprf{\vnbproofpart{\texttt{c-int-add}}}
%
\vnbprf{\textbf{Proof.} \texttt{theorem-\allowbreak{}library/c-\allowbreak{}int.scm} --- 27 recorded steps.}
%
\vnbprf{\textbf{Lemmas.} \texttt{antiderivative-endpoints}, \texttt{antiderivative-map-in-fun}, \texttt{fun-apply-type-c}, \texttt{rr-sub-in-rr}, \texttt{antiderivative-add}, \texttt{c-int-value}.}
%
\vnbprf{\textbf{Oracles.} \texttt{crs}, \texttt{ineq}, \texttt{arith}.  \textbf{Bill:} \texttt{proven modulo 0}.}
%
\vnbprf{\vnbproofpart{\texttt{antiderivative-add}}}
%
\vnbprf{\textbf{Proof.} \texttt{theorem-\allowbreak{}library/antiderivative.scm} --- 80 recorded steps.}
%
\vnbprf{\textbf{Lemmas.} \texttt{sum-lam-in-fun}, \texttt{sum-continuous-at}, \texttt{fun-apply-type-c}, \texttt{deriv-sum}.}
%
\vnbprf{\textbf{Oracles.} \texttt{ineq}, \texttt{crs}, \texttt{arith}.  \textbf{Bill:} \texttt{proven modulo 0}.}
%
\end{proof}
%
\noindent\textbf{Deviation from the notes.} Restated 2026-09-22 for functions ON [a, b] (the first name(s)), for the PIECEWISE primitive is-pw-antiderivative: continuous on [a, b], with derivative phi(t) at every t of (a, b) outside a finite set (Dieudonne 8.7 with a finite exceptional set; the notes' Definition 4.6 is the case of an empty exceptional set, which the library does not single out yet). The remaining names are the older forms for functions TOTAL on RR (is-antiderivative, c-int); they remain in the library.
%
\renewcommand{\thethm}{4.16}
\begin{prop}[{Differentiating a sequence whose derivatives converge uniformly}]\label{thm:primitive-uniform-limit}
%
\leavevmode
\vnbpart{\texttt{primitive-uniform-limit}}
\vnbstmt{\begin{array}{@{}l@{}} \forall a, b \in \mathbb{R}, m, c \in (\mathbb{N} \to (\operatorname{ccint}(a, b) \to \mathbb{R})), f \in (\operatorname{ccint}(a, b) \to \mathbb{R}). \\ \quad (a < b) \wedge \\ \quad \forall k \in \mathbb{N}.\; \operatorname{is\text{-}primitive}(m(k), c(k), a, b) \wedge \\ \quad \forall s. \\ \quad \quad \operatorname{pos\text{-}rr}(s) \Rightarrow  \\ \quad \quad \quad \exists n \in \mathbb{N}. \forall k \in \mathbb{N}, x \in \operatorname{ccint}(a, b). \\ \quad \quad \quad \quad ((n \leq k) \Rightarrow (\lvert (c(k)(x) - f(x)) \rvert \leq s)) \wedge \\ \quad \forall k \in \mathbb{N}.\; (m(k)(a) = 0) \Rightarrow \\ \quad \quad \exists g. \\ \quad \quad \quad \operatorname{is\text{-}primitive}(g, f, a, b) \wedge \\ \quad \quad \quad \forall x \in \operatorname{ccint}(a, b). \\ \quad \quad \quad \quad \operatorname{converges\text{-}to}(\operatorname{rr\text{-}ms}, (\lambda k \in \mathbb{N}.\; m(k)(x)), g(x)) \end{array}}
\vnbpart{\texttt{antiderivative-family-uniform-limit}}
\vnbstmt{\begin{array}{@{}l@{}} \forall m, c \in (\mathbb{N} \to (\mathbb{R} \to \mathbb{R})), i \in (\mathbb{R} \to \mathbb{R}), a, b \in \mathbb{R}. \\ \quad (a < b) \wedge \\ \quad \forall k \in \mathbb{N}. \\ \quad \quad (\operatorname{is\text{-}antiderivative}(m(k), c(k), a, b) \wedge (m(k)(a) = 0)) \wedge \\ \quad \forall s. \\ \quad \quad \operatorname{pos\text{-}rr}(s) \Rightarrow  \\ \quad \quad \quad \exists n \in \mathbb{N}. \forall k \in \mathbb{N}, x \in \mathbb{R}. \\ \quad \quad \quad \quad ((n \leq k) \wedge (x \in \operatorname{ccint}(a, b))) \Rightarrow  \\ \quad \quad \quad \quad \quad (\lvert (i(x) - c(k)(x)) \rvert \leq s) \Rightarrow \\ \quad \quad \exists f \in (\mathbb{R} \to \mathbb{R}). \\ \quad \quad \quad \operatorname{is\text{-}antiderivative}(f, i, a, b) \wedge \\ \quad \quad \quad \forall s. \\ \quad \quad \quad \quad \operatorname{pos\text{-}rr}(s) \Rightarrow  \\ \quad \quad \quad \quad \quad \exists n \in \mathbb{N}. \forall k \in \mathbb{N}, x \in \mathbb{R}. \\ \quad \quad \quad \quad \quad \quad ((n \leq k) \wedge (x \in \operatorname{ccint}(a, b))) \Rightarrow  \\ \quad \quad \quad \quad \quad \quad \quad (\lvert (f(x) - m(k)(x)) \rvert \leq s) \end{array}}
\vnbpart{\texttt{antiderivable-uniform-limit}}
\vnbstmt{\begin{array}{@{}l@{}} \forall m \in (\mathbb{N} \to (\mathbb{R} \to \mathbb{R})), i \in (\mathbb{R} \to \mathbb{R}), a, b \in \mathbb{R}. \\ \quad \forall k \in \mathbb{N}.\; \operatorname{is\text{-}antiderivable}(m(k), a, b) \wedge \\ \quad \forall s. \\ \quad \quad \operatorname{pos\text{-}rr}(s) \Rightarrow  \\ \quad \quad \quad \exists n \in \mathbb{N}. \forall k \in \mathbb{N}, x \in \mathbb{R}. \\ \quad \quad \quad \quad ((n \leq k) \wedge (x \in \operatorname{ccint}(a, b))) \Rightarrow  \\ \quad \quad \quad \quad \quad (\lvert (i(x) - m(k)(x)) \rvert \leq s) \Rightarrow \\ \quad \quad \operatorname{is\text{-}antiderivable}(i, a, b) \end{array}}
\vnbpart{\texttt{diff-at-ptwise-limit}}
\vnbstmt{\begin{array}{@{}l@{}} \forall m, a \in (\mathbb{N} \to (\mathbb{R} \to \mathbb{R})), f, i \in (\mathbb{R} \to \mathbb{R}), t, o \in \mathbb{R}. \\ \quad (0 < o) \wedge \\ \quad \operatorname{is\text{-}continuous\text{-}at}(\operatorname{rr\text{-}ms}, \operatorname{rr\text{-}ms}, i, t) \wedge \\ \quad \forall k \in \mathbb{N}, y \in \mathbb{R}. \\ \quad \quad (\lvert (y - t) \rvert \leq o) \Rightarrow  \\ \quad \quad \quad ((m(k)(y) - m(k)(t)) = (a(k)(y) \cdot (y - t))) \wedge \\ \quad \forall y \in \mathbb{R}. \\ \quad \quad (\lvert (y - t) \rvert \leq o) \Rightarrow  \\ \quad \quad \quad \operatorname{converges\text{-}to}(\operatorname{rr\text{-}ms}, (\lambda b \in \mathbb{N}.\; m(b)(y)), f(y)) \wedge \\ \quad \forall y \in \mathbb{R}. \\ \quad \quad (\lvert (y - t) \rvert \leq o) \Rightarrow  \\ \quad \quad \quad \operatorname{converges\text{-}to}(\operatorname{rr\text{-}ms}, (\lambda b \in \mathbb{N}.\; a(b)(y)), i(y)) \Rightarrow \\ \quad \quad \operatorname{is\text{-}diff\text{-}at}(f, t, i(t)) \end{array}}
%
\end{prop}
%
\begin{proof}
%
The proofs: the file that holds each, the lemmas it cites, its oracles and its bill.
%
\vnbprf{\vnbproofpart{\texttt{primitive-uniform-limit}}}
%
\vnbprf{\textbf{Proof.} \texttt{theorem-\allowbreak{}library/primitive-\allowbreak{}uniform-\allowbreak{}limit.scm} --- 387 recorded steps.}
%
\vnbprf{\textbf{Lemmas.} \texttt{ccint-subset-rr}, \texttt{primitive-family-limit}, \texttt{rr-pos-halvable}, \texttt{rr-lt-of-pos-rr}, \texttt{rr-pos-rr-in-rr}, \texttt{fun-apply-type-c}, \texttt{rr-sub-in-rr}, \texttt{rr-abs-le-parts}, \texttt{rr-abs-le-of-parts}, \texttt{restrict-in-fun}, \texttt{fun-domain-in-set}, \texttt{countable-union-nn}, \texttt{primitive-exceptional-set}, \texttt{subset-def}, \texttt{countable-subset}, \texttt{cont-on-of-total}, \texttt{subset-mem-fwd}, \texttt{restrict-apply}, \texttt{ooint-elt-in-rr}, \texttt{ooint-subset-ccint}, \texttt{nn-is-set}, \texttt{rr-converges-to-abs}, \texttt{primitive-family-deriv-at}, \texttt{has-deriv-at-iff-diff-at}, \texttt{ooint-inner-radius}, \texttt{has-deriv-at-local}.}
%
\vnbprf{\textbf{Oracles.} \texttt{ineq}, \texttt{crs}, \texttt{arith}.  \textbf{Bill:} \texttt{proven modulo 0}.}
%
\vnbprf{\vnbproofpart{\texttt{antiderivative-family-uniform-limit}}}
%
\vnbprf{\textbf{Proof.} \texttt{theorem-\allowbreak{}library/antiderivable-\allowbreak{}uniform-\allowbreak{}limit.scm} --- 595 recorded steps.}
%
\vnbprf{\textbf{Lemmas.} \texttt{rr-zero-in}, \texttt{rr-one-in}, \texttt{rr-zero-lt-one}, \texttt{rr-lt-implies-le}, \texttt{rr-add-closed}, \texttt{rr-sub-in-rr}, \texttt{uniform-cauchy-limit}, \texttt{rr-scale-eps}, \texttt{fun-apply-type-c}, \texttt{rr-abs-closed}, \texttt{rr-abs-triangle-c}, \texttt{rr-abs-sub-sym}, \texttt{rr-mul-closed}, \texttt{antiderivative-pair-close}, \texttt{rr-le-trans}, \texttt{nn-is-set}, \texttt{clamp-compose-lam-in-fun}, \texttt{clamp-in-ccint}, \texttt{clamp-in-rr}, \texttt{clamp-compose-continuous-at}, \texttt{uniform-limit-continuous-at}, \texttt{ccint-membership}, \texttt{rr-min-pos}, \texttt{rr-le-abs}, \texttt{rr-neg-abs-le}, \texttt{diff-at-witness-family}, \texttt{clamp-fixes}, \texttt{caratheodory-witness-pair-bound}, \texttt{rr-converges-to-abs}, \texttt{uniform-witness-diff-at}.}
%
\vnbprf{\textbf{Oracles.} \texttt{ineq}, \texttt{crs}, \texttt{arith}.  \textbf{Bill:} \texttt{proven modulo 0}.}
%
\vnbprf{\vnbproofpart{\texttt{antiderivable-uniform-limit}}}
%
\vnbprf{\textbf{Proof.} \texttt{theorem-\allowbreak{}library/antiderivable-\allowbreak{}uniform-\allowbreak{}limit.scm} --- 19 recorded steps.}
%
\vnbprf{\textbf{Lemmas.} \texttt{antiderivable-family-normalized}, \texttt{nn-zero-in}, \texttt{antiderivative-endpoints}, \texttt{antiderivative-family-uniform-limit}.}
%
\vnbprf{\textbf{Oracles.} \texttt{crs}, \texttt{arith}, \texttt{ineq}.  \textbf{Bill:} \texttt{proven modulo 0}.}
%
\vnbprf{\vnbproofpart{\texttt{diff-at-ptwise-limit}}}
%
\vnbprf{\textbf{Proof.} \texttt{theorem-\allowbreak{}library/diff-\allowbreak{}at-\allowbreak{}local.scm} --- 129 recorded steps.}
%
\vnbprf{\textbf{Lemmas.} \texttt{rr-zero-in}, \texttt{rr-lt-implies-le}, \texttt{rr-leq-reflexive}, \texttt{rr-abs-of-nonneg}, \texttt{fun-apply-type-c}, \texttt{rr-sub-in-rr}, \texttt{rr-mul-closed}, \texttt{nn-is-set}, \texttt{rr-limit-sub}, \texttt{rr-limit-scale}, \texttt{rr-limit-unique}, \texttt{diff-at-local}.}
%
\vnbprf{\textbf{Oracles.} \texttt{crs}, \texttt{ineq}, \texttt{arith}.  \textbf{Bill:} \texttt{proven modulo 0}.}
%
\end{proof}
%
\noindent\textbf{Deviation from the notes.} UPDATED 2026-09-23 (batch 22): primitive-uniform-limit is Dieudonne (8.6.4) on [a, b] with COUNTABLE exceptional sets -- g\_k a primitive of f\_k, the f\_k uniformly convergent to f on [a, b], g\_k(a) = 0 -- concluding a primitive of f to which the g\_k converge; it is still the NORMALISED form (g\_k(a) = 0 rather than the notes' convergent f\_k(a)), hence partial. EARLIER NOTE: The library version NORMALISES the sequence: it demands fam(k)(a) = 0 for EVERY k, where the notes assume only that \{f\_k(a)\} is convergent, so the notes' hypothesis (2) is replaced by a stronger and non-equivalent one. Hypothesis (1) is replaced too: instead of f\_k' being defined and CONTINUOUS on [a,b], the library assumes is-antiderivative(fam(k), phifam(k), a, b), which gives differentiability only on the OPEN interval (a,b) and says nothing about continuity of phifam(k). Of the four conclusions of the notes, it delivers the existence of f = lim f\_k (as a uniform limit on ccint(a, b)) and is-antiderivative(f, psi, a, b); it does NOT state that psi is continuous, does not state that f is differentiable at the ENDPOINTS (the notes explicitly extend the conclusion to them via one-sided derivatives), and does not state f\_k' --$>$ f' uniformly. All functions are in fun(rr, rr), TOTAL on RR, uniform convergence is written inline over ccint(a, b), and a $<$ b is strict.
%
\renewcommand{\thethm}{4.17}
\begin{cor}[{The antiderivable functions are closed under uniform limits}]\label{thm:antiderivable-uniform-limit}
%
\vnbstmt{\begin{array}{@{}l@{}} \forall m \in (\mathbb{N} \to (\mathbb{R} \to \mathbb{R})), i \in (\mathbb{R} \to \mathbb{R}), a, b \in \mathbb{R}. \\ \quad \forall k \in \mathbb{N}.\; \operatorname{is\text{-}antiderivable}(m(k), a, b) \wedge \\ \quad \forall s. \\ \quad \quad \operatorname{pos\text{-}rr}(s) \Rightarrow  \\ \quad \quad \quad \exists n \in \mathbb{N}. \forall k \in \mathbb{N}, x \in \mathbb{R}. \\ \quad \quad \quad \quad ((n \leq k) \wedge (x \in \operatorname{ccint}(a, b))) \Rightarrow  \\ \quad \quad \quad \quad \quad (\lvert (i(x) - m(k)(x)) \rvert \leq s) \Rightarrow \\ \quad \quad \operatorname{is\text{-}antiderivable}(i, a, b) \end{array}}
%
\end{cor}
%
\begin{proof}
%
The proofs: the file that holds each, the lemmas it cites, its oracles and its bill.
%
\vnbprf{\textbf{Proof.} \texttt{theorem-\allowbreak{}library/antiderivable-\allowbreak{}uniform-\allowbreak{}limit.scm} --- 19 recorded steps.}
%
\vnbprf{\textbf{Lemmas.} \texttt{antiderivable-family-normalized}, \texttt{nn-zero-in}, \texttt{antiderivative-endpoints}, \texttt{antiderivative-family-uniform-limit}.}
%
\vnbprf{\textbf{Oracles.} \texttt{crs}, \texttt{arith}, \texttt{ineq}.  \textbf{Bill:} \texttt{proven modulo 0}.}
%
\end{proof}
%
\noindent\textbf{Deviation from the notes.} Uniform convergence is spelled out inline (for every eps $>$ 0 there is n\_ in nn with abs(psi(x\_) - (phifam(k\_))(x\_)) $<$= eps for all k\_ $>$= n\_ and all x\_ in ccint(a, b)) rather than through a named predicate, and is required only on ccint(a, b). The sequence and its limit are in fun(rr, rr), TOTAL on RR, and a $<$ b is strict, so the degenerate interval is excluded.
%
\renewcommand{\thethm}{5.2}
\begin{thm}[{Bernstein approximation of a continuous function on [0,1]}]\label{thm:bernstein-uniform-approximation}
%
\vnbstmt{\begin{array}{@{}l@{}} \forall f, s. \\ \quad (f \in (\mathbb{R} \to \mathbb{R})) \wedge \\ \quad \forall x \in \operatorname{ccint}(0, 1). \\ \quad \quad \operatorname{is\text{-}continuous\text{-}at}(\operatorname{rr\text{-}ms}, \operatorname{rr\text{-}ms}, f, x) \wedge \\ \quad \operatorname{pos\text{-}rr}(s) \Rightarrow \\ \quad \quad \exists p \in \mathbb{N}. \forall n, a. \\ \quad \quad \quad ((n \in \mathbb{N}) \wedge ((p \leq n) \wedge (a \in \operatorname{ccint}(0, 1)))) \Rightarrow  \\ \quad \quad \quad \quad (\lvert (f(a) - \operatorname{bernstein\text{-}poly}(f, n, a)) \rvert \leq s) \end{array}}
%
\end{thm}
%
\begin{proof}
%
The proofs: the file that holds each, the lemmas it cites, its oracles and its bill.
%
\vnbprf{\textbf{Proof.} \texttt{theorem-\allowbreak{}library/bernstein-\allowbreak{}density.scm} --- 311 recorded steps.}
%
\vnbprf{\textbf{Lemmas.} \texttt{rr-pos-halvable}, \texttt{rr-le-ne-lt}, \texttt{neq-sym}, \texttt{rr-zero-in}, \texttt{rr-one-in}, \texttt{continuous-uniformly-continuous-on-ccint}, \texttt{continuous-abs-bounded-on-ccint}, \texttt{rr-pos-ne-zero}, \texttt{rr-recip-closed}, \texttt{rr-recip-pos}, \texttt{rr-recip-inverse}, \texttt{rr-lt-implies-le}, \texttt{rr-add-in-rr}, \texttt{rr-add-nonneg}, \texttt{rr-sq-nonneg}, \texttt{rr-mul-in-rr}, \texttt{rr-leq-mul-nonneg}, \texttt{nn-unbounded-in-rr}, \texttt{nn-in-rr}, \texttt{rr-le-lt-trans}, \texttt{rr-lt-le-trans}, \texttt{rr-le-scale-nonneg}, \texttt{rr-leq-reflexive}, \texttt{rr-le-trans}, \texttt{eq-sym}, \texttt{bernstein-approx-bound}.}
%
\vnbprf{\textbf{Oracles.} \texttt{arith}, \texttt{crs}, \texttt{ineq}, \texttt{sos}.  \textbf{Bill:} \texttt{proven modulo 0}.}
%
\end{proof}
%
\noindent\textbf{Deviation from the notes.} f is required to be in fun(rr, rr), i.e. TOTAL on RR, with continuity assumed only at the points of ccint(0, 1) and there TWO-SIDED, including at 0 and at 1; the notes have f continuous on [0,1]. The conclusion is given as a single eps/N estimate -- eps is a PARAMETER of the theorem and the claim is forsome cap\_ in nn, for all n\_ $>$= cap\_ and all x\_ in ccint(0,1), abs(f\_(x\_) - bernstein-poly(f\_, n\_, x\_)) $<$= eps\_ -- rather than as the statement that f is the uniform limit of the sequence B\_n. bernstein-poly(f, n, x) is the functoid series-partial-sum(l\_ |--$>$ f(l\_/n) * (bernstein-basis(n,x))(l\_), succ(n)), so the notes' sum from 0 to n appears as a partial sum of length n+1.
%
\renewcommand{\thethm}{5.4}
\begin{prop}[{The two Bernstein estimates}]\label{thm:bernstein-partition}
%
\leavevmode
\vnbpart{\texttt{bernstein-partition}}
\vnbstmt{\begin{array}{@{}l@{}} \forall n \in \mathbb{N}, x \in \mathbb{R}. \\ \quad \operatorname{series\text{-}partial\text{-}sum}(\operatorname{bernstein\text{-}basis}(n, x), (n + 1)) \\ \quad =\; 1 \end{array}}
\vnbpart{\texttt{bernstein-variance-bound}}
\vnbstmt{\begin{array}{@{}l@{}} \forall n, x, h. \\ \quad (n \in \mathbb{N}) \wedge \\ \quad \neg (n = 0) \wedge \\ \quad (x \in \mathbb{R}) \wedge \\ \quad \forall k \in \mathbb{N}. \\ \quad \quad h(k) \\ \quad \quad =\; (((k - (n \cdot x)) \cdot (k - (n \cdot x))) \cdot \operatorname{bernstein\text{-}basis}(n, x)(k)) \Rightarrow \\ \quad \quad ((/(1, n) \cdot \operatorname{series\text{-}partial\text{-}sum}(h, (n + 1))) \leq 1) \end{array}}
%
\end{prop}
%
\begin{proof}
%
The proofs: the file that holds each, the lemmas it cites, its oracles and its bill.
%
\vnbprf{\vnbproofpart{\texttt{bernstein-partition}}}
%
\vnbprf{\textbf{Proof.} \texttt{theorem-\allowbreak{}library/bernstein-\allowbreak{}basis.scm} --- 48 recorded steps.}
%
\vnbprf{\textbf{Lemmas.} \texttt{rr-is-normed-field}, \texttt{rr-one-in}, \texttt{rr-sub-in-rr}, \texttt{nn-succ-closed}, \texttt{rr-scalar-ring-carr}, \texttt{series-partial-sum-is-ring-sum}, \texttt{binomial-sum-value}, \texttt{rr-scalar-ring-add}, \texttt{rr-scalar-ring-one}, \texttt{lam-slot-add-apply}, \texttt{rr-add-closed}, \texttt{ring-power-of-one}.}
%
\vnbprf{\textbf{Oracles.} \texttt{crs}, \texttt{arith}, \texttt{ineq}.  \textbf{Bill:} \texttt{proven modulo 0}.}
%
\vnbprf{\vnbproofpart{\texttt{bernstein-variance-bound}}}
%
\vnbprf{\textbf{Proof.} \texttt{theorem-\allowbreak{}library/bernstein-\allowbreak{}moments.scm} --- 26 recorded steps.}
%
\vnbprf{\textbf{Lemmas.} \texttt{nn-in-rr}, \texttt{bernstein-variance}, \texttt{rr-recip-closed}, \texttt{rr-recip-inverse}.}
%
\vnbprf{\textbf{Oracles.} \texttt{crs}, \texttt{sos}, \texttt{arith}, \texttt{ineq}.  \textbf{Bill:} \texttt{proven modulo 0}.}
%
\end{proof}
%
\noindent\textbf{Deviation from the notes.} Both identities are stated with series-partial-sum(..., succ(n)) in place of the notes' sum from l = 0 to n, and bernstein-basis(n, x) is a function on ZZ (bernstein-basis-in-fun), so the out-of-range terms are zero by the separate bernstein-basis-above and bernstein-basis-null rather than by the range of the sum. The second estimate's summand cannot be written inline: it must be supplied as an auxiliary sequence h\_ pinned pointwise by forall k\_ in nn, h\_(k\_) = ((k\_ - n\_ x\_)\textasciicircum{}2) (bernstein-basis(n\_, x\_))(k\_). bernstein-variance-bound additionally requires n\_ /= 0, which the notes leave implicit in the factor 1/n.
%
\renewcommand{\thethm}{5.5}
\begin{lem}[{The four Bernstein identities}]\label{thm:calculus-5.5}
%
\leavevmode
\vnbpart{\texttt{bernstein-partition}}
\vnbstmt{\begin{array}{@{}l@{}} \forall n \in \mathbb{N}, x \in \mathbb{R}. \\ \quad \operatorname{series\text{-}partial\text{-}sum}(\operatorname{bernstein\text{-}basis}(n, x), (n + 1)) \\ \quad =\; 1 \end{array}}
\vnbpart{\texttt{bernstein-moment-1}}
\vnbstmt{\begin{array}{@{}l@{}} \forall n \in \mathbb{N}, x, h. \\ \quad (x \in \mathbb{R}) \wedge \\ \quad \forall k \in \mathbb{N}. \\ \quad \quad (h(k) = (k \cdot \operatorname{bernstein\text{-}basis}(n, x)(k))) \Rightarrow \\ \quad \quad (\operatorname{series\text{-}partial\text{-}sum}(h, (n + 1)) = (n \cdot x)) \end{array}}
\vnbpart{\texttt{bernstein-moment-2}}
\vnbstmt{\begin{array}{@{}l@{}} \forall n \in \mathbb{N}, x, h. \\ \quad (x \in \mathbb{R}) \wedge \\ \quad \forall k \in \mathbb{N}. \\ \quad \quad (h(k) = ((k \cdot (k - 1)) \cdot \operatorname{bernstein\text{-}basis}(n, x)(k))) \Rightarrow \\ \quad \quad \operatorname{series\text{-}partial\text{-}sum}(h, (n + 1)) \\ \quad \quad =\; ((n \cdot (n - 1)) \cdot (x \cdot x)) \end{array}}
\vnbpart{\texttt{bernstein-variance}}
\vnbstmt{\begin{array}{@{}l@{}} \forall n, x, h. \\ \quad (n \in \mathbb{N}) \wedge \\ \quad (x \in \mathbb{R}) \wedge \\ \quad \forall k \in \mathbb{N}. \\ \quad \quad h(k) \\ \quad \quad =\; (((k - (n \cdot x)) \cdot (k - (n \cdot x))) \cdot \operatorname{bernstein\text{-}basis}(n, x)(k)) \Rightarrow \\ \quad \quad (\operatorname{series\text{-}partial\text{-}sum}(h, (n + 1)) = ((n \cdot x) \cdot (1 - x))) \end{array}}
%
\end{lem}
%
\begin{proof}
%
The proofs: the file that holds each, the lemmas it cites, its oracles and its bill.
%
\vnbprf{\vnbproofpart{\texttt{bernstein-partition}}}
%
\vnbprf{\textbf{Proof.} \texttt{theorem-\allowbreak{}library/bernstein-\allowbreak{}basis.scm} --- 48 recorded steps.}
%
\vnbprf{\textbf{Lemmas.} \texttt{rr-is-normed-field}, \texttt{rr-one-in}, \texttt{rr-sub-in-rr}, \texttt{nn-succ-closed}, \texttt{rr-scalar-ring-carr}, \texttt{series-partial-sum-is-ring-sum}, \texttt{binomial-sum-value}, \texttt{rr-scalar-ring-add}, \texttt{rr-scalar-ring-one}, \texttt{lam-slot-add-apply}, \texttt{rr-add-closed}, \texttt{ring-power-of-one}.}
%
\vnbprf{\textbf{Oracles.} \texttt{crs}, \texttt{arith}, \texttt{ineq}.  \textbf{Bill:} \texttt{proven modulo 0}.}
%
\vnbprf{\vnbproofpart{\texttt{bernstein-moment-1}}}
%
\vnbprf{\textbf{Proof.} \texttt{theorem-\allowbreak{}library/bernstein-\allowbreak{}moments.scm} --- 190 recorded steps.}
%
\vnbprf{\textbf{Lemmas.} \texttt{nn-zero-in}, \texttt{zz-zero-in}, \texttt{rr-zero-in}, \texttt{bernstein-basis-ptwise-in-rr}, \texttt{series-partial-sum-zero}, \texttt{rr-mul-closed}, \texttt{series-partial-sum-succ}, \texttt{bt-succ-in-nn}, \texttt{nn-subset-zz}, \texttt{rr-one-in}, \texttt{rr-sub-in-rr}, \texttt{nn-in-rr}, \texttt{zz-in-rr}, \texttt{nn-succ-plus-one}, \texttt{bernstein-weighted-step}, \texttt{bernstein-basis-in-fun}, \texttt{series-partial-sum-add-ptwise}, \texttt{bernstein-partition}.}
%
\vnbprf{\textbf{Oracles.} \texttt{crs}, \texttt{arith}, \texttt{ineq}.  \textbf{Bill:} \texttt{proven modulo 0}.}
%
\vnbprf{\vnbproofpart{\texttt{bernstein-moment-2}}}
%
\vnbprf{\textbf{Proof.} \texttt{theorem-\allowbreak{}library/bernstein-\allowbreak{}moments.scm} --- 355 recorded steps.}
%
\vnbprf{\textbf{Lemmas.} \texttt{nn-zero-in}, \texttt{zz-zero-in}, \texttt{rr-zero-in}, \texttt{rr-one-in}, \texttt{bernstein-basis-ptwise-in-rr}, \texttt{series-partial-sum-zero}, \texttt{rr-sub-in-rr}, \texttt{rr-mul-closed}, \texttt{series-partial-sum-succ}, \texttt{bt-succ-in-nn}, \texttt{nn-subset-zz}, \texttt{nn-in-rr}, \texttt{zz-in-rr}, \texttt{rr-add-closed}, \texttt{nn-succ-plus-one}, \texttt{bernstein-weighted-step}, \texttt{series-partial-sum-add-ptwise}, \texttt{bernstein-moment-1}.}
%
\vnbprf{\textbf{Oracles.} \texttt{crs}, \texttt{arith}, \texttt{ineq}.  \textbf{Bill:} \texttt{proven modulo 0}.}
%
\vnbprf{\vnbproofpart{\texttt{bernstein-variance}}}
%
\vnbprf{\textbf{Proof.} \texttt{theorem-\allowbreak{}library/bernstein-\allowbreak{}moments.scm} --- 179 recorded steps.}
%
\vnbprf{\textbf{Lemmas.} \texttt{nn-in-rr}, \texttt{rr-one-in}, \texttt{nn-subset-zz}, \texttt{bt-succ-in-nn}, \texttt{rr-mul-in-rr}, \texttt{rr-sub-in-rr}, \texttt{bernstein-basis-ptwise-in-rr}, \texttt{rr-add-in-rr}, \texttt{series-partial-sum-add-ptwise}, \texttt{series-partial-sum-scale-ptwise}, \texttt{bernstein-basis-in-fun}, \texttt{bernstein-moment-2}, \texttt{bernstein-moment-1}, \texttt{bernstein-partition}.}
%
\vnbprf{\textbf{Oracles.} \texttt{crs}, \texttt{arith}, \texttt{ineq}.  \textbf{Bill:} \texttt{proven modulo 0}.}
%
\end{proof}
%
\noindent\textbf{Deviation from the notes.} Each of the last three identities is stated through an auxiliary sequence h\_ pinned pointwise (forall k\_ in nn, h\_(k\_) = ...) and the sum is series-partial-sum(h\_, succ(n\_)), because the library has no inline finite-sum term with a bound variable; bernstein-basis(n, x) is a function on ZZ, so the range of the notes' sum is enforced by bernstein-basis-above and bernstein-basis-null rather than by the summation. The notes derive the identities from the Binomial Theorem by differentiating in x; the library's route is bernstein-basis-succ and induction, so the proof does not follow the notes (binomial-theorem exists independently).
%
\renewcommand{\thethm}{5.7}
\begin{cor}[{Weierstrass approximation on an arbitrary interval}]\label{thm:bernstein-uniform-approximation-ccint}
%
\vnbstmt{\begin{array}{@{}l@{}} \forall n, a, b, p. \\ \quad (n \in (\mathbb{R} \to \mathbb{R})) \wedge \\ \quad (a \in \mathbb{R}) \wedge \\ \quad (b \in \mathbb{R}) \wedge \\ \quad (a < b) \wedge \\ \quad \forall c \in \operatorname{ccint}(a, b). \\ \quad \quad \operatorname{is\text{-}continuous\text{-}at}(\operatorname{rr\text{-}ms}, \operatorname{rr\text{-}ms}, n, c) \wedge \\ \quad \operatorname{pos\text{-}rr}(p) \Rightarrow \\ \quad \quad \exists \mathit{cap} \in \mathbb{N}. \forall d \in \mathbb{N}, c \in \operatorname{ccint}(a, b). \\ \quad \quad \quad (\mathit{cap} \leq d) \Rightarrow  \\ \quad \quad \quad \quad \lvert (n(c) - \operatorname{bernstein\text{-}poly}(\operatorname{compose}(n, (\lambda x \in \mathbb{R}.\; (a + ((b - a) \cdot x)))), d, ((c - a) \cdot \operatorname{recip}((b - a))))) \rvert \\ \quad \quad \quad \quad \leq\; p \end{array}}
%
\end{cor}
%
\begin{proof}
%
The proofs: the file that holds each, the lemmas it cites, its oracles and its bill.
%
\vnbprf{\textbf{Proof.} \texttt{theorem-\allowbreak{}library/bernstein-\allowbreak{}ccint.scm} --- 129 recorded steps.}
%
\vnbprf{\textbf{Lemmas.} \texttt{rr-sub-in-rr}, \texttt{rr-is-set}, \texttt{affine-lam-in-fun}, \texttt{compose-type}, \texttt{ccint-parts}, \texttt{affine-continuous-at}, \texttt{ccint-affine-in}, \texttt{compose-continuous-at}, \texttt{bernstein-uniform-approximation}, \texttt{ccint-coaffine-in}, \texttt{ccint-affine-inverse}, \texttt{fun-apply-type-c}, \texttt{compose-apply}.}
%
\vnbprf{\textbf{Oracles.} \texttt{crs}, \texttt{arith}, \texttt{ineq}, \texttt{sos}.  \textbf{Bill:} \texttt{proven modulo 0}.}
%
\end{proof}
%
\noindent\textbf{Deviation from the notes.} The library exhibits ONE SPECIFIC approximant -- the transported Bernstein polynomial bernstein-poly(compose(fn\_, x |--$>$ a + (b-a)*x), n\_, (x\_ - a) * recip(b - a)) -- and bounds abs(fn\_(x\_) - that) by eps. It never says that this term IS a polynomial, and there is no statement anywhere of the form the notes make, that any continuous function on [a,b] is a uniform limit of a sequence of POLYNOMIALS: there is no predicate for being a polynomial function on RR to connect the two. Further: fn\_ in fun(rr, rr), TOTAL on RR; continuity assumed two-sided at every x\_ in ccint(a, b); a $<$ b strictly, so the degenerate interval is excluded; and eps is a parameter with the claim in eps/N form rather than the notes' uniform-limit form.
%
\section{Linear algebra}
%
\renewcommand{\thethm}{3.3}
\begin{lem}[{Action of the matrix unit E\_n[k,l] on the right}]\label{thm:matunit-col-shift}
%
\vnbstmt{\begin{array}{@{}l@{}} \forall a, m, n, p, k, l.; \forall i \in \operatorname{interval}(1, m), c \in \operatorname{interval}(1, n). \\ \quad \operatorname{is\text{-}ring}(a) \wedge \\ \quad (p \in \operatorname{mat}(m, n, \operatorname{carr}(a))) \wedge \\ \quad (k \in \operatorname{interval}(1, n)) \wedge \\ \quad (l \in \operatorname{interval}(1, n)) \Rightarrow \\ \quad \quad \operatorname{entry}(\operatorname{matmul}(a, p, \operatorname{matunit}(a, n, k, l)), i, c) \\ \quad \quad =\; \left(\text{if }(c = l)\text{ then }\operatorname{entry}(p, i, k)\text{ else }\operatorname{zero}(a)\right) \end{array}}
%
\end{lem}
%
\begin{proof}
%
The proofs: the file that holds each, the lemmas it cites, its oracles and its bill.
%
\vnbprf{\textbf{Proof.} \texttt{theorem-\allowbreak{}library/matunit-\allowbreak{}shift-\allowbreak{}proof.scm} --- 80 recorded steps.}
%
\vnbprf{\textbf{Lemmas.} \texttt{mat-cols-in-nn}, \texttt{interval-elt-in-nn}, \texttt{interval-lo}, \texttt{interval-hi}, \texttt{nn-one-in}, \texttt{nn-le-trans-guarded}, \texttt{matunit-type}, \texttt{matmul-entry}, \texttt{interval-in-set}, \texttt{mat-rows-in-nn}, \texttt{interval-card-in-nn}, \texttt{matunit-summand-type}, \texttt{matunit-entry-off-row}, \texttt{entry-in-carrier}, \texttt{ring-mul-zero-right}, \texttt{ras-id}, \texttt{finsum-single-support}, \texttt{matunit-entry-k-row}.}
%
\vnbprf{\textbf{Oracles.} \texttt{arith}, \texttt{ineq}.  \textbf{Bill:} \texttt{proven modulo 0}.}
%
\end{proof}
%
\noindent\textbf{Deviation from the notes.} Generality GAINED: the library proves it over an arbitrary ring (is-ring), while the notes' section fixes R commutative.  The library states the conclusion entrywise -- entry(matmul(a,p,matunit(a,n,k,l)),i,c) = if(c = l, entry(p,i,k), zero(a)) -- where the notes name the resulting matrix A'; same content.  No k /= l guard on either side.
%
\renewcommand{\thethm}{3.5}
\begin{prop}[{The elementary column operations}]\label{thm:elem-f-action}
%
\leavevmode
\vnbpart{\texttt{elem-f-action}}
\vnbstmt{\begin{array}{@{}l@{}} \forall a, m, n, p, k, l.; \forall i \in \operatorname{interval}(1, m), c \in \operatorname{interval}(1, n). \\ \quad \operatorname{is\text{-}ring}(a) \wedge \\ \quad (p \in \operatorname{mat}(m, n, \operatorname{carr}(a))) \wedge \\ \quad (k \in \operatorname{interval}(1, n)) \wedge \\ \quad (l \in \operatorname{interval}(1, n)) \wedge \\ \quad \neg (k = l) \Rightarrow \\ \quad \quad \operatorname{entry}(\operatorname{matmul}(a, p, \operatorname{elem\text{-}f}(a, n, k, l)), i, c) \\ \quad \quad =\; \left(\text{if }(c = k)\text{ then }\operatorname{entry}(p, i, l)\text{ else }\left(\text{if }(c = l)\text{ then }\operatorname{entry}(p, i, k)\text{ else }\operatorname{entry}(p, i, c)\right)\right) \end{array}}
\vnbpart{\texttt{elem-g-action}}
\vnbstmt{\begin{array}{@{}l@{}} \forall a, m, n, p, r, k, l.; \forall i \in \operatorname{interval}(1, m), c \in \operatorname{interval}(1, n). \\ \quad \operatorname{is\text{-}ring}(a) \wedge \\ \quad (p \in \operatorname{mat}(m, n, \operatorname{carr}(a))) \wedge \\ \quad (r \in \operatorname{carr}(a)) \wedge \\ \quad (k \in \operatorname{interval}(1, n)) \wedge \\ \quad (l \in \operatorname{interval}(1, n)) \wedge \\ \quad \neg (k = l) \Rightarrow \\ \quad \quad \operatorname{entry}(\operatorname{matmul}(a, p, \operatorname{elem\text{-}g}(a, n, r, k, l)), i, c) \\ \quad \quad =\; \left(\text{if }(c = l)\text{ then }\operatorname{add}(a)(\operatorname{entry}(p, i, l), \operatorname{mul}(a)(\operatorname{entry}(p, i, k), r))\text{ else }\operatorname{entry}(p, i, c)\right) \end{array}}
\vnbpart{\texttt{elem-h-action}}
\vnbstmt{\begin{array}{@{}l@{}} \forall a, m, n, p, r, k.; \forall i \in \operatorname{interval}(1, m), c \in \operatorname{interval}(1, n). \\ \quad \operatorname{is\text{-}ring}(a) \wedge \\ \quad (p \in \operatorname{mat}(m, n, \operatorname{carr}(a))) \wedge \\ \quad (r \in \operatorname{carr}(a)) \wedge \\ \quad (k \in \operatorname{interval}(1, n)) \Rightarrow \\ \quad \quad \operatorname{entry}(\operatorname{matmul}(a, p, \operatorname{elem\text{-}h}(a, n, r, k)), i, c) \\ \quad \quad =\; \left(\text{if }(c = k)\text{ then }\operatorname{mul}(a)(\operatorname{entry}(p, i, k), r)\text{ else }\operatorname{entry}(p, i, c)\right) \end{array}}
%
\end{prop}
%
\begin{proof}
%
The proofs: the file that holds each, the lemmas it cites, its oracles and its bill.
%
\vnbprf{\vnbproofpart{\texttt{elem-f-action}}}
%
\vnbprf{\textbf{Proof.} \texttt{theorem-\allowbreak{}library/elem-\allowbreak{}actions-\allowbreak{}proof.scm} --- 140 recorded steps.}
%
\vnbprf{\textbf{Lemmas.} \texttt{mat-cols-in-nn}, \texttt{interval-elt-in-nn}, \texttt{interval-lo}, \texttt{interval-hi}, \texttt{nn-one-in}, \texttt{nn-le-trans-guarded}, \texttt{elem-f-type}, \texttt{matmul-entry}, \texttt{interval-in-set}, \texttt{mat-rows-in-nn}, \texttt{interval-card-in-nn}, \texttt{matprod-summand-type}, \texttt{elem-f-ck-off}, \texttt{entry-in-carrier}, \texttt{ring-mul-zero-right}, \texttt{ras-id}, \texttt{finsum-single-support}, \texttt{elem-f-ck-at}, \texttt{elem-f-cl-off}, \texttt{elem-f-cl-at}, \texttt{elem-f-co-off}, \texttt{elem-f-co-at}.}
%
\vnbprf{\textbf{Oracles.} \texttt{arith}, \texttt{ineq}.  \textbf{Bill:} \texttt{proven modulo 0}.}
%
\vnbprf{\vnbproofpart{\texttt{elem-g-action}}}
%
\vnbprf{\textbf{Proof.} \texttt{theorem-\allowbreak{}library/elem-\allowbreak{}actions-\allowbreak{}proof.scm} --- 102 recorded steps.}
%
\vnbprf{\textbf{Lemmas.} \texttt{mat-cols-in-nn}, \texttt{interval-elt-in-nn}, \texttt{interval-lo}, \texttt{interval-hi}, \texttt{nn-one-in}, \texttt{nn-le-trans-guarded}, \texttt{elem-g-type}, \texttt{matmul-entry}, \texttt{interval-in-set}, \texttt{mat-rows-in-nn}, \texttt{interval-card-in-nn}, \texttt{matprod-summand-type}, \texttt{eq-sym}, \texttt{elem-g-entry-l-vanish}, \texttt{entry-in-carrier}, \texttt{ring-mul-zero-right}, \texttt{ras-id}, \texttt{finsum-two-support}, \texttt{elem-g-entry-l-at-l}, \texttt{elem-g-entry-l-at-k}, \texttt{ras-op}, \texttt{elem-g-entry-off-l-off-diag}, \texttt{finsum-single-support}, \texttt{elem-g-entry-off-l-diag}.}
%
\vnbprf{\textbf{Oracles.} \texttt{arith}, \texttt{ineq}.  \textbf{Bill:} \texttt{proven modulo 0}.}
%
\vnbprf{\vnbproofpart{\texttt{elem-h-action}}}
%
\vnbprf{\textbf{Proof.} \texttt{theorem-\allowbreak{}library/elem-\allowbreak{}actions-\allowbreak{}proof.scm} --- 79 recorded steps.}
%
\vnbprf{\textbf{Lemmas.} \texttt{mat-cols-in-nn}, \texttt{interval-elt-in-nn}, \texttt{interval-lo}, \texttt{interval-hi}, \texttt{nn-one-in}, \texttt{nn-le-trans-guarded}, \texttt{elem-h-type}, \texttt{matmul-entry}, \texttt{interval-in-set}, \texttt{mat-rows-in-nn}, \texttt{interval-card-in-nn}, \texttt{matprod-summand-type}, \texttt{elem-h-entry-off-diag}, \texttt{entry-in-carrier}, \texttt{ring-mul-zero-right}, \texttt{ras-id}, \texttt{finsum-single-support}, \texttt{elem-h-entry-diag}.}
%
\vnbprf{\textbf{Oracles.} \texttt{arith}, \texttt{ineq}.  \textbf{Bill:} \texttt{proven modulo 0}.}
%
\end{proof}
%
\noindent\textbf{Deviation from the notes.} Generality GAINED: stated over an arbitrary ring (is-ring), not a commutative ring, and elem-h-action asks only r in carr(a), NOT that r be invertible (the notes build invertibility of r into the definition of H).  Split into three entrywise theorems on matmul(a,p,elem-x(...)) instead of one prose sentence; elem-f-action and elem-g-action carry the notes' own k /= l guard.  The row-operation mirrors (elem-f-row-action, elem-g-row-action, elem-h-row-action) are proven as well, which the notes only define by transposition.
%
\renewcommand{\thethm}{3.6}
\begin{cor}[{The elementary column matrices are invertible, with explicit inverses}]\label{thm:elem-f-inverse}
%
\leavevmode
\vnbpart{\texttt{elem-f-inverse}}
\vnbstmt{\begin{array}{@{}l@{}} \forall a, n, k, l. \\ \quad \operatorname{is\text{-}ring}(a) \wedge \\ \quad (n \in \mathbb{N}) \wedge \\ \quad (k \in \operatorname{interval}(1, n)) \wedge \\ \quad (l \in \operatorname{interval}(1, n)) \wedge \\ \quad \neg (k = l) \Rightarrow \\ \quad \quad \operatorname{matmul}(a, \operatorname{elem\text{-}f}(a, n, k, l), \operatorname{elem\text{-}f}(a, n, l, k)) \\ \quad \quad =\; \operatorname{identmat}(a, n) \end{array}}
\vnbpart{\texttt{elem-g-inverse}}
\vnbstmt{\begin{array}{@{}l@{}} \forall a, n, r, k, l. \\ \quad \operatorname{is\text{-}ring}(a) \wedge \\ \quad (n \in \mathbb{N}) \wedge \\ \quad (r \in \operatorname{carr}(a)) \wedge \\ \quad (k \in \operatorname{interval}(1, n)) \wedge \\ \quad (l \in \operatorname{interval}(1, n)) \wedge \\ \quad \neg (k = l) \Rightarrow \\ \quad \quad \operatorname{matmul}(a, \operatorname{elem\text{-}g}(a, n, r, k, l), \operatorname{elem\text{-}g}(a, n, \operatorname{neg}(a)(r), k, l)) \\ \quad \quad =\; \operatorname{identmat}(a, n) \end{array}}
\vnbpart{\texttt{elem-h-inverse}}
\vnbstmt{\begin{array}{@{}l@{}} \forall a, n, r, s, k. \\ \quad \operatorname{is\text{-}ring}(a) \wedge \\ \quad (n \in \mathbb{N}) \wedge \\ \quad (r \in \operatorname{carr}(a)) \wedge \\ \quad (s \in \operatorname{carr}(a)) \wedge \\ \quad (\operatorname{mul}(a)(r, s) = \operatorname{one}(a)) \wedge \\ \quad (k \in \operatorname{interval}(1, n)) \Rightarrow \\ \quad \quad \operatorname{matmul}(a, \operatorname{elem\text{-}h}(a, n, r, k), \operatorname{elem\text{-}h}(a, n, s, k)) \\ \quad \quad =\; \operatorname{identmat}(a, n) \end{array}}
\vnbpart{\texttt{elem-f-invertible}}
\vnbstmt{\begin{array}{@{}l@{}} \forall a, n, k, l. \\ \quad \operatorname{is\text{-}ring}(a) \wedge \\ \quad (n \in \mathbb{N}) \wedge \\ \quad (k \in \operatorname{interval}(1, n)) \wedge \\ \quad (l \in \operatorname{interval}(1, n)) \wedge \\ \quad \neg (k = l) \Rightarrow \\ \quad \quad \operatorname{is\text{-}invertible\text{-}mat}(a, n, \operatorname{elem\text{-}f}(a, n, k, l)) \end{array}}
\vnbpart{\texttt{elem-g-invertible}}
\vnbstmt{\begin{array}{@{}l@{}} \forall a, n, r, k, l. \\ \quad \operatorname{is\text{-}ring}(a) \wedge \\ \quad (n \in \mathbb{N}) \wedge \\ \quad (r \in \operatorname{carr}(a)) \wedge \\ \quad (k \in \operatorname{interval}(1, n)) \wedge \\ \quad (l \in \operatorname{interval}(1, n)) \wedge \\ \quad \neg (k = l) \Rightarrow \\ \quad \quad \operatorname{is\text{-}invertible\text{-}mat}(a, n, \operatorname{elem\text{-}g}(a, n, r, k, l)) \end{array}}
%
\end{cor}
%
\begin{proof}
%
The proofs: the file that holds each, the lemmas it cites, its oracles and its bill.
%
\vnbprf{\vnbproofpart{\texttt{elem-f-inverse}}}
%
\vnbprf{\textbf{Proof.} \texttt{theorem-\allowbreak{}library/elem-\allowbreak{}inverses-\allowbreak{}proof.scm} --- 93 recorded steps.}
%
\vnbprf{\textbf{Lemmas.} \texttt{eq-sym}, \texttt{elem-f-type}, \texttt{matmul-type}, \texttt{identmat-type}, \texttt{elem-f-action}, \texttt{entry-of-identmat}, \texttt{elem-f-col-k}, \texttt{elem-f-col-l}, \texttt{elem-f-col-other}, \texttt{matrix-entry-extensionality}.}
%
\vnbprf{\textbf{Oracles.} \texttt{arith}, \texttt{ineq}.  \textbf{Bill:} \texttt{proven modulo 0}.}
%
\vnbprf{\vnbproofpart{\texttt{elem-g-inverse}}}
%
\vnbprf{\textbf{Proof.} \texttt{theorem-\allowbreak{}library/elem-\allowbreak{}inverses-\allowbreak{}proof.scm} --- 157 recorded steps.}
%
\vnbprf{\textbf{Lemmas.} \texttt{ring-neg-in-carr}, \texttt{elem-g-type}, \texttt{matmul-type}, \texttt{identmat-type}, \texttt{elem-g-action}, \texttt{entry-of-identmat}, \texttt{elem-g-col-l}, \texttt{elem-g-col-k}, \texttt{eq-sym}, \texttt{ring-mul-zero-left}, \texttt{ring-one-in}, \texttt{ring-add-right-id}, \texttt{ring-add-right-inv}, \texttt{elem-g-col-other}, \texttt{matrix-entry-extensionality}.}
%
\vnbprf{\textbf{Oracles.} \texttt{arith}, \texttt{ineq}.  \textbf{Bill:} \texttt{proven modulo 0}.}
%
\vnbprf{\vnbproofpart{\texttt{elem-h-inverse}}}
%
\vnbprf{\textbf{Proof.} \texttt{theorem-\allowbreak{}library/elem-\allowbreak{}inverses-\allowbreak{}proof.scm} --- 130 recorded steps.}
%
\vnbprf{\textbf{Lemmas.} \texttt{elem-h-type}, \texttt{matmul-type}, \texttt{identmat-type}, \texttt{elem-h-action}, \texttt{entry-of-identmat}, \texttt{elem-h-entry-diag}, \texttt{elem-h-entry-off-diag}, \texttt{ring-mul-zero-left}, \texttt{matrix-entry-extensionality}.}
%
\vnbprf{\textbf{Oracles.} \texttt{arith}, \texttt{ineq}.  \textbf{Bill:} \texttt{proven modulo 0}.}
%
\vnbprf{\vnbproofpart{\texttt{elem-f-invertible}}}
%
\vnbprf{\textbf{Proof.} \texttt{theorem-\allowbreak{}library/mat-\allowbreak{}equiv-\allowbreak{}proof.scm} --- 31 recorded steps.}
%
\vnbprf{\textbf{Lemmas.} \texttt{elem-f-type}, \texttt{elem-f-inverse}, \texttt{eq-sym}.}
%
\vnbprf{\textbf{Oracles.} \texttt{arith}, \texttt{ineq}.  \textbf{Bill:} \texttt{proven modulo 0}.}
%
\vnbprf{\vnbproofpart{\texttt{elem-g-invertible}}}
%
\vnbprf{\textbf{Proof.} \texttt{theorem-\allowbreak{}library/mat-\allowbreak{}equiv-\allowbreak{}proof.scm} --- 34 recorded steps.}
%
\vnbprf{\textbf{Lemmas.} \texttt{ring-neg-in-carr}, \texttt{elem-g-type}, \texttt{elem-g-inverse}, \texttt{ring-neg-neg}, \texttt{eq-sym}, \texttt{elem-g-param-cong}.}
%
\vnbprf{\textbf{Oracles.} \texttt{arith}, \texttt{ineq}.  \textbf{Bill:} \texttt{proven modulo 0}.}
%
\end{proof}
%
\noindent\textbf{Deviation from the notes.} Two differences.  (1) F and G are given is-invertible-mat (elem-f-invertible, elem-g-invertible), but there is NO elem-h-invertible: for H the library installs only the product equation matmul(a,elem-h(a,n,r,k),elem-h(a,n,s,k)) = identmat(a,n) under the hypothesis mul(a)(r,s) = one(a), so the inverse is a SUPPLIED s, not r\textasciicircum{}\{-1\} -- the library has no ring-inverse operation.  (2) Generality gained: is-ring rather than a commutative ring.  The notes' G inverse formula G[-r,k,l] and the k /= l guard match elem-g-inverse exactly.
%
\renewcommand{\thethm}{3.8}
\begin{lem}[{Column reduction of the first row over a field}]\label{thm:clear-first-row}
%
\leavevmode
\vnbpart{\texttt{clear-first-row}}
\vnbstmt{\begin{array}{@{}l@{}} \forall a, m, n.; \forall p \in \operatorname{mat}(m, n, \operatorname{carr}(a)). \\ \quad \operatorname{is\text{-}euclidean\text{-}ring}(a) \wedge \\ \quad (m \in \mathbb{N}) \wedge \\ \quad (n \in \mathbb{N}) \wedge \\ \quad (1 \in \operatorname{interval}(1, m)) \wedge \\ \quad (1 \in \operatorname{interval}(1, n)) \wedge \\ \quad \neg (\operatorname{entry}(p, 1, 1) = \operatorname{zero}(a)) \wedge \\ \quad \forall c, i, b. \\ \quad \quad \operatorname{mat\text{-}equiv}(a, m, n, p, c) \wedge \\ \quad \quad (i \in \operatorname{interval}(1, m)) \wedge \\ \quad \quad (b \in \operatorname{interval}(1, n)) \wedge \\ \quad \quad \neg (\operatorname{entry}(c, i, b) = \operatorname{zero}(a)) \Rightarrow \\ \quad \quad \quad \operatorname{gauge}(a)(\operatorname{entry}(p, 1, 1)) \\ \quad \quad \quad \leq\; \operatorname{gauge}(a)(\operatorname{entry}(c, i, b)) \Rightarrow \\ \quad \quad \exists q. \\ \quad \quad \quad \operatorname{mat\text{-}equiv}(a, m, n, p, q) \wedge \\ \quad \quad \quad (\operatorname{entry}(q, 1, 1) = \operatorname{entry}(p, 1, 1)) \wedge \\ \quad \quad \quad \forall j \in \operatorname{interval}(1, n). \\ \quad \quad \quad \quad \neg (j = 1) \Rightarrow  \\ \quad \quad \quad \quad \quad (\operatorname{entry}(q, 1, j) = \operatorname{zero}(a)) \end{array}}
\vnbpart{\texttt{class-min-pivot}}
\vnbstmt{\begin{array}{@{}l@{}} \forall a, m, n.; \forall p \in \operatorname{mat}(m, n, \operatorname{carr}(a)). \\ \quad \operatorname{is\text{-}euclidean\text{-}ring}(a) \wedge \\ \quad (1 \in \operatorname{interval}(1, m)) \wedge \\ \quad (1 \in \operatorname{interval}(1, n)) \wedge \\ \quad \exists d \in \operatorname{interval}(1, m), f \in \operatorname{interval}(1, n). \\ \quad \quad \neg (\operatorname{entry}(p, d, f) = \operatorname{zero}(a)) \Rightarrow \\ \quad \quad \exists b. \\ \quad \quad \quad \operatorname{mat\text{-}equiv}(a, m, n, p, b) \wedge \\ \quad \quad \quad \neg (\operatorname{entry}(b, 1, 1) = \operatorname{zero}(a)) \wedge \\ \quad \quad \quad \forall c, i, j. \\ \quad \quad \quad \quad \operatorname{mat\text{-}equiv}(a, m, n, p, c) \wedge \\ \quad \quad \quad \quad (i \in \operatorname{interval}(1, m)) \wedge \\ \quad \quad \quad \quad (j \in \operatorname{interval}(1, n)) \wedge \\ \quad \quad \quad \quad \neg (\operatorname{entry}(c, i, j) = \operatorname{zero}(a)) \Rightarrow \\ \quad \quad \quad \quad \quad \operatorname{gauge}(a)(\operatorname{entry}(b, 1, 1)) \\ \quad \quad \quad \quad \quad \leq\; \operatorname{gauge}(a)(\operatorname{entry}(c, i, j)) \end{array}}
%
\end{lem}
%
\begin{proof}
%
The proofs: the file that holds each, the lemmas it cites, its oracles and its bill.
%
\vnbprf{\vnbproofpart{\texttt{clear-first-row}}}
%
\vnbprf{\textbf{Proof.} \texttt{theorem-\allowbreak{}library/clear-\allowbreak{}first-\allowbreak{}row-\allowbreak{}proof.scm} --- 28 recorded steps.}
%
\vnbprf{\textbf{Lemmas.} \texttt{clear-row-upto}, \texttt{interval-hi}.}
%
\vnbprf{\textbf{Oracles.} \texttt{arith}, \texttt{ineq}, \texttt{crs}.  \textbf{Bill:} \texttt{proven modulo 0}.}
%
\vnbprf{\vnbproofpart{\texttt{class-min-pivot}}}
%
\vnbprf{\textbf{Proof.} \texttt{theorem-\allowbreak{}library/class-\allowbreak{}min-\allowbreak{}pivot-\allowbreak{}proof.scm} --- 99 recorded steps.}
%
\vnbprf{\textbf{Lemmas.} \texttt{euclidean-ring-is-integral-domain}, \texttt{integral-domain-is-commutative-ring}, \texttt{commutative-ring-is-ring}, \texttt{entry-in-carrier}, \texttt{gauge-is-degree}, \texttt{fun-apply-type-c}, \texttt{mat-equiv-refl}, \texttt{swap-to-corner-gen}, \texttt{mat-equiv-trans}, \texttt{mat-equiv-target-is-mat}.}
%
\vnbprf{\textbf{Oracles.} \texttt{arith}, \texttt{ineq}.  \textbf{Bill:} \texttt{proven modulo 0}.}
%
\end{proof}
%
\noindent\textbf{Deviation from the notes.} Three losses and one gain.  The library works over a EUCLIDEAN RING, not a field (gain), and therefore (a) cannot scale the pivot: clear-first-row concludes entry(q,1,1) = entry(p,1,1) and never normalises the corner to 0 or 1, which is precisely the notes' conclusion b\_11 in \{0,1\}; (b) needs the pivot to be nonzero AND to MINIMISE the gauge over the whole equivalence class (that hypothesis is supplied by class-min-pivot), where the notes need no such thing; (c) works modulo MAT-EQUIV, which is two-sided (an invertible U on the left and V on the right), not the notes' column equivalence.  Extra guards: m, n in nn and 1 in interval(1,m), 1 in interval(1,n).
%
\renewcommand{\thethm}{3.10}
\begin{prop}[{Column reduction to a step matrix with zeros above the diagonal}]\label{thm:smith-diagonalization}
%
\vnbstmt{\begin{array}{@{}l@{}} \forall a.; \forall k \in \mathbb{N}, n, p. \\ \quad \operatorname{is\text{-}euclidean\text{-}ring}(a) \wedge \\ \quad (n \in \mathbb{N}) \wedge \\ \quad (p \in \operatorname{mat}(k, n, \operatorname{carr}(a))) \Rightarrow \\ \quad \quad \exists d. \\ \quad \quad \quad \operatorname{mat\text{-}equiv}(a, k, n, p, d) \wedge \\ \quad \quad \quad \operatorname{is\text{-}diagonal}(a, k, n, d) \end{array}}
%
\end{prop}
%
\begin{proof}
%
The proofs: the file that holds each, the lemmas it cites, its oracles and its bill.
%
\vnbprf{\textbf{Proof.} \texttt{theorem-\allowbreak{}library/smith-\allowbreak{}diagonalization-\allowbreak{}proof.scm} --- 162 recorded steps.}
%
\vnbprf{\textbf{Lemmas.} \texttt{euclidean-ring-is-integral-domain}, \texttt{integral-domain-is-commutative-ring}, \texttt{commutative-ring-is-ring}, \texttt{mat-equiv-refl}, \texttt{interval-lo}, \texttt{interval-hi}, \texttt{interval-elt-in-nn}, \texttt{nn-not-le-zero-pos}, \texttt{nn-one-in}, \texttt{nn-le-trans-guarded}, \texttt{nn-pos-is-succ}, \texttt{eq-sym}, \texttt{clear-pivot-cross}, \texttt{mat-equiv-target-is-mat}, \texttt{submat-type}, \texttt{one-in-interval}, \texttt{entry-in-carrier}, \texttt{bordering}, \texttt{border-is-diagonal}, \texttt{border-type}, \texttt{mat-equiv-trans}.}
%
\vnbprf{\textbf{Oracles.} \texttt{arith}, \texttt{ineq}, \texttt{crs}.  \textbf{Bill:} \texttt{proven modulo 0}.}
%
\end{proof}
%
\noindent\textbf{Deviation from the notes.} Neither statement implies the other.  smith-diagonalization reaches a FULLY DIAGONAL d (is-diagonal: every off-diagonal entry zero) over a euclidean ring, which is stronger than the notes' zeros above the diagonal; but it reaches it through MAT-EQUIV -- row AND column operations -- where the notes use COLUMN operations only, and it does not normalise the diagonal entries to 0 or 1 (that step needs a field).  The library has no column-only equivalence relation at all.
%
\renewcommand{\thethm}{3.12}
\begin{cor}[{Over a field, invertible iff column equivalent to the identity}]\label{thm:linear-algebra-3.12}
%
\textit{If R is a field, A in M\_n(R) is invertible if and only if A is column equivalent to the identity matrix.}
%
Not in the library.
%
\end{cor}
%
\begin{proof}
%
The proofs: the file that holds each, the lemmas it cites, its oracles and its bill.
%
\vnbprf{There is no library statement, and so no proof.}
%
\end{proof}
%
\noindent\textbf{Deviation from the notes.} Nothing in the library relates invertibility to equivalence with the identity matrix, and there is no column-only equivalence to state it with.  Nearest: the definitions is-invertible-mat and mat-equiv, plus identmat-invertible and product-of-invertibles-is-invertible.
%
\renewcommand{\thethm}{3.13}
\begin{cor}[{The elementary column matrices generate the group of invertible matrices}]\label{thm:linear-algebra-3.13}
%
\textit{If R is a field, the n x n elementary column matrices generate the group of invertible matrices in M\_n(R).}
%
Not in the library.
%
\end{cor}
%
\begin{proof}
%
The proofs: the file that holds each, the lemmas it cites, its oracles and its bill.
%
\vnbprf{There is no library statement, and so no proof.}
%
\end{proof}
%
\noindent\textbf{Deviation from the notes.} The library never forms the group of invertible n x n matrices (mat-ring is a ring, and no unit group is carved out of it) and states no generation result for it.  Nearest: elem-f-invertible, elem-g-invertible, product-of-invertibles-is-invertible.
%
\renewcommand{\thethm}{3.15}
\begin{prop}[{A permutation of the indeterminates is a ring automorphism of R[x\_1,...,x\_n]}]\label{thm:linear-algebra-3.15}
%
\textit{For a permutation sigma of the indeterminates, p |-$>$ sigma . p is an isomorphism of the ring R[x\_1, x\_2, ..., x\_n].}
%
Not in the library.
%
\end{prop}
%
\begin{proof}
%
The proofs: the file that holds each, the lemmas it cites, its oracles and its bill.
%
\vnbprf{There is no library statement, and so no proof.}
%
\end{proof}
%
\noindent\textbf{Deviation from the notes.} The library has the monoid algebra monalg(a,m) and, on top of it, poly(a) = monalg over the additive monoid of NN, i.e. ONE-VARIABLE polynomials only; there is no R[x\_1,...,x\_n].  permutations(n) is defined (injection(ord-segment(n), ord-segment(n))) but acts on nothing.  No sign of a permutation, hence none of the sign machinery under this proposition either.
%
\renewcommand{\thethm}{3.18}
\begin{prop}[{det is multilinear and alternating, and det of the transpose equals det}]\label{thm:det-row-linear}
%
\leavevmode
\vnbpart{\texttt{det-row-linear}}
\vnbstmt{\begin{array}{@{}l@{}} \forall r, n, a, b, c, p, d. \\ \quad \operatorname{is\text{-}commutative\text{-}ring}(r) \wedge \\ \quad (n \in \mathbb{N}) \wedge \\ \quad (a \in \operatorname{mat}(n, n, \operatorname{carr}(r))) \wedge \\ \quad (b \in \operatorname{mat}(n, n, \operatorname{carr}(r))) \wedge \\ \quad (c \in \operatorname{mat}(n, n, \operatorname{carr}(r))) \wedge \\ \quad (p \in \operatorname{interval}(1, n)) \wedge \\ \quad (d \in \operatorname{carr}(r)) \wedge \\ \quad \forall i, k \in \operatorname{interval}(1, n). \\ \quad \quad \neg (i = p) \Rightarrow  \\ \quad \quad \quad (\operatorname{entry}(b, i, k) = \operatorname{entry}(a, i, k)) \wedge \\ \quad \forall i, k \in \operatorname{interval}(1, n). \\ \quad \quad \neg (i = p) \Rightarrow  \\ \quad \quad \quad (\operatorname{entry}(c, i, k) = \operatorname{entry}(a, i, k)) \wedge \\ \quad \forall k \in \operatorname{interval}(1, n). \\ \quad \quad \operatorname{entry}(c, p, k) \\ \quad \quad =\; \operatorname{add}(r)(\operatorname{mul}(r)(d, \operatorname{entry}(a, p, k)), \operatorname{entry}(b, p, k)) \Rightarrow \\ \quad \quad \operatorname{det}(r, n, c) \\ \quad \quad =\; \operatorname{add}(r)(\operatorname{mul}(r)(d, \operatorname{det}(r, n, a)), \operatorname{det}(r, n, b)) \end{array}}
\vnbpart{\texttt{det-col-linear}}
\vnbstmt{\begin{array}{@{}l@{}} \forall r, n, a, b, c, p, d. \\ \quad \operatorname{is\text{-}commutative\text{-}ring}(r) \wedge \\ \quad (n \in \mathbb{N}) \wedge \\ \quad (a \in \operatorname{mat}(n, n, \operatorname{carr}(r))) \wedge \\ \quad (b \in \operatorname{mat}(n, n, \operatorname{carr}(r))) \wedge \\ \quad (c \in \operatorname{mat}(n, n, \operatorname{carr}(r))) \wedge \\ \quad (p \in \operatorname{interval}(1, n)) \wedge \\ \quad (d \in \operatorname{carr}(r)) \wedge \\ \quad \forall j, i \in \operatorname{interval}(1, n). \\ \quad \quad \neg (j = p) \Rightarrow  \\ \quad \quad \quad (\operatorname{entry}(b, i, j) = \operatorname{entry}(a, i, j)) \wedge \\ \quad \forall j, i \in \operatorname{interval}(1, n). \\ \quad \quad \neg (j = p) \Rightarrow  \\ \quad \quad \quad (\operatorname{entry}(c, i, j) = \operatorname{entry}(a, i, j)) \wedge \\ \quad \forall i \in \operatorname{interval}(1, n). \\ \quad \quad \operatorname{entry}(c, i, p) \\ \quad \quad =\; \operatorname{add}(r)(\operatorname{mul}(r)(d, \operatorname{entry}(a, i, p)), \operatorname{entry}(b, i, p)) \Rightarrow \\ \quad \quad \operatorname{det}(r, n, c) \\ \quad \quad =\; \operatorname{add}(r)(\operatorname{mul}(r)(d, \operatorname{det}(r, n, a)), \operatorname{det}(r, n, b)) \end{array}}
\vnbpart{\texttt{det-swap-rows}}
\vnbstmt{\begin{array}{@{}l@{}} \forall r, n, a, b, p, q. \\ \quad \operatorname{is\text{-}commutative\text{-}ring}(r) \wedge \\ \quad (n \in \mathbb{N}) \wedge \\ \quad (a \in \operatorname{mat}(n, n, \operatorname{carr}(r))) \wedge \\ \quad (b \in \operatorname{mat}(n, n, \operatorname{carr}(r))) \wedge \\ \quad (p \in \operatorname{interval}(1, n)) \wedge \\ \quad (q \in \operatorname{interval}(1, n)) \wedge \\ \quad \neg (p = q) \wedge \\ \quad \forall i, k \in \operatorname{interval}(1, n). \\ \quad \quad (\neg (i = p) \wedge \neg (i = q)) \Rightarrow  \\ \quad \quad \quad (\operatorname{entry}(b, i, k) = \operatorname{entry}(a, i, k)) \wedge \\ \quad \forall k \in \operatorname{interval}(1, n). \\ \quad \quad (\operatorname{entry}(b, p, k) = \operatorname{entry}(a, q, k)) \wedge \\ \quad \forall k \in \operatorname{interval}(1, n). \\ \quad \quad (\operatorname{entry}(b, q, k) = \operatorname{entry}(a, p, k)) \Rightarrow \\ \quad \quad (\operatorname{det}(r, n, b) = \operatorname{neg}(r)(\operatorname{det}(r, n, a))) \end{array}}
\vnbpart{\texttt{det-swap-cols}}
\vnbstmt{\begin{array}{@{}l@{}} \forall r, n, a, b, p, q. \\ \quad \operatorname{is\text{-}commutative\text{-}ring}(r) \wedge \\ \quad (n \in \mathbb{N}) \wedge \\ \quad (a \in \operatorname{mat}(n, n, \operatorname{carr}(r))) \wedge \\ \quad (b \in \operatorname{mat}(n, n, \operatorname{carr}(r))) \wedge \\ \quad (p \in \operatorname{interval}(1, n)) \wedge \\ \quad (q \in \operatorname{interval}(1, n)) \wedge \\ \quad \neg (p = q) \wedge \\ \quad \forall j, i \in \operatorname{interval}(1, n). \\ \quad \quad (\neg (j = p) \wedge \neg (j = q)) \Rightarrow  \\ \quad \quad \quad (\operatorname{entry}(b, i, j) = \operatorname{entry}(a, i, j)) \wedge \\ \quad \forall i \in \operatorname{interval}(1, n). \\ \quad \quad (\operatorname{entry}(b, i, p) = \operatorname{entry}(a, i, q)) \wedge \\ \quad \forall i \in \operatorname{interval}(1, n). \\ \quad \quad (\operatorname{entry}(b, i, q) = \operatorname{entry}(a, i, p)) \Rightarrow \\ \quad \quad (\operatorname{det}(r, n, b) = \operatorname{neg}(r)(\operatorname{det}(r, n, a))) \end{array}}
\vnbpart{\texttt{det-alternating-rows}}
\vnbstmt{\begin{array}{@{}l@{}} \forall r, n, a, p, q. \\ \quad \operatorname{is\text{-}commutative\text{-}ring}(r) \wedge \\ \quad (n \in \mathbb{N}) \wedge \\ \quad (a \in \operatorname{mat}(n, n, \operatorname{carr}(r))) \wedge \\ \quad (p \in \operatorname{interval}(1, n)) \wedge \\ \quad (q \in \operatorname{interval}(1, n)) \wedge \\ \quad \neg (p = q) \wedge \\ \quad \forall k \in \operatorname{interval}(1, n). \\ \quad \quad (\operatorname{entry}(a, p, k) = \operatorname{entry}(a, q, k)) \Rightarrow \\ \quad \quad (\operatorname{det}(r, n, a) = \operatorname{zero}(r)) \end{array}}
\vnbpart{\texttt{det-equal-cols-zero}}
\vnbstmt{\begin{array}{@{}l@{}} \forall r, n, a, p, q. \\ \quad \operatorname{is\text{-}commutative\text{-}ring}(r) \wedge \\ \quad (n \in \mathbb{N}) \wedge \\ \quad (a \in \operatorname{mat}(n, n, \operatorname{carr}(r))) \wedge \\ \quad (p \in \operatorname{interval}(1, n)) \wedge \\ \quad (q \in \operatorname{interval}(1, n)) \wedge \\ \quad \neg (p = q) \wedge \\ \quad \forall i \in \operatorname{interval}(1, n). \\ \quad \quad (\operatorname{entry}(a, i, p) = \operatorname{entry}(a, i, q)) \Rightarrow \\ \quad \quad (\operatorname{det}(r, n, a) = \operatorname{zero}(r)) \end{array}}
\vnbpart{\texttt{det-transpose}}
\vnbstmt{\begin{array}{@{}l@{}} \forall r, n, a. \\ \quad \operatorname{is\text{-}commutative\text{-}ring}(r) \wedge \\ \quad (n \in \mathbb{N}) \wedge \\ \quad (a \in \operatorname{mat}(n, n, \operatorname{carr}(r))) \Rightarrow \\ \quad \quad (\operatorname{det}(r, n, \operatorname{transpose}(a, n, n)) = \operatorname{det}(r, n, a)) \end{array}}
%
\end{prop}
%
\begin{proof}
%
The proofs: the file that holds each, the lemmas it cites, its oracles and its bill.
%
\vnbprf{\vnbproofpart{\texttt{det-row-linear}}}
%
\vnbprf{\textbf{Proof.} \texttt{theorem-\allowbreak{}library/det-\allowbreak{}rows.scm} --- 13 recorded steps.}
%
\vnbprf{\textbf{Lemmas.} \texttt{det-row-linear-ind}.}
%
\vnbprf{\textbf{Oracles.} \texttt{ineq}, \texttt{arith}, \texttt{crs}.  \textbf{Bill:} \texttt{proven modulo 0}.}
%
\vnbprf{\vnbproofpart{\texttt{det-col-linear}}}
%
\vnbprf{\textbf{Proof.} \texttt{theorem-\allowbreak{}library/det-\allowbreak{}transpose.scm} --- 280 recorded steps.}
%
\vnbprf{\textbf{Lemmas.} \texttt{commutative-ring-is-ring}, \texttt{dtt-col1-ind}, \texttt{interval-elt-in-nn}, \texttt{interval-lo}, \texttt{interval-hi}, \texttt{nn-one-in}, \texttt{nn-in-rr}, \texttt{nn-zero-le}, \texttt{interval-mem-intro}, \texttt{neq-sym}, \texttt{elem-f-type}, \texttt{matmul-type}, \texttt{det-in-carrier}, \texttt{det-elem-f-col}, \texttt{elem-f-action}, \texttt{dtt-ring-neg-lin}, \texttt{dtt-ring-neg-neg}.}
%
\vnbprf{\textbf{Oracles.} \texttt{ineq}, \texttt{arith}, \texttt{crs}.  \textbf{Bill:} \texttt{proven modulo 0}.}
%
\vnbprf{\vnbproofpart{\texttt{det-swap-rows}}}
%
\vnbprf{\textbf{Proof.} \texttt{theorem-\allowbreak{}library/det-\allowbreak{}rows.scm} --- 70 recorded steps.}
%
\vnbprf{\textbf{Lemmas.} \texttt{interval-elt-in-nn}, \texttt{interval-lo}, \texttt{interval-hi}, \texttt{nn-in-rr}, \texttt{nn-zero-le}, \texttt{rr-zero-in}, \texttt{rr-lt-irrefl}, \texttt{nn-nonzero-is-succ}, \texttt{dtr-swap-s}.}
%
\vnbprf{\textbf{Oracles.} \texttt{ineq}, \texttt{arith}, \texttt{crs}.  \textbf{Bill:} \texttt{proven modulo 0}.}
%
\vnbprf{\vnbproofpart{\texttt{det-swap-cols}}}
%
\vnbprf{\textbf{Proof.} \texttt{theorem-\allowbreak{}library/det-\allowbreak{}transpose.scm} --- 116 recorded steps.}
%
\vnbprf{\textbf{Lemmas.} \texttt{commutative-ring-is-ring}, \texttt{elem-f-type}, \texttt{matmul-type}, \texttt{det-elem-f-col}, \texttt{elem-f-action}, \texttt{matrix-entry-extensionality}.}
%
\vnbprf{\textbf{Oracles.} \texttt{arith}, \texttt{ineq}, \texttt{crs}.  \textbf{Bill:} \texttt{proven modulo 0}.}
%
\vnbprf{\vnbproofpart{\texttt{det-alternating-rows}}}
%
\vnbprf{\textbf{Proof.} \texttt{theorem-\allowbreak{}library/det-\allowbreak{}rows.scm} --- 10 recorded steps.}
%
\vnbprf{\textbf{Lemmas.} \texttt{det-alternating-ind}.}
%
\vnbprf{\textbf{Oracles.} \texttt{ineq}, \texttt{arith}, \texttt{crs}.  \textbf{Bill:} \texttt{proven modulo 0}.}
%
\vnbprf{\vnbproofpart{\texttt{det-equal-cols-zero}}}
%
\vnbprf{\textbf{Proof.} \texttt{theorem-\allowbreak{}library/det-\allowbreak{}transpose.scm} --- 45 recorded steps.}
%
\vnbprf{\textbf{Lemmas.} \texttt{commutative-ring-is-ring}, \texttt{ring-one-in}, \texttt{ring-neg-in-carr}, \texttt{elem-g-type}, \texttt{matmul-type}, \texttt{det-elem-g-col}, \texttt{dtt-zero-col}, \texttt{elem-g-action}, \texttt{entry-in-carrier}, \texttt{dtt-ring-x-xneg}.}
%
\vnbprf{\textbf{Oracles.} \texttt{arith}, \texttt{ineq}, \texttt{crs}.  \textbf{Bill:} \texttt{proven modulo 0}.}
%
\vnbprf{\vnbproofpart{\texttt{det-transpose}}}
%
\vnbprf{\textbf{Proof.} \texttt{theorem-\allowbreak{}library/det-\allowbreak{}transpose.scm} --- 21 recorded steps.}
%
\vnbprf{\textbf{Lemmas.} \texttt{commutative-ring-is-ring}, \texttt{identmat-type}, \texttt{dtt-tr-fun}, \texttt{dtt-tr-lin}, \texttt{dtt-tr-alt}, \texttt{det-alternating-form}, \texttt{transpose-identity}, \texttt{det-identity-ind}, \texttt{det-in-carrier}, \texttt{dtt-ring-mul-one}.}
%
\vnbprf{\textbf{Oracles.} \texttt{arith}, \texttt{ineq}, \texttt{crs}.  \textbf{Bill:} \texttt{proven modulo 0}.}
%
\end{proof}
%
\noindent\textbf{Deviation from the notes.} UPDATED 2026-09-25 (batch 32): PROVEN in full over any commutative ring: TRANSPOSE is defined entrywise (structure-library/transpose.scm, the user's decision), det-col-linear, det-swap-cols and det-equal-cols-zero are the column clauses, and det-transpose (det(A\textasciicircum{}t) = det A) follows from the uniqueness theorem det-alternating-form since X |-$>$ det(X\textasciicircum{}t) is alternating n-linear in the rows. UPDATED 2026-09-24 (batch 25): the ROW clauses are PROVEN over any commutative ring (theorem-library/det-rows.scm): det-row-linear (linear in each row), det-swap-rows (interchanging two rows negates det), and det-alternating-rows (two equal rows give 0) is now a THEOREM, its support retired; the proof from the first-row cofactor definition needs a second level of expansion (rows 1 and 2 equal: a double sum cancelling in pairs, no division by 2, so characteristic 2 is covered). The COLUMN clauses and det(A\textasciicircum{}t) = det A wait on a TRANSPOSE, which is a definition (proposed in BATCH25-REPORTS.md, not added). EARLIER NOTE: Only one of the three claims exists, and in a different form.  det-alternating-rows is ASSERTED (warrant reference) and says that two EQUAL rows force det = zero(r) -- not that interchanging two rows changes the sign.  Linearity in each row and in each column is absent.  det(A\textasciicircum{}t) = det(A) is absent and unstatable: the library has NO matrix transpose operation.  Ring hypothesis agrees (is-commutative-ring).  Note that the library's det is DEFINED by cofactor expansion along the first row, not by the notes' sum over permutations, so even the definitions are only provably equal in principle.
%
\renewcommand{\thethm}{3.19}
\begin{cor}[{Determinant under an elementary column operation}]\label{thm:det-elem-f-col}
%
\leavevmode
\vnbpart{\texttt{det-elem-f-col}}
\vnbstmt{\begin{array}{@{}l@{}} \forall r, n, a, k, l. \\ \quad \operatorname{is\text{-}commutative\text{-}ring}(r) \wedge \\ \quad (n \in \mathbb{N}) \wedge \\ \quad (a \in \operatorname{mat}(n, n, \operatorname{carr}(r))) \wedge \\ \quad (k \in \operatorname{interval}(1, n)) \wedge \\ \quad (l \in \operatorname{interval}(1, n)) \wedge \\ \quad \neg (k = l) \Rightarrow \\ \quad \quad \operatorname{det}(r, n, \operatorname{matmul}(r, a, \operatorname{elem\text{-}f}(r, n, k, l))) \\ \quad \quad =\; \operatorname{neg}(r)(\operatorname{det}(r, n, a)) \end{array}}
\vnbpart{\texttt{det-elem-g-col}}
\vnbstmt{\begin{array}{@{}l@{}} \forall r, n, a, c, k, l. \\ \quad \operatorname{is\text{-}commutative\text{-}ring}(r) \wedge \\ \quad (n \in \mathbb{N}) \wedge \\ \quad (a \in \operatorname{mat}(n, n, \operatorname{carr}(r))) \wedge \\ \quad (c \in \operatorname{carr}(r)) \wedge \\ \quad (k \in \operatorname{interval}(1, n)) \wedge \\ \quad (l \in \operatorname{interval}(1, n)) \wedge \\ \quad \neg (k = l) \Rightarrow \\ \quad \quad \operatorname{det}(r, n, \operatorname{matmul}(r, a, \operatorname{elem\text{-}g}(r, n, c, k, l))) \\ \quad \quad =\; \operatorname{det}(r, n, a) \end{array}}
\vnbpart{\texttt{det-elem-h-col}}
\vnbstmt{\begin{array}{@{}l@{}} \forall r, n, a, c, k. \\ \quad \operatorname{is\text{-}commutative\text{-}ring}(r) \wedge \\ \quad (n \in \mathbb{N}) \wedge \\ \quad (a \in \operatorname{mat}(n, n, \operatorname{carr}(r))) \wedge \\ \quad (c \in \operatorname{carr}(r)) \wedge \\ \quad (k \in \operatorname{interval}(1, n)) \Rightarrow \\ \quad \quad \operatorname{det}(r, n, \operatorname{matmul}(r, a, \operatorname{elem\text{-}h}(r, n, c, k))) \\ \quad \quad =\; \operatorname{mul}(r)(c, \operatorname{det}(r, n, a)) \end{array}}
\vnbpart{\texttt{det-mul-elem-f-col}}
\vnbstmt{\begin{array}{@{}l@{}} \forall r, n, a, k, l. \\ \quad \operatorname{is\text{-}commutative\text{-}ring}(r) \wedge \\ \quad (n \in \mathbb{N}) \wedge \\ \quad (a \in \operatorname{mat}(n, n, \operatorname{carr}(r))) \wedge \\ \quad (k \in \operatorname{interval}(1, n)) \wedge \\ \quad (l \in \operatorname{interval}(1, n)) \wedge \\ \quad \neg (k = l) \Rightarrow \\ \quad \quad \operatorname{det}(r, n, \operatorname{matmul}(r, a, \operatorname{elem\text{-}f}(r, n, k, l))) \\ \quad \quad =\; \operatorname{mul}(r)(\operatorname{det}(r, n, a), \operatorname{det}(r, n, \operatorname{elem\text{-}f}(r, n, k, l))) \end{array}}
\vnbpart{\texttt{det-mul-elem-g-col}}
\vnbstmt{\begin{array}{@{}l@{}} \forall r, n, a, c, k, l. \\ \quad \operatorname{is\text{-}commutative\text{-}ring}(r) \wedge \\ \quad (n \in \mathbb{N}) \wedge \\ \quad (a \in \operatorname{mat}(n, n, \operatorname{carr}(r))) \wedge \\ \quad (c \in \operatorname{carr}(r)) \wedge \\ \quad (k \in \operatorname{interval}(1, n)) \wedge \\ \quad (l \in \operatorname{interval}(1, n)) \wedge \\ \quad \neg (k = l) \Rightarrow \\ \quad \quad \operatorname{det}(r, n, \operatorname{matmul}(r, a, \operatorname{elem\text{-}g}(r, n, c, k, l))) \\ \quad \quad =\; \operatorname{mul}(r)(\operatorname{det}(r, n, a), \operatorname{det}(r, n, \operatorname{elem\text{-}g}(r, n, c, k, l))) \end{array}}
\vnbpart{\texttt{det-mul-elem-h-col}}
\vnbstmt{\begin{array}{@{}l@{}} \forall r, n, a, c, k. \\ \quad \operatorname{is\text{-}commutative\text{-}ring}(r) \wedge \\ \quad (n \in \mathbb{N}) \wedge \\ \quad (a \in \operatorname{mat}(n, n, \operatorname{carr}(r))) \wedge \\ \quad (c \in \operatorname{carr}(r)) \wedge \\ \quad (k \in \operatorname{interval}(1, n)) \Rightarrow \\ \quad \quad \operatorname{det}(r, n, \operatorname{matmul}(r, a, \operatorname{elem\text{-}h}(r, n, c, k))) \\ \quad \quad =\; \operatorname{mul}(r)(\operatorname{det}(r, n, a), \operatorname{det}(r, n, \operatorname{elem\text{-}h}(r, n, c, k))) \end{array}}
%
\end{cor}
%
\begin{proof}
%
The proofs: the file that holds each, the lemmas it cites, its oracles and its bill.
%
\vnbprf{\vnbproofpart{\texttt{det-elem-f-col}}}
%
\vnbprf{\textbf{Proof.} \texttt{theorem-\allowbreak{}library/det-\allowbreak{}transpose.scm} --- 15 recorded steps.}
%
\vnbprf{\textbf{Lemmas.} \texttt{commutative-ring-is-ring}, \texttt{det-in-carrier}, \texttt{det-mul-elem-f-col}, \texttt{det-elem-f}, \texttt{dtt-ring-mul-neg1}.}
%
\vnbprf{\textbf{Oracles.} \texttt{ineq}, \texttt{arith}, \texttt{crs}.  \textbf{Bill:} \texttt{proven modulo 0}.}
%
\vnbprf{\vnbproofpart{\texttt{det-elem-g-col}}}
%
\vnbprf{\textbf{Proof.} \texttt{theorem-\allowbreak{}library/det-\allowbreak{}transpose.scm} --- 16 recorded steps.}
%
\vnbprf{\textbf{Lemmas.} \texttt{commutative-ring-is-ring}, \texttt{det-in-carrier}, \texttt{det-mul-elem-g-col}, \texttt{det-elem-g}, \texttt{dtt-ring-mul-one}.}
%
\vnbprf{\textbf{Oracles.} \texttt{ineq}, \texttt{arith}, \texttt{crs}.  \textbf{Bill:} \texttt{proven modulo 0}.}
%
\vnbprf{\vnbproofpart{\texttt{det-elem-h-col}}}
%
\vnbprf{\textbf{Proof.} \texttt{theorem-\allowbreak{}library/det-\allowbreak{}transpose.scm} --- 14 recorded steps.}
%
\vnbprf{\textbf{Lemmas.} \texttt{commutative-ring-is-ring}, \texttt{det-in-carrier}, \texttt{det-mul-elem-h-col}, \texttt{det-elem-h}, \texttt{commutative-ring-mul-comm}.}
%
\vnbprf{\textbf{Oracles.} \texttt{ineq}, \texttt{arith}, \texttt{crs}.  \textbf{Bill:} \texttt{proven modulo 0}.}
%
\vnbprf{\vnbproofpart{\texttt{det-mul-elem-f-col}}}
%
\vnbprf{\textbf{Proof.} \texttt{theorem-\allowbreak{}library/det-\allowbreak{}transpose.scm} --- 11 recorded steps.}
%
\vnbprf{\textbf{Lemmas.} \texttt{commutative-ring-is-ring}, \texttt{elem-f-type}, \texttt{det-multiplicative}.}
%
\vnbprf{\textbf{Oracles.} \texttt{arith}, \texttt{ineq}, \texttt{crs}.  \textbf{Bill:} \texttt{proven modulo 0}.}
%
\vnbprf{\vnbproofpart{\texttt{det-mul-elem-g-col}}}
%
\vnbprf{\textbf{Proof.} \texttt{theorem-\allowbreak{}library/det-\allowbreak{}transpose.scm} --- 12 recorded steps.}
%
\vnbprf{\textbf{Lemmas.} \texttt{commutative-ring-is-ring}, \texttt{elem-g-type}, \texttt{det-multiplicative}.}
%
\vnbprf{\textbf{Oracles.} \texttt{arith}, \texttt{ineq}, \texttt{crs}.  \textbf{Bill:} \texttt{proven modulo 0}.}
%
\vnbprf{\vnbproofpart{\texttt{det-mul-elem-h-col}}}
%
\vnbprf{\textbf{Proof.} \texttt{theorem-\allowbreak{}library/det-\allowbreak{}transpose.scm} --- 10 recorded steps.}
%
\vnbprf{\textbf{Lemmas.} \texttt{commutative-ring-is-ring}, \texttt{elem-h-type}, \texttt{det-multiplicative}.}
%
\vnbprf{\textbf{Oracles.} \texttt{arith}, \texttt{ineq}, \texttt{crs}.  \textbf{Bill:} \texttt{proven modulo 0}.}
%
\end{proof}
%
\noindent\textbf{Deviation from the notes.} UPDATED 2026-09-25 (batch 32): PROVEN (theorem-library/det-transpose.scm), each in one line from det-multiplicative and the determinants of the elementary matrices; the H clause is stated for every scalar r of the ring, not only a unit. The determinant of an elementary matrix is computed nowhere in the library -- there is no det of elem-f, elem-g or elem-h, and no det of a diagonal matrix from which det H could be read.  The asserted det-multiplicative would reduce these to those three values, which are missing.
%
\renewcommand{\thethm}{3.20}
\begin{cor}[{det(AV) = det A det V for invertible V over a field}]\label{thm:det-mul-invertible}
%
\leavevmode
\vnbpart{\texttt{det-mul-invertible}}
\vnbstmt{\begin{array}{@{}l@{}} \forall r, n, a, v. \\ \quad \operatorname{is\text{-}field\text{-}ring}(r) \wedge \\ \quad (n \in \mathbb{N}) \wedge \\ \quad (a \in \operatorname{mat}(n, n, \operatorname{carr}(r))) \wedge \\ \quad \operatorname{is\text{-}invertible\text{-}mat}(r, n, v) \Rightarrow \\ \quad \quad \operatorname{det}(r, n, \operatorname{matmul}(r, a, v)) \\ \quad \quad =\; \operatorname{mul}(r)(\operatorname{det}(r, n, a), \operatorname{det}(r, n, v)) \end{array}}
\vnbpart{\texttt{det-multiplicative}}
\vnbstmt{\begin{array}{@{}l@{}} \forall r, n, p, q. \\ \quad \operatorname{is\text{-}commutative\text{-}ring}(r) \wedge \\ \quad (n \in \mathbb{N}) \wedge \\ \quad (p \in \operatorname{mat}(n, n, \operatorname{carr}(r))) \wedge \\ \quad (q \in \operatorname{mat}(n, n, \operatorname{carr}(r))) \Rightarrow \\ \quad \quad \operatorname{det}(r, n, \operatorname{matmul}(r, p, q)) \\ \quad \quad =\; \operatorname{mul}(r)(\operatorname{det}(r, n, p), \operatorname{det}(r, n, q)) \end{array}}
%
\end{cor}
%
\begin{proof}
%
The proofs: the file that holds each, the lemmas it cites, its oracles and its bill.
%
\vnbprf{\vnbproofpart{\texttt{det-mul-invertible}}}
%
\vnbprf{\textbf{Proof.} \texttt{theorem-\allowbreak{}library/det-\allowbreak{}multiplicative.scm} --- 15 recorded steps.}
%
\vnbprf{\textbf{Lemmas.} \texttt{invertible-mat-is-mat}, \texttt{det-multiplicative}.}
%
\vnbprf{\textbf{Oracles.} \texttt{arith}, \texttt{ineq}, \texttt{crs}.  \textbf{Bill:} \texttt{proven modulo 0}.}
%
\vnbprf{\vnbproofpart{\texttt{det-multiplicative}}}
%
\vnbprf{\textbf{Proof.} \texttt{theorem-\allowbreak{}library/det-\allowbreak{}multiplicative.scm} --- 21 recorded steps.}
%
\vnbprf{\textbf{Lemmas.} \texttt{commutative-ring-is-ring}, \texttt{identmat-type}, \texttt{daf-det-rm-fun}, \texttt{daf-det-rm-lin}, \texttt{daf-det-rm-alt}, \texttt{det-alternating-form}, \texttt{identmat-left-identity}, \texttt{det-in-carrier}, \texttt{ring-carrier-closed-mul}.}
%
\vnbprf{\textbf{Oracles.} \texttt{arith}, \texttt{ineq}, \texttt{crs}.  \textbf{Bill:} \texttt{proven modulo 0}.}
%
\end{proof}
%
\noindent\textbf{Deviation from the notes.} UPDATED 2026-09-24 (batch 26-B): PROVEN as a corollary of 3.27 (det-mul-invertible: over a field, V invertible, det(AV) = det A det V). EARLIER NOTE: UPDATED 2026-09-24 (batch 25): det(EA) = det E det A is PROVEN for each elementary E (theorem-library/det-mul.scm), but the tree has no theorem that an invertible matrix over a field is a product of elementary matrices (and no finite product of matrices to state it with), so this corollary and det-multiplicative remain asserted. The one lemma that proves 3.27 for EVERY commutative ring, and 3.20 with it, is Hoffman-Kunze 5.3 Theorem 2: an alternating n-linear D satisfies D(A) = det(A) D(I); with row linearity and alternation now proven it is about one file. EARLIER NOTE: The library asserts the UNRESTRICTED multiplicativity -- any commutative ring, any two square matrices of the same size, no invertibility and no field -- which subsumes this corollary.  It carries warrant reference, i.e. it is trusted, not proven.
%
\renewcommand{\thethm}{3.22}
\begin{prop}[{Expansion by minors along a row and along a column}]\label{thm:det-expand-row}
%
\leavevmode
\vnbpart{\texttt{det-expand-row}}
\vnbstmt{\begin{array}{@{}l@{}} \forall r, n, a, i. \\ \quad \operatorname{is\text{-}commutative\text{-}ring}(r) \wedge \\ \quad (n \in \mathbb{N}) \wedge \\ \quad (a \in \operatorname{mat}((n + 1), (n + 1), \operatorname{carr}(r))) \wedge \\ \quad (i \in \operatorname{interval}(1, (n + 1))) \Rightarrow \\ \quad \quad \operatorname{det}(r, (n + 1), a) \\ \quad \quad =\; \operatorname{finsum}(\operatorname{ring\text{-}additive\text{-}ag}(r), \\ \quad \quad \quad (\lambda j \in \operatorname{interval}(1, (n + 1)).\; \\ \quad \quad \quad \quad \operatorname{mul}(r)(\operatorname{mpow}(\operatorname{ring\text{-}multiplicative\text{-}monoid}(r), \operatorname{neg}(r)(\operatorname{one}(r)), (i + j)), \\ \quad \quad \quad \quad \quad \operatorname{mul}(r)(\operatorname{entry}(a, i, j), \operatorname{det}(r, n, \operatorname{minor}(a, i, j, n))))), \\ \quad \quad \quad \operatorname{interval}(1, (n + 1))) \end{array}}
\vnbpart{\texttt{det-expand-col}}
\vnbstmt{\begin{array}{@{}l@{}} \forall r, n, a, j. \\ \quad \operatorname{is\text{-}commutative\text{-}ring}(r) \wedge \\ \quad (n \in \mathbb{N}) \wedge \\ \quad (a \in \operatorname{mat}((n + 1), (n + 1), \operatorname{carr}(r))) \wedge \\ \quad (j \in \operatorname{interval}(1, (n + 1))) \Rightarrow \\ \quad \quad \operatorname{det}(r, (n + 1), a) \\ \quad \quad =\; \operatorname{finsum}(\operatorname{ring\text{-}additive\text{-}ag}(r), \\ \quad \quad \quad (\lambda e \in \operatorname{interval}(1, (n + 1)).\; \\ \quad \quad \quad \quad \operatorname{mul}(r)(\operatorname{mpow}(\operatorname{ring\text{-}multiplicative\text{-}monoid}(r), \operatorname{neg}(r)(\operatorname{one}(r)), (e + j)), \\ \quad \quad \quad \quad \quad \operatorname{mul}(r)(\operatorname{entry}(a, e, j), \operatorname{det}(r, n, \operatorname{minor}(a, e, j, n))))), \\ \quad \quad \quad \operatorname{interval}(1, (n + 1))) \end{array}}
\vnbpart{\texttt{det-cofactor}}
\vnbstmt{\begin{array}{@{}l@{}} \forall r, n, a. \\ \quad \operatorname{is\text{-}ring}(r) \wedge \\ \quad (n \in \mathbb{N}) \wedge \\ \quad (a \in \operatorname{mat}((n + 1), (n + 1), \operatorname{carr}(r))) \Rightarrow \\ \quad \quad (\operatorname{det}(r, (n + 1), a) \simeq \operatorname{finsum}(\operatorname{ring\text{-}additive\text{-}ag}(r), (\lambda j \in \operatorname{interval}(1, (n + 1)).\; \operatorname{mul}(r)(\operatorname{mpow}(\operatorname{ring\text{-}multiplicative\text{-}monoid}(r), \operatorname{neg}(r)(\operatorname{one}(r)), (j + 1)), \operatorname{mul}(r)(\operatorname{entry}(a, 1, j), \operatorname{det}(r, n, \operatorname{minor}(a, 1, j, n))))), \operatorname{interval}(1, (n + 1)))) \end{array}}
%
\end{prop}
%
\begin{proof}
%
The proofs: the file that holds each, the lemmas it cites, its oracles and its bill.
%
\vnbprf{\vnbproofpart{\texttt{det-expand-row}}}
%
\vnbprf{\textbf{Proof.} \texttt{theorem-\allowbreak{}library/det-\allowbreak{}rows.scm} --- 10 recorded steps.}
%
\vnbprf{\textbf{Lemmas.} \texttt{interval-elt-in-nn}, \texttt{interval-lo}, \texttt{interval-hi}, \texttt{det-expand-row-ind}.}
%
\vnbprf{\textbf{Oracles.} \texttt{arith}, \texttt{ineq}, \texttt{crs}.  \textbf{Bill:} \texttt{proven modulo 0}.}
%
\vnbprf{\vnbproofpart{\texttt{det-expand-col}}}
%
\vnbprf{\textbf{Proof.} \texttt{theorem-\allowbreak{}library/det-\allowbreak{}transpose.scm} --- 102 recorded steps.}
%
\vnbprf{\textbf{Lemmas.} \texttt{commutative-ring-is-ring}, \texttt{nn-succ-closed}, \texttt{nn-one-in}, \texttt{transpose-in-mat}, \texttt{det-transpose}, \texttt{det-expand-row}, \texttt{ras-carr}, \texttt{interval-in-set}, \texttt{interval-card-in-nn}, \texttt{interval-elt-in-nn}, \texttt{interval-lo}, \texttt{interval-hi}, \texttt{finsum-congruence}, \texttt{nn-add-closed}, \texttt{dtr-sign-type}, \texttt{entry-in-carrier}, \texttt{minor-type}, \texttt{det-in-carrier}, \texttt{ring-carrier-closed-mul}, \texttt{transpose-entry}, \texttt{dtt-minor-transpose}, \texttt{nn-add-comm}.}
%
\vnbprf{\textbf{Oracles.} \texttt{arith}, \texttt{ineq}, \texttt{crs}.  \textbf{Bill:} \texttt{proven modulo 0}.}
%
\vnbprf{\vnbproofpart{\texttt{det-cofactor}}}
%
\vnbprf{Definitional; no proof in the library.}
%
\end{proof}
%
\noindent\textbf{Deviation from the notes.} UPDATED 2026-09-25 (batch 32): the COLUMN half is PROVEN, det-expand-col (for each j, det A = sum\_i (-1)\textasciicircum{}(i+j) a\_ij det A[i,j]), by transposing det-expand-row with (minor of A\textasciicircum{}t at (j,i)) = (minor of A at (i,j))\textasciicircum{}t. UPDATED 2026-09-24 (batch 25): the ROW half is PROVEN: det-expand-row expands along any row i with the sign (-1)\textasciicircum{}(i+j). The column half waits on the transpose. EARLIER NOTE: The library DEFINES det by the cofactor expansion along the FIRST ROW: det-cofactor is the defining equation for det(r, succ(n), a), not a theorem.  So the notes' statement is available only at i = 1, and only as a definition; the expansion along an arbitrary row i, and the column expansion, are nowhere stated.  The library's minor(s,r,c,n) is the notes' A[i,j].  The sign (-1)\textasciicircum{}\{i+j\} appears in det-cofactor as mpow of neg(one(r)) to succ(j), the i = 1 specialisation.
%
\renewcommand{\thethm}{3.23}
\begin{cor}[{The determinant of a triangular matrix is the product of its diagonal entries}]\label{thm:linear-algebra-3.23}
%
\textit{If A is lower triangular then det A is the product of the diagonal entries; likewise for upper triangular A.}
%
Not in the library.
%
\end{cor}
%
\begin{proof}
%
The proofs: the file that holds each, the lemmas it cites, its oracles and its bill.
%
\vnbprf{There is no library statement, and so no proof.}
%
\end{proof}
%
\noindent\textbf{Deviation from the notes.} The library defines is-diagonal (all off-diagonal entries zero) and border/submat, but computes no determinant from any of them, and has no notion of triangular matrix at all.
%
\renewcommand{\thethm}{3.25}
\begin{prop}[{The formal adjoint satisfies adj(A) A = det(A) I}]\label{thm:linear-algebra-3.25}
%
\textit{For any n x n matrix A, the formal adjoint with entries (-1)\textasciicircum{}\{i+j\} det(A[j,i]) satisfies adj(A) A = det(A) I; hence A\textasciicircum{}\{-1\} = det(A)\textasciicircum{}\{-1\} adj(A).}
%
Not in the library.
%
\end{prop}
%
\begin{proof}
%
The proofs: the file that holds each, the lemmas it cites, its oracles and its bill.
%
\vnbprf{There is no library statement, and so no proof.}
%
\end{proof}
%
\noindent\textbf{Deviation from the notes.} No adjoint (adjugate) operator, no cofactor matrix, no Cramer rule.  minor and det exist separately; nothing assembles the cofactors into a matrix.
%
\renewcommand{\thethm}{3.26}
\begin{cor}[{Over a field, invertible iff the determinant is nonzero}]\label{thm:linear-algebra-3.26}
%
\textit{If R is a field, A in M\_n(R) is invertible if and only if det A /= 0.}
%
Not in the library.
%
\end{cor}
%
\begin{proof}
%
The proofs: the file that holds each, the lemmas it cites, its oracles and its bill.
%
\vnbprf{There is no library statement, and so no proof.}
%
\end{proof}
%
\noindent\textbf{Deviation from the notes.} Nothing in the library connects is-invertible-mat with det in either direction.
%
\renewcommand{\thethm}{3.27}
\begin{prop}[{det is multiplicative}]\label{thm:det-multiplicative}
%
\leavevmode
\vnbpart{\texttt{det-multiplicative}}
\vnbstmt{\begin{array}{@{}l@{}} \forall r, n, p, q. \\ \quad \operatorname{is\text{-}commutative\text{-}ring}(r) \wedge \\ \quad (n \in \mathbb{N}) \wedge \\ \quad (p \in \operatorname{mat}(n, n, \operatorname{carr}(r))) \wedge \\ \quad (q \in \operatorname{mat}(n, n, \operatorname{carr}(r))) \Rightarrow \\ \quad \quad \operatorname{det}(r, n, \operatorname{matmul}(r, p, q)) \\ \quad \quad =\; \operatorname{mul}(r)(\operatorname{det}(r, n, p), \operatorname{det}(r, n, q)) \end{array}}
\vnbpart{\texttt{det-alternating-form}}
\vnbstmt{\begin{array}{@{}l@{}} \forall r, n, d, a. \\ \quad \operatorname{is\text{-}commutative\text{-}ring}(r) \wedge \\ \quad (n \in \mathbb{N}) \wedge \\ \quad (d \in (\operatorname{mat}(n, n, \operatorname{carr}(r)) \to \operatorname{carr}(r))) \wedge \\ \quad \forall \mathit{dfa}, b, c, p, s. \\ \quad \quad (\mathit{dfa} \in \operatorname{mat}(n, n, \operatorname{carr}(r))) \wedge \\ \quad \quad (b \in \operatorname{mat}(n, n, \operatorname{carr}(r))) \wedge \\ \quad \quad (c \in \operatorname{mat}(n, n, \operatorname{carr}(r))) \wedge \\ \quad \quad (p \in \operatorname{interval}(1, n)) \wedge \\ \quad \quad (s \in \operatorname{carr}(r)) \wedge \\ \quad \quad \forall i, k \in \operatorname{interval}(1, n). \\ \quad \quad \quad \neg (i = p) \Rightarrow  \\ \quad \quad \quad \quad (\operatorname{entry}(b, i, k) = \operatorname{entry}(\mathit{dfa}, i, k)) \wedge \\ \quad \quad \forall i, k \in \operatorname{interval}(1, n). \\ \quad \quad \quad \neg (i = p) \Rightarrow  \\ \quad \quad \quad \quad (\operatorname{entry}(c, i, k) = \operatorname{entry}(\mathit{dfa}, i, k)) \wedge \\ \quad \quad \forall k \in \operatorname{interval}(1, n). \\ \quad \quad \quad \operatorname{entry}(c, p, k) \\ \quad \quad \quad =\; \operatorname{add}(r)(\operatorname{mul}(r)(s, \operatorname{entry}(\mathit{dfa}, p, k)), \operatorname{entry}(b, p, k)) \Rightarrow \\ \quad \quad \quad (d(c) = \operatorname{add}(r)(\operatorname{mul}(r)(s, d(\mathit{dfa})), d(b))) \wedge \\ \quad \forall \mathit{dfa}, p, q. \\ \quad \quad (\mathit{dfa} \in \operatorname{mat}(n, n, \operatorname{carr}(r))) \wedge \\ \quad \quad (p \in \operatorname{interval}(1, n)) \wedge \\ \quad \quad (q \in \operatorname{interval}(1, n)) \wedge \\ \quad \quad \neg (p = q) \wedge \\ \quad \quad \forall k \in \operatorname{interval}(1, n). \\ \quad \quad \quad (\operatorname{entry}(\mathit{dfa}, p, k) = \operatorname{entry}(\mathit{dfa}, q, k)) \Rightarrow \\ \quad \quad \quad (d(\mathit{dfa}) = \operatorname{zero}(r)) \wedge \\ \quad (a \in \operatorname{mat}(n, n, \operatorname{carr}(r))) \Rightarrow \\ \quad \quad (d(a) = \operatorname{mul}(r)(\operatorname{det}(r, n, a), d(\operatorname{identmat}(r, n)))) \end{array}}
%
\end{prop}
%
\begin{proof}
%
The proofs: the file that holds each, the lemmas it cites, its oracles and its bill.
%
\vnbprf{\vnbproofpart{\texttt{det-multiplicative}}}
%
\vnbprf{\textbf{Proof.} \texttt{theorem-\allowbreak{}library/det-\allowbreak{}multiplicative.scm} --- 21 recorded steps.}
%
\vnbprf{\textbf{Lemmas.} \texttt{commutative-ring-is-ring}, \texttt{identmat-type}, \texttt{daf-det-rm-fun}, \texttt{daf-det-rm-lin}, \texttt{daf-det-rm-alt}, \texttt{det-alternating-form}, \texttt{identmat-left-identity}, \texttt{det-in-carrier}, \texttt{ring-carrier-closed-mul}.}
%
\vnbprf{\textbf{Oracles.} \texttt{arith}, \texttt{ineq}, \texttt{crs}.  \textbf{Bill:} \texttt{proven modulo 0}.}
%
\vnbprf{\vnbproofpart{\texttt{det-alternating-form}}}
%
\vnbprf{\textbf{Proof.} \texttt{theorem-\allowbreak{}library/det-\allowbreak{}alternating-\allowbreak{}form.scm} --- 9 recorded steps.}
%
\vnbprf{\textbf{Lemmas.} \texttt{det-alternating-form-ind}.}
%
\vnbprf{\textbf{Oracles.} \texttt{ineq}, \texttt{arith}, \texttt{crs}.  \textbf{Bill:} \texttt{proven modulo 0}.}
%
\end{proof}
%
\noindent\textbf{Deviation from the notes.} UPDATED 2026-09-24 (batch 26-B): PROVEN for every commutative ring, theorem-library/det-multiplicative.scm: X |-$>$ det(XQ) is linear in each row and zero on two equal rows, and Hoffman-Kunze 5.3 Theorem 2 (det-alternating-form: such a D is det(A) D(I)) is proven by induction on n without permutations or signs; the support is retired. EARLIER NOTE: Same statement, same hypothesis (is-commutative-ring, both matrices n x n), but ASSERTED with warrant reference -- it is trusted surface, never proven.  It is one of only two determinant facts on the debt bundle (with det-alternating-rows).
%
\renewcommand{\thethm}{3.29}
\begin{prop}[{Behaviour of the sigma-th order subdeterminants under F and G operations}]\label{thm:linear-algebra-3.29}
%
\textit{If B comes from A by an F row or column operation then sdet\_sigma B = sdet\_sigma A up to sign; if by a G operation then every element of sdet\_sigma B has the form a + s a' with a, a' in sdet\_sigma A and s = 0 or s = +/- r.}
%
Not in the library.
%
\end{prop}
%
\begin{proof}
%
The proofs: the file that holds each, the lemmas it cites, its oracles and its bill.
%
\vnbprf{There is no library statement, and so no proof.}
%
\end{proof}
%
\noindent\textbf{Deviation from the notes.} The library has no submatrix obtained by deleting SEVERAL rows and columns (minor deletes exactly one of each; submat and block cut a leading or trailing rectangle), no set minor\_sigma A, and no sdet\_sigma A.  Nothing about how determinants of such minors move under elementary operations.
%
\renewcommand{\thethm}{3.31}
\begin{cor}[{The gcd of the sigma-th subdeterminants is an invariant over Z}]\label{thm:linear-algebra-3.31}
%
\textit{If A, B in M\_\{m x n\}(Z) and A -$>$ B by elementary F, G row and column operations over Z, then gcd sdet\_sigma A = +/- gcd sdet\_sigma B.}
%
Not in the library.
%
\end{cor}
%
\begin{proof}
%
The proofs: the file that holds each, the lemmas it cites, its oracles and its bill.
%
\vnbprf{There is no library statement, and so no proof.}
%
\end{proof}
%
\noindent\textbf{Deviation from the notes.} Depends on 3.29, which is absent.  The library has zz-is-gcd and zz-bezout for a PAIR of integers, but no gcd of a set of integers and no subdeterminants.
%
\renewcommand{\thethm}{3.36}
\begin{prop}[{Integer normal form: existence and uniqueness of the diagonal form}]\label{thm:smith-normal-form}
%
\leavevmode
\vnbpart{\texttt{smith-normal-form}}
\vnbstmt{\begin{array}{@{}l@{}} \forall a.; \forall k \in \mathbb{N}, n, p. \\ \quad \operatorname{is\text{-}euclidean\text{-}ring}(a) \wedge \\ \quad (n \in \mathbb{N}) \wedge \\ \quad (p \in \operatorname{mat}(k, n, \operatorname{carr}(a))) \Rightarrow \\ \quad \quad \exists d.; \exists b \in \mathbb{N}. \\ \quad \quad \quad \operatorname{mat\text{-}equiv}(a, k, n, p, d) \wedge \\ \quad \quad \quad \operatorname{smith\text{-}staircase}(a, k, n, d, b) \end{array}}
\vnbpart{\texttt{smith-diagonalization}}
\vnbstmt{\begin{array}{@{}l@{}} \forall a.; \forall k \in \mathbb{N}, n, p. \\ \quad \operatorname{is\text{-}euclidean\text{-}ring}(a) \wedge \\ \quad (n \in \mathbb{N}) \wedge \\ \quad (p \in \operatorname{mat}(k, n, \operatorname{carr}(a))) \Rightarrow \\ \quad \quad \exists d. \\ \quad \quad \quad \operatorname{mat\text{-}equiv}(a, k, n, p, d) \wedge \\ \quad \quad \quad \operatorname{is\text{-}diagonal}(a, k, n, d) \end{array}}
%
\end{prop}
%
\begin{proof}
%
The proofs: the file that holds each, the lemmas it cites, its oracles and its bill.
%
\vnbprf{\vnbproofpart{\texttt{smith-normal-form}}}
%
\vnbprf{\textbf{Proof.} \texttt{theorem-\allowbreak{}library/smith-\allowbreak{}staircase-\allowbreak{}proof.scm} --- 317 recorded steps.}
%
\vnbprf{\textbf{Lemmas.} \texttt{euclidean-ring-is-integral-domain}, \texttt{integral-domain-is-commutative-ring}, \texttt{commutative-ring-is-ring}, \texttt{nn-zero-in}, \texttt{mat-equiv-refl}, \texttt{nn-zero-le}, \texttt{interval-lo}, \texttt{interval-hi}, \texttt{interval-elt-in-nn}, \texttt{nn-not-le-zero-pos}, \texttt{nn-succ-closed}, \texttt{nn-one-in}, \texttt{nn-le-trans-guarded}, \texttt{nn-pos-is-succ}, \texttt{eq-sym}, \texttt{clear-pivot-cross}, \texttt{mat-equiv-target-is-mat}, \texttt{submat-type}, \texttt{one-in-interval}, \texttt{entry-in-carrier}, \texttt{bordering}, \texttt{border-type}, \texttt{mat-equiv-trans}, \texttt{border-staircase}.}
%
\vnbprf{\textbf{Oracles.} \texttt{arith}, \texttt{ineq}, \texttt{crs}.  \textbf{Bill:} \texttt{proven modulo 0}.}
%
\vnbprf{\vnbproofpart{\texttt{smith-diagonalization}}}
%
\vnbprf{\textbf{Proof.} \texttt{theorem-\allowbreak{}library/smith-\allowbreak{}diagonalization-\allowbreak{}proof.scm} --- 162 recorded steps.}
%
\vnbprf{\textbf{Lemmas.} \texttt{euclidean-ring-is-integral-domain}, \texttt{integral-domain-is-commutative-ring}, \texttt{commutative-ring-is-ring}, \texttt{mat-equiv-refl}, \texttt{interval-lo}, \texttt{interval-hi}, \texttt{interval-elt-in-nn}, \texttt{nn-not-le-zero-pos}, \texttt{nn-one-in}, \texttt{nn-le-trans-guarded}, \texttt{nn-pos-is-succ}, \texttt{eq-sym}, \texttt{clear-pivot-cross}, \texttt{mat-equiv-target-is-mat}, \texttt{submat-type}, \texttt{one-in-interval}, \texttt{entry-in-carrier}, \texttt{bordering}, \texttt{border-is-diagonal}, \texttt{border-type}, \texttt{mat-equiv-trans}.}
%
\vnbprf{\textbf{Oracles.} \texttt{arith}, \texttt{ineq}, \texttt{crs}.  \textbf{Bill:} \texttt{proven modulo 0}.}
%
\end{proof}
%
\noindent\textbf{Deviation from the notes.} The heaviest gap in the chapter.  smith-normal-form gives, over a EUCLIDEAN RING (Z qualifies, zz-is-euclidean-ring), the EXISTENCE of a mat-equiv d that is is-diagonal and whose nonzero diagonal entries form an initial segment of length kk (smith-staircase).  It does NOT give the divisibility chain b\_ii divides b\_jj, does NOT give non-negativity of the entries, and does NOT give UNIQUENESS -- the notes' whole second half.  Moreover the notes' relation \textasciitilde{} is generated by integral elementary operations, whereas mat-equiv is d = U c V for SOME invertible U, V; the two are never shown to coincide.
%
\renewcommand{\thethm}{3.37}
\begin{cor}[{The reducing matrices U and V are products of elementary matrices}]\label{thm:linear-algebra-3.37}
%
\vnbstmt{\begin{array}{@{}l@{}} \forall a.; \forall k \in \mathbb{N}, n, p. \\ \quad \operatorname{is\text{-}euclidean\text{-}ring}(a) \wedge \\ \quad (n \in \mathbb{N}) \wedge \\ \quad (p \in \operatorname{mat}(k, n, \operatorname{carr}(a))) \Rightarrow \\ \quad \quad \exists d.; \exists b \in \mathbb{N}. \\ \quad \quad \quad \operatorname{mat\text{-}equiv}(a, k, n, p, d) \wedge \\ \quad \quad \quad \operatorname{smith\text{-}staircase}(a, k, n, d, b) \end{array}}
%
\end{cor}
%
\begin{proof}
%
The proofs: the file that holds each, the lemmas it cites, its oracles and its bill.
%
\vnbprf{\textbf{Proof.} \texttt{theorem-\allowbreak{}library/smith-\allowbreak{}staircase-\allowbreak{}proof.scm} --- 317 recorded steps.}
%
\vnbprf{\textbf{Lemmas.} \texttt{euclidean-ring-is-integral-domain}, \texttt{integral-domain-is-commutative-ring}, \texttt{commutative-ring-is-ring}, \texttt{nn-zero-in}, \texttt{mat-equiv-refl}, \texttt{nn-zero-le}, \texttt{interval-lo}, \texttt{interval-hi}, \texttt{interval-elt-in-nn}, \texttt{nn-not-le-zero-pos}, \texttt{nn-succ-closed}, \texttt{nn-one-in}, \texttt{nn-le-trans-guarded}, \texttt{nn-pos-is-succ}, \texttt{eq-sym}, \texttt{clear-pivot-cross}, \texttt{mat-equiv-target-is-mat}, \texttt{submat-type}, \texttt{one-in-interval}, \texttt{entry-in-carrier}, \texttt{bordering}, \texttt{border-type}, \texttt{mat-equiv-trans}, \texttt{border-staircase}.}
%
\vnbprf{\textbf{Oracles.} \texttt{arith}, \texttt{ineq}, \texttt{crs}.  \textbf{Bill:} \texttt{proven modulo 0}.}
%
\end{proof}
%
\noindent\textbf{Deviation from the notes.} Unfolding mat-equiv inside smith-normal-form does produce invertible U (m x m) and V (n x n) with d = U A V, so part (1) holds as far as 3.36 goes.  Part (2) -- that U and V are PRODUCTS OF ELEMENTARY matrices -- is never established: mat-equiv is defined by existence of invertible factors, not by a sequence of elementary operations, and no theorem bridges the two.  The normal form reached also lacks 3.36's divisibility and non-negativity.
%
\renewcommand{\thethm}{3.38}
\begin{cor}[{An invertible integer matrix is a product of elementary matrices}]\label{thm:linear-algebra-3.38}
%
\textit{U in M\_m(Z) is invertible over Z if and only if it is a product U V with U a product of column elementary matrices and V a product of row elementary matrices.}
%
Not in the library.
%
\end{cor}
%
\begin{proof}
%
The proofs: the file that holds each, the lemmas it cites, its oracles and its bill.
%
\vnbprf{There is no library statement, and so no proof.}
%
\end{proof}
%
\noindent\textbf{Deviation from the notes.} No factorisation of an invertible matrix into elementary matrices anywhere in the library, in either direction.  The easy direction's ingredients exist (elem-f-invertible, elem-g-invertible, product-of-invertibles-is-invertible) but are never assembled.
%
\renewcommand{\thethm}{3.40}
\begin{lem}[{An invertible change of sequence preserves generation and freeness}]\label{thm:generates-transport}
%
\leavevmode
\vnbpart{\texttt{generates-transport}}
\vnbstmt{\begin{array}{@{}l@{}} \forall d, n, u, m. \\ \quad \operatorname{is\text{-}module}(d) \wedge \\ \quad (u \in \operatorname{mat}(n, 1, \operatorname{vec}(d))) \wedge \\ \quad \operatorname{is\text{-}invertible\text{-}mat}(\operatorname{scal}(d), n, m) \wedge \\ \quad \operatorname{generates}(d, n, u) \Rightarrow \\ \quad \quad \operatorname{generates}(d, n, \operatorname{matact}(d, m, u)) \end{array}}
\vnbpart{\texttt{free-transport}}
\vnbstmt{\begin{array}{@{}l@{}} \forall d, m, v, a. \\ \quad \operatorname{is\text{-}module}(d) \wedge \\ \quad (v \in \operatorname{mat}(m, 1, \operatorname{vec}(d))) \wedge \\ \quad \operatorname{is\text{-}invertible\text{-}mat}(\operatorname{scal}(d), m, a) \wedge \\ \quad \operatorname{rel\text{-}free}(d, m, v) \Rightarrow \\ \quad \quad \operatorname{rel\text{-}free}(d, m, \operatorname{matact}(d, a, v)) \end{array}}
\vnbpart{\texttt{spans-transport}}
\vnbstmt{\begin{array}{@{}l@{}} \forall d, n, u, m, a. \\ \quad \operatorname{is\text{-}module}(d) \wedge \\ \quad (u \in \operatorname{mat}(n, 1, \operatorname{vec}(d))) \wedge \\ \quad \operatorname{is\text{-}invertible\text{-}mat}(\operatorname{scal}(d), n, m) \wedge \\ \quad \operatorname{is\text{-}submodule}(d, a) \wedge \\ \quad \operatorname{spans}(d, n, u, a) \Rightarrow \\ \quad \quad \operatorname{spans}(d, n, \operatorname{matact}(d, m, u), a) \end{array}}
%
\end{lem}
%
\begin{proof}
%
The proofs: the file that holds each, the lemmas it cites, its oracles and its bill.
%
\vnbprf{\vnbproofpart{\texttt{generates-transport}}}
%
\vnbprf{\textbf{Proof.} \texttt{theorem-\allowbreak{}library/mod-\allowbreak{}basis-\allowbreak{}proof.scm} --- 114 recorded steps.}
%
\vnbprf{\textbf{Lemmas.} \texttt{mat-rows-in-nn}, \texttt{matact-type}, \texttt{lincomb-empty}, \texttt{nn-nonzero-is-succ}, \texttt{nn-one-le-succ}, \texttt{matmul-type}, \texttt{matact-lincomb}, \texttt{matact-assoc}, \texttt{eq-sym}, \texttt{matmul-assoc}, \texttt{identmat-right-identity}.}
%
\vnbprf{\textbf{Oracles.} \texttt{arith}, \texttt{ineq}.  \textbf{Bill:} \texttt{proven modulo 0}.}
%
\vnbprf{\vnbproofpart{\texttt{free-transport}}}
%
\vnbprf{\textbf{Proof.} \texttt{theorem-\allowbreak{}library/mod-\allowbreak{}basis-\allowbreak{}proof.scm} --- 141 recorded steps.}
%
\vnbprf{\textbf{Lemmas.} \texttt{one-in-interval-1}, \texttt{mat-rows-in-nn}, \texttt{interval-in-set}, \texttt{interval-card-in-nn}, \texttt{interval-elt-in-nn}, \texttt{interval-lo}, \texttt{interval-hi}, \texttt{nn-one-in}, \texttt{nn-le-trans-guarded}, \texttt{nn-nonzero-is-succ}, \texttt{nn-one-le-succ}, \texttt{matmul-type}, \texttt{matact-type}, \texttt{matact-lincomb}, \texttt{matact-assoc}, \texttt{matmul-assoc}, \texttt{identmat-right-identity}, \texttt{eq-sym}, \texttt{matmul-entry}, \texttt{matprod-summand-type}, \texttt{entry-in-carrier}, \texttt{ring-mul-zero-left}, \texttt{ras-id}, \texttt{finsum-all-id}.}
%
\vnbprf{\textbf{Oracles.} \texttt{arith}, \texttt{ineq}.  \textbf{Bill:} \texttt{proven modulo 0}.}
%
\vnbprf{\vnbproofpart{\texttt{spans-transport}}}
%
\vnbprf{\textbf{Proof.} \texttt{theorem-\allowbreak{}library/spans-\allowbreak{}transport-\allowbreak{}proof.scm} --- 148 recorded steps.}
%
\vnbprf{\textbf{Lemmas.} \texttt{one-in-interval-1}, \texttt{mat-rows-in-nn}, \texttt{interval-in-set}, \texttt{interval-card-in-nn}, \texttt{interval-elt-in-nn}, \texttt{interval-lo}, \texttt{interval-hi}, \texttt{nn-one-in}, \texttt{nn-le-trans-guarded}, \texttt{matact-entry}, \texttt{matact-summand-type}, \texttt{entry-in-carrier}, \texttt{submodule-finsum-closed}, \texttt{matact-type}, \texttt{lincomb-empty}, \texttt{nn-nonzero-is-succ}, \texttt{nn-one-le-succ}, \texttt{matmul-type}, \texttt{matact-lincomb}, \texttt{matact-assoc}, \texttt{eq-sym}, \texttt{matmul-assoc}, \texttt{identmat-right-identity}.}
%
\vnbprf{\textbf{Oracles.} \texttt{arith}, \texttt{ineq}.  \textbf{Bill:} \texttt{proven modulo 0}.}
%
\end{proof}
%
\noindent\textbf{Deviation from the notes.} The library's two statements are ONE-WAY implications, not the notes' iff: generates(md,n,u) implies generates(md,n,matact(md,pm,u)), and rel-free(md,m,v) implies rel-free(md,m,matact(md,qm,v)).  The converses (which follow by applying the same theorems to the inverse matrix) are not installed.  Setting: a module md over any ring with the transport matrix invertible over scal(md); the notes' E is an ABELIAN GROUP and U is over Z, and the library has NO abelian-group-as-ZZ-module view (zz-act exists as an integer action, but no module structure is built from it), so the notes' own case is not an instance of the library statement.  spans-transport is the submodule version, not in the notes.
%
\renewcommand{\thethm}{3.41}
\begin{prop}[{A relation-free sequence is no longer than a generating sequence}]\label{thm:free-length-le-generators}
%
\vnbstmt{\begin{array}{@{}l@{}} \forall d, n, m, u, v. \\ \quad \operatorname{is\text{-}module}(d) \wedge \\ \quad \operatorname{is\text{-}euclidean\text{-}ring}(\operatorname{scal}(d)) \wedge \\ \quad (n \in \mathbb{N}) \wedge \\ \quad (m \in \mathbb{N}) \wedge \\ \quad (u \in \operatorname{mat}(n, 1, \operatorname{vec}(d))) \wedge \\ \quad (v \in \operatorname{mat}(m, 1, \operatorname{vec}(d))) \wedge \\ \quad \operatorname{generates}(d, n, u) \wedge \\ \quad \operatorname{rel\text{-}free}(d, m, v) \Rightarrow \\ \quad \quad (m \leq n) \end{array}}
%
\end{prop}
%
\begin{proof}
%
The proofs: the file that holds each, the lemmas it cites, its oracles and its bill.
%
\vnbprf{\textbf{Proof.} \texttt{theorem-\allowbreak{}library/rank-\allowbreak{}bound-\allowbreak{}proof.scm} --- 356 recorded steps.}
%
\vnbprf{\textbf{Lemmas.} \texttt{euclidean-ring-is-integral-domain}, \texttt{integral-domain-nontrivial}, \texttt{one-in-interval-1}, \texttt{interval-in-set}, \texttt{interval-card-in-nn}, \texttt{nn-nonzero-is-succ}, \texttt{nn-one-le-succ}, \texttt{nn-one-in}, \texttt{nn-le-refl}, \texttt{interval-mem-intro}, \texttt{entry-in-carrier}, \texttt{lincomb-empty}, \texttt{unitrow-type}, \texttt{lincomb-unfold}, \texttt{matact-summand-type}, \texttt{unitrow-entry-off}, \texttt{module-zero-act}, \texttt{mvag-id}, \texttt{finsum-single-support}, \texttt{unitrow-entry-at}, \texttt{eq-sym}, \texttt{eq-trans}, \texttt{generates-coeff-matrix}, \texttt{smith-diagonalization}, \texttt{mat-equiv-target-is-mat}, \texttt{invertible-mat-is-mat}, \texttt{matact-type}, \texttt{free-transport}, \texttt{matmul-type}, \texttt{matact-assoc}, \texttt{matmul-assoc}, \texttt{identmat-right-identity}, \texttt{nn-not-le-succ-le}, \texttt{nn-succ-closed}, \texttt{matact-entry}, \texttt{interval-elt-in-nn}, \texttt{interval-hi}, \texttt{nn-le-imp-neq-succ}, \texttt{neq-sym}, \texttt{diagonal-off-entry}, \texttt{finsum-all-id}.}
%
\vnbprf{\textbf{Oracles.} \texttt{arith}, \texttt{ineq}, \texttt{crs}.  \textbf{Bill:} \texttt{proven modulo 0}.}
%
\end{proof}
%
\noindent\textbf{Deviation from the notes.} The first sentence is proven, for a module whose SCALAR RING IS EUCLIDEAN (the notes' Z case is only an instance once the missing abelian-group-as-ZZ-module view exists).  The second sentence -- that all free generating sequences of an abelian group have the same LENGTH, the fact that makes rank well defined -- is NOT stated anywhere, although it follows by two applications of the first.  Guards the notes do not carry: n, m in nn, and both sequences are column matrices u in mat(n,1,vec(md)), v in mat(m,1,vec(md)).
%
\renewcommand{\thethm}{3.43}
\begin{cor}[{GL\_n(Z) acts freely and transitively on the free generating sequences}]\label{thm:linear-algebra-3.43}
%
\textit{If E is a free abelian group of rank n, the action (A, u) |-$>$ A . u of M\_n\textasciicircum{}x(Z) on free generating sequences of E is free and transitive.}
%
Not in the library.
%
\end{cor}
%
\begin{proof}
%
The proofs: the file that holds each, the lemmas it cites, its oracles and its bill.
%
\vnbprf{There is no library statement, and so no proof.}
%
\end{proof}
%
\noindent\textbf{Deviation from the notes.} Only the one-way transport exists (generates-transport, free-transport).  matact is never presented as a GROUP ACTION, and neither freeness nor transitivity of it is stated; the library does not form the group M\_n\textasciicircum{}x(Z) at all.
%
\renewcommand{\thethm}{3.44}
\begin{prop}[{Adapted basis for a subgroup, with the divisibility chain}]\label{thm:submodule-of-free-is-free}
%
\vnbstmt{\begin{array}{@{}l@{}} \forall d.; \forall n \in \mathbb{N}, u \in \operatorname{mat}(n, 1, \operatorname{vec}(d)), m. \\ \quad \operatorname{is\text{-}module}(d) \wedge \\ \quad \operatorname{is\text{-}euclidean\text{-}ring}(\operatorname{scal}(d)) \wedge \\ \quad \operatorname{generates}(d, n, u) \wedge \\ \quad \operatorname{rel\text{-}free}(d, n, u) \wedge \\ \quad \operatorname{is\text{-}submodule}(d, m) \Rightarrow \\ \quad \quad \exists k. \\ \quad \quad \quad (k \in \mathbb{N}) \wedge \\ \quad \quad \quad (k \leq n) \wedge \\ \quad \quad \quad \exists w \in \operatorname{mat}(k, 1, \operatorname{vec}(d)). \\ \quad \quad \quad \quad (\operatorname{spans}(d, k, w, m) \wedge \operatorname{rel\text{-}free}(d, k, w)) \end{array}}
%
\end{prop}
%
\begin{proof}
%
The proofs: the file that holds each, the lemmas it cites, its oracles and its bill.
%
\vnbprf{Asserted (warrant: reference); no proof in the library.}
%
\end{proof}
%
\noindent\textbf{Deviation from the notes.} Much weaker.  submodule-of-free-is-free (ASSERTED, warrant reference) gives only that the submodule has SOME family w of length k $<$= n that spans it and is rel-free.  It gives no ADAPTED BASIS e\_1,...,e\_n OF E, no coefficients b\_i, hence no divisibility chain b\_i divides b\_\{i+1\}, no non-negativity, and no claim that the b\_i are invariants of F.  Setting is a module over a euclidean ring rather than a free abelian group.
%
\renewcommand{\thethm}{3.45}
\begin{cor}[{The norm N\_F(E): the absolute determinant is basis-independent}]\label{thm:linear-algebra-3.45}
%
\textit{If F is contained in E, with free generating sets v, u and v = A . u, then |det A| does not depend on the choice of v and u.}
%
Not in the library.
%
\end{cor}
%
\begin{proof}
%
The proofs: the file that holds each, the lemmas it cites, its oracles and its bill.
%
\vnbprf{There is no library statement, and so no proof.}
%
\end{proof}
%
\noindent\textbf{Deviation from the notes.} No determinant of a change-of-basis matrix is ever taken, and no index or norm invariant of a submodule is defined.  Depends on 3.44's adapted basis and on determinant theory the library does not have.
%
\renewcommand{\thethm}{3.46}
\begin{cor}[{A subgroup of a free abelian group of rank n is free of rank at most n}]\label{thm:linear-algebra-3.46}
%
\vnbstmt{\begin{array}{@{}l@{}} \forall d.; \forall n \in \mathbb{N}, u \in \operatorname{mat}(n, 1, \operatorname{vec}(d)), m. \\ \quad \operatorname{is\text{-}module}(d) \wedge \\ \quad \operatorname{is\text{-}euclidean\text{-}ring}(\operatorname{scal}(d)) \wedge \\ \quad \operatorname{generates}(d, n, u) \wedge \\ \quad \operatorname{rel\text{-}free}(d, n, u) \wedge \\ \quad \operatorname{is\text{-}submodule}(d, m) \Rightarrow \\ \quad \quad \exists k. \\ \quad \quad \quad (k \in \mathbb{N}) \wedge \\ \quad \quad \quad (k \leq n) \wedge \\ \quad \quad \quad \exists w \in \operatorname{mat}(k, 1, \operatorname{vec}(d)). \\ \quad \quad \quad \quad (\operatorname{spans}(d, k, w, m) \wedge \operatorname{rel\text{-}free}(d, k, w)) \end{array}}
%
\end{cor}
%
\begin{proof}
%
The proofs: the file that holds each, the lemmas it cites, its oracles and its bill.
%
\vnbprf{Asserted (warrant: reference); no proof in the library.}
%
\end{proof}
%
\noindent\textbf{Deviation from the notes.} ASSERTED with warrant reference -- trusted, not proven, and one of only 32 reference-tier supports in the library.  The library form is for a submodule sm of a module md over a EUCLIDEAN scalar ring, with a generating AND rel-free column family u of length n; the conclusion is a k $<$= n and a w in mat(k,1,vec(md)) with spans(md,k,w,sm) and rel-free(md,k,w) -- the notes' free of rank $<$= n, granted the missing abelian-group-as-ZZ-module view.  Its own harder input, 3.47, IS proven, so this is assembly debt, not mathematical debt.
%
\renewcommand{\thethm}{3.47}
\begin{cor}[{A subgroup of an n-generated group is generated by at most n elements}]\label{thm:submodule-fg}
%
\leavevmode
\vnbpart{\texttt{submodule-fg}}
\vnbstmt{\begin{array}{@{}l@{}} \forall d.; \forall n \in \mathbb{N}, u \in \operatorname{mat}(n, 1, \operatorname{vec}(d)), m. \\ \quad \operatorname{is\text{-}module}(d) \wedge \\ \quad \operatorname{is\text{-}euclidean\text{-}ring}(\operatorname{scal}(d)) \wedge \\ \quad \operatorname{generates}(d, n, u) \wedge \\ \quad \operatorname{is\text{-}submodule}(d, m) \Rightarrow \\ \quad \quad \exists k. \\ \quad \quad \quad (k \in \mathbb{N}) \wedge \\ \quad \quad \quad (k \leq n) \wedge \\ \quad \quad \quad \exists w \in \operatorname{mat}(k, 1, \operatorname{vec}(d)). \\ \quad \quad \quad \quad \operatorname{spans}(d, k, w, m) \end{array}}
\vnbpart{\texttt{spans-submodule-fg}}
\vnbstmt{\begin{array}{@{}l@{}} \forall d.; \forall n \in \mathbb{N}, u \in \operatorname{mat}(n, 1, \operatorname{vec}(d)), m, a. \\ \quad \operatorname{is\text{-}module}(d) \wedge \\ \quad \operatorname{is\text{-}euclidean\text{-}ring}(\operatorname{scal}(d)) \wedge \\ \quad \operatorname{is\text{-}submodule}(d, m) \wedge \\ \quad \operatorname{spans}(d, n, u, m) \wedge \\ \quad \operatorname{is\text{-}submodule}(d, a) \wedge \\ \quad (a \subseteq m) \Rightarrow \\ \quad \quad \exists k. \\ \quad \quad \quad (k \in \mathbb{N}) \wedge \\ \quad \quad \quad (k \leq n) \wedge \\ \quad \quad \quad \exists w \in \operatorname{mat}(k, 1, \operatorname{vec}(d)). \\ \quad \quad \quad \quad \operatorname{spans}(d, k, w, a) \end{array}}
%
\end{cor}
%
\begin{proof}
%
The proofs: the file that holds each, the lemmas it cites, its oracles and its bill.
%
\vnbprf{\vnbproofpart{\texttt{submodule-fg}}}
%
\vnbprf{\textbf{Proof.} \texttt{theorem-\allowbreak{}library/submodule-\allowbreak{}fg-\allowbreak{}proof.scm} --- 19 recorded steps.}
%
\vnbprf{\textbf{Lemmas.} \texttt{whole-module-is-submodule}, \texttt{generates-implies-spans-vec}, \texttt{spans-submodule-fg}.}
%
\vnbprf{\textbf{Oracles.} \texttt{arith}, \texttt{ineq}, \texttt{crs}.  \textbf{Bill:} \texttt{proven modulo 0}.}
%
\vnbprf{\vnbproofpart{\texttt{spans-submodule-fg}}}
%
\vnbprf{\textbf{Proof.} \texttt{theorem-\allowbreak{}library/spans-\allowbreak{}submodule-\allowbreak{}fg-\allowbreak{}proof.scm} --- 9 recorded steps.}
%
\vnbprf{\textbf{Lemmas.} \texttt{spans-fg-base}, \texttt{spans-fg-step}.}
%
\vnbprf{\textbf{Oracles.} \texttt{arith}, \texttt{ineq}, \texttt{crs}.  \textbf{Bill:} \texttt{proven modulo 0}.}
%
\end{proof}
%
\noindent\textbf{Deviation from the notes.} This is the library's headline linear-algebra result, proven and billing modulo 0 since 2026-09-17.  Differences: it is stated for a MODULE over a euclidean ring (the notes' abelian group is not an instance, the abelian-group-as-ZZ-module view being absent), and its conclusion is a column family w in mat(k,1,vec(md)) with k $<$= n that SPANS the submodule.  spans-submodule-fg is the strictly more general form the descent actually needs: a submodule of a submodule spanned by n elements is spanned by k $<$= n.  The notes' own proof of this corollary (through Z\textasciicircum{}n and 3.44) is CIRCULAR -- Smith form consumes a finite generating set for F and cannot produce one -- so the library proves it directly, by descent on the ideal of last coefficients; the \textbackslash{}iffalse'd induction at tex:1651 is load-bearing in the notes.
%
\renewcommand{\thethm}{3.48}
\begin{prop}[{Uniqueness of the principal sequence of a subgroup}]\label{thm:linear-algebra-3.48}
%
\textit{The principal sequence b of a subgroup F of E is uniquely determined.}
%
Not in the library.
%
\end{prop}
%
\begin{proof}
%
The proofs: the file that holds each, the lemmas it cites, its oracles and its bill.
%
\vnbprf{There is no library statement, and so no proof.}
%
\end{proof}
%
\noindent\textbf{Deviation from the notes.} The uniqueness half of the Smith theory is nowhere in the library: smith-normal-form is pure existence, and neither the principal sequence nor the notion of invariant factor is defined.  It is exactly the part of 3.36 that the library dropped.
%
\renewcommand{\thethm}{3.49}
\begin{lem}[{The quotient E/F for a diagonal subgroup}]\label{thm:linear-algebra-3.49}
%
\textit{If E is freely generated by e\_1,...,e\_n, the a\_1,...,a\_k are non-zero and F is generated by a\_1 e\_1,...,a\_k e\_k, then E/F is isomorphic to Z/(a\_1) + ... + Z/(a\_k) + Z\textasciicircum{}\{n-k\}.}
%
Not in the library.
%
\end{lem}
%
\begin{proof}
%
The proofs: the file that holds each, the lemmas it cites, its oracles and its bill.
%
\vnbprf{There is no library statement, and so no proof.}
%
\end{proof}
%
\noindent\textbf{Deviation from the notes.} The library has no quotient MODULE and no quotient group -- only the quotient of a SETOID (quotient, class, proj, quotient-universal) and ringoid-quotient-ring.  It also has no direct sum of modules and no cyclic group Z/(a).  is-hom-module-def exists as a homomorphism predicate, but nothing builds the map of this lemma.
%
\renewcommand{\thethm}{3.50}
\begin{prop}[{Structure of a finitely generated abelian group}]\label{thm:linear-algebra-3.50}
%
\textit{If E is a finitely generated abelian group there are w\_1,...,w\_l in E and non-negative a\_1,...,a\_l with a\_i dividing a\_j for i $<$= j such that (x\_1,...,x\_l) |-$>$ sum x\_i w\_i induces an isomorphism Z/(a\_1) + ... + Z/(a\_l) -$>$ E; the sequence a\_1,...,a\_l is unique up to sign.}
%
Not in the library.
%
\end{prop}
%
\begin{proof}
%
The proofs: the file that holds each, the lemmas it cites, its oracles and its bill.
%
\vnbprf{There is no library statement, and so no proof.}
%
\end{proof}
%
\noindent\textbf{Deviation from the notes.} Absent for the reasons listed under 3.49 (no quotient module, no direct sum, no cyclic group) and under 3.48 (no uniqueness).  The nearest library assets are submodule-fg and smith-normal-form, which supply the existence half of the input but nothing of the conclusion.
%
\renewcommand{\thethm}{3.51}
\begin{cor}[{Every finitely generated abelian group is a direct sum of cyclic groups}]\label{thm:linear-algebra-3.51}
%
\textit{Every finitely generated abelian group is a finite direct sum Z/(m\_1) + ... + Z/(m\_l) with m\_k $>$ 0 and m\_k dividing m\_\{k+1\}; the representation is unique.}
%
Not in the library.
%
\end{cor}
%
\begin{proof}
%
The proofs: the file that holds each, the lemmas it cites, its oracles and its bill.
%
\vnbprf{There is no library statement, and so no proof.}
%
\end{proof}
%
\noindent\textbf{Deviation from the notes.} Absent; it is the restatement of 3.50, and the library has neither the direct sum nor the cyclic groups nor the uniqueness.
%
\renewcommand{\thethm}{3.52}
\begin{cor}[{The cardinality of E/F is the norm N\_E(F)}]\label{thm:linear-algebra-3.52}
%
\textit{If E is a finitely generated free abelian group, F a subgroup and N\_E(F) /= 0, then card(E/F) = N\_E(F).}
%
Not in the library.
%
\end{cor}
%
\begin{proof}
%
The proofs: the file that holds each, the lemmas it cites, its oracles and its bill.
%
\vnbprf{There is no library statement, and so no proof.}
%
\end{proof}
%
\noindent\textbf{Deviation from the notes.} Depends on 3.45 (the norm) and 3.49 (the quotient), both absent.  CARD is defined in the library since 2026-09-20, but the quotient E/F it would be applied to does not exist.
%
\section{Complex analysis}
%
\renewcommand{\thethm}{1.2}
\begin{prop}[{Uniqueness of the least upper bound}]\label{thm:esup-unique}
%
\vnbstmt{\begin{array}{@{}l@{}} \forall s, b_{1}, b_{2}. \\ \quad (b_{1} \in \operatorname{rr\text{-}pos\text{-}star}) \wedge \\ \quad \forall x \in s.\; (x \leq b_{1}) \wedge \\ \quad \forall b \in \operatorname{rr\text{-}pos\text{-}star}. \\ \quad \quad (\forall x \in s.\; (x \leq b) \Rightarrow (b_{1} \leq b)) \wedge \\ \quad (b_{2} \in \operatorname{rr\text{-}pos\text{-}star}) \wedge \\ \quad \forall x \in s.\; (x \leq b_{2}) \wedge \\ \quad \forall b \in \operatorname{rr\text{-}pos\text{-}star}. \\ \quad \quad (\forall x \in s.\; (x \leq b) \Rightarrow (b_{2} \leq b)) \Rightarrow \\ \quad \quad (b_{1} = b_{2}) \end{array}}
%
\end{prop}
%
\begin{proof}
%
The proofs: the file that holds each, the lemmas it cites, its oracles and its bill.
%
\vnbprf{\textbf{Proof.} \texttt{theorem-\allowbreak{}library/rake-\allowbreak{}esup-\allowbreak{}defined.scm} --- 43 recorded steps.}
%
\vnbprf{\textbf{Lemmas.} \texttt{rr-pos-star-le-antisymm}.}
%
\vnbprf{\textbf{Oracles.} none.  \textbf{Bill:} \texttt{proven modulo 0}.}
%
\end{proof}
%
\noindent\textbf{Deviation from the notes.} The library states uniqueness of the least upper bound only on the EXTENDED non-negative reals: esup-unique quantifies over s\_ subset RR-POS-STAR with both candidate bounds in RR-POS-STAR. For a subset of RR the supremum is the primitive constant SUP, characterised by the three AXIOMS rr-sup-in / rr-sup-upper / rr-sup-least (warrant: reference), each guarded by non-emptiness of the set and boundedness above; uniqueness of a least upper bound of a subset of RR is nowhere stated.
%
\renewcommand{\thethm}{1.3}
\begin{prop}[{C as 2x2 real matrices}]\label{thm:complex-analysis-1.3}
%
\textit{Phi : a + b i |-$>$ [[a, b], [-b, a]] is an algebraic isomorphism from C onto the set of 2x2 real matrices of that form.}
%
Not in the library.
%
\end{prop}
%
\begin{proof}
%
The proofs: the file that holds each, the lemmas it cites, its oracles and its bill.
%
\vnbprf{There is no library statement, and so no proof.}
%
\end{proof}
%
\noindent\textbf{Deviation from the notes.} Nothing about matrix representations of C. The nearest is the R-linear identification of C with R x R: cc-coords-is-bijection and cc-of-pair-is-bijection (proven), with cc-coords-add and cc-coords-real-mul; there is no set of 2x2 matrices of the form [[a,b],[-b,a]] and no statement that any map from C is an algebra (ring) isomorphism onto it.
%
\renewcommand{\thethm}{1.5}
\begin{prop}[{Unions and finite intersections of open sets}]\label{thm:complex-analysis-1.5}
%
\leavevmode
\vnbpart{\texttt{union-of-opens-open}}
\vnbstmt{\begin{array}{@{}l@{}} \forall s, a, g. \\ \quad \operatorname{is\text{-}metric\text{-}space}(s) \wedge \\ \quad \forall i \in a.\; \operatorname{is\text{-}open}(s, g(i)) \Rightarrow \\ \quad \quad \operatorname{is\text{-}open}(s, \operatorname{big\text{-}union}(i, a, g(i))) \end{array}}
\vnbpart{\texttt{inter-of-opens-open}}
\vnbstmt{\begin{array}{@{}l@{}} \forall s, u, w. \\ \quad \operatorname{is\text{-}metric\text{-}space}(s) \wedge \\ \quad \operatorname{is\text{-}open}(s, u) \wedge \\ \quad \operatorname{is\text{-}open}(s, w) \Rightarrow \\ \quad \quad \operatorname{is\text{-}open}(s, (u \cap w)) \end{array}}
%
\end{prop}
%
\begin{proof}
%
The proofs: the file that holds each, the lemmas it cites, its oracles and its bill.
%
\vnbprf{\vnbproofpart{\texttt{union-of-opens-open}}}
%
\vnbprf{\textbf{Proof.} \texttt{theorem-\allowbreak{}library/rake-\allowbreak{}open-\allowbreak{}sets.scm} --- 43 recorded steps.}
%
\vnbprf{\textbf{Lemmas.} \texttt{subset-def}, \texttt{subset-mem-fwd}.}
%
\vnbprf{\textbf{Oracles.} none.  \textbf{Bill:} \texttt{proven modulo 0}.}
%
\vnbprf{\vnbproofpart{\texttt{inter-of-opens-open}}}
%
\vnbprf{\textbf{Proof.} \texttt{theorem-\allowbreak{}library/rake-\allowbreak{}open-\allowbreak{}sets.scm} --- 114 recorded steps.}
%
\vnbprf{\textbf{Lemmas.} \texttt{subset-def}, \texttt{intersection-membership}, \texttt{subset-mem-fwd}, \texttt{rr-pos-rr-in-rr}, \texttt{rr-lt-of-pos-rr}, \texttt{rr-min-pos}, \texttt{rr-pos-rr-of-lt}, \texttt{ball-sep-unfold}, \texttt{metric-dist-real}.}
%
\vnbprf{\textbf{Oracles.} \texttt{ineq}.  \textbf{Bill:} \texttt{proven modulo 0}.}
%
\end{proof}
%
\noindent\textbf{Deviation from the notes.} The union half matches (union-of-opens-open: an indexed family g over an index set a, conclusion is-open(s, big-union(i, a, g(i)))). The intersection half is only BINARY: inter-of-opens-open takes two open sets u, w. The notes' claim for the intersection of FINITELY many open sets is not stated anywhere.
%
\renewcommand{\thethm}{1.7}
\begin{prop}[{Monotone convergence for real sequences}]\label{thm:complex-analysis-1.7}
%
\vnbstmt{\begin{array}{@{}l@{}} \forall f \in (\mathbb{N} \to \mathbb{R}). \\ \quad \forall k \in \mathbb{N}.\; (f(k) \leq f((k + 1))) \wedge \\ \quad \exists d \in \mathbb{R}.\; \forall k \in \mathbb{N}.\; (f(k) \leq d) \Rightarrow \\ \quad \quad \operatorname{converges}(\operatorname{rr\text{-}ms}, f) \end{array}}
%
\end{prop}
%
\begin{proof}
%
The proofs: the file that holds each, the lemmas it cites, its oracles and its bill.
%
\vnbprf{\textbf{Proof.} \texttt{theorem-\allowbreak{}library/monotone-\allowbreak{}convergence-\allowbreak{}proof.scm} --- 123 recorded steps.}
%
\vnbprf{\textbf{Lemmas.} \texttt{nn-zero-in}, \texttt{fun-apply-type-c}, \texttt{subset-def}, \texttt{rr-sup-in}, \texttt{rr-sup-upper}, \texttt{rr-is-metric-space}, \texttt{rr-sup-approx}, \texttt{nn-monotone-step-implies-le}, \texttt{rr-sub-in-rr}, \texttt{rr-ms-dist}, \texttt{rr-abs-bound}.}
%
\vnbprf{\textbf{Oracles.} \texttt{ineq}, \texttt{arith}, \texttt{crs}.  \textbf{Bill:} \texttt{proven modulo 0}.}
%
\end{proof}
%
\noindent\textbf{Deviation from the notes.} The library proves only that the sequence CONVERGES; it nowhere identifies the limit with sup\_n x\_n, which is the whole content of the notes' statement. Only the non-decreasing / bounded-above case is stated; the non-increasing / bounded-below mirror (lim = inf) is absent. Monotonicity is stated in the one-step form forall k in NN. f(k) $<$= f(succ(k)), and boundedness as forsome bnd in RR. forall k. f(k) $<$= bnd rather than through SUP.
%
\renewcommand{\thethm}{1.9}
\begin{prop}[{Closed sets and limits of sequences}]\label{thm:complex-analysis-1.9}
%
\leavevmode
\vnbpart{\texttt{closed-limit-in}}
\vnbstmt{\begin{array}{@{}l@{}} \forall s, a.; \forall f \in (\mathbb{N} \to a), l. \\ \quad \operatorname{is\text{-}metric\text{-}space}(s) \wedge \\ \quad \operatorname{is\text{-}closed}(s, a) \wedge \\ \quad \operatorname{converges\text{-}to}(s, f, l) \Rightarrow \\ \quad \quad (l \in a) \end{array}}
\vnbpart{\texttt{closed-set-limit-in}}
\vnbstmt{\begin{array}{@{}l@{}} \forall s, a, f, l. \\ \quad \operatorname{is\text{-}closed}(s, a) \wedge \\ \quad \operatorname{converges\text{-}to}(s, f, l) \wedge \\ \quad \forall n \in \mathbb{N}.\; (f(n) \in a) \Rightarrow \\ \quad \quad (l \in a) \end{array}}
%
\end{prop}
%
\begin{proof}
%
The proofs: the file that holds each, the lemmas it cites, its oracles and its bill.
%
\vnbprf{\vnbproofpart{\texttt{closed-limit-in}}}
%
\vnbprf{\textbf{Proof.} \texttt{theorem-\allowbreak{}library/metric-\allowbreak{}subspace-\allowbreak{}laws.scm} --- 105 recorded steps.}
%
\vnbprf{\textbf{Lemmas.} \texttt{complement-in-membership}, \texttt{rr-pos-halvable}, \texttt{rr-lt-of-pos-rr}, \texttt{rr-pos-rr-in-rr}, \texttt{nn-in-rr}, \texttt{rr-leq-reflexive}, \texttt{fun-apply-type-c}, \texttt{subset-mem-fwd}, \texttt{metric-dist-real}, \texttt{metric-sym}, \texttt{ball-mem-from-le}.}
%
\vnbprf{\textbf{Oracles.} \texttt{ineq}, \texttt{crs}.  \textbf{Bill:} \texttt{proven modulo 0}.}
%
\vnbprf{\vnbproofpart{\texttt{closed-set-limit-in}}}
%
\vnbprf{\textbf{Proof.} \texttt{theorem-\allowbreak{}library/metric-\allowbreak{}closure-\allowbreak{}laws.scm} --- 93 recorded steps.}
%
\vnbprf{\textbf{Lemmas.} \texttt{complement-in-membership}, \texttt{rr-pos-rr-in-rr}, \texttt{rr-lt-of-pos-rr}, \texttt{ball-is-set}, \texttt{rr-pos-halvable}, \texttt{nn-le-refl}, \texttt{subset-mem-fwd}, \texttt{metric-dist-real}, \texttt{metric-sym}, \texttt{ball-mem-from-le}.}
%
\vnbprf{\textbf{Oracles.} \texttt{ineq}, \texttt{crs}.  \textbf{Bill:} \texttt{proven modulo 0}.}
%
\end{proof}
%
\noindent\textbf{Deviation from the notes.} Only the FORWARD direction is in the library (A closed implies every convergent sequence of elements of A has its limit in A), in two forms: closed-limit-in with f in fun(NN, a), and closed-set-limit-in with the pointwise membership forall n in NN. f(n) in a. The CONVERSE -- sequentially closed implies closed -- is not stated, so the notes' `if and only if' is not available, nor the corollary drawn from it in the notes (a set is closed iff it contains all its accumulation points).
%
\renewcommand{\thethm}{1.11}
\begin{prop}[{Characterisations of compactness}]\label{thm:complex-analysis-1.11}
%
\leavevmode
\vnbpart{\texttt{compact-iff-fip}}
\vnbstmt{\begin{array}{@{}l@{}} \forall s. \\ \quad \operatorname{is\text{-}metric\text{-}space}(s) \Rightarrow  \\ \quad \quad (\operatorname{is\text{-}compact}(s) \Leftrightarrow \forall c.\; (\operatorname{has\text{-}fip}(s, c) \Rightarrow \exists p.\; (p \in \operatorname{intersection\text{-}of}(c)))) \end{array}}
\vnbpart{\texttt{compact-iff-cluster-point}}
\vnbstmt{\begin{array}{@{}l@{}} \forall s. \\ \quad \operatorname{is\text{-}metric\text{-}space}(s) \Rightarrow  \\ \quad \quad (\operatorname{is\text{-}compact}(s) \Leftrightarrow \forall f \in (\mathbb{N} \to \operatorname{pts}(s)).\; \exists x.\; \operatorname{cluster\text{-}point}(s, f, x)) \end{array}}
\vnbpart{\texttt{compact-iff-tb-complete}}
\vnbstmt{\begin{array}{@{}l@{}} \forall s. \\ \quad \operatorname{is\text{-}metric\text{-}space}(s) \Rightarrow  \\ \quad \quad (\operatorname{is\text{-}compact}(s) \Leftrightarrow (\operatorname{totally\text{-}bounded}(s) \wedge \operatorname{is\text{-}complete}(s))) \end{array}}
%
\end{prop}
%
\begin{proof}
%
The proofs: the file that holds each, the lemmas it cites, its oracles and its bill.
%
\vnbprf{\vnbproofpart{\texttt{compact-iff-fip}}}
%
\vnbprf{\textbf{Proof.} \texttt{theorem-\allowbreak{}library/rake-\allowbreak{}compact-\allowbreak{}fip.scm} --- 6 recorded steps.}
%
\vnbprf{\textbf{Lemmas.} \texttt{compact-implies-fip}, \texttt{fip-implies-compact}.}
%
\vnbprf{\textbf{Oracles.} \texttt{ineq}, \texttt{arith}.  \textbf{Bill:} \texttt{proven modulo 0}.}
%
\vnbprf{\vnbproofpart{\texttt{compact-iff-cluster-point}}}
%
\vnbprf{\textbf{Proof.} \texttt{theorem-\allowbreak{}library/rake-\allowbreak{}compact-\allowbreak{}equivalences.scm} --- 58 recorded steps.}
%
\vnbprf{\textbf{Lemmas.} \texttt{compact-seq-has-cluster}, \texttt{cluster-point-has-convergent-subseq}, \texttt{seq-compact-implies-compact}.}
%
\vnbprf{\textbf{Oracles.} \texttt{ineq}, \texttt{arith}, \texttt{crs}.  \textbf{Bill:} \texttt{proven modulo 0}.}
%
\vnbprf{\vnbproofpart{\texttt{compact-iff-tb-complete}}}
%
\vnbprf{\textbf{Proof.} \texttt{theorem-\allowbreak{}library/rake-\allowbreak{}compact-\allowbreak{}equivalences.scm} --- 68 recorded steps.}
%
\vnbprf{\textbf{Lemmas.} \texttt{compact-implies-totally-bounded}, \texttt{compact-implies-complete}, \texttt{totally-bounded-has-cauchy-subsequence}, \texttt{subseq-is-fun}, \texttt{seq-compact-implies-compact}.}
%
\vnbprf{\textbf{Oracles.} \texttt{ineq}, \texttt{arith}, \texttt{crs}.  \textbf{Bill:} \texttt{proven modulo 0}.}
%
\end{proof}
%
\noindent\textbf{Deviation from the notes.} All four conditions are present and each equivalence is proven, but as three separate IFF theorems against (1) rather than one four-way equivalence. In (2) the library's HAS-FIP carries two inhabitedness guards the notes do not: the family c must be non-empty (forsome a0. a0 in c), and only NON-EMPTY finite subfamilies are required to have a non-empty intersection. (3) is stated with CLUSTER-POINT defined by: for every eps $>$ 0 and every m in NN there is n $>$= m with dist(f(n), x) $<$ eps -- the notes define a cluster point by the set \{n : x\_n in B(y,r)\} being infinite.
%
\renewcommand{\thethm}{1.15}
\begin{thm}[{C is an algebraically closed field}]\label{thm:complex-analysis-1.15}
%
\textit{C is an algebraically closed field: every polynomial equation a\_n z\textasciicircum{}n + ... + a\_0 = 0 with n $>$= 1 and a\_n /= 0 has a solution in C.}
%
Not in the library.
%
\end{thm}
%
\begin{proof}
%
The proofs: the file that holds each, the lemmas it cites, its oracles and its bill.
%
\vnbprf{There is no library statement, and so no proof.}
%
\end{proof}
%
\noindent\textbf{Deviation from the notes.} Nothing about roots of polynomials over CC. The library has no polynomial with complex coefficients as a function on CC, no n-th roots of a complex number, and no statement that a non-constant polynomial has a zero. The nearest facts are the algebraic ones about CC: cc-is-normed-field (proven) -- note that IS-NORMED-FIELD has no MUL-INV slot and does NOT require multiplicative inverses, so it is not the field property -- together with the flat CC arithmetic axioms cc-recip-closed and cc-recip-inverse, which do supply inverses of non-zero complex numbers.
%
\renewcommand{\thethm}{1.16}
\begin{lem}[{Convergence in the extended plane}]\label{thm:complex-analysis-1.16}
%
\textit{For a sequence \{x\_n\} in C: it converges to x in C in the topology of C-hat iff it converges to x in the Euclidean metric; it converges to infinity iff |x\_n| -$>$ +infinity.}
%
Not in the library.
%
\end{lem}
%
\begin{proof}
%
The proofs: the file that holds each, the lemmas it cites, its oracles and its bill.
%
\vnbprf{There is no library statement, and so no proof.}
%
\end{proof}
%
\noindent\textbf{Deviation from the notes.} The library has no one-point compactification of C: no C-hat, no topology on C union \{infinity\}, no |x\_n| -$>$ +infinity. Only convergence in the Euclidean metric space CC-MS is available (converges-to(cc-ms, f, l), with cc-converges-iff-coords and cc-complete), which is exactly the first half of the notes' `if and only if' but not the statement about the extended topology. The extended line RR-STAR / POS-INF exists but is a different object.
%
\renewcommand{\thethm}{1.18}
\begin{prop}[{Linear maps on C\textasciicircum{}2 and linear fractional transformations}]\label{thm:complex-analysis-1.18}
%
\textit{An invertible complex linear map S on C\textasciicircum{}2 with matrix [[a,b],[c,d]] factors through a transformation L\_S on the projective plane P; under the identification of P with C-hat, S |-$>$ L\_S is a surjective homomorphism from the group of invertible complex linear maps onto the set of linear fractional transformations, with kernel the scalar multiples of the identity.}
%
Not in the library.
%
\end{prop}
%
\begin{proof}
%
The proofs: the file that holds each, the lemmas it cites, its oracles and its bill.
%
\vnbprf{There is no library statement, and so no proof.}
%
\end{proof}
%
\noindent\textbf{Deviation from the notes.} Nothing near. The library has no complex projective space P, no quotient of C\textasciicircum{}2 minus 0 by scalars, and no linear fractional transformations at all (grep for `projective', `fractional', `moebius' over THEOREMS.md returns nothing). C\textasciicircum{}2 itself is not a library object.
%
\renewcommand{\thethm}{1.19}
\begin{cor}[{The linear fractional transformations form a group}]\label{thm:complex-analysis-1.19}
%
\textit{The set of linear fractional transformations is a group.}
%
Not in the library.
%
\end{cor}
%
\begin{proof}
%
The proofs: the file that holds each, the lemmas it cites, its oracles and its bill.
%
\vnbprf{There is no library statement, and so no proof.}
%
\end{proof}
%
\noindent\textbf{Deviation from the notes.} Nothing near: no linear fractional transformations in the library.
%
\renewcommand{\thethm}{1.20}
\begin{cor}[{SL(2,C) onto the linear fractional transformations}]\label{thm:complex-analysis-1.20}
%
\textit{S |-$>$ L\_S is a surjective homomorphism from the group of linear maps on C\textasciicircum{}2 with determinant 1 onto the set of linear fractional transformations, with kernel \{-I, I\}.}
%
Not in the library.
%
\end{cor}
%
\begin{proof}
%
The proofs: the file that holds each, the lemmas it cites, its oracles and its bill.
%
\vnbprf{There is no library statement, and so no proof.}
%
\end{proof}
%
\noindent\textbf{Deviation from the notes.} Nothing near: no linear fractional transformations, no determinant-1 group of complex 2x2 matrices. The library's DET is for matrices over a commutative ring and is unconnected to any of this.
%
\renewcommand{\thethm}{1.21}
\begin{prop}[{The Cayley transform}]\label{thm:complex-analysis-1.21}
%
\textit{The linear fractional transformation C z = (z - i)/(z + i) maps R bijectively onto T minus \{1\}, and maps the upper half plane H+ bijectively onto the open unit disc D.}
%
Not in the library.
%
\end{prop}
%
\begin{proof}
%
The proofs: the file that holds each, the lemmas it cites, its oracles and its bill.
%
\vnbprf{There is no library statement, and so no proof.}
%
\end{proof}
%
\noindent\textbf{Deviation from the notes.} Nothing near: no Cayley transform, no upper half plane, no unit circle T, no unit disc as a named object.
%
\renewcommand{\thethm}{1.22}
\begin{prop}[{The Hopf map is continuous and onto S\textasciicircum{}2}]\label{thm:complex-analysis-1.22}
%
\textit{pi(u1, u2) = (2 u1 conj(u2), |u1|\textasciicircum{}2 - |u2|\textasciicircum{}2) is a continuous mapping C\textasciicircum{}2 -$>$ C x R, and maps the sphere S\textasciicircum{}3 of C\textasciicircum{}2 onto the unit sphere S\textasciicircum{}2 of R\textasciicircum{}2 x R.}
%
Not in the library.
%
\end{prop}
%
\begin{proof}
%
The proofs: the file that holds each, the lemmas it cites, its oracles and its bill.
%
\vnbprf{There is no library statement, and so no proof.}
%
\end{proof}
%
\noindent\textbf{Deviation from the notes.} Nothing near: no C\textasciicircum{}2, no spheres S\textasciicircum{}3 or S\textasciicircum{}2, no map pi(u1,u2) = (2 u1 conj(u2), |u1|\textasciicircum{}2 - |u2|\textasciicircum{}2). The library does have CONJUGATE and MAGNITUDE on CC (cc-conjugate-mul, cc-magnitude-mul, ...), the ingredients only.
%
\renewcommand{\thethm}{1.24}
\begin{cor}[{The Hopf fibration factors through a bijection S\textasciicircum{}2 -$>$ P}]\label{thm:complex-analysis-1.24}
%
\textit{The mapping pi factors through a bijection S\textasciicircum{}2 -$>$ P.}
%
Not in the library.
%
\end{cor}
%
\begin{proof}
%
The proofs: the file that holds each, the lemmas it cites, its oracles and its bill.
%
\vnbprf{There is no library statement, and so no proof.}
%
\end{proof}
%
\noindent\textbf{Deviation from the notes.} Nothing near: no projective plane, no spheres.
%
\renewcommand{\thethm}{2.2}
\begin{prop}[{Differentiable implies continuous}]\label{thm:diff-on-implies-continuous}
%
\vnbstmt{\begin{array}{@{}l@{}} \forall k, u, f, a, l. \\ \quad \operatorname{is\text{-}diff\text{-}on}(k, u, f, a, l) \Rightarrow  \\ \quad \quad \operatorname{is\text{-}continuous\text{-}at}(\operatorname{subspace\text{-}ms}(\operatorname{nf\text{-}metric\text{-}space}(k), u), \operatorname{nf\text{-}metric\text{-}space}(k), f, a) \end{array}}
%
\end{prop}
%
\begin{proof}
%
The proofs: the file that holds each, the lemmas it cites, its oracles and its bill.
%
\vnbprf{\textbf{Proof.} \texttt{theorem-\allowbreak{}library/diff-\allowbreak{}on-\allowbreak{}laws.scm} --- 201 recorded steps.}
%
\vnbprf{\textbf{Lemmas.} \texttt{diff-on-normed-field}, \texttt{diff-on-open}, \texttt{diff-on-in-fun}, \texttt{diff-on-pt-in}, \texttt{diff-on-deriv-in-carr}, \texttt{nf-open-subset-carr}, \texttt{nf-metric-space-is-metric-space}, \texttt{subspace-is-metric-space}, \texttt{subspace-pts}, \texttt{rr-pos-rr-in-rr}, \texttt{rr-one-in}, \texttt{nf-norm-in-rr}, \texttt{nf-norm-nonneg}, \texttt{rr-add-in-rr}, \texttt{rr-pos-rr-of-lt}, \texttt{rr-scale-eps}, \texttt{rr-lt-of-pos-rr}, \texttt{rr-min-pos}, \texttt{subset-mem-fwd}, \texttt{fun-apply-type-c}, \texttt{nf-neg-in-carr}, \texttt{nf-add-in-carr}, \texttt{metric-dist-real}, \texttt{subspace-dist}, \texttt{metric-sym}, \texttt{nf-norm-le-add}, \texttt{nf-norm-mult}, \texttt{rr-mul-le-right}, \texttt{rr-mul-in-rr}.}
%
\vnbprf{\textbf{Oracles.} \texttt{ineq}, \texttt{crs}, \texttt{arith}.  \textbf{Bill:} \texttt{proven modulo 0}.}
%
\end{proof}
%
\noindent\textbf{Deviation from the notes.} Stated for an arbitrary normed field K rather than for CC (more general; the CC instance is got by taking K = CC-NORMED-FIELD, but no CC-specialised corollary is installed). The conclusion is continuity at a in the SUBSPACE metric on U -- is-continuous-at(subspace-ms(nf-metric-space(k), u), nf-metric-space(k), f, a) -- not continuity of a map defined near a in C; the two agree because IS-DIFF-ON forces U open, but the library does not say so.
%
\renewcommand{\thethm}{2.3}
\begin{prop}[{Sum and product rules}]\label{thm:diff-on-sum}
%
\leavevmode
\vnbpart{\texttt{diff-on-sum}}
\vnbstmt{\begin{array}{@{}l@{}} \forall k, u, f, g, a, l, m. \\ \quad \operatorname{is\text{-}normed\text{-}field}(k) \wedge \\ \quad \operatorname{is\text{-}diff\text{-}on}(k, u, f, a, l) \wedge \\ \quad \operatorname{is\text{-}diff\text{-}on}(k, u, g, a, m) \Rightarrow \\ \quad \quad \operatorname{is\text{-}diff\text{-}on}(k, u, (\lambda x \in u.\; \operatorname{add}(k)(f(x), g(x))), a, \operatorname{add}(k)(l, m)) \end{array}}
\vnbpart{\texttt{diff-on-product}}
\vnbstmt{\begin{array}{@{}l@{}} \forall k, u, f, g, a, l, m. \\ \quad \operatorname{is\text{-}normed\text{-}field}(k) \wedge \\ \quad \operatorname{is\text{-}diff\text{-}on}(k, u, f, a, l) \wedge \\ \quad \operatorname{is\text{-}diff\text{-}on}(k, u, g, a, m) \Rightarrow \\ \quad \quad \operatorname{is\text{-}diff\text{-}on}(k, u, (\lambda x \in u.\; \operatorname{mul}(k)(f(x), g(x))), a, \\ \quad \quad \quad \operatorname{add}(k)(\operatorname{mul}(k)(l, g(a)), \operatorname{mul}(k)(m, f(a)))) \end{array}}
\vnbpart{\texttt{holomorphic-sum}}
\vnbstmt{\begin{array}{@{}l@{}} \forall u, f, g, m. \\ \quad \operatorname{holomorphic\text{-}on}(u, f) \wedge \\ \quad \operatorname{holomorphic\text{-}on}(u, g) \wedge \\ \quad (m \in (u \to \mathbb{C})) \wedge \\ \quad \forall x \in u.\; (m(x) \simeq (f(x) + g(x))) \Rightarrow \\ \quad \quad \operatorname{holomorphic\text{-}on}(u, m) \end{array}}
\vnbpart{\texttt{holomorphic-product}}
\vnbstmt{\begin{array}{@{}l@{}} \forall u, f, g, m. \\ \quad \operatorname{holomorphic\text{-}on}(u, f) \wedge \\ \quad \operatorname{holomorphic\text{-}on}(u, g) \wedge \\ \quad (m \in (u \to \mathbb{C})) \wedge \\ \quad \forall x \in u.\; (m(x) \simeq (f(x) \cdot g(x))) \Rightarrow \\ \quad \quad \operatorname{holomorphic\text{-}on}(u, m) \end{array}}
%
\end{prop}
%
\begin{proof}
%
The proofs: the file that holds each, the lemmas it cites, its oracles and its bill.
%
\vnbprf{\vnbproofpart{\texttt{diff-on-sum}}}
%
\vnbprf{\textbf{Proof.} \texttt{theorem-\allowbreak{}library/diff-\allowbreak{}on-\allowbreak{}laws-\allowbreak{}2.scm} --- 173 recorded steps.}
%
\vnbprf{\textbf{Lemmas.} \texttt{diff-on-open}, \texttt{diff-on-pt-in}, \texttt{diff-on-in-fun}, \texttt{diff-on-deriv-in-carr}, \texttt{nf-metric-space-is-metric-space}, \texttt{nf-open-subset-carr}, \texttt{subspace-is-metric-space}, \texttt{subspace-pts}, \texttt{nf-open-is-set}, \texttt{subset-mem-fwd}, \texttt{nf-add-in-carr}, \texttt{nf-lam-sum-in-fun}, \texttt{nf-sum-continuous-on}, \texttt{fun-apply-type-c}, \texttt{commutative-ring-is-ring}, \texttt{normed-field-ring-view-carr}, \texttt{normed-field-ring-view-add}, \texttt{normed-field-ring-view-mul}, \texttt{normed-field-ring-view-neg}, \texttt{normed-field-ring-view-zero}, \texttt{normed-field-ring-view-one}.}
%
\vnbprf{\textbf{Oracles.} \texttt{crs}, \texttt{ineq}, \texttt{arith}.  \textbf{Bill:} \texttt{proven modulo 0}.}
%
\vnbprf{\vnbproofpart{\texttt{diff-on-product}}}
%
\vnbprf{\textbf{Proof.} \texttt{theorem-\allowbreak{}library/diff-\allowbreak{}on-\allowbreak{}laws-\allowbreak{}2.scm} --- 198 recorded steps.}
%
\vnbprf{\textbf{Lemmas.} \texttt{diff-on-open}, \texttt{diff-on-pt-in}, \texttt{diff-on-in-fun}, \texttt{diff-on-deriv-in-carr}, \texttt{diff-on-implies-continuous}, \texttt{nf-metric-space-is-metric-space}, \texttt{nf-open-subset-carr}, \texttt{subspace-is-metric-space}, \texttt{subspace-pts}, \texttt{nf-open-is-set}, \texttt{subset-mem-fwd}, \texttt{fun-apply-type-c}, \texttt{nf-mul-in-carr}, \texttt{nf-add-in-carr}, \texttt{nf-lam-prod-in-fun}, \texttt{nf-lam-const-z-in-fun}, \texttt{ms-const-continuous-on}, \texttt{nf-lam-prod-v-in-fun}, \texttt{nf-product-continuous-v}, \texttt{nf-lam-sum-in-fun}, \texttt{nf-sum-continuous-on}, \texttt{commutative-ring-is-ring}, \texttt{normed-field-ring-view-carr}, \texttt{normed-field-ring-view-add}, \texttt{normed-field-ring-view-mul}, \texttt{normed-field-ring-view-neg}, \texttt{normed-field-ring-view-zero}, \texttt{normed-field-ring-view-one}.}
%
\vnbprf{\textbf{Oracles.} \texttt{crs}, \texttt{ineq}, \texttt{arith}.  \textbf{Bill:} \texttt{proven modulo 0}.}
%
\vnbprf{\vnbproofpart{\texttt{holomorphic-sum}}}
%
\vnbprf{\textbf{Proof.} \texttt{theorem-\allowbreak{}library/holomorphic-\allowbreak{}basics-\allowbreak{}2.scm} --- 49 recorded steps.}
%
\vnbprf{\textbf{Lemmas.} \texttt{cc-is-normed-field}, \texttt{holomorphic-on-open}, \texttt{holomorphic-on-in-fun}, \texttt{holomorphic-on-diff}, \texttt{cc-nf-carr}, \texttt{diff-on-sum}, \texttt{fun-apply-type-c}, \texttt{cc-nf-add-apply}, \texttt{diff-on-transfer-ptwise-eq}.}
%
\vnbprf{\textbf{Oracles.} \texttt{arith}, \texttt{crs}, \texttt{ineq}.  \textbf{Bill:} \texttt{proven modulo 0}.}
%
\vnbprf{\vnbproofpart{\texttt{holomorphic-product}}}
%
\vnbprf{\textbf{Proof.} \texttt{theorem-\allowbreak{}library/holomorphic-\allowbreak{}basics-\allowbreak{}2.scm} --- 49 recorded steps.}
%
\vnbprf{\textbf{Lemmas.} \texttt{cc-is-normed-field}, \texttt{holomorphic-on-open}, \texttt{holomorphic-on-in-fun}, \texttt{holomorphic-on-diff}, \texttt{cc-nf-carr}, \texttt{diff-on-product}, \texttt{fun-apply-type-c}, \texttt{cc-nf-mul-apply}, \texttt{diff-on-transfer-ptwise-eq}.}
%
\vnbprf{\textbf{Oracles.} \texttt{arith}, \texttt{crs}, \texttt{ineq}.  \textbf{Bill:} \texttt{proven modulo 0}.}
%
\end{proof}
%
\noindent\textbf{Deviation from the notes.} Two separate developments, neither of which is exactly the notes' statement. diff-on-sum / diff-on-product give the derivative at a, but only for the LITERAL function VNB-LAMBDA x in U. (add K)(f(x), g(x)) (resp. the product); a pointwise-equal representative of f + g or f . g has to be transported by diff-on-transfer-ptwise-eq. holomorphic-sum / holomorphic-product do take an arbitrary sm in FUN(U, CC) that agrees pointwise, but they conclude only HOLOMORPHIC-ON(U, sm) and say nothing about the value of the derivative at a. Both diff-on laws are stated for a general normed field K.
%
\renewcommand{\thethm}{2.4}
\begin{prop}[{Chain Rule}]\label{thm:diff-on-chain}
%
\vnbstmt{\begin{array}{@{}l@{}} \forall k, u, s, f, g, a, l, m. \\ \quad \operatorname{is\text{-}normed\text{-}field}(k) \wedge \\ \quad \operatorname{is\text{-}diff\text{-}on}(k, u, f, a, l) \wedge \\ \quad \operatorname{is\text{-}diff\text{-}on}(k, s, g, f(a), m) \wedge \\ \quad \forall y \in u.\; (f(y) \in s) \Rightarrow \\ \quad \quad \operatorname{is\text{-}diff\text{-}on}(k, u, (\lambda x \in u.\; g(f(x))), a, \operatorname{mul}(k)(m, l)) \end{array}}
%
\end{prop}
%
\begin{proof}
%
The proofs: the file that holds each, the lemmas it cites, its oracles and its bill.
%
\vnbprf{\textbf{Proof.} \texttt{theorem-\allowbreak{}library/diff-\allowbreak{}on-\allowbreak{}laws-\allowbreak{}2.scm} --- 164 recorded steps.}
%
\vnbprf{\textbf{Lemmas.} \texttt{diff-on-open}, \texttt{diff-on-pt-in}, \texttt{diff-on-in-fun}, \texttt{diff-on-deriv-in-carr}, \texttt{diff-on-implies-continuous}, \texttt{nf-metric-space-is-metric-space}, \texttt{nf-open-subset-carr}, \texttt{subspace-is-metric-space}, \texttt{subspace-pts}, \texttt{nf-open-is-set}, \texttt{subset-mem-fwd}, \texttt{nf-mul-in-carr}, \texttt{nf-lam-comp-z-in-fun}, \texttt{ms-corestrict-continuous}, \texttt{ms-compose-continuous-on}, \texttt{nf-lam-prod-v-in-fun}, \texttt{nf-product-continuous-v}, \texttt{fun-apply-type-c}, \texttt{commutative-ring-is-ring}, \texttt{normed-field-ring-view-carr}, \texttt{normed-field-ring-view-add}, \texttt{normed-field-ring-view-mul}, \texttt{normed-field-ring-view-neg}, \texttt{normed-field-ring-view-zero}, \texttt{normed-field-ring-view-one}.}
%
\vnbprf{\textbf{Oracles.} \texttt{crs}, \texttt{ineq}, \texttt{arith}.  \textbf{Bill:} \texttt{proven modulo 0}.}
%
\end{proof}
%
\noindent\textbf{Deviation from the notes.} The library adds hypotheses the notes leave implicit: both U and the target set S must be open (carried by IS-DIFF-ON), and forall y in U. f(y) in S is an explicit antecedent (the notes only require g to be differentiable at f(a)). The composite is the literal VNB-LAMBDA x in U. g(f(x)), so a pointwise-equal representative needs diff-on-transfer-ptwise-eq; the derivative is written (mul K)(m, l) in the normed field K. Stated for a general K, not specialised to CC.
%
\renewcommand{\thethm}{2.5}
\begin{lem}[{Root Test}]\label{thm:root-test-converges}
%
\leavevmode
\vnbpart{\texttt{root-test-converges}}
\vnbstmt{\begin{array}{@{}l@{}} \forall h \in (\mathbb{N} \to \mathbb{R}). \\ \quad \forall k \in \mathbb{N}.\; (0 \leq h(k)) \wedge \\ \quad (\operatorname{elimsup}((\lambda j \in \mathbb{N}.\; {h(j)}^{{j}^{\ast\!-1}})) < 1) \Rightarrow \\ \quad \quad \operatorname{series\text{-}converges}(h) \end{array}}
\vnbpart{\texttt{root-test-diverges}}
\vnbstmt{\begin{array}{@{}l@{}} \forall h \in (\mathbb{N} \to \mathbb{R}). \\ \quad \forall k \in \mathbb{N}.\; (0 \leq h(k)) \wedge \\ \quad (1 < \operatorname{elimsup}((\lambda j \in \mathbb{N}.\; {h(j)}^{{j}^{\ast\!-1}}))) \Rightarrow \\ \quad \quad \neg \operatorname{series\text{-}converges}(h) \end{array}}
%
\end{lem}
%
\begin{proof}
%
The proofs: the file that holds each, the lemmas it cites, its oracles and its bill.
%
\vnbprf{\vnbproofpart{\texttt{root-test-converges}}}
%
\vnbprf{\textbf{Proof.} \texttt{theorem-\allowbreak{}library/limsup-\allowbreak{}tests.scm} --- 148 recorded steps.}
%
\vnbprf{\textbf{Lemmas.} \texttt{fun-apply-type-c}, \texttt{nn-in-rr}, \texttt{recip-star-in-rr}, \texttt{recip-star-nn-nonneg}, \texttt{rpow-star-nonneg-in-rr}, \texttt{rpow-star-nonneg-nonneg}, \texttt{rr-one-in}, \texttt{elimsup-in-rr-pos-star}, \texttt{rr-pos-star-nonneg}, \texttt{rr-pos-star-le-real-in-rr}, \texttt{rr-add-in-rr}, \texttt{rr-mul-in-rr}, \texttt{elimsup-eventually-below}, \texttt{nn-succ-closed}, \texttt{nn-le-succ}, \texttt{nn-one-le-succ}, \texttt{nn-le-trans-guarded}, \texttt{power-real-closed}, \texttt{rr-power-nonneg}, \texttt{rpow-star-root-power}, \texttt{rr-power-mono-base}, \texttt{series-eventually-geometric-converges}.}
%
\vnbprf{\textbf{Oracles.} \texttt{arith}, \texttt{ineq}, \texttt{crs}, \texttt{sos}.  \textbf{Bill:} \texttt{proven modulo 0}.}
%
\vnbprf{\vnbproofpart{\texttt{root-test-diverges}}}
%
\vnbprf{\textbf{Proof.} \texttt{theorem-\allowbreak{}library/limsup-\allowbreak{}tests.scm} --- 129 recorded steps.}
%
\vnbprf{\textbf{Lemmas.} \texttt{fun-apply-type-c}, \texttt{nn-in-rr}, \texttt{recip-star-in-rr}, \texttt{recip-star-nn-nonneg}, \texttt{rpow-star-nonneg-in-rr}, \texttt{rpow-star-nonneg-nonneg}, \texttt{rr-one-in}, \texttt{elimsup-frequently-above}, \texttt{nn-succ-closed}, \texttt{nn-le-succ}, \texttt{nn-one-le-succ}, \texttt{nn-le-trans-guarded}, \texttt{recip-star-nn-pos}, \texttt{rpow-star-zero-base}, \texttt{rpow-star-root-power}, \texttt{rr-power-mono-base}, \texttt{rr-power-one}, \texttt{power-real-closed}, \texttt{rr-zero-lt-one}, \texttt{series-frequently-large-diverges}.}
%
\vnbprf{\textbf{Oracles.} \texttt{ineq}, \texttt{arith}, \texttt{crs}, \texttt{sos}.  \textbf{Bill:} \texttt{proven modulo 0}.}
%
\end{proof}
%
\noindent\textbf{Deviation from the notes.} UPDATED 2026-09-24 (batch 25): PROVEN as the notes state it, both halves, for a non-negative real series, with the lim sup taken in [0, +inf] (ELIMSUP; every lim sup the notes take is of a non-negative sequence), theorem-library/limsup-tests.scm. EARLIER NOTE: No root test in any form. The library has no lim sup for real sequences: ELIMSUP / ELIMINF exist only for sequences into RR-POS-STAR = [0, +inf] and carry no theory beyond the definition of ECONVERGES-TO. The nearest convergence criteria for a non-negative real series are comparison-test (termwise domination by a convergent series) and geometric-series-converges; the divergence half has no analogue at all.
%
\renewcommand{\thethm}{2.7}
\begin{prop}[{Radius of convergence by the root test}]\label{thm:cps-radius-formula}
%
\vnbmissing{cps-radius-formula}
%
\end{prop}
%
\begin{proof}
%
The proofs: the file that holds each, the lemmas it cites, its oracles and its bill.
%
\vnbprf{No theorem named \texttt{cps-radius-formula} is installed in the library.}
%
\end{proof}
%
\noindent\textbf{Deviation from the notes.} UPDATED 2026-09-24 (batch 25): PROVEN in the notes' three cases (lim sup 0 gives +inf, +inf gives 0, finite positive gives the reciprocal), for the DEFINED radius CPS-RADIUS(cf) = the sup in [0, +inf] of the r $>$= 0 with sum |a\_k| r\textasciicircum{}k convergent (structure-library/cc-power-series.scm, quoting the notes p. 25; defining the radius by the formula itself would make 2.7 a tautology). EARLIER NOTE: The library has no notion of RADIUS OF CONVERGENCE at all, and no complex power series. Its power-series layer (structure-library/power-series.scm) is REAL and pointwise only: PS-PARTIAL-SUM(coef, x, k), PS-CONVERGES-AT(coef, x) and PS-ABSOLUTELY-CONVERGES-AT for coef in FUN(NN, RR) and x in RR, with the bridges ps-converges-as-series and ps-partial-sum-as-series. There is no lim sup of |a\_k|\textasciicircum{}(1/k) to state the formula with.
%
\renewcommand{\thethm}{2.8}
\begin{lem}[{Normal convergence inside the radius}]\label{thm:cps-normally-convergent-inside}
%
\leavevmode
\vnbpart{\texttt{cps-normally-convergent-inside}}
\vnbmissing{cps-normally-convergent-inside}
\vnbpart{\texttt{cps-converges-inside}}
\vnbmissing{cps-converges-inside}
%
\end{lem}
%
\begin{proof}
%
The proofs: the file that holds each, the lemmas it cites, its oracles and its bill.
%
\vnbprf{\vnbproofpart{\texttt{cps-normally-convergent-inside}}}
%
\vnbprf{No theorem named \texttt{cps-normally-convergent-inside} is installed in the library.}
%
\vnbprf{\vnbproofpart{\texttt{cps-converges-inside}}}
%
\vnbprf{No theorem named \texttt{cps-converges-inside} is installed in the library.}
%
\end{proof}
%
\noindent\textbf{Deviation from the notes.} UPDATED 2026-09-24 (batch 25): PROVEN: normal convergence (IS-NORMALLY-CONVERGENT, the notes p. 26, with CC-SUP-NORM = formula (29) allowed to be +inf) on the closed ball of radius r $<$ R, hence absolute, pointwise and uniform convergence there (M-test), theorem-library/cc-power-series-laws.scm. EARLIER NOTE: No radius of convergence, no supremum norm on complex-valued functions, and no notion of normal convergence (the Weierstrass M-test is not in the library). The nearest machinery is uniform convergence of sequences of REAL functions on a metric space (CONVERGES-UNIFORMLY, unif-cauchy-has-uniform-limit, uniform-limit-continuous) and absolute convergence of a real series (series-abs-converges, ps-absolute-implies-convergent); none of it is applied to a power series on a closed ball.
%
\renewcommand{\thethm}{2.9}
\begin{lem}[{Polynomially bounded factors do not shrink the radius}]\label{thm:cps-radius-poly-factor}
%
\vnbmissing{cps-radius-poly-factor}
%
\end{lem}
%
\begin{proof}
%
The proofs: the file that holds each, the lemmas it cites, its oracles and its bill.
%
\vnbprf{No theorem named \texttt{cps-radius-poly-factor} is installed in the library.}
%
\end{proof}
%
\noindent\textbf{Deviation from the notes.} UPDATED 2026-09-24 (batch 25): PROVEN from the definition of the radius (the notes' route needs lim n\textasciicircum{}(1/n) = 1, absent; the bound k\textasciicircum{}m q\textasciicircum{}k -$>$ 0 for 0 $<$= q $<$ 1 goes through the m-th root of q instead). EARLIER NOTE: Nothing near: no radius of convergence, no comparison of two power series by their coefficients.
%
\renewcommand{\thethm}{2.10}
\begin{prop}[{A power series is holomorphic and differentiates termwise}]\label{thm:cps-holomorphic-ball}
%
\leavevmode
\vnbpart{\texttt{cps-holomorphic-ball}}
\vnbstmt{\begin{array}{@{}l@{}} \forall f \in (\mathbb{N} \to \mathbb{C}), r \in \mathbb{R}. \\ \quad ((0 < r) \wedge (r < \operatorname{cps\text{-}radius}(f))) \Rightarrow  \\ \quad \quad \operatorname{holomorphic\text{-}on}(\operatorname{ball}(\operatorname{cc\text{-}ms}, 0, r), \\ \quad \quad \quad (\lambda z \in \operatorname{ball}(\operatorname{cc\text{-}ms}, 0, r).\; \\ \quad \quad \quad \quad \operatorname{cc\text{-}series\text{-}limit}((\lambda k \in \mathbb{N}.\; (f(k) \cdot \operatorname{power}(z, k)))))) \wedge \\ \quad \quad \operatorname{cps\text{-}radius}((\lambda n \in \mathbb{N}.\; ((n + 1) \cdot f((n + 1))))) \\ \quad \quad =\; \operatorname{cps\text{-}radius}(f) \wedge \\ \quad \quad \forall a \in \operatorname{ball}(\operatorname{cc\text{-}ms}, 0, r). \\ \quad \quad \quad \operatorname{is\text{-}diff\text{-}on}(\operatorname{cc\text{-}normed\text{-}field}, \operatorname{ball}(\operatorname{cc\text{-}ms}, 0, r), \\ \quad \quad \quad \quad (\lambda z \in \operatorname{ball}(\operatorname{cc\text{-}ms}, 0, r).\; \\ \quad \quad \quad \quad \quad \operatorname{cc\text{-}series\text{-}limit}((\lambda k \in \mathbb{N}.\; (f(k) \cdot \operatorname{power}(z, k))))), \\ \quad \quad \quad \quad a, \\ \quad \quad \quad \quad \operatorname{cc\text{-}series\text{-}limit}((\lambda k \in \mathbb{N}.\; (((k + 1) \cdot f((k + 1))) \cdot \operatorname{power}(a, k))))) \end{array}}
\vnbpart{\texttt{cps-radius-derived}}
\vnbstmt{\begin{array}{@{}l@{}} \forall f, a \in (\mathbb{N} \to \mathbb{C}). \\ \quad \forall j \in \mathbb{N}.\; (a(j) = ((j + 1) \cdot f((j + 1)))) \Rightarrow  \\ \quad \quad (\operatorname{cps\text{-}radius}(a) = \operatorname{cps\text{-}radius}(f)) \end{array}}
\vnbpart{\texttt{cps-diff-on-ball}}
\vnbstmt{\begin{array}{@{}l@{}} \forall f, a \in (\mathbb{N} \to \mathbb{C}), r \in \mathbb{R}, \mathit{cpda} \in \operatorname{ball}(\operatorname{cc\text{-}ms}, 0, r). \\ \quad \forall j \in \mathbb{N}.\; (a(j) = ((j + 1) \cdot f((j + 1)))) \wedge \\ \quad (0 < r) \wedge \\ \quad (r < \operatorname{cps\text{-}radius}(f)) \Rightarrow \\ \quad \quad \operatorname{is\text{-}diff\text{-}on}(\operatorname{cc\text{-}normed\text{-}field}, \operatorname{ball}(\operatorname{cc\text{-}ms}, 0, r), \\ \quad \quad \quad (\lambda z \in \operatorname{ball}(\operatorname{cc\text{-}ms}, 0, r).\; \\ \quad \quad \quad \quad \operatorname{cc\text{-}series\text{-}limit}((\lambda k \in \mathbb{N}.\; (f(k) \cdot \operatorname{power}(z, k))))), \\ \quad \quad \quad \mathit{cpda}, \\ \quad \quad \quad \operatorname{cc\text{-}series\text{-}limit}((\lambda k \in \mathbb{N}.\; (a(k) \cdot \operatorname{power}(\mathit{cpda}, k))))) \end{array}}
%
\end{prop}
%
\begin{proof}
%
The proofs: the file that holds each, the lemmas it cites, its oracles and its bill.
%
\vnbprf{\vnbproofpart{\texttt{cps-holomorphic-ball}}}
%
\vnbprf{\textbf{Proof.} \texttt{theorem-\allowbreak{}library/cc-\allowbreak{}power-\allowbreak{}series-\allowbreak{}deriv.scm} --- 283 recorded steps.}
%
\vnbprf{\textbf{Lemmas.} \texttt{cc-is-metric-space}, \texttt{cc-zero-in}, \texttt{rr-lt-implies-le}, \texttt{rr-pos-ne-zero}, \texttt{neq-sym}, \texttt{ball-is-open}, \texttt{cc-ms-open-iff}, \texttt{ball-is-set}, \texttt{ball-in-closed-ball}, \texttt{nn-is-set}, \texttt{nn-succ-closed}, \texttt{fun-apply-type-c}, \texttt{nn-in-rr}, \texttt{rr-subset-cc}, \texttt{cc-mul-closed}, \texttt{cps-radius-derived}, \texttt{cps-diff-on-ball}, \texttt{subset-mem-fwd}, \texttt{cc-closed-ball-origin}, \texttt{power-typing-nonneg}, \texttt{cps-converges-closed-ball}, \texttt{cc-series-limit-in-cc}, \texttt{rr-one-in}, \texttt{rr-zero-in}, \texttt{cc-series-limit-linear}.}
%
\vnbprf{\textbf{Oracles.} \texttt{crs}, \texttt{arith}, \texttt{ineq}, \texttt{sos}.  \textbf{Bill:} \texttt{proven modulo 0}.}
%
\vnbprf{\vnbproofpart{\texttt{cps-radius-derived}}}
%
\vnbprf{\textbf{Proof.} \texttt{theorem-\allowbreak{}library/cc-\allowbreak{}power-\allowbreak{}series-\allowbreak{}deriv.scm} --- 417 recorded steps.}
%
\vnbprf{\textbf{Lemmas.} \texttt{nn-is-set}, \texttt{nn-succ-closed}, \texttt{fun-apply-type-c}, \texttt{nn-in-rr}, \texttt{rr-subset-cc}, \texttt{cc-mul-closed}, \texttt{cps-radius-in}, \texttt{rr-one-in}, \texttt{nn-one-in}, \texttt{nn-zero-le}, \texttt{rr-magnitude-is-abs}, \texttt{rr-abs-of-nonneg}, \texttt{nn-one-is-succ-zero}, \texttt{nn-zero-in}, \texttt{power-succ}, \texttt{power-zero}, \texttt{cps-radius-factor}, \texttt{cc-magnitude-closed}, \texttt{cps-shift-conv}, \texttt{rr-lt-implies-le}, \texttt{cc-magnitude-nonneg}, \texttt{power-real-closed}, \texttt{rr-power-nonneg}, \texttt{rr-mul-in-rr}, \texttt{rr-leq-mul-nonneg}, \texttt{cc-magnitude-mul}, \texttt{nn-one-le-succ}, \texttt{rr-le-scale-nonneg}, \texttt{rr-mul-le-right}, \texttt{comparison-test}, \texttt{cps-radius-mono}, \texttt{rr-pos-star-le-trans}, \texttt{rr-pos-star-le-antisymm}.}
%
\vnbprf{\textbf{Oracles.} \texttt{ineq}, \texttt{crs}, \texttt{arith}, \texttt{sos}.  \textbf{Bill:} \texttt{proven modulo 0}.}
%
\vnbprf{\vnbproofpart{\texttt{cps-diff-on-ball}}}
%
\vnbprf{\textbf{Proof.} \texttt{theorem-\allowbreak{}library/cc-\allowbreak{}power-\allowbreak{}series-\allowbreak{}deriv.scm} --- 1170 recorded steps.}
%
\vnbprf{\textbf{Lemmas.} \texttt{cc-is-metric-space}, \texttt{cc-zero-in}, \texttt{rr-lt-implies-le}, \texttt{rr-pos-ne-zero}, \texttt{neq-sym}, \texttt{ball-is-open}, \texttt{cc-ms-open-iff}, \texttt{ball-is-set}, \texttt{ball-in-closed-ball}, \texttt{cps-radius-derived}, \texttt{subset-mem-fwd}, \texttt{cc-closed-ball-origin}, \texttt{nn-is-set}, \texttt{fun-apply-type-c}, \texttt{power-typing-nonneg}, \texttt{cc-mul-closed}, \texttt{cps-converges-closed-ball}, \texttt{cc-series-limit-in-cc}, \texttt{cps-radius-approx}, \texttt{nn-in-rr}, \texttt{rr-subset-cc}, \texttt{rr-one-in}, \texttt{nn-two-is-succ-succ-zero}, \texttt{nn-zero-in}, \texttt{nn-succ-closed}, \texttt{rr-mul-in-rr}, \texttt{nn-zero-le}, \texttt{rr-leq-mul-nonneg}, \texttt{rr-magnitude-is-abs}, \texttt{rr-abs-of-nonneg}, \texttt{power-succ}, \texttt{power-zero}, \texttt{cps-factor-conv}, \texttt{cc-magnitude-closed}, \texttt{cc-magnitude-nonneg}, \texttt{power-real-closed}, \texttt{rr-power-nonneg}, \texttt{series-shift-converges-iff}, \texttt{series-partial-sum-bounded}, \texttt{series-partial-sum-nonneg}, \texttt{series-partial-sum-in-rr}, \texttt{rr-mul-pos}, \texttt{rr-recip-closed}, \texttt{rr-recip-inverse}, \texttt{rr-recip-pos}, \texttt{cc-sub-in-cc}, \texttt{cc-partial-error-limit}, \texttt{cc-series-partial-sum-in-cc}, \texttt{cps-partial-error-bound}, \texttt{rr-le-scale-nonneg}, \texttt{rr-mul-le-cancel-pos}, \texttt{cc-limit-magnitude-le}, \texttt{cc-diff-on-quadratic}.}
%
\vnbprf{\textbf{Oracles.} \texttt{crs}, \texttt{ineq}, \texttt{arith}, \texttt{sos}.  \textbf{Bill:} \texttt{proven modulo 0}.}
%
\end{proof}
%
\noindent\textbf{Deviation from the notes.} UPDATED 2026-09-24 (batch 26-A, theorem-library/cc-power-series-deriv.scm): PROVEN with one convention -- the library has no ball of radius +inf, so the statement is made on B(0, r) for every real 0 $<$ r $<$ R (which covers B(0, R) for finite R and CC for R = +inf): f is holomorphic on B(0, r), the derived series sum (k+1) a\_(k+1) z\textasciicircum{}k has the same radius, and f'(u) is its sum. The proof avoids the notes' double sum and lim n\textasciicircum{}(1/n) = 1 (absent): a one-step recurrence gives the second-order estimate, a radius-comparison lemma and Lemma 2.9 with b\_k = k give the radius. EARLIER NOTE: Nothing near, in either the complex or the real case. The library's power series are real and pointwise (PS-CONVERGES-AT); it has no theorem that the sum function of a power series is differentiable, no termwise-differentiated series, and no statement that the two series share a radius. HOLOMORPHIC-ON is available (holomorphic-basics.scm) but the only holomorphic functions exhibited are constants, the identity, powers, sums and products (holomorphic-const, holomorphic-identity, holomorphic-power, holomorphic-sum, holomorphic-product).
%
\renewcommand{\thethm}{2.12}
\begin{prop}[{Principle of Analytic Continuation}]\label{thm:complex-analysis-2.12}
%
\textit{Let U be a connected open set. If f is analytic in U and the set of zeros of f in U has an accumulation point in U, then f is identically 0 on U.}
%
Not in the library.
%
\end{prop}
%
\begin{proof}
%
The proofs: the file that holds each, the lemmas it cites, its oracles and its bill.
%
\vnbprf{There is no library statement, and so no proof.}
%
\end{proof}
%
\noindent\textbf{Deviation from the notes.} Nothing near. The library has no notion of an ANALYTIC function (no local power-series representation), and no connectedness of a subset of a metric space -- IS-CONNECTED does not exist in THEOREMS.md. Accumulation point is available only indirectly through CLOSURE / CLUSTER-POINT.
%
\renewcommand{\thethm}{2.13}
\begin{cor}[{Finitely many zeros in a closed ball}]\label{thm:complex-analysis-2.13}
%
\textit{Let U be a connected open set. If f is analytic in U and not identically 0, then f has only finitely many zeros within any closed ball B contained in U.}
%
Not in the library.
%
\end{cor}
%
\begin{proof}
%
The proofs: the file that holds each, the lemmas it cites, its oracles and its bill.
%
\vnbprf{There is no library statement, and so no proof.}
%
\end{proof}
%
\noindent\textbf{Deviation from the notes.} Nothing near: no analytic functions, no connectedness. CLOSED-BALL and card(...) in NN (finiteness) both exist, so only the statement is missing, not the vocabulary for its conclusion.
%
\renewcommand{\thethm}{2.14}
\begin{prop}[{Taylor coefficients of an analytic function}]\label{thm:taylor-coefficients-at}
%
\leavevmode
\vnbpart{\texttt{taylor-coefficients-at}}
\vnbstmt{\begin{array}{@{}l@{}} \forall f, u.; \forall c \in (\mathbb{N} \to \mathbb{C}), a \in \mathbb{C}, r \in \mathbb{R}, k \in \mathbb{N}. \\ \quad (f \in (u \to \mathbb{C})) \wedge \\ \quad (0 < r) \wedge \\ \quad (r < \operatorname{cps\text{-}radius}(c)) \wedge \\ \quad (\operatorname{ball}(\operatorname{cc\text{-}ms}, a, r) \subseteq u) \wedge \\ \quad \forall z \in \operatorname{ball}(\operatorname{cc\text{-}ms}, a, r). \\ \quad \quad f(z) \\ \quad \quad =\; \operatorname{cc\text{-}series\text{-}limit}((\lambda \mathit{cpsk} \in \mathbb{N}.\; \\ \quad \quad \quad \quad (c(\mathit{cpsk}) \cdot \operatorname{power}((z - a), \mathit{cpsk})))) \Rightarrow \\ \quad \quad \operatorname{restrict}(f, \operatorname{ball}(\operatorname{cc\text{-}ms}, a, r))^{(k)}(a) \\ \quad \quad =\; (\operatorname{factorial}(k) \cdot c(k)) \end{array}}
\vnbpart{\texttt{taylor-coefficients}}
\vnbstmt{\begin{array}{@{}l@{}} \forall f \in (\mathbb{N} \to \mathbb{C}), r \in \mathbb{R}, k \in \mathbb{N}. \\ \quad ((0 < r) \wedge (r < \operatorname{cps\text{-}radius}(f))) \Rightarrow  \\ \quad \quad (\lambda z \in \operatorname{ball}(\operatorname{cc\text{-}ms}, 0, r).\; \operatorname{cc\text{-}series\text{-}limit}((\lambda \mathit{cpsk} \in \mathbb{N}.\; (f(\mathit{cpsk}) \cdot \operatorname{power}(z, \mathit{cpsk})))))^{(k)}(0) \\ \quad \quad =\; (\operatorname{factorial}(k) \cdot f(k)) \end{array}}
\vnbpart{\texttt{taylor-coefficient-formula}}
\vnbstmt{\begin{array}{@{}l@{}} \forall f \in (\mathbb{N} \to \mathbb{C}), r \in \mathbb{R}, k \in \mathbb{N}. \\ \quad ((0 < r) \wedge (r < \operatorname{cps\text{-}radius}(f))) \Rightarrow  \\ \quad \quad f(k) \\ \quad \quad =\; ((\lambda z \in \operatorname{ball}(\operatorname{cc\text{-}ms}, 0, r).\; \operatorname{cc\text{-}series\text{-}limit}((\lambda \mathit{cpsk} \in \mathbb{N}.\; (f(\mathit{cpsk}) \cdot \operatorname{power}(z, \mathit{cpsk})))))^{(k)}(0) \cdot \operatorname{recip}(\operatorname{factorial}(k))) \end{array}}
\vnbpart{\texttt{cps-translate-nth-deriv}}
\vnbstmt{\begin{array}{@{}l@{}} \forall c \in (\mathbb{N} \to \mathbb{C}), a \in \mathbb{C}, r \in \mathbb{R}, k \in \mathbb{N}, g \in (\mathbb{N} \to \mathbb{C}). \\ \quad (0 < r) \wedge \\ \quad (r < \operatorname{cps\text{-}radius}(c)) \wedge \\ \quad \forall n \in \mathbb{N}. \\ \quad \quad g(n) \\ \quad \quad =\; ((\operatorname{factorial}((n + k)) \cdot \operatorname{recip}(\operatorname{factorial}(n))) \cdot c((n + k))) \Rightarrow \\ \quad \quad (\lambda z \in \operatorname{ball}(\operatorname{cc\text{-}ms}, a, r).\; \operatorname{cc\text{-}series\text{-}limit}((\lambda \mathit{cpsk} \in \mathbb{N}.\; (c(\mathit{cpsk}) \cdot \operatorname{power}((z - a), \mathit{cpsk})))))^{(k)} \\ \quad \quad =\; (\lambda z \in \operatorname{ball}(\operatorname{cc\text{-}ms}, a, r).\; \\ \quad \quad \quad \operatorname{cc\text{-}series\text{-}limit}((\lambda \mathit{cpsk} \in \mathbb{N}.\; \\ \quad \quad \quad \quad \quad (g(\mathit{cpsk}) \cdot \operatorname{power}((z - a), \mathit{cpsk}))))) \end{array}}
\vnbpart{\texttt{analytic-holomorphic-ball}}
\vnbstmt{\begin{array}{@{}l@{}} \forall f, u, a. \\ \quad f \text{ is analytic at } a \Rightarrow  \\ \quad \quad \exists r \in \mathbb{R}. \\ \quad \quad \quad (0 < r) \wedge \\ \quad \quad \quad (\operatorname{ball}(\operatorname{cc\text{-}ms}, a, r) \subseteq u) \wedge \\ \quad \quad \quad \operatorname{holomorphic\text{-}on}(\operatorname{ball}(\operatorname{cc\text{-}ms}, a, r), \operatorname{restrict}(f, \operatorname{ball}(\operatorname{cc\text{-}ms}, a, r))) \end{array}}
\vnbpart{\texttt{cps-taylor-coefficients}}
\vnbstmt{\begin{array}{@{}l@{}} \forall f \in (\mathbb{N} \to \mathbb{C}), k \in \mathbb{N}, g, h \in (\mathbb{N} \to \mathbb{C}), r \in \mathbb{R}. \\ \quad \forall n \in \mathbb{N}. \\ \quad \quad g(n) \\ \quad \quad =\; ((\operatorname{factorial}((n + k)) \cdot \operatorname{recip}(\operatorname{factorial}(n))) \cdot f((n + k))) \wedge \\ \quad \forall n \in \mathbb{N}. \\ \quad \quad h(n) \\ \quad \quad =\; ((\operatorname{factorial}((n + (k + 1))) \cdot \operatorname{recip}(\operatorname{factorial}(n))) \cdot f((n + (k + 1)))) \wedge \\ \quad (0 < r) \wedge \\ \quad (r < \operatorname{cps\text{-}radius}(f)) \Rightarrow \\ \quad \quad \operatorname{holomorphic\text{-}on}(\operatorname{ball}(\operatorname{cc\text{-}ms}, 0, r), \\ \quad \quad \quad (\lambda z \in \operatorname{ball}(\operatorname{cc\text{-}ms}, 0, r).\; \\ \quad \quad \quad \quad \operatorname{cc\text{-}series\text{-}limit}((\lambda \mathit{cpsk} \in \mathbb{N}.\; \\ \quad \quad \quad \quad \quad \quad (g(\mathit{cpsk}) \cdot \operatorname{power}(z, \mathit{cpsk})))))) \wedge \\ \quad \quad (\lambda z \in \operatorname{ball}(\operatorname{cc\text{-}ms}, 0, r).\; \operatorname{cc\text{-}series\text{-}limit}((\lambda \mathit{cpsk} \in \mathbb{N}.\; (g(\mathit{cpsk}) \cdot \operatorname{power}(z, \mathit{cpsk})))))(0) \\ \quad \quad =\; (\operatorname{factorial}(k) \cdot f(k)) \wedge \\ \quad \quad \forall a \in \operatorname{ball}(\operatorname{cc\text{-}ms}, 0, r). \\ \quad \quad \quad \operatorname{deriv\text{-}on}(\operatorname{cc\text{-}normed\text{-}field}, \operatorname{ball}(\operatorname{cc\text{-}ms}, 0, r), \\ \quad \quad \quad \quad (\lambda z \in \operatorname{ball}(\operatorname{cc\text{-}ms}, 0, r).\; \\ \quad \quad \quad \quad \quad \operatorname{cc\text{-}series\text{-}limit}((\lambda \mathit{cpsk} \in \mathbb{N}.\; \\ \quad \quad \quad \quad \quad \quad \quad (g(\mathit{cpsk}) \cdot \operatorname{power}(z, \mathit{cpsk}))))), \\ \quad \quad \quad \quad a) \\ \quad \quad \quad =\; (\lambda z \in \operatorname{ball}(\operatorname{cc\text{-}ms}, 0, r).\; \operatorname{cc\text{-}series\text{-}limit}((\lambda \mathit{cpsk} \in \mathbb{N}.\; (h(\mathit{cpsk}) \cdot \operatorname{power}(z, \mathit{cpsk})))))(a) \end{array}}
%
\end{prop}
%
\begin{proof}
%
The proofs: the file that holds each, the lemmas it cites, its oracles and its bill.
%
\vnbprf{\vnbproofpart{\texttt{taylor-coefficients-at}}}
%
\vnbprf{\textbf{Proof.} \texttt{theorem-\allowbreak{}library/analytic-\allowbreak{}log-\allowbreak{}laws.scm} --- 329 recorded steps.}
%
\vnbprf{\textbf{Lemmas.} \texttt{cc-is-metric-space}, \texttt{cc-is-normed-field}, \texttt{cc-zero-in}, \texttt{nn-zero-in}, \texttt{rr-zero-in}, \texttt{rr-lt-irrefl}, \texttt{ball-is-set}, \texttt{ball-is-open}, \texttt{cc-ms-open-iff}, \texttt{ball-center-in}, \texttt{alg-shift-in-ball}, \texttt{cc-sub-in-cc}, \texttt{cps-holomorphic-ball}, \texttt{holomorphic-on-in-fun}, \texttt{fun-apply-type-c}, \texttt{restrict-in-fun}, \texttt{fun-codomain-iff}, \texttt{restrict-apply}, \texttt{fun-domain-extensionality}, \texttt{nn-add-closed}, \texttt{factorial-real-pos}, \texttt{recip-factorial-in-rr}, \texttt{rr-mul-in-rr}, \texttt{rr-subset-cc}, \texttt{cc-mul-closed}, \texttt{nn-is-set}, \texttt{cps-translate-nth-deriv}, \texttt{nn-succ-closed}, \texttt{cps-taylor-coefficients}.}
%
\vnbprf{\textbf{Oracles.} \texttt{ineq}, \texttt{crs}, \texttt{arith}, \texttt{sos}.  \textbf{Bill:} \texttt{proven modulo 0}.}
%
\vnbprf{\vnbproofpart{\texttt{taylor-coefficients}}}
%
\vnbprf{\textbf{Proof.} \texttt{theorem-\allowbreak{}library/nth-\allowbreak{}deriv-\allowbreak{}laws.scm} --- 178 recorded steps.}
%
\vnbprf{\textbf{Lemmas.} \texttt{cc-is-metric-space}, \texttt{cc-zero-in}, \texttt{nn-zero-in}, \texttt{rr-zero-in}, \texttt{ball-is-set}, \texttt{nn-add-closed}, \texttt{factorial-real-pos}, \texttt{recip-factorial-in-rr}, \texttt{rr-mul-in-rr}, \texttt{rr-subset-cc}, \texttt{fun-apply-type-c}, \texttt{cc-mul-closed}, \texttt{nn-is-set}, \texttt{cps-nth-deriv}, \texttt{nn-succ-closed}, \texttt{cps-taylor-coefficients}.}
%
\vnbprf{\textbf{Oracles.} \texttt{arith}, \texttt{crs}, \texttt{ineq}, \texttt{sos}.  \textbf{Bill:} \texttt{proven modulo 0}.}
%
\vnbprf{\vnbproofpart{\texttt{taylor-coefficient-formula}}}
%
\vnbprf{\textbf{Proof.} \texttt{theorem-\allowbreak{}library/nth-\allowbreak{}deriv-\allowbreak{}laws.scm} --- 99 recorded steps.}
%
\vnbprf{\textbf{Lemmas.} \texttt{cc-is-metric-space}, \texttt{cc-zero-in}, \texttt{nn-zero-in}, \texttt{rr-zero-in}, \texttt{ball-is-set}, \texttt{nn-add-closed}, \texttt{factorial-real-pos}, \texttt{recip-factorial-in-rr}, \texttt{rr-mul-in-rr}, \texttt{rr-subset-cc}, \texttt{fun-apply-type-c}, \texttt{cc-mul-closed}, \texttt{nn-is-set}, \texttt{cps-nth-deriv}, \texttt{cps-coefficient-formula}.}
%
\vnbprf{\textbf{Oracles.} \texttt{arith}, \texttt{crs}, \texttt{ineq}, \texttt{sos}.  \textbf{Bill:} \texttt{proven modulo 0}.}
%
\vnbprf{\vnbproofpart{\texttt{cps-translate-nth-deriv}}}
%
\vnbprf{\textbf{Proof.} \texttt{theorem-\allowbreak{}library/analytic-\allowbreak{}log-\allowbreak{}laws.scm} --- 319 recorded steps.}
%
\vnbprf{\textbf{Lemmas.} \texttt{cc-is-metric-space}, \texttt{cc-is-normed-field}, \texttt{cc-zero-in}, \texttt{nn-zero-in}, \texttt{rr-zero-in}, \texttt{rr-lt-irrefl}, \texttt{ball-is-set}, \texttt{ball-is-open}, \texttt{cc-ms-open-iff}, \texttt{fun-codomain-iff}, \texttt{fun-domain-extensionality}, \texttt{nn-add-zero}, \texttt{rr-recip-factorial}, \texttt{fun-apply-type-c}, \texttt{nn-succ-closed}, \texttt{nn-add-closed}, \texttt{factorial-real-pos}, \texttt{recip-factorial-in-rr}, \texttt{rr-mul-in-rr}, \texttt{rr-subset-cc}, \texttt{cc-mul-closed}, \texttt{nn-is-set}, \texttt{cps-kth-derived-radius}, \texttt{cps-kth-derived-step}, \texttt{alg-shift-in-ball}, \texttt{cc-sub-in-cc}, \texttt{cps-holomorphic-ball}, \texttt{holomorphic-on-in-fun}, \texttt{cps-translate-diff}, \texttt{cc-deriv-on-of-is-diff-on}.}
%
\vnbprf{\textbf{Oracles.} \texttt{ineq}, \texttt{crs}, \texttt{arith}, \texttt{sos}.  \textbf{Bill:} \texttt{proven modulo 0}.}
%
\vnbprf{\vnbproofpart{\texttt{analytic-holomorphic-ball}}}
%
\vnbprf{\textbf{Proof.} \texttt{theorem-\allowbreak{}library/analytic-\allowbreak{}log-\allowbreak{}laws.scm} --- 43 recorded steps.}
%
\vnbprf{\textbf{Lemmas.} \texttt{cc-is-metric-space}, \texttt{subset-mem-fwd}, \texttt{analytic-rep-holomorphic}.}
%
\vnbprf{\textbf{Oracles.} \texttt{arith}, \texttt{crs}, \texttt{ineq}, \texttt{sos}.  \textbf{Bill:} \texttt{proven modulo 0}.}
%
\vnbprf{\vnbproofpart{\texttt{cps-taylor-coefficients}}}
%
\vnbprf{\textbf{Proof.} \texttt{theorem-\allowbreak{}library/cps-\allowbreak{}taylor-\allowbreak{}coefficients.scm} --- 155 recorded steps.}
%
\vnbprf{\textbf{Lemmas.} \texttt{cc-is-metric-space}, \texttt{cc-zero-in}, \texttt{nn-zero-in}, \texttt{rr-zero-in}, \texttt{rr-lt-implies-le}, \texttt{rr-pos-ne-zero}, \texttt{neq-sym}, \texttt{ball-is-set}, \texttt{ball-center-in}, \texttt{cps-kth-derived-radius}, \texttt{cps-holomorphic-ball}, \texttt{cps-kth-derived-step}, \texttt{cps-diff-on-ball}, \texttt{cc-is-normed-field}, \texttt{cc-nf-small-elements}, \texttt{diff-on-unique}, \texttt{cps-sum-at-zero}, \texttt{nn-add-comm}, \texttt{nn-add-zero}, \texttt{nn-one-is-succ-zero}, \texttt{rr-recip-one}, \texttt{factorial-real-pos}, \texttt{rr-subset-cc}, \texttt{fun-apply-type-c}.}
%
\vnbprf{\textbf{Oracles.} \texttt{crs}, \texttt{arith}, \texttt{ineq}, \texttt{sos}.  \textbf{Bill:} \texttt{proven modulo 0}.}
%
\end{proof}
%
\noindent\textbf{Deviation from the notes.} PROVEN 2026-09-27 (batch 38) AT A CENTRE: IS-ANALYTIC-AT(f, U, a) is DEFINED (structure-library/analytic-log.scm, the notes' 2.11: U open, f in fun(U, cc), a power series of positive radius on a ball B(a, r) inside U), the power-series layer is translated from 0 to a (cps-translate-*), and taylor-coefficients-at states f\textasciicircum{}(k)(a) = k! a\_k with f\textasciicircum{}(k) the k-th NTH-DERIV-ON of f restricted to B(a, r). Earlier: UPDATED 2026-09-26 (batch 36): NTH-DERIV-ON(K, U, f, k) is DEFINED by NN-recursion (structure-library/nth-deriv.scm) and 2.14 is stated LITERALLY -- taylor-coefficients: the k-th derivative of the sum function at 0 is k! a\_k, for the power series at 0 on every ball B(0, r), r $<$ radius; taylor-coefficient-formula is the coefficient form. Still partial for one reason only: the notes' 'f analytic at a' (no IS-ANALYTIC predicate, no translation to a centre a). Earlier: Proven 2026-09-24 (batch 27-A) for a function GIVEN as a power series at 0 (the notes' proof reduces to this case by translation; the tree has no IS-ANALYTIC predicate, so the hypothesis 'f analytic at a' is not stated). The tree has no iterated derivative over CC: f\textasciicircum{}(k) is read as F\_k, the sum function of the k-th derived series, whose coefficients (n+k)!/n! a\_(n+k) are given by their values. For every real 0 $<$ r $<$ R: F\_k is holomorphic on B(0, r), DERIV-ON(F\_k) = F\_(k+1) on the ball, and F\_k(0) = k! a\_k; hence a\_k = F\_k(0)/k! (the second name).
%
\renewcommand{\thethm}{2.15}
\begin{lem}[{Ratio Test}]\label{thm:ratio-test-converges-limsup}
%
\leavevmode
\vnbpart{\texttt{ratio-test-converges-limsup}}
\vnbstmt{\begin{array}{@{}l@{}} \forall h \in (\mathbb{N} \to \mathbb{R}). \\ \quad \forall k \in \mathbb{N}.\; (0 < h(k)) \wedge \\ \quad \operatorname{elimsup}((\lambda j \in \mathbb{N}.\; (h((j + 1)) \cdot \operatorname{recip}(h(j))))) \\ \quad <\; 1 \Rightarrow \\ \quad \quad \operatorname{series\text{-}converges}(h) \end{array}}
\vnbpart{\texttt{ratio-test-diverges-liminf}}
\vnbstmt{\begin{array}{@{}l@{}} \forall h \in (\mathbb{N} \to \mathbb{R}). \\ \quad \forall k \in \mathbb{N}.\; (0 < h(k)) \wedge \\ \quad 1 \\ \quad <\; \operatorname{eliminf}((\lambda j \in \mathbb{N}.\; (h((j + 1)) \cdot \operatorname{recip}(h(j))))) \Rightarrow \\ \quad \quad \neg \operatorname{series\text{-}converges}(h) \end{array}}
\vnbpart{\texttt{ratio-test-converges}}
\vnbstmt{\begin{array}{@{}l@{}} \forall f \in (\mathbb{N} \to \mathbb{R}), l, x \in \mathbb{R}. \\ \quad (0 \leq l) \wedge \\ \quad \forall n \in \mathbb{N}.\; \neg (f(n) = 0) \wedge \\ \quad \operatorname{ps\text{-}ratio\text{-}limit}(f, l) \wedge \\ \quad ((\lvert x \rvert \cdot l) < 1) \Rightarrow \\ \quad \quad \operatorname{ps\text{-}converges\text{-}at}(f, x) \end{array}}
%
\end{lem}
%
\begin{proof}
%
The proofs: the file that holds each, the lemmas it cites, its oracles and its bill.
%
\vnbprf{\vnbproofpart{\texttt{ratio-test-converges-limsup}}}
%
\vnbprf{\textbf{Proof.} \texttt{theorem-\allowbreak{}library/limsup-\allowbreak{}tests.scm} --- 431 recorded steps.}
%
\vnbprf{\textbf{Lemmas.} \texttt{nn-succ-closed}, \texttt{fun-apply-type-c}, \texttt{rr-recip-closed}, \texttt{rr-recip-pos}, \texttt{rr-mul-in-rr}, \texttt{rr-mul-pos}, \texttt{rr-one-in}, \texttt{elimsup-in-rr-pos-star}, \texttt{rr-pos-star-nonneg}, \texttt{rr-pos-star-le-real-in-rr}, \texttt{rr-add-in-rr}, \texttt{elimsup-eventually-below}, \texttt{power-real-closed}, \texttt{rr-power-pos}, \texttt{rr-recip-cancel-right}, \texttt{nn-le-zero-is-zero}, \texttt{nn-le-succ-cases}, \texttt{rr-mul-le-right}, \texttt{rr-le-scale-nonneg}, \texttt{rr-subset-cc}, \texttt{power-succ}, \texttt{series-eventually-geometric-converges}.}
%
\vnbprf{\textbf{Oracles.} \texttt{arith}, \texttt{ineq}, \texttt{crs}.  \textbf{Bill:} \texttt{proven modulo 0}.}
%
\vnbprf{\vnbproofpart{\texttt{ratio-test-diverges-liminf}}}
%
\vnbprf{\textbf{Proof.} \texttt{theorem-\allowbreak{}library/limsup-\allowbreak{}tests.scm} --- 319 recorded steps.}
%
\vnbprf{\textbf{Lemmas.} \texttt{nn-succ-closed}, \texttt{fun-apply-type-c}, \texttt{rr-recip-closed}, \texttt{rr-recip-pos}, \texttt{rr-mul-in-rr}, \texttt{rr-mul-pos}, \texttt{rr-one-in}, \texttt{eliminf-eventually-above}, \texttt{nn-le-zero-is-zero}, \texttt{nn-le-succ-cases}, \texttt{rr-recip-cancel-right}, \texttt{rr-mul-le-right}, \texttt{nn-add-closed}, \texttt{nn-le-add-right}, \texttt{nn-le-add-left}, \texttt{series-frequently-large-diverges}.}
%
\vnbprf{\textbf{Oracles.} \texttt{ineq}, \texttt{crs}, \texttt{arith}.  \textbf{Bill:} \texttt{proven modulo 0}.}
%
\vnbprf{\vnbproofpart{\texttt{ratio-test-converges}}}
%
\vnbprf{\textbf{Proof.} \texttt{theorem-\allowbreak{}library/ratio-\allowbreak{}test-\allowbreak{}proof.scm} --- 911 recorded steps.}
%
\vnbprf{\textbf{Lemmas.} \texttt{ps-converges-as-series}, \texttt{rr-zero-in}, \texttt{rr-one-in}, \texttt{nn-zero-in}, \texttt{nn-is-set}, \texttt{fun-apply-type-c}, \texttt{power-real-closed}, \texttt{rr-mul-in-rr}, \texttt{rr-abs-closed}, \texttt{rr-abs-nonneg}, \texttt{nn-one-in}, \texttt{series-eventually-geometric-converges}, \texttt{nn-zero-or-succ}, \texttt{nn-in-rr}, \texttt{rr-lt-irrefl}, \texttt{nn-succ-closed}, \texttt{power-zero-base}, \texttt{rr-mul-zero}, \texttt{rr-abs-zero-value}, \texttt{rr-abs-zero}, \texttt{rr-pos-ne-zero}, \texttt{rr-power-pos}, \texttt{rr-mul-pos}, \texttt{ps-abs-term}, \texttt{ratio-test-converges-limsup}, \texttt{rr-recip-closed}, \texttt{rr-recip-pos}, \texttt{rr-limit-scale}, \texttt{rr-recip-cancel-right}, \texttt{rr-subset-cc}, \texttt{power-succ}, \texttt{rr-abs-mult}, \texttt{rr-mul-recip-cancel}, \texttt{rr-sub-in-rr}, \texttt{rr-add-in-rr}, \texttt{rr-limit-tail-le}, \texttt{elimsup-eventually-le}, \texttt{elimsup-in-rr-pos-star}, \texttt{rr-pos-star-le-real-in-rr}, \texttt{series-abs-converges}.}
%
\vnbprf{\textbf{Oracles.} \texttt{ineq}, \texttt{crs}, \texttt{arith}.  \textbf{Bill:} \texttt{proven modulo 0}.}
%
\end{proof}
%
\noindent\textbf{Deviation from the notes.} UPDATED 2026-09-24 (batch 25): the first half PROVEN as stated (lim sup of the ratios $<$ 1). THE SECOND HALF IS FALSE AS PRINTED: lim sup theta\_(k+1)/theta\_k $>$ 1 does not imply divergence (theta\_(2j) = 4\textasciicircum{}-j, theta\_(2j+1) = 2 * 4\textasciicircum{}-j has ratios alternating 2 and 1/8, lim sup 2, and sums to 4); the library states and proves the divergence half with lim INF $>$ 1. The older asserted ratio-test-converges (a real power series with a LIMIT of ratios) follows from the new one with two short bridges and remains asserted for now. EARLIER NOTE: A PSS entry (warrant: informal), not proven. The library's ratio test is not the notes' statement about a bare series of positive reals: it is about a REAL POWER SERIES, coef in FUN(NN, RR) with every coef(n) non-zero, and the ratio hypothesis is PS-RATIO-LIMIT(coef, l) -- the ratios |coef(n+1)/coef(n)| CONVERGE to l, a limit, not a lim sup -- with l in RR, 0 $<$= l; the conclusion is PS-CONVERGES-AT(coef, x) for every real x with |x| l $<$ 1. The divergence half (lim sup of the ratios $>$ 1 implies divergence) is absent.
%
\renewcommand{\thethm}{2.16}
\begin{lem}[{The exponential series has infinite radius}]\label{thm:cc-exp-radius}
%
\leavevmode
\vnbpart{\texttt{cc-exp-radius}}
\vnbstmt{\begin{array}{@{}l@{}} \forall f \in (\mathbb{N} \to \mathbb{C}). \\ \quad \forall k \in \mathbb{N}.\; (f(k) = \operatorname{recip}(\operatorname{factorial}(k))) \Rightarrow  \\ \quad \quad (\operatorname{cps\text{-}radius}(f) = \operatorname{pos\text{-}inf}) \end{array}}
\vnbpart{\texttt{cps-radius-pos-inf}}
\vnbstmt{\begin{array}{@{}l@{}} \forall f \in (\mathbb{N} \to \mathbb{C}). \\ \quad \forall t \in \mathbb{R}. \\ \quad \quad (0 < t) \Rightarrow  \\ \quad \quad \quad \operatorname{series\text{-}converges}((\lambda k \in \mathbb{N}.\; (\operatorname{magnitude}(f(k)) \cdot \operatorname{power}(t, k)))) \Rightarrow \\ \quad \quad (\operatorname{cps\text{-}radius}(f) = \operatorname{pos\text{-}inf}) \end{array}}
\vnbpart{\texttt{cc-cos-radius}}
\vnbstmt{\begin{array}{@{}l@{}} \forall f \in (\mathbb{N} \to \mathbb{C}). \\ \quad \forall k \in \mathbb{N}. \\ \quad \quad f(k) \\ \quad \quad =\; (\operatorname{power}(-1, k) \cdot \operatorname{recip}(\operatorname{factorial}((2 \cdot k)))) \Rightarrow \\ \quad \quad (\operatorname{cps\text{-}radius}(f) = \operatorname{pos\text{-}inf}) \end{array}}
\vnbpart{\texttt{cc-sin-radius}}
\vnbstmt{\begin{array}{@{}l@{}} \forall f \in (\mathbb{N} \to \mathbb{C}). \\ \quad \forall k \in \mathbb{N}. \\ \quad \quad f(k) \\ \quad \quad =\; (\operatorname{power}(-1, k) \cdot \operatorname{recip}(\operatorname{factorial}(((2 \cdot k) + 1)))) \Rightarrow \\ \quad \quad (\operatorname{cps\text{-}radius}(f) = \operatorname{pos\text{-}inf}) \end{array}}
%
\end{lem}
%
\begin{proof}
%
The proofs: the file that holds each, the lemmas it cites, its oracles and its bill.
%
\vnbprf{\vnbproofpart{\texttt{cc-exp-radius}}}
%
\vnbprf{\textbf{Proof.} \texttt{theorem-\allowbreak{}library/cc-\allowbreak{}exp.scm} --- 83 recorded steps.}
%
\vnbprf{\textbf{Lemmas.} \texttt{cex-exp-series-real-converges}, \texttt{nn-is-set}, \texttt{fun-apply-type-c}, \texttt{cc-magnitude-closed}, \texttt{power-real-closed}, \texttt{rr-mul-in-rr}, \texttt{recip-factorial-in-rr}, \texttt{fun-codomain-iff}, \texttt{factorial-real-pos}, \texttt{rr-recip-pos}, \texttt{rr-magnitude-is-abs}, \texttt{rr-abs-of-nonneg}, \texttt{fun-domain-extensionality}, \texttt{cps-radius-pos-inf}.}
%
\vnbprf{\textbf{Oracles.} \texttt{ineq}, \texttt{crs}, \texttt{arith}.  \textbf{Bill:} \texttt{proven modulo 0}.}
%
\vnbprf{\vnbproofpart{\texttt{cps-radius-pos-inf}}}
%
\vnbprf{\textbf{Proof.} \texttt{theorem-\allowbreak{}library/cc-\allowbreak{}exp.scm} --- 61 recorded steps.}
%
\vnbprf{\textbf{Lemmas.} \texttt{subset-def}, \texttt{esup-in}, \texttt{esup-upper}, \texttt{rr-one-in}, \texttt{rr-add-in-rr}, \texttt{rr-zero-in}, \texttt{rr-lt-irrefl}.}
%
\vnbprf{\textbf{Oracles.} \texttt{ineq}.  \textbf{Bill:} \texttt{proven modulo 0}.}
%
\vnbprf{\vnbproofpart{\texttt{cc-cos-radius}}}
%
\vnbprf{\textbf{Proof.} \texttt{theorem-\allowbreak{}library/cc-\allowbreak{}exp.scm} --- 118 recorded steps.}
%
\vnbprf{\textbf{Lemmas.} \texttt{cex-exp-series-real-converges}, \texttt{nn-is-set}, \texttt{fun-apply-type-c}, \texttt{cc-magnitude-closed}, \texttt{power-real-closed}, \texttt{rr-mul-in-rr}, \texttt{recip-factorial-in-rr}, \texttt{comparison-test}, \texttt{nn-mul-closed}, \texttt{factorial-real-pos}, \texttt{rr-recip-pos}, \texttt{rr-power-pos}, \texttt{factorial-mono}, \texttt{nn-add-closed}, \texttt{nn-in-rr}, \texttt{nn-zero-le}, \texttt{rr-recip-antitone}, \texttt{rr-mul-le-right}, \texttt{rr-mul-pos}, \texttt{cex-sign-magnitude}, \texttt{cps-radius-pos-inf}.}
%
\vnbprf{\textbf{Oracles.} \texttt{arith}, \texttt{ineq}, \texttt{crs}.  \textbf{Bill:} \texttt{proven modulo 0}.}
%
\vnbprf{\vnbproofpart{\texttt{cc-sin-radius}}}
%
\vnbprf{\textbf{Proof.} \texttt{theorem-\allowbreak{}library/cc-\allowbreak{}exp.scm} --- 129 recorded steps.}
%
\vnbprf{\textbf{Lemmas.} \texttt{cex-exp-series-real-converges}, \texttt{nn-is-set}, \texttt{fun-apply-type-c}, \texttt{cc-magnitude-closed}, \texttt{power-real-closed}, \texttt{rr-mul-in-rr}, \texttt{recip-factorial-in-rr}, \texttt{comparison-test}, \texttt{nn-mul-closed}, \texttt{nn-succ-closed}, \texttt{factorial-real-pos}, \texttt{rr-recip-pos}, \texttt{rr-power-pos}, \texttt{factorial-mono}, \texttt{nn-add-closed}, \texttt{nn-in-rr}, \texttt{nn-zero-le}, \texttt{nn-succ-plus-one}, \texttt{rr-recip-antitone}, \texttt{rr-mul-le-right}, \texttt{rr-mul-pos}, \texttt{cex-sign-magnitude}, \texttt{cps-radius-pos-inf}.}
%
\vnbprf{\textbf{Oracles.} \texttt{arith}, \texttt{ineq}, \texttt{crs}.  \textbf{Bill:} \texttt{proven modulo 0}.}
%
\end{proof}
%
\noindent\textbf{Deviation from the notes.} UPDATED 2026-09-25 (batch 31-A, the user's decision: exp, cos, sin DEFINED by their power series, structure-library/cc-elementary.scm): PROVEN, cc-exp-radius (cps-radius of the coefficients 1/k! is pos-inf, by the notes' ratio test), with the same for the cos and sin series. Nothing near: no radius of convergence, and the series sum\_k z\textasciicircum{}k / k! is nowhere formed. FACTORIAL exists (factorial-real-pos, recip-factorial-in-rr) and is used inside the Taylor polynomial, but there is no theorem about the convergence of the exponential series, real or complex.
%
\renewcommand{\thethm}{2.17}
\begin{prop}[{exp is a continuous homomorphism into C*}]\label{thm:cc-exp-add}
%
\leavevmode
\vnbpart{\texttt{cc-exp-add}}
\vnbstmt{\forall z, w \in \mathbb{C}.\; (\exp (z + w) = (\exp z \cdot \exp w))}
\vnbpart{\texttt{cc-exp-zero}}
\vnbstmt{(\exp 0 = 1)}
\vnbpart{\texttt{cc-exp-neg}}
\vnbstmt{\forall z \in \mathbb{C}.\; (\exp -z = \operatorname{recip}(\exp z))}
\vnbpart{\texttt{cc-exp-nonzero}}
\vnbstmt{\forall z \in \mathbb{C}.\; \neg (\exp z = 0)}
\vnbpart{\texttt{cc-exp-maps-to-nonzero}}
\vnbstmt{\begin{array}{@{}l@{}} \forall z \in \mathbb{C}. \\ \quad (\exp z \in \operatorname{difference}(\mathbb{C}, \operatorname{singleton}(0))) \end{array}}
\vnbpart{\texttt{cc-exp-continuous}}
\vnbstmt{\begin{array}{@{}l@{}} \forall z \in \mathbb{C}. \\ \quad \operatorname{is\text{-}continuous\text{-}at}(\operatorname{subspace\text{-}ms}(\operatorname{nf\text{-}metric\text{-}space}(\operatorname{cc\text{-}normed\text{-}field}), \mathbb{C}), \\ \quad \quad \operatorname{nf\text{-}metric\text{-}space}(\operatorname{cc\text{-}normed\text{-}field}), (\lambda y \in \mathbb{C}.\; \exp y), z) \end{array}}
\vnbpart{\texttt{cc-exp-deriv}}
\vnbstmt{\begin{array}{@{}l@{}} \forall z \in \mathbb{C}. \\ \quad \operatorname{is\text{-}diff\text{-}on}(\operatorname{cc\text{-}normed\text{-}field}, \mathbb{C}, (\lambda y \in \mathbb{C}.\; \exp y), z, \exp z) \end{array}}
\vnbpart{\texttt{cc-exp-holomorphic}}
\vnbstmt{\operatorname{holomorphic\text{-}on}(\mathbb{C}, (\lambda y \in \mathbb{C}.\; \exp y))}
%
\end{prop}
%
\begin{proof}
%
The proofs: the file that holds each, the lemmas it cites, its oracles and its bill.
%
\vnbprf{\vnbproofpart{\texttt{cc-exp-add}}}
%
\vnbprf{\textbf{Proof.} \texttt{theorem-\allowbreak{}library/cc-\allowbreak{}exp.scm} --- 32 recorded steps.}
%
\vnbprf{\textbf{Lemmas.} \texttt{cc-add-closed}, \texttt{cc-exp-shift}, \texttt{cc-sub-in-cc}, \texttt{cc-exp-converges}, \texttt{cc-mul-closed}.}
%
\vnbprf{\textbf{Oracles.} \texttt{crs}, \texttt{ineq}, \texttt{arith}, \texttt{sos}.  \textbf{Bill:} \texttt{proven modulo 0}.}
%
\vnbprf{\vnbproofpart{\texttt{cc-exp-zero}}}
%
\vnbprf{\textbf{Proof.} \texttt{theorem-\allowbreak{}library/cc-\allowbreak{}exp.scm} --- 55 recorded steps.}
%
\vnbprf{\textbf{Lemmas.} \texttt{cc-zero-in}, \texttt{nn-zero-in}, \texttt{nn-is-set}, \texttt{recip-factorial-in-rr}, \texttt{rr-subset-cc}, \texttt{nn-succ-closed}, \texttt{tgd-recip-factorial-succ}, \texttt{cc-exp-radius}, \texttt{cc-exp-as-power-series}, \texttt{cps-sum-at-zero}, \texttt{nn-one-is-succ-zero}, \texttt{rr-recip-one}, \texttt{rr-one-in}.}
%
\vnbprf{\textbf{Oracles.} \texttt{arith}, \texttt{ineq}, \texttt{crs}, \texttt{sos}.  \textbf{Bill:} \texttt{proven modulo 0}.}
%
\vnbprf{\vnbproofpart{\texttt{cc-exp-neg}}}
%
\vnbprf{\textbf{Proof.} \texttt{theorem-\allowbreak{}library/cc-\allowbreak{}exp.scm} --- 54 recorded steps.}
%
\vnbprf{\textbf{Lemmas.} \texttt{cc-zero-in}, \texttt{cc-exp-shift}, \texttt{cc-neg-closed}, \texttt{cc-sub-in-cc}, \texttt{cc-exp-zero}, \texttt{cc-one-in}, \texttt{cc-exp-converges}, \texttt{cc-exp-nonzero}, \texttt{cc-recip-closed}, \texttt{cc-recip-inverse}, \texttt{cc-one-mul}.}
%
\vnbprf{\textbf{Oracles.} \texttt{crs}, \texttt{ineq}, \texttt{arith}, \texttt{sos}.  \textbf{Bill:} \texttt{proven modulo 0}.}
%
\vnbprf{\vnbproofpart{\texttt{cc-exp-nonzero}}}
%
\vnbprf{\textbf{Proof.} \texttt{theorem-\allowbreak{}library/cc-\allowbreak{}exp.scm} --- 40 recorded steps.}
%
\vnbprf{\textbf{Lemmas.} \texttt{cc-zero-in}, \texttt{cc-exp-shift}, \texttt{cc-neg-closed}, \texttt{cc-sub-in-cc}, \texttt{cc-exp-zero}, \texttt{cc-one-in}, \texttt{cc-exp-converges}, \texttt{rr-zero-in}, \texttt{rr-lt-irrefl}.}
%
\vnbprf{\textbf{Oracles.} \texttt{crs}, \texttt{ineq}, \texttt{arith}, \texttt{sos}.  \textbf{Bill:} \texttt{proven modulo 0}.}
%
\vnbprf{\vnbproofpart{\texttt{cc-exp-maps-to-nonzero}}}
%
\vnbprf{\textbf{Proof.} \texttt{theorem-\allowbreak{}library/cc-\allowbreak{}exp.scm} --- 10 recorded steps.}
%
\vnbprf{\textbf{Lemmas.} \texttt{cc-exp-converges}, \texttt{cc-exp-nonzero}, \texttt{cc-nonzero-membership}.}
%
\vnbprf{\textbf{Oracles.} \texttt{arith}, \texttt{ineq}, \texttt{crs}, \texttt{sos}.  \textbf{Bill:} \texttt{proven modulo 0}.}
%
\vnbprf{\vnbproofpart{\texttt{cc-exp-continuous}}}
%
\vnbprf{\textbf{Proof.} \texttt{theorem-\allowbreak{}library/cc-\allowbreak{}exp.scm} --- 6 recorded steps.}
%
\vnbprf{\textbf{Lemmas.} \texttt{cc-exp-converges}, \texttt{cc-exp-deriv}, \texttt{diff-on-implies-continuous}.}
%
\vnbprf{\textbf{Oracles.} \texttt{arith}, \texttt{ineq}, \texttt{crs}, \texttt{sos}.  \textbf{Bill:} \texttt{proven modulo 0}.}
%
\vnbprf{\vnbproofpart{\texttt{cc-exp-deriv}}}
%
\vnbprf{\textbf{Proof.} \texttt{theorem-\allowbreak{}library/cc-\allowbreak{}exp.scm} --- 149 recorded steps.}
%
\vnbprf{\textbf{Lemmas.} \texttt{cc-is-normed-field}, \texttt{nf-cc-is-metric-space}, \texttt{cc-nf-carr}, \texttt{nn-is-set}, \texttt{recip-factorial-in-rr}, \texttt{rr-subset-cc}, \texttt{nn-succ-closed}, \texttt{tgd-recip-factorial-succ}, \texttt{cc-exp-radius}, \texttt{cc-exp-as-power-series}, \texttt{cps-entire-diff}, \texttt{cc-is-set}, \texttt{cc-exp-converges}, \texttt{cps-entire-converges}, \texttt{fun-apply-type-c}, \texttt{power-typing-nonneg}, \texttt{cc-mul-closed}, \texttt{cc-series-limit-in-cc}, \texttt{diff-on-transfer-ptwise-eq}.}
%
\vnbprf{\textbf{Oracles.} \texttt{arith}, \texttt{crs}, \texttt{ineq}, \texttt{sos}.  \textbf{Bill:} \texttt{proven modulo 0}.}
%
\vnbprf{\vnbproofpart{\texttt{cc-exp-holomorphic}}}
%
\vnbprf{\textbf{Proof.} \texttt{theorem-\allowbreak{}library/cc-\allowbreak{}exp.scm} --- 38 recorded steps.}
%
\vnbprf{\textbf{Lemmas.} \texttt{cc-is-normed-field}, \texttt{nf-cc-is-metric-space}, \texttt{cc-nf-carr}, \texttt{cc-is-set}, \texttt{cc-exp-converges}, \texttt{carrier-is-open}, \texttt{cc-exp-deriv}.}
%
\vnbprf{\textbf{Oracles.} \texttt{arith}, \texttt{crs}, \texttt{ineq}, \texttt{sos}.  \textbf{Bill:} \texttt{proven modulo 0}.}
%
\end{proof}
%
\noindent\textbf{Deviation from the notes.} UPDATED 2026-09-25 (batch 31-A): PROVEN for the COMPLEX exponential defined by its series (theorem-library/cc-exp.scm): exp(z + w) = exp z exp w (cc-exp-add, by the zero-derivative route along a segment: exp(z) exp(w - z) is constant, no Cauchy product needed), exp 0 = 1, exp(-z) = 1/exp z, exp z /= 0, continuity and exp' = exp on CC. The tree has no group-homomorphism predicate, so 2.17 is these equations. The real exponential R-EXP (inverse of the integral LOG) is a separate object; the bridge 'exp restricted to RR is R-EXP' is not yet stated. There is no complex exponential in the library. R-EXP is the REAL exponential, and it is not defined by the power series: LOG is defined as the oriented integral of 1/x from 1 (log-unfold) and R-EXP is its inverse, so the identification of R-EXP with the restriction of the notes' exp is not a library fact. What is proven is the real restriction of the statement: r-exp-add (exp(u+v) = exp(u) exp(v) for real u, v), r-exp-pos (0 $<$ exp(y), i.e. the values land in the multiplicative group of the positive reals) and r-exp-continuous-at (continuity at each real point, in RR-MS). Nothing about C or C* = C minus \{0\}.
%
\renewcommand{\thethm}{2.18}
\begin{lem}[{The kernel of exp is 2 pi i Z}]\label{thm:cc-exp-one-iff}
%
\leavevmode
\vnbpart{\texttt{cc-exp-one-iff}}
\vnbstmt{\begin{array}{@{}l@{}} \forall z \in \mathbb{C}. \\ \quad ((\exp z = 1) \Leftrightarrow \exists k \in \mathbb{Z}.\; (z = (k \cdot ((2 \cdot \pi) \cdot +i)))) \end{array}}
\vnbpart{\texttt{pi-spec}}
\vnbstmt{\begin{array}{@{}l@{}} (\pi \in \mathbb{R}) \wedge \\ (0 < \pi) \wedge \\ \forall z \in \mathbb{C}. \\ \quad ((\exp z = 1) \Leftrightarrow \exists k \in \mathbb{Z}.\; (z = (k \cdot ((2 \cdot \pi) \cdot +i)))) \end{array}}
\vnbpart{\texttt{cc-exp-kernel-generator}}
\vnbstmt{\begin{array}{@{}l@{}} \exists g \in \mathbb{C}. \\ \quad (0 < \operatorname{imag\text{-}part}(g)) \wedge \\ \quad \forall z \in \mathbb{C}. \\ \quad \quad ((\exp z = 1) \Leftrightarrow \exists k \in \mathbb{Z}.\; (z = (k \cdot g))) \end{array}}
\vnbpart{\texttt{cc-exp-kernel-generator-unique}}
\vnbstmt{\begin{array}{@{}l@{}} \forall g \in \mathbb{C}. \\ \quad (0 < \operatorname{imag\text{-}part}(g)) \wedge \\ \quad \forall z \in \mathbb{C}. \\ \quad \quad ((\exp z = 1) \Leftrightarrow \exists k \in \mathbb{Z}.\; (z = (k \cdot g))) \Rightarrow \\ \quad \quad (g = ((2 \cdot \pi) \cdot +i)) \end{array}}
\vnbpart{\texttt{cc-exp-two-pi-i}}
\vnbstmt{(\exp ((2 \cdot \pi) \cdot +i) = 1)}
\vnbpart{\texttt{cc-exp-i-pi}}
\vnbstmt{(\exp (+i \cdot \pi) = -1)}
%
\end{lem}
%
\begin{proof}
%
The proofs: the file that holds each, the lemmas it cites, its oracles and its bill.
%
\vnbprf{\vnbproofpart{\texttt{cc-exp-one-iff}}}
%
\vnbprf{\textbf{Proof.} \texttt{theorem-\allowbreak{}library/cc-\allowbreak{}exp-\allowbreak{}kernel.scm} --- 4 recorded steps.}
%
\vnbprf{\textbf{Lemmas.} \texttt{pi-spec}.}
%
\vnbprf{\textbf{Oracles.} \texttt{arith}, \texttt{crs}, \texttt{ineq}, \texttt{sos}.  \textbf{Bill:} \texttt{proven modulo 0}.}
%
\vnbprf{\vnbproofpart{\texttt{pi-spec}}}
%
\vnbprf{\textbf{Proof.} \texttt{theorem-\allowbreak{}library/cc-\allowbreak{}exp-\allowbreak{}kernel.scm} --- 17 recorded steps.}
%
\vnbprf{\textbf{Lemmas.} \texttt{cc-pi-exists}, \texttt{cc-pi-unique}.}
%
\vnbprf{\textbf{Oracles.} \texttt{arith}, \texttt{crs}, \texttt{ineq}, \texttt{sos}.  \textbf{Bill:} \texttt{proven modulo 0}.}
%
\vnbprf{\vnbproofpart{\texttt{cc-exp-kernel-generator}}}
%
\vnbprf{\textbf{Proof.} \texttt{theorem-\allowbreak{}library/cc-\allowbreak{}exp-\allowbreak{}kernel.scm} --- 46 recorded steps.}
%
\vnbprf{\textbf{Lemmas.} \texttt{pi-pos}, \texttt{rr-subset-cc}, \texttt{rr-mul-in-rr}, \texttt{cc-i-in}, \texttt{cc-mul-closed}, \texttt{imag-part-in-rr}, \texttt{cc-im-real-mul}, \texttt{cc-exp-one-iff}.}
%
\vnbprf{\textbf{Oracles.} \texttt{arith}, \texttt{crs}, \texttt{ineq}, \texttt{sos}.  \textbf{Bill:} \texttt{proven modulo 0}.}
%
\vnbprf{\vnbproofpart{\texttt{cc-exp-kernel-generator-unique}}}
%
\vnbprf{\textbf{Proof.} \texttt{theorem-\allowbreak{}library/cc-\allowbreak{}exp-\allowbreak{}kernel.scm} --- 86 recorded steps.}
%
\vnbprf{\textbf{Lemmas.} \texttt{pi-pos}, \texttt{rr-subset-cc}, \texttt{nn-one-in}, \texttt{nn-subset-zz}, \texttt{cc-exp-kernel-on-imaginary-axis}, \texttt{imag-part-in-rr}, \texttt{rr-mul-in-rr}, \texttt{cc-i-in}, \texttt{cc-mul-closed}, \texttt{pi-spec}, \texttt{cc-pi-unique}.}
%
\vnbprf{\textbf{Oracles.} \texttt{crs}, \texttt{arith}, \texttt{ineq}, \texttt{sos}.  \textbf{Bill:} \texttt{proven modulo 0}.}
%
\vnbprf{\vnbproofpart{\texttt{cc-exp-two-pi-i}}}
%
\vnbprf{\textbf{Proof.} \texttt{theorem-\allowbreak{}library/cc-\allowbreak{}exp-\allowbreak{}kernel.scm} --- 35 recorded steps.}
%
\vnbprf{\textbf{Lemmas.} \texttt{pi-pos}, \texttt{rr-subset-cc}, \texttt{rr-mul-in-rr}, \texttt{cc-i-in}, \texttt{cc-mul-closed}, \texttt{cc-exp-one-iff}, \texttt{nn-one-in}, \texttt{nn-subset-zz}.}
%
\vnbprf{\textbf{Oracles.} \texttt{arith}, \texttt{crs}, \texttt{ineq}, \texttt{sos}.  \textbf{Bill:} \texttt{proven modulo 0}.}
%
\vnbprf{\vnbproofpart{\texttt{cc-exp-i-pi}}}
%
\vnbprf{\textbf{Proof.} \texttt{theorem-\allowbreak{}library/cc-\allowbreak{}exp-\allowbreak{}kernel.scm} --- 278 recorded steps.}
%
\vnbprf{\textbf{Lemmas.} \texttt{pi-pos}, \texttt{rr-subset-cc}, \texttt{cc-i-in}, \texttt{cc-mul-closed}, \texttt{cc-exp-converges}, \texttt{cc-add-closed}, \texttt{rr-mul-in-rr}, \texttt{cc-exp-two-pi-i}, \texttt{cc-exp-add}, \texttt{cc-exp-one-iff}, \texttt{zz-in-rr}, \texttt{imag-part-in-rr}, \texttt{cc-im-real-mul}, \texttt{rr-zero-in}, \texttt{rr-lt-trichotomy}, \texttt{rr-mul-le-right}, \texttt{rr-lt-irrefl}, \texttt{zz-discrete}, \texttt{rr-one-in}, \texttt{cc-one-in}, \texttt{cc-sub-in-cc}, \texttt{cc-mul-nonzero}.}
%
\vnbprf{\textbf{Oracles.} \texttt{arith}, \texttt{crs}, \texttt{ineq}, \texttt{sos}.  \textbf{Bill:} \texttt{proven modulo 0}.}
%
\end{proof}
%
\noindent\textbf{Deviation from the notes.} UPDATED 2026-09-25 (batch 31-B, the user's decision on PI): PROVEN, cc-exp-one-iff: exp z = 1 iff z = k (2 pi i) for some integer k. PI is DEFINED (structure-library/cc-pi.scm) as the positive real p with kernel(exp) = ZZ (2 p i), i.e. the generator of the kernel in the upper half plane divided by 2i; existence (the kernel is a subgroup of i RR, discrete, non-trivial since cos has a zero in (0,2), hence cyclic) and uniqueness are proven before the IOTA is formed. The notes use pi without defining it. Nothing near: no complex exponential, no pi (the symbol pi occurs in THEOREMS.md only as a bound variable in pigeonhole-infinite), no trigonometric functions. The real fragment of the statement -- exp t = 1 for real t only when t = 0 -- follows from r-exp-strictly-increasing and r-exp-zero but is not itself installed.
%
\renewcommand{\thethm}{2.19}
\begin{prop}[{exp is a bijection on a horizontal strip}]\label{thm:cc-exp-bijection-strip}
%
\leavevmode
\vnbpart{\texttt{cc-exp-bijection-strip}}
\vnbstmt{\begin{array}{@{}l@{}} (\lambda z \in \{q \in \mathbb{C} : ((-\pi \leq \operatorname{imag\text{-}part}(q)) \wedge (\operatorname{imag\text{-}part}(q) < \pi))\}.\; \\ \quad \exp z) \\ \in\; \mathrm{Bij}(\{q \in \mathbb{C} : ((-\pi \leq \operatorname{imag\text{-}part}(q)) \wedge (\operatorname{imag\text{-}part}(q) < \pi))\}, \operatorname{difference}(\mathbb{C}, \operatorname{singleton}(0))) \end{array}}
\vnbpart{\texttt{cc-exp-bijection-open-strip}}
\vnbstmt{\begin{array}{@{}l@{}} (\lambda z \in \{q \in \mathbb{C} : ((-\pi < \operatorname{imag\text{-}part}(q)) \wedge (\operatorname{imag\text{-}part}(q) < \pi))\}.\; \\ \quad \exp z) \\ \in\; \mathrm{Bij}(\{q \in \mathbb{C} : ((-\pi < \operatorname{imag\text{-}part}(q)) \wedge (\operatorname{imag\text{-}part}(q) < \pi))\}, \{u \in \mathbb{C} : \neg ((\operatorname{imag\text{-}part}(u) = 0) \wedge (\operatorname{real\text{-}part}(u) \leq 0))\}) \end{array}}
\vnbpart{\texttt{cc-exp-surjective-strip}}
\vnbstmt{\begin{array}{@{}l@{}} \forall w \in \mathbb{C}. \\ \quad \neg (w = 0) \Rightarrow  \\ \quad \quad \exists z \in \mathbb{C}. \\ \quad \quad \quad (-\pi \leq \operatorname{imag\text{-}part}(z)) \wedge \\ \quad \quad \quad (\operatorname{imag\text{-}part}(z) < \pi) \wedge \\ \quad \quad \quad (\exp z = w) \end{array}}
\vnbpart{\texttt{cc-exp-injective-strip}}
\vnbstmt{\begin{array}{@{}l@{}} \forall z, w \in \mathbb{C}. \\ \quad (-\pi \leq \operatorname{imag\text{-}part}(z)) \wedge \\ \quad (\operatorname{imag\text{-}part}(z) < \pi) \wedge \\ \quad (-\pi \leq \operatorname{imag\text{-}part}(w)) \wedge \\ \quad (\operatorname{imag\text{-}part}(w) < \pi) \wedge \\ \quad (\exp z = \exp w) \Rightarrow \\ \quad \quad (z = w) \end{array}}
\vnbpart{\texttt{cc-exp-open-strip-onto}}
\vnbstmt{\begin{array}{@{}l@{}} \forall w \in \mathbb{C}. \\ \quad \neg ((\operatorname{imag\text{-}part}(w) = 0) \wedge (\operatorname{real\text{-}part}(w) \leq 0)) \Rightarrow  \\ \quad \quad \exists z \in \mathbb{C}. \\ \quad \quad \quad (-\pi < \operatorname{imag\text{-}part}(z)) \wedge \\ \quad \quad \quad (\operatorname{imag\text{-}part}(z) < \pi) \wedge \\ \quad \quad \quad (\exp z = w) \end{array}}
\vnbpart{\texttt{cc-exp-open-strip-avoids-negatives}}
\vnbstmt{\begin{array}{@{}l@{}} \forall z \in \mathbb{C}. \\ \quad ((-\pi < \operatorname{imag\text{-}part}(z)) \wedge (\operatorname{imag\text{-}part}(z) < \pi)) \Rightarrow  \\ \quad \quad \neg ((\operatorname{imag\text{-}part}(\exp z) = 0) \wedge (\operatorname{real\text{-}part}(\exp z) \leq 0)) \end{array}}
\vnbpart{\texttt{cc-polar-form}}
\vnbstmt{\begin{array}{@{}l@{}} \forall w \in \mathbb{C}. \\ \quad \neg (w = 0) \Rightarrow  \\ \quad \quad \exists r, h \in \mathbb{R}. \\ \quad \quad \quad ((0 < r) \wedge ((-\pi \leq h) \wedge ((h < \pi) \wedge (w = (r \cdot \exp (+i \cdot h)))))) \end{array}}
\vnbpart{\texttt{cc-exp-real-is-r-exp}}
\vnbstmt{\forall t \in \mathbb{R}.\; (\exp t = \exp\left(t\right))}
%
\end{prop}
%
\begin{proof}
%
The proofs: the file that holds each, the lemmas it cites, its oracles and its bill.
%
\vnbprf{\vnbproofpart{\texttt{cc-exp-bijection-strip}}}
%
\vnbprf{\textbf{Proof.} \texttt{theorem-\allowbreak{}library/cc-\allowbreak{}exp-\allowbreak{}surjective.scm} --- 94 recorded steps.}
%
\vnbprf{\textbf{Lemmas.} \texttt{pi-pos}, \texttt{rr-subset-cc}, \texttt{cc-exp-maps-to-nonzero}, \texttt{cc-is-set}, \texttt{rr-zero-in}, \texttt{cc-exp-injective-strip}, \texttt{difference-membership}, \texttt{singleton-membership}, \texttt{membership-implies-sethood}, \texttt{cc-exp-surjective-strip}.}
%
\vnbprf{\textbf{Oracles.} \texttt{arith}, \texttt{crs}, \texttt{ineq}, \texttt{sos}.  \textbf{Bill:} \texttt{proven modulo 0}.}
%
\vnbprf{\vnbproofpart{\texttt{cc-exp-bijection-open-strip}}}
%
\vnbprf{\textbf{Proof.} \texttt{theorem-\allowbreak{}library/cc-\allowbreak{}exp-\allowbreak{}surjective.scm} --- 94 recorded steps.}
%
\vnbprf{\textbf{Lemmas.} \texttt{pi-pos}, \texttt{rr-subset-cc}, \texttt{cc-exp-converges}, \texttt{cc-exp-open-strip-avoids-negatives}, \texttt{cc-is-set}, \texttt{rr-zero-in}, \texttt{imag-part-in-rr}, \texttt{cc-exp-injective-strip}, \texttt{cc-exp-open-strip-onto}.}
%
\vnbprf{\textbf{Oracles.} \texttt{ineq}, \texttt{arith}, \texttt{crs}, \texttt{sos}.  \textbf{Bill:} \texttt{proven modulo 0}.}
%
\vnbprf{\vnbproofpart{\texttt{cc-exp-surjective-strip}}}
%
\vnbprf{\textbf{Proof.} \texttt{theorem-\allowbreak{}library/cc-\allowbreak{}exp-\allowbreak{}surjective.scm} --- 98 recorded steps.}
%
\vnbprf{\textbf{Lemmas.} \texttt{pi-pos}, \texttt{rr-subset-cc}, \texttt{cc-polar-form}, \texttt{cc-exp-real-surjective-pos}, \texttt{cc-i-in}, \texttt{cc-mul-closed}, \texttt{cc-add-closed}, \texttt{cc-re-im-of}, \texttt{cc-exp-converges}, \texttt{cc-exp-add}.}
%
\vnbprf{\textbf{Oracles.} \texttt{crs}, \texttt{arith}, \texttt{ineq}, \texttt{sos}.  \textbf{Bill:} \texttt{proven modulo 0}.}
%
\vnbprf{\vnbproofpart{\texttt{cc-exp-injective-strip}}}
%
\vnbprf{\textbf{Proof.} \texttt{theorem-\allowbreak{}library/cc-\allowbreak{}exp-\allowbreak{}kernel.scm} --- 188 recorded steps.}
%
\vnbprf{\textbf{Lemmas.} \texttt{pi-pos}, \texttt{rr-subset-cc}, \texttt{cc-sub-in-cc}, \texttt{cc-neg-closed}, \texttt{cc-add-closed}, \texttt{cc-exp-converges}, \texttt{cc-zero-in}, \texttt{cc-exp-shift}, \texttt{cc-exp-zero}, \texttt{cc-one-in}, \texttt{cc-exp-add}, \texttt{cc-exp-one-iff}, \texttt{zz-in-rr}, \texttt{rr-mul-in-rr}, \texttt{cc-i-in}, \texttt{cc-mul-closed}, \texttt{imag-part-in-rr}, \texttt{cc-im-real-mul}, \texttt{cc-im-sub}, \texttt{rr-zero-in}, \texttt{rr-one-in}, \texttt{rr-lt-trichotomy}, \texttt{zz-neg-closed}, \texttt{rr-neg-closed}, \texttt{zz-discrete}, \texttt{rr-mul-le-right}, \texttt{rr-lt-irrefl}.}
%
\vnbprf{\textbf{Oracles.} \texttt{crs}, \texttt{arith}, \texttt{ineq}, \texttt{sos}.  \textbf{Bill:} \texttt{proven modulo 0}.}
%
\vnbprf{\vnbproofpart{\texttt{cc-exp-open-strip-onto}}}
%
\vnbprf{\textbf{Proof.} \texttt{theorem-\allowbreak{}library/cc-\allowbreak{}exp-\allowbreak{}surjective.scm} --- 170 recorded steps.}
%
\vnbprf{\textbf{Lemmas.} \texttt{pi-pos}, \texttt{rr-subset-cc}, \texttt{rr-zero-in}, \texttt{cc-exp-surjective-strip}, \texttt{imag-part-in-rr}, \texttt{real-part-in-rr}, \texttt{cc-i-in}, \texttt{cc-mul-closed}, \texttt{cc-add-closed}, \texttt{cc-exp-converges}, \texttt{cc-exp-real-valued}, \texttt{rr-neg-closed}, \texttt{cc-exp-real}, \texttt{cc-neg-closed}, \texttt{cc-exp-i-pi}, \texttt{cc-exp-neg}, \texttt{cc-re-im-decompose}, \texttt{cc-exp-add}, \texttt{cc-re-im-of}, \texttt{rr-le-ne-lt}.}
%
\vnbprf{\textbf{Oracles.} \texttt{arith}, \texttt{crs}, \texttt{ineq}, \texttt{sos}.  \textbf{Bill:} \texttt{proven modulo 0}.}
%
\vnbprf{\vnbproofpart{\texttt{cc-exp-open-strip-avoids-negatives}}}
%
\vnbprf{\textbf{Proof.} \texttt{theorem-\allowbreak{}library/cc-\allowbreak{}exp-\allowbreak{}surjective.scm} --- 178 recorded steps.}
%
\vnbprf{\textbf{Lemmas.} \texttt{pi-pos}, \texttt{rr-subset-cc}, \texttt{rr-zero-in}, \texttt{cc-exp-converges}, \texttt{real-part-in-rr}, \texttt{imag-part-in-rr}, \texttt{rr-neg-closed}, \texttt{cc-exp-nonzero}, \texttt{cc-re-im-decompose}, \texttt{cc-i-in}, \texttt{cc-mul-closed}, \texttt{cc-add-closed}, \texttt{rr-le-ne-lt}, \texttt{cc-exp-real-surjective-pos}, \texttt{cc-neg-closed}, \texttt{cc-exp-i-pi}, \texttt{cc-exp-neg}, \texttt{cc-exp-add}, \texttt{cc-re-im-of}, \texttt{cc-exp-injective-strip}, \texttt{rr-lt-irrefl}.}
%
\vnbprf{\textbf{Oracles.} \texttt{crs}, \texttt{ineq}, \texttt{arith}, \texttt{sos}.  \textbf{Bill:} \texttt{proven modulo 0}.}
%
\vnbprf{\vnbproofpart{\texttt{cc-polar-form}}}
%
\vnbprf{\textbf{Proof.} \texttt{theorem-\allowbreak{}library/cc-\allowbreak{}exp-\allowbreak{}surjective.scm} --- 241 recorded steps.}
%
\vnbprf{\textbf{Lemmas.} \texttt{pi-pos}, \texttt{rr-subset-cc}, \texttt{rr-zero-in}, \texttt{cc-magnitude-closed}, \texttt{real-part-in-rr}, \texttt{imag-part-in-rr}, \texttt{cc-magnitude-nonneg}, \texttt{cc-magnitude-zero-iff}, \texttt{rr-le-ne-lt}, \texttt{rr-recip-closed}, \texttt{cc-recip-inverse}, \texttt{rr-mul-in-rr}, \texttt{cc-magnitude-sq}, \texttt{nn-two-is-succ-succ-zero}, \texttt{nn-zero-in}, \texttt{nn-succ-closed}, \texttt{power-succ}, \texttt{power-zero}, \texttt{cc-circle-arg}, \texttt{cc-exp-i}, \texttt{cc-re-im-decompose}, \texttt{cc-i-in}, \texttt{cc-mul-closed}, \texttt{cc-exp-converges}, \texttt{cc-add-closed}.}
%
\vnbprf{\textbf{Oracles.} \texttt{ineq}, \texttt{crs}, \texttt{arith}, \texttt{sos}.  \textbf{Bill:} \texttt{proven modulo 0}.}
%
\vnbprf{\vnbproofpart{\texttt{cc-exp-real-is-r-exp}}}
%
\vnbprf{\textbf{Proof.} \texttt{theorem-\allowbreak{}library/cc-\allowbreak{}exp-\allowbreak{}surjective.scm} --- 286 recorded steps.}
%
\vnbprf{\textbf{Lemmas.} \texttt{rr-zero-in}, \texttt{rr-subset-cc}, \texttt{cc-exp-real-valued}, \texttt{cc-exp-converges}, \texttt{cc-exp-real}, \texttt{cxs-log-exp-deriv}, \texttt{diff-at-in-fun}, \texttt{rr-lt-trichotomy}, \texttt{deriv-zero-implies-constant}, \texttt{ccint-membership}, \texttt{diff-implies-continuous}, \texttt{fun-apply-type-c}, \texttt{log-in-rr}, \texttt{rr-sub-in-rr}, \texttt{cc-exp-zero}, \texttt{log-one}, \texttt{r-exp-char}, \texttt{r-exp-zero}.}
%
\vnbprf{\textbf{Oracles.} \texttt{ineq}, \texttt{arith}, \texttt{crs}, \texttt{sos}.  \textbf{Bill:} \texttt{proven modulo 0}.}
%
\end{proof}
%
\noindent\textbf{Deviation from the notes.} PROVEN 2026-09-26 (batch 34) as the notes state it: cc-exp-bijection-strip says the lambda z in S. exp z is in bijection(S, cc minus \{0\}) for S the SEP \{-pi $<$= Im $<$ pi\}; cc-exp-bijection-open-strip the same for the open strip onto cc minus (-inf, 0]. Surjectivity goes through the polar form (cc-polar-form, cc-circle-arg by the IVT on cos) and the bridge cc-exp-real-is-r-exp (exp restricted to RR is the real exponential R-EXP, via the derivative of log(exp t) - t). Earlier: UPDATED 2026-09-25 (batch 31-B): the INJECTIVITY half is PROVEN for the complex exponential on the strip -pi $<$= Im z $<$ pi (cc-exp-injective-strip, by Lemma 2.18); surjectivity onto CC minus 0 and the open-strip statement wait on the polar form (exp unbounded on RR, IVT for the modulus and the argument). Only the real restriction, and under the library's own definition of R-EXP as the inverse of the integral LOG rather than as a power series. r-exp-log and log-r-exp say that R-EXP and LOG are mutually inverse between RR and the positive reals, log-surjective and r-exp-pos give surjectivity onto (0, +inf), and log-injective-pos / r-exp-strictly-increasing give injectivity. Nothing about the strips S = \{-pi $<$= Im z $<$ pi\} or S' = \{-pi $<$ Im z $<$ pi\}, about C minus \{0\} or about C minus (-inf, 0]; there is no pi and no complex exponential to state them with.
%
\renewcommand{\thethm}{2.20}
\begin{lem}[{The logarithm series has radius 1}]\label{thm:cps-log-radius-one}
%
\leavevmode
\vnbpart{\texttt{cps-log-radius-one}}
\vnbstmt{\begin{array}{@{}l@{}} \forall c \in (\mathbb{N} \to \mathbb{C}). \\ \quad (c(0) = 0) \wedge \\ \quad \forall j \in \mathbb{N}. \\ \quad \quad (c((j + 1)) = (\operatorname{power}(-1, j) \cdot \operatorname{recip}((j + 1)))) \Rightarrow \\ \quad \quad (\operatorname{cps\text{-}radius}(c) = 1) \end{array}}
\vnbpart{\texttt{cps-log-series}}
\vnbstmt{\begin{array}{@{}l@{}} \forall c \in (\mathbb{N} \to \mathbb{C}), r \in \mathbb{R}, z \in \operatorname{ball}(\operatorname{cc\text{-}ms}, 0, r). \\ \quad (c(0) = 0) \wedge \\ \quad \forall j \in \mathbb{N}. \\ \quad \quad (c((j + 1)) = (\operatorname{power}(-1, j) \cdot \operatorname{recip}((j + 1)))) \wedge \\ \quad (0 < r) \wedge \\ \quad (r < 1) \Rightarrow \\ \quad \quad \operatorname{is\text{-}diff\text{-}on}(\operatorname{cc\text{-}normed\text{-}field}, \operatorname{ball}(\operatorname{cc\text{-}ms}, 0, r), \\ \quad \quad \quad (\lambda \mathit{cpdz} \in \operatorname{ball}(\operatorname{cc\text{-}ms}, 0, r).\; \\ \quad \quad \quad \quad \operatorname{cc\text{-}series\text{-}limit}((\lambda k \in \mathbb{N}.\; (c(k) \cdot \operatorname{power}(\mathit{cpdz}, k))))), \\ \quad \quad \quad z, \operatorname{recip}((1 + z))) \end{array}}
\vnbpart{\texttt{cps-log-radius}}
\vnbstmt{\begin{array}{@{}l@{}} \forall c \in (\mathbb{N} \to \mathbb{C}), r \in \mathbb{R}. \\ \quad (c(0) = 0) \wedge \\ \quad \forall j \in \mathbb{N}. \\ \quad \quad (c((j + 1)) = (\operatorname{power}(-1, j) \cdot \operatorname{recip}((j + 1)))) \wedge \\ \quad (0 \leq r) \wedge \\ \quad (r < 1) \Rightarrow \\ \quad \quad (r < \operatorname{cps\text{-}radius}(c)) \end{array}}
\vnbpart{\texttt{cps-log-radius-le-one}}
\vnbstmt{\begin{array}{@{}l@{}} \forall c \in (\mathbb{N} \to \mathbb{C}), r \in \mathbb{R}. \\ \quad (c(0) = 0) \wedge \\ \quad \forall j \in \mathbb{N}. \\ \quad \quad (c((j + 1)) = (\operatorname{power}(-1, j) \cdot \operatorname{recip}((j + 1)))) \wedge \\ \quad (1 < r) \Rightarrow \\ \quad \quad \neg (r < \operatorname{cps\text{-}radius}(c)) \end{array}}
\vnbpart{\texttt{cc-series-terms-small}}
\vnbstmt{\begin{array}{@{}l@{}} \forall f \in (\mathbb{N} \to \mathbb{C}), e. \\ \quad (\operatorname{cc\text{-}series\text{-}converges}(f) \wedge \operatorname{pos\text{-}rr}(e)) \Rightarrow  \\ \quad \quad \exists n \in \mathbb{N}. \forall m \in \mathbb{N}. \\ \quad \quad \quad ((n \leq m) \Rightarrow (\operatorname{magnitude}(f(m)) \leq e)) \end{array}}
\vnbpart{\texttt{cc-geometric-alternating}}
\vnbstmt{\begin{array}{@{}l@{}} \forall g \in (\mathbb{N} \to \mathbb{C}), z \in \mathbb{C}. \\ \quad \forall j \in \mathbb{N}.\; (g(j) = \operatorname{power}(-1, j)) \wedge \\ \quad (\operatorname{magnitude}(z) < 1) \Rightarrow \\ \quad \quad \operatorname{cc\text{-}series\text{-}converges}((\lambda k \in \mathbb{N}.\; (g(k) \cdot \operatorname{power}(z, k)))) \wedge \\ \quad \quad \operatorname{cc\text{-}series\text{-}limit}((\lambda k \in \mathbb{N}.\; (g(k) \cdot \operatorname{power}(z, k)))) \\ \quad \quad =\; \operatorname{recip}((1 + z)) \end{array}}
%
\end{lem}
%
\begin{proof}
%
The proofs: the file that holds each, the lemmas it cites, its oracles and its bill.
%
\vnbprf{\vnbproofpart{\texttt{cps-log-radius-one}}}
%
\vnbprf{\textbf{Proof.} \texttt{theorem-\allowbreak{}library/cc-\allowbreak{}log-\allowbreak{}holomorphic.scm} --- 83 recorded steps.}
%
\vnbprf{\textbf{Lemmas.} \texttt{rr-zero-in}, \texttt{rr-one-in}, \texttt{cps-radius-in}, \texttt{rr-add-in-rr}, \texttt{rr-mul-in-rr}, \texttt{rr-lt-trichotomy}, \texttt{cps-log-radius}, \texttt{rr-lt-irrefl}, \texttt{cps-log-radius-le-one}, \texttt{rr-subset-rr-star}, \texttt{subset-mem-fwd}, \texttt{pos-inf-upper-bound}, \texttt{pos-inf-not-in-rr}.}
%
\vnbprf{\textbf{Oracles.} \texttt{arith}, \texttt{ineq}, \texttt{crs}, \texttt{sos}.  \textbf{Bill:} \texttt{proven modulo 0}.}
%
\vnbprf{\vnbproofpart{\texttt{cps-log-series}}}
%
\vnbprf{\textbf{Proof.} \texttt{theorem-\allowbreak{}library/analytic-\allowbreak{}log-\allowbreak{}laws.scm} --- 207 recorded steps.}
%
\vnbprf{\textbf{Lemmas.} \texttt{cc-is-metric-space}, \texttt{cc-is-normed-field}, \texttt{cc-zero-in}, \texttt{rr-one-in}, \texttt{rr-zero-in}, \texttt{rr-lt-irrefl}, \texttt{cps-log-radius}, \texttt{nn-succ-closed}, \texttt{nn-in-rr}, \texttt{rr-subset-cc}, \texttt{fun-apply-type-c}, \texttt{cc-mul-closed}, \texttt{nn-is-set}, \texttt{nn-one-le-succ}, \texttt{rr-recip-closed}, \texttt{power-real-closed}, \texttt{rr-recip-inverse}, \texttt{ball-membership}, \texttt{cc-sub-in-cc}, \texttt{cc-magnitude-closed}, \texttt{cc-ms-dist}, \texttt{cc-magnitude-neg}, \texttt{cps-diff-on-ball}, \texttt{cc-geometric-alternating}.}
%
\vnbprf{\textbf{Oracles.} \texttt{ineq}, \texttt{arith}, \texttt{crs}, \texttt{sos}.  \textbf{Bill:} \texttt{proven modulo 0}.}
%
\vnbprf{\vnbproofpart{\texttt{cps-log-radius}}}
%
\vnbprf{\textbf{Proof.} \texttt{theorem-\allowbreak{}library/analytic-\allowbreak{}log-\allowbreak{}laws.scm} --- 268 recorded steps.}
%
\vnbprf{\textbf{Lemmas.} \texttt{rr-zero-in}, \texttt{rr-one-in}, \texttt{rr-add-in-rr}, \texttt{rr-mul-in-rr}, \texttt{fun-apply-type-c}, \texttt{cc-magnitude-closed}, \texttt{power-real-closed}, \texttt{nn-is-set}, \texttt{geometric-series-converges}, \texttt{cc-magnitude-zero}, \texttt{nn-nonzero-is-succ}, \texttt{nn-succ-closed}, \texttt{nn-in-rr}, \texttt{nn-one-le-succ}, \texttt{rr-lt-irrefl}, \texttt{rr-recip-closed}, \texttt{rr-subset-cc}, \texttt{cc-magnitude-mul}, \texttt{cc-magnitude-power}, \texttt{rr-magnitude-is-abs}, \texttt{rr-power-one}, \texttt{rr-recip-pos}, \texttt{rr-abs-of-nonneg}, \texttt{rr-mul-le-right}, \texttt{rr-recip-inverse}, \texttt{cc-magnitude-nonneg}, \texttt{rr-power-nonneg}, \texttt{comparison-test}, \texttt{subset-def}, \texttt{esup-upper}, \texttt{esup-in}, \texttt{cps-radius-unfold}, \texttt{rr-pos-star-below-pos-inf}, \texttt{pos-inf-not-in-rr}.}
%
\vnbprf{\textbf{Oracles.} \texttt{arith}, \texttt{ineq}, \texttt{crs}.  \textbf{Bill:} \texttt{proven modulo 0}.}
%
\vnbprf{\vnbproofpart{\texttt{cps-log-radius-le-one}}}
%
\vnbprf{\textbf{Proof.} \texttt{theorem-\allowbreak{}library/cc-\allowbreak{}log-\allowbreak{}holomorphic.scm} --- 284 recorded steps.}
%
\vnbprf{\textbf{Lemmas.} \texttt{cc-is-metric-space}, \texttt{cc-is-normed-field}, \texttt{cc-zero-in}, \texttt{rr-zero-in}, \texttt{rr-one-in}, \texttt{nn-zero-in}, \texttt{rr-subset-cc}, \texttt{cc-ms-dist}, \texttt{rr-sub-in-rr}, \texttt{closed-ball-membership}, \texttt{cc-magnitude-neg}, \texttt{rr-magnitude-is-abs}, \texttt{rr-abs-of-nonneg}, \texttt{cps-converges-closed-ball}, \texttt{fun-apply-type-c}, \texttt{power-real-closed}, \texttt{cc-mul-closed}, \texttt{nn-is-set}, \texttt{rr-mul-in-rr}, \texttt{cc-series-terms-small}, \texttt{nn-succ-closed}, \texttt{nn-le-succ}, \texttt{nn-in-rr}, \texttt{nn-one-le-succ}, \texttt{rr-lt-irrefl}, \texttt{rr-recip-closed}, \texttt{rr-recip-pos}, \texttt{rr-recip-inverse}, \texttt{cc-magnitude-mul}, \texttt{cc-magnitude-power}, \texttt{rr-power-one}, \texttt{rr-bernoulli}, \texttt{rr-add-in-rr}, \texttt{rr-le-scale-nonneg}, \texttt{cc-magnitude-closed}.}
%
\vnbprf{\textbf{Oracles.} \texttt{ineq}, \texttt{crs}, \texttt{arith}, \texttt{sos}.  \textbf{Bill:} \texttt{proven modulo 0}.}
%
\vnbprf{\vnbproofpart{\texttt{cc-series-terms-small}}}
%
\vnbprf{\textbf{Proof.} \texttt{theorem-\allowbreak{}library/cc-\allowbreak{}log-\allowbreak{}holomorphic.scm} --- 142 recorded steps.}
%
\vnbprf{\textbf{Lemmas.} \texttt{cc-is-metric-space}, \texttt{cc-is-normed-field}, \texttt{cc-zero-in}, \texttt{rr-zero-in}, \texttt{rr-pos-rr-in-rr}, \texttt{rr-lt-of-pos-rr}, \texttt{rr-mul-in-rr}, \texttt{cc-series-limit-in-cc}, \texttt{cc-series-limit-converges-to}, \texttt{nn-succ-closed}, \texttt{nn-le-succ}, \texttt{nn-le-trans-guarded}, \texttt{fun-apply-type-c}, \texttt{cc-sub-in-cc}, \texttt{cc-series-partial-sum-in-cc}, \texttt{cc-series-partial-sum-succ}, \texttt{cc-ms-dist}, \texttt{cc-magnitude-closed}, \texttt{cc-magnitude-le-add}, \texttt{cc-neg-closed}, \texttt{cc-magnitude-neg}.}
%
\vnbprf{\textbf{Oracles.} \texttt{arith}, \texttt{ineq}, \texttt{crs}.  \textbf{Bill:} \texttt{proven modulo 0}.}
%
\vnbprf{\vnbproofpart{\texttt{cc-geometric-alternating}}}
%
\vnbprf{\textbf{Proof.} \texttt{theorem-\allowbreak{}library/analytic-\allowbreak{}log-\allowbreak{}laws.scm} --- 375 recorded steps.}
%
\vnbprf{\textbf{Lemmas.} \texttt{cc-is-metric-space}, \texttt{cc-is-normed-field}, \texttt{cc-zero-in}, \texttt{nn-zero-in}, \texttt{cc-magnitude-closed}, \texttt{cc-magnitude-nonneg}, \texttt{rr-subset-cc}, \texttt{cc-one-in}, \texttt{cc-add-closed}, \texttt{rr-magnitude-is-abs}, \texttt{rr-zero-in}, \texttt{rr-lt-irrefl}, \texttt{cc-recip-closed}, \texttt{fun-apply-type-c}, \texttt{power-typing-nonneg}, \texttt{cc-mul-closed}, \texttt{nn-is-set}, \texttt{cc-series-partial-sum-zero}, \texttt{power-zero}, \texttt{cc-series-partial-sum-in-cc}, \texttt{power-real-closed}, \texttt{nn-succ-closed}, \texttt{cc-series-partial-sum-succ}, \texttt{power-succ}, \texttt{cc-recip-inverse}, \texttt{rr-mul-in-rr}, \texttt{rr-abs-of-nonneg}, \texttt{power-tends-to-zero}, \texttt{rr-null-scale}, \texttt{cc-magnitude-neg}, \texttt{cc-magnitude-mul}, \texttt{cc-magnitude-power}, \texttt{rr-power-one}, \texttt{cpd-cc-squeeze}, \texttt{cc-series-limit-in-cc}, \texttt{cc-series-limit-converges-to}, \texttt{cc-limit-unique}.}
%
\vnbprf{\textbf{Oracles.} \texttt{arith}, \texttt{crs}, \texttt{ineq}.  \textbf{Bill:} \texttt{proven modulo 0}.}
%
\end{proof}
%
\noindent\textbf{Deviation from the notes.} PROVEN 2026-09-27 (batches 38 and 38-B): cps-radius of the series is exactly 1 (cps-log-radius-one; the upper half through the new general lemma cc-series-terms-small, the terms of a convergent complex series tend to 0, and Bernoulli), and the derivative is recip(1 + z) on every B(0, r), r $<$ 1. Earlier: PARTIAL 2026-09-27 (batch 38): the series is formed (as a lambda inside the theorems), its radius is at least 1 (cps-log-radius) and its derivative is recip(1 + z) on every B(0, r), r $<$ 1 (cps-log-series, by 2.10 termwise and the alternating geometric series cc-geometric-alternating). The other half, radius at most 1, is batch 38-B's item (4). Earlier: Nothing near: no radius of convergence, and the series sum\_k (-1)\textasciicircum{}(k-1) z\textasciicircum{}k / k is nowhere formed. The library's derivative of the logarithm is log-deriv (is-diff-at of x |-$>$ log(x) at c $>$ 0 with value recip-star(c)), proved from the integral definition, not from a series.
%
\renewcommand{\thethm}{2.21}
\begin{prop}[{The logarithm series agrees with ln on (0,2)}]\label{thm:log0-is-ln}
%
\leavevmode
\vnbpart{\texttt{log0-is-ln}}
\vnbstmt{\begin{array}{@{}l@{}} \forall c \in (\mathbb{N} \to \mathbb{C}), t \in \mathbb{R}. \\ \quad (c(0) = 0) \wedge \\ \quad \forall j \in \mathbb{N}. \\ \quad \quad (c((j + 1)) = (\operatorname{power}(-1, j) \cdot \operatorname{recip}((j + 1)))) \wedge \\ \quad (0 < t) \wedge \\ \quad (t < 2) \Rightarrow \\ \quad \quad \operatorname{cc\text{-}series\text{-}limit}((\lambda k \in \mathbb{N}.\; (c(k) \cdot \operatorname{power}((t - 1), k)))) \\ \quad \quad =\; \log\left(t\right) \end{array}}
\vnbpart{\texttt{log0-exp}}
\vnbstmt{\begin{array}{@{}l@{}} \forall c \in (\mathbb{N} \to \mathbb{C}), r \in \mathbb{R}, z \in \operatorname{ball}(\operatorname{cc\text{-}ms}, 1, r). \\ \quad (c(0) = 0) \wedge \\ \quad \forall j \in \mathbb{N}. \\ \quad \quad (c((j + 1)) = (\operatorname{power}(-1, j) \cdot \operatorname{recip}((j + 1)))) \wedge \\ \quad (0 < r) \wedge \\ \quad (r < 1) \Rightarrow \\ \quad \quad \exp \operatorname{cc\text{-}series\text{-}limit}((\lambda k \in \mathbb{N}.\; (c(k) \cdot \operatorname{power}((z - 1), k)))) \\ \quad \quad =\; z \end{array}}
\vnbpart{\texttt{log0-is-log-on-ball}}
\vnbstmt{\begin{array}{@{}l@{}} \forall c \in (\mathbb{N} \to \mathbb{C}), r \in \mathbb{R}, z \in \operatorname{ball}(\operatorname{cc\text{-}ms}, 1, r). \\ \quad (c(0) = 0) \wedge \\ \quad \forall j \in \mathbb{N}. \\ \quad \quad (c((j + 1)) = (\operatorname{power}(-1, j) \cdot \operatorname{recip}((j + 1)))) \wedge \\ \quad (0 < r) \wedge \\ \quad ((r \cdot (1 + \pi)) \leq \pi) \Rightarrow \\ \quad \quad \log\left(z\right) \\ \quad \quad =\; \operatorname{cc\text{-}series\text{-}limit}((\lambda k \in \mathbb{N}.\; (c(k) \cdot \operatorname{power}((z - 1), k)))) \end{array}}
\vnbpart{\texttt{cc-log-holomorphic-near-one}}
\vnbstmt{\begin{array}{@{}l@{}} \forall r \in \mathbb{R}. \\ \quad ((0 < r) \wedge ((r \cdot (1 + \pi)) \leq \pi)) \Rightarrow  \\ \quad \quad \operatorname{holomorphic\text{-}on}(\operatorname{ball}(\operatorname{cc\text{-}ms}, 1, r), (\lambda z \in \operatorname{ball}(\operatorname{cc\text{-}ms}, 1, r).\; \log\left(z\right))) \end{array}}
%
\end{prop}
%
\begin{proof}
%
The proofs: the file that holds each, the lemmas it cites, its oracles and its bill.
%
\vnbprf{\vnbproofpart{\texttt{log0-is-ln}}}
%
\vnbprf{\textbf{Proof.} \texttt{theorem-\allowbreak{}library/analytic-\allowbreak{}log-\allowbreak{}laws.scm} --- 126 recorded steps.}
%
\vnbprf{\textbf{Lemmas.} \texttt{cc-is-metric-space}, \texttt{cc-is-normed-field}, \texttt{cc-zero-in}, \texttt{cc-one-in}, \texttt{rr-one-in}, \texttt{rr-zero-in}, \texttt{rr-subset-cc}, \texttt{rr-sub-in-rr}, \texttt{rr-abs-closed}, \texttt{rr-abs-nonneg}, \texttt{rr-abs-lt-of-parts}, \texttt{rr-add-in-rr}, \texttt{rr-mul-in-rr}, \texttt{cc-ms-dist}, \texttt{rr-magnitude-is-abs}, \texttt{rr-abs-neg}, \texttt{ball-mem-from-le}, \texttt{log0-real}, \texttt{log0-exp}, \texttt{pi-pos}, \texttt{cc-log-char}, \texttt{cc-log-real}.}
%
\vnbprf{\textbf{Oracles.} \texttt{ineq}, \texttt{arith}, \texttt{crs}, \texttt{sos}.  \textbf{Bill:} \texttt{proven modulo 0}.}
%
\vnbprf{\vnbproofpart{\texttt{log0-exp}}}
%
\vnbprf{\textbf{Proof.} \texttt{theorem-\allowbreak{}library/analytic-\allowbreak{}log-\allowbreak{}laws.scm} --- 1023 recorded steps.}
%
\vnbprf{\textbf{Lemmas.} \texttt{cc-is-metric-space}, \texttt{cc-is-normed-field}, \texttt{cc-zero-in}, \texttt{cc-one-in}, \texttt{rr-one-in}, \texttt{rr-zero-in}, \texttt{nn-zero-in}, \texttt{cc-nf-carr}, \texttt{rr-lt-irrefl}, \texttt{cps-log-radius}, \texttt{ball-is-set}, \texttt{ball-is-open}, \texttt{cc-ms-open-iff}, \texttt{ball-center-in}, \texttt{alg-shift-in-ball}, \texttt{cc-sub-in-cc}, \texttt{cps-holomorphic-ball}, \texttt{holomorphic-on-in-fun}, \texttt{fun-apply-type-c}, \texttt{ball-membership}, \texttt{cc-magnitude-closed}, \texttt{cc-ms-dist}, \texttt{cc-magnitude-neg}, \texttt{rr-magnitude-is-abs}, \texttt{cc-exp-converges}, \texttt{cc-recip-closed}, \texttt{cc-mul-closed}, \texttt{cps-log-deriv-at-one}, \texttt{cc-exp-deriv}, \texttt{diff-on-chain}, \texttt{diff-on-identity}, \texttt{cc-mul-nonzero}, \texttt{diff-on-recip-cc}, \texttt{cc-neg-closed}, \texttt{diff-on-product}, \texttt{cc-nf-mul-apply}, \texttt{diff-on-transfer-ptwise-eq}, \texttt{cc-nf-add-apply}, \texttt{cc-recip-mul}, \texttt{alg-zero-deriv-ball}, \texttt{cps-sum-at-zero}, \texttt{cc-exp-zero}, \texttt{cc-recip-inverse}.}
%
\vnbprf{\textbf{Oracles.} \texttt{ineq}, \texttt{crs}, \texttt{arith}, \texttt{sos}.  \textbf{Bill:} \texttt{proven modulo 0}.}
%
\vnbprf{\vnbproofpart{\texttt{log0-is-log-on-ball}}}
%
\vnbprf{\textbf{Proof.} \texttt{theorem-\allowbreak{}library/analytic-\allowbreak{}log-\allowbreak{}laws.scm} --- 324 recorded steps.}
%
\vnbprf{\textbf{Lemmas.} \texttt{cc-is-metric-space}, \texttt{cc-is-normed-field}, \texttt{cc-zero-in}, \texttt{cc-one-in}, \texttt{rr-one-in}, \texttt{rr-zero-in}, \texttt{pi-pos}, \texttt{rr-subset-cc}, \texttt{rr-add-in-rr}, \texttt{rr-mul-in-rr}, \texttt{rr-lt-trichotomy}, \texttt{rr-mul-le-right}, \texttt{rr-lt-irrefl}, \texttt{cps-log-radius}, \texttt{ball-membership}, \texttt{cc-sub-in-cc}, \texttt{cc-magnitude-closed}, \texttt{cc-ms-dist}, \texttt{cc-magnitude-neg}, \texttt{alg-shift-in-ball}, \texttt{cps-holomorphic-ball}, \texttt{holomorphic-on-in-fun}, \texttt{fun-apply-type-c}, \texttt{imag-part-in-rr}, \texttt{log0-bound}, \texttt{cc-abs-im-le-magnitude}, \texttt{rr-sub-in-rr}, \texttt{cc-magnitude-nonneg}, \texttt{rr-recip-closed}, \texttt{rr-recip-pos}, \texttt{rr-recip-inverse}, \texttt{rr-abs-closed}, \texttt{rr-abs-le-parts}, \texttt{rr-lt-scale-pos}, \texttt{log0-exp}, \texttt{cc-log-char}.}
%
\vnbprf{\textbf{Oracles.} \texttt{crs}, \texttt{ineq}, \texttt{arith}, \texttt{sos}.  \textbf{Bill:} \texttt{proven modulo 0}.}
%
\vnbprf{\vnbproofpart{\texttt{cc-log-holomorphic-near-one}}}
%
\vnbprf{\textbf{Proof.} \texttt{theorem-\allowbreak{}library/analytic-\allowbreak{}log-\allowbreak{}laws.scm} --- 315 recorded steps.}
%
\vnbprf{\textbf{Lemmas.} \texttt{cc-is-metric-space}, \texttt{cc-is-normed-field}, \texttt{cc-zero-in}, \texttt{cc-one-in}, \texttt{rr-one-in}, \texttt{rr-zero-in}, \texttt{nn-zero-in}, \texttt{pi-pos}, \texttt{rr-subset-cc}, \texttt{rr-add-in-rr}, \texttt{rr-mul-in-rr}, \texttt{rr-lt-trichotomy}, \texttt{rr-mul-le-right}, \texttt{rr-lt-irrefl}, \texttt{ball-is-set}, \texttt{ball-is-open}, \texttt{cc-ms-open-iff}, \texttt{nn-in-rr}, \texttt{recip-star-in-rr}, \texttt{power-real-closed}, \texttt{cc-neg-closed}, \texttt{cc-mul-closed}, \texttt{nn-is-set}, \texttt{recip-star-zero}, \texttt{nn-succ-closed}, \texttt{nn-one-le-succ}, \texttt{rr-recip-closed}, \texttt{power-succ}, \texttt{recip-star-value}, \texttt{cps-log-radius}, \texttt{cps-translate-holomorphic}, \texttt{log0-is-log-on-ball}, \texttt{alg-shift-in-ball}, \texttt{cc-sub-in-cc}, \texttt{cps-holomorphic-ball}, \texttt{holomorphic-on-in-fun}, \texttt{fun-apply-type-c}, \texttt{holomorphic-transfer-ptwise-eq}.}
%
\vnbprf{\textbf{Oracles.} \texttt{crs}, \texttt{ineq}, \texttt{arith}, \texttt{sos}.  \textbf{Bill:} \texttt{proven modulo 0}.}
%
\end{proof}
%
\noindent\textbf{Deviation from the notes.} PROVEN 2026-09-27 (batch 38) in the notes' wording: log0-is-ln says the series (39) at a real t in (0, 2) equals LOG(t), by the notes' argument (equal derivatives, equal value at 1). Also exp(log0 z) = z on every B(1, r), r $<$ 1 (log0-exp), and log0 IS the principal logarithm CC-LOG on B(1, pi/(1+pi)) (log0-is-log-on-ball: the radius where |log0 z| $<$ pi keeps the value in the strip; the full B(1, 1) waits for the identity theorem 3.35). Earlier: Nothing near: the power series log\_0 z = sum\_\{k$>$=1\} (-1)\textasciicircum{}(k-1) (z-1)\textasciicircum{}k / k does not exist in the library. LOG itself is present and is exactly the notes' ln (log-unfold: LOG(x) = the oriented integral of 1/z from 1 to x, with log-one, log-mul, log-deriv, log-strictly-increasing proven), so only the series and its identification with LOG are missing.
%
\renewcommand{\thethm}{2.3.3.1}
\begin{defn}[{The principal branch of the logarithm}]\label{thm:cc-log-spec}
%
\leavevmode
\vnbpart{\texttt{cc-log-spec}}
\vnbstmt{\begin{array}{@{}l@{}} \forall w \in \mathbb{C}. \\ \quad \neg ((\operatorname{imag\text{-}part}(w) = 0) \wedge (\operatorname{real\text{-}part}(w) \leq 0)) \Rightarrow  \\ \quad \quad (\log\left(w\right) \in \mathbb{C}) \wedge \\ \quad \quad (-\pi < \operatorname{imag\text{-}part}(\log\left(w\right))) \wedge \\ \quad \quad (\operatorname{imag\text{-}part}(\log\left(w\right)) < \pi) \wedge \\ \quad \quad (\exp \log\left(w\right) = w) \end{array}}
\vnbpart{\texttt{cc-exp-log}}
\vnbstmt{\begin{array}{@{}l@{}} \forall w \in \mathbb{C}. \\ \quad \neg ((\operatorname{imag\text{-}part}(w) = 0) \wedge (\operatorname{real\text{-}part}(w) \leq 0)) \Rightarrow  \\ \quad \quad (\exp \log\left(w\right) = w) \end{array}}
\vnbpart{\texttt{cc-log-exp}}
\vnbstmt{\begin{array}{@{}l@{}} \forall z \in \mathbb{C}. \\ \quad ((-\pi < \operatorname{imag\text{-}part}(z)) \wedge (\operatorname{imag\text{-}part}(z) < \pi)) \Rightarrow  \\ \quad \quad (\log\left(\exp z\right) = z) \end{array}}
\vnbpart{\texttt{cc-log-char}}
\vnbstmt{\begin{array}{@{}l@{}} \forall w, z \in \mathbb{C}. \\ \quad (-\pi < \operatorname{imag\text{-}part}(z)) \wedge \\ \quad (\operatorname{imag\text{-}part}(z) < \pi) \wedge \\ \quad (\exp z = w) \Rightarrow \\ \quad \quad (\log\left(w\right) = z) \end{array}}
\vnbpart{\texttt{cc-log-unique}}
\vnbstmt{\begin{array}{@{}l@{}} \forall a, b \in \mathbb{C}. \\ \quad (-\pi < \operatorname{imag\text{-}part}(a)) \wedge \\ \quad (\operatorname{imag\text{-}part}(a) < \pi) \wedge \\ \quad (-\pi < \operatorname{imag\text{-}part}(b)) \wedge \\ \quad (\operatorname{imag\text{-}part}(b) < \pi) \wedge \\ \quad (\exp a = \exp b) \Rightarrow \\ \quad \quad (a = b) \end{array}}
\vnbpart{\texttt{cc-log-real}}
\vnbstmt{\begin{array}{@{}l@{}} \forall t \in \mathbb{R}. \\ \quad ((0 < t) \Rightarrow (\log\left(t\right) = \log\left(t\right))) \end{array}}
\vnbpart{\texttt{cc-log-one}}
\vnbstmt{(\log\left(1\right) = 0)}
\vnbpart{\texttt{cc-log-polar}}
\vnbstmt{\begin{array}{@{}l@{}} \forall r, h \in \mathbb{R}. \\ \quad ((0 < r) \wedge ((-\pi < h) \wedge (h < \pi))) \Rightarrow  \\ \quad \quad (\log\left((r \cdot \exp (+i \cdot h))\right) = (\log\left(r\right) + (+i \cdot h))) \end{array}}
\vnbpart{\texttt{cc-log-scale}}
\vnbstmt{\begin{array}{@{}l@{}} \forall s \in \mathbb{R}, w \in \mathbb{C}. \\ \quad (0 < s) \wedge \\ \quad \neg ((\operatorname{imag\text{-}part}(w) = 0) \wedge (\operatorname{real\text{-}part}(w) \leq 0)) \Rightarrow \\ \quad \quad (\log\left((s \cdot w)\right) = (\log\left(s\right) + \log\left(w\right))) \end{array}}
\vnbpart{\texttt{cc-log-in-slit}}
\vnbstmt{\begin{array}{@{}l@{}} \forall w \in \mathbb{C}. \\ \quad \neg ((\operatorname{imag\text{-}part}(w) = 0) \wedge (\operatorname{real\text{-}part}(w) \leq 0)) \Rightarrow  \\ \quad \quad (\log\left(w\right) \in \mathbb{C}) \wedge \\ \quad \quad (-\pi < \operatorname{imag\text{-}part}(\log\left(w\right))) \wedge \\ \quad \quad (\operatorname{imag\text{-}part}(\log\left(w\right)) < \pi) \end{array}}
%
\end{defn}
%
\begin{proof}
%
The proofs: the file that holds each, the lemmas it cites, its oracles and its bill.
%
\vnbprf{\vnbproofpart{\texttt{cc-log-spec}}}
%
\vnbprf{\textbf{Proof.} \texttt{theorem-\allowbreak{}library/analytic-\allowbreak{}log-\allowbreak{}laws.scm} --- 34 recorded steps.}
%
\vnbprf{\textbf{Lemmas.} \texttt{cc-exp-open-strip-onto}, \texttt{cc-exp-converges}, \texttt{cc-log-unique}.}
%
\vnbprf{\textbf{Oracles.} \texttt{arith}, \texttt{crs}, \texttt{ineq}, \texttt{sos}.  \textbf{Bill:} \texttt{proven modulo 0}.}
%
\vnbprf{\vnbproofpart{\texttt{cc-exp-log}}}
%
\vnbprf{\textbf{Proof.} \texttt{theorem-\allowbreak{}library/analytic-\allowbreak{}log-\allowbreak{}laws.scm} --- 7 recorded steps.}
%
\vnbprf{\textbf{Lemmas.} \texttt{cc-log-spec}.}
%
\vnbprf{\textbf{Oracles.} \texttt{arith}, \texttt{crs}, \texttt{ineq}, \texttt{sos}.  \textbf{Bill:} \texttt{proven modulo 0}.}
%
\vnbprf{\vnbproofpart{\texttt{cc-log-exp}}}
%
\vnbprf{\textbf{Proof.} \texttt{theorem-\allowbreak{}library/analytic-\allowbreak{}log-\allowbreak{}laws.scm} --- 11 recorded steps.}
%
\vnbprf{\textbf{Lemmas.} \texttt{cc-exp-converges}, \texttt{cc-log-char}.}
%
\vnbprf{\textbf{Oracles.} \texttt{arith}, \texttt{ineq}, \texttt{crs}, \texttt{sos}.  \textbf{Bill:} \texttt{proven modulo 0}.}
%
\vnbprf{\vnbproofpart{\texttt{cc-log-char}}}
%
\vnbprf{\textbf{Proof.} \texttt{theorem-\allowbreak{}library/analytic-\allowbreak{}log-\allowbreak{}laws.scm} --- 24 recorded steps.}
%
\vnbprf{\textbf{Lemmas.} \texttt{cc-exp-open-strip-avoids-negatives}, \texttt{cc-log-spec}, \texttt{cc-exp-converges}, \texttt{cc-log-unique}.}
%
\vnbprf{\textbf{Oracles.} \texttt{crs}, \texttt{ineq}, \texttt{arith}, \texttt{sos}.  \textbf{Bill:} \texttt{proven modulo 0}.}
%
\vnbprf{\vnbproofpart{\texttt{cc-log-unique}}}
%
\vnbprf{\textbf{Proof.} \texttt{theorem-\allowbreak{}library/analytic-\allowbreak{}log-\allowbreak{}laws.scm} --- 21 recorded steps.}
%
\vnbprf{\textbf{Lemmas.} \texttt{pi-pos}, \texttt{rr-subset-cc}, \texttt{imag-part-in-rr}, \texttt{cc-exp-injective-strip}.}
%
\vnbprf{\textbf{Oracles.} \texttt{ineq}, \texttt{arith}, \texttt{crs}, \texttt{sos}.  \textbf{Bill:} \texttt{proven modulo 0}.}
%
\vnbprf{\vnbproofpart{\texttt{cc-log-real}}}
%
\vnbprf{\textbf{Proof.} \texttt{theorem-\allowbreak{}library/analytic-\allowbreak{}log-\allowbreak{}laws.scm} --- 53 recorded steps.}
%
\vnbprf{\textbf{Lemmas.} \texttt{pi-pos}, \texttt{rr-subset-cc}, \texttt{log-in-rr}, \texttt{cc-exp-real-is-r-exp}, \texttt{r-exp-log}, \texttt{rr-zero-in}, \texttt{cc-add-closed}, \texttt{cc-re-im-of}, \texttt{cc-log-char}.}
%
\vnbprf{\textbf{Oracles.} \texttt{arith}, \texttt{crs}, \texttt{ineq}, \texttt{sos}.  \textbf{Bill:} \texttt{proven modulo 0}.}
%
\vnbprf{\vnbproofpart{\texttt{cc-log-one}}}
%
\vnbprf{\textbf{Proof.} \texttt{theorem-\allowbreak{}library/analytic-\allowbreak{}log-\allowbreak{}laws.scm} --- 11 recorded steps.}
%
\vnbprf{\textbf{Lemmas.} \texttt{rr-one-in}, \texttt{cc-log-real}, \texttt{log-one}, \texttt{rr-zero-in}.}
%
\vnbprf{\textbf{Oracles.} \texttt{ineq}, \texttt{arith}, \texttt{crs}, \texttt{sos}.  \textbf{Bill:} \texttt{proven modulo 0}.}
%
\vnbprf{\vnbproofpart{\texttt{cc-log-polar}}}
%
\vnbprf{\textbf{Proof.} \texttt{theorem-\allowbreak{}library/analytic-\allowbreak{}log-\allowbreak{}laws.scm} --- 74 recorded steps.}
%
\vnbprf{\textbf{Lemmas.} \texttt{pi-pos}, \texttt{rr-subset-cc}, \texttt{log-in-rr}, \texttt{cc-exp-real-is-r-exp}, \texttt{r-exp-log}, \texttt{cc-i-in}, \texttt{cc-mul-closed}, \texttt{cc-add-closed}, \texttt{cc-exp-converges}, \texttt{cc-re-im-of}, \texttt{cc-exp-add}, \texttt{cc-log-char}.}
%
\vnbprf{\textbf{Oracles.} \texttt{crs}, \texttt{arith}, \texttt{ineq}, \texttt{sos}.  \textbf{Bill:} \texttt{proven modulo 0}.}
%
\vnbprf{\vnbproofpart{\texttt{cc-log-scale}}}
%
\vnbprf{\textbf{Proof.} \texttt{theorem-\allowbreak{}library/analytic-\allowbreak{}log-\allowbreak{}laws.scm} --- 90 recorded steps.}
%
\vnbprf{\textbf{Lemmas.} \texttt{cc-log-spec}, \texttt{log-in-rr}, \texttt{rr-subset-cc}, \texttt{cc-exp-real-is-r-exp}, \texttt{r-exp-log}, \texttt{cc-add-closed}, \texttt{cc-mul-closed}, \texttt{rr-zero-in}, \texttt{cc-re-im-of}, \texttt{imag-part-in-rr}, \texttt{cc-im-add}, \texttt{cc-exp-add}, \texttt{cc-log-char}.}
%
\vnbprf{\textbf{Oracles.} \texttt{arith}, \texttt{crs}, \texttt{ineq}, \texttt{sos}.  \textbf{Bill:} \texttt{proven modulo 0}.}
%
\vnbprf{\vnbproofpart{\texttt{cc-log-in-slit}}}
%
\vnbprf{\textbf{Proof.} \texttt{theorem-\allowbreak{}library/analytic-\allowbreak{}log-\allowbreak{}laws.scm} --- 15 recorded steps.}
%
\vnbprf{\textbf{Lemmas.} \texttt{cc-log-spec}.}
%
\vnbprf{\textbf{Oracles.} \texttt{arith}, \texttt{crs}, \texttt{ineq}, \texttt{sos}.  \textbf{Bill:} \texttt{proven modulo 0}.}
%
\end{proof}
%
\noindent\textbf{Deviation from the notes.} DEFINED 2026-09-27 (batch 38): CC-LOG(w) is the IOTA of the z with -pi $<$ Im z $<$ pi and exp z = w, for w in the slit plane \{u in cc : not(Im u = 0 and Re u $<$= 0)\} (structure-library/analytic-log.scm); existence is cc-exp-open-strip-onto (batch 34), uniqueness cc-log-unique, the characterisation cc-log-spec by iota-d. Laws: exp(log w) = w, log(exp z) = z on the strip, log = LOG on the positive reals (through the bridge cc-exp-real-is-r-exp), log 1 = 0, the polar form (40) log(r exp(i th)) = LOG r + i th, log(s w) = LOG s + log w for real s $>$ 0.
%
\renewcommand{\thethm}{2.3.3.1-hol}
\begin{prop}[{log is holomorphic on C minus (-inf, 0]}]\label{thm:cc-log-holomorphic}
%
\leavevmode
\vnbpart{\texttt{cc-log-holomorphic}}
\vnbstmt{\begin{array}{@{}l@{}} \operatorname{holomorphic\text{-}on}(\{u \in \mathbb{C} : \neg ((\operatorname{imag\text{-}part}(u) = 0) \wedge (\operatorname{real\text{-}part}(u) \leq 0))\}, \\ \quad (\lambda w \in \{u \in \mathbb{C} : \neg ((\operatorname{imag\text{-}part}(u) = 0) \wedge (\operatorname{real\text{-}part}(u) \leq 0))\}.\; \\ \quad \quad \log\left(w\right))) \end{array}}
\vnbpart{\texttt{cc-log-deriv}}
\vnbstmt{\begin{array}{@{}l@{}} \forall b \in \mathbb{C}. \\ \quad \neg ((\operatorname{imag\text{-}part}(b) = 0) \wedge (\operatorname{real\text{-}part}(b) \leq 0)) \Rightarrow  \\ \quad \quad \operatorname{is\text{-}diff\text{-}on}(\operatorname{cc\text{-}normed\text{-}field}, \\ \quad \quad \quad \{u \in \mathbb{C} : \neg ((\operatorname{imag\text{-}part}(u) = 0) \wedge (\operatorname{real\text{-}part}(u) \leq 0))\}, \\ \quad \quad \quad (\lambda w \in \{u \in \mathbb{C} : \neg ((\operatorname{imag\text{-}part}(u) = 0) \wedge (\operatorname{real\text{-}part}(u) \leq 0))\}.\; \\ \quad \quad \quad \quad \log\left(w\right)), \\ \quad \quad \quad b, \operatorname{recip}(b)) \end{array}}
\vnbpart{\texttt{cc-log-holomorphic-at}}
\vnbstmt{\begin{array}{@{}l@{}} \forall a \in \mathbb{C}. \\ \quad \neg ((\operatorname{imag\text{-}part}(a) = 0) \wedge (\operatorname{real\text{-}part}(a) \leq 0)) \Rightarrow  \\ \quad \quad \exists p \in \mathbb{R}. \\ \quad \quad \quad (0 < p) \wedge \\ \quad \quad \quad \operatorname{ball}(\operatorname{cc\text{-}ms}, a, p) \\ \quad \quad \quad \subseteq\; \{u \in \mathbb{C} : \neg ((\operatorname{imag\text{-}part}(u) = 0) \wedge (\operatorname{real\text{-}part}(u) \leq 0))\} \wedge \\ \quad \quad \quad \operatorname{holomorphic\text{-}on}(\operatorname{ball}(\operatorname{cc\text{-}ms}, a, p), (\lambda x \in \operatorname{ball}(\operatorname{cc\text{-}ms}, a, p).\; \log\left(x\right))) \end{array}}
\vnbpart{\texttt{cc-log-im-near-one}}
\vnbstmt{\begin{array}{@{}l@{}} \forall d, r \in \mathbb{R}, z \in \operatorname{ball}(\operatorname{cc\text{-}ms}, 1, r). \\ \quad ((0 < d) \wedge ((0 < r) \wedge (((r \cdot (1 + \pi)) \leq \pi) \wedge ((r \cdot (1 + d)) \leq d)))) \Rightarrow  \\ \quad \quad \neg ((\operatorname{imag\text{-}part}(z) = 0) \wedge (\operatorname{real\text{-}part}(z) \leq 0)) \wedge \\ \quad \quad (-d < \operatorname{imag\text{-}part}(\log\left(z\right))) \wedge \\ \quad \quad (\operatorname{imag\text{-}part}(\log\left(z\right)) < d) \end{array}}
\vnbpart{\texttt{cc-log-rotate}}
\vnbstmt{\begin{array}{@{}l@{}} \forall w \in \mathbb{C}. \\ \quad (0 < \operatorname{real\text{-}part}(w)) \Rightarrow  \\ \quad \quad (\log\left((+i \cdot w)\right) = (\log\left(w\right) + (+i \cdot (1/2 \cdot \pi)))) \end{array}}
\vnbpart{\texttt{cc-log-rotate-neg}}
\vnbstmt{\begin{array}{@{}l@{}} \forall w \in \mathbb{C}. \\ \quad (0 < \operatorname{real\text{-}part}(w)) \Rightarrow  \\ \quad \quad (\log\left((-+i \cdot w)\right) = (\log\left(w\right) - (+i \cdot (1/2 \cdot \pi)))) \end{array}}
\vnbpart{\texttt{cc-log-mul}}
\vnbstmt{\begin{array}{@{}l@{}} \forall u, v \in \mathbb{C}. \\ \quad \neg ((\operatorname{imag\text{-}part}(u) = 0) \wedge (\operatorname{real\text{-}part}(u) \leq 0)) \wedge \\ \quad \neg ((\operatorname{imag\text{-}part}(v) = 0) \wedge (\operatorname{real\text{-}part}(v) \leq 0)) \wedge \\ \quad -\pi \\ \quad <\; (\operatorname{imag\text{-}part}(\log\left(u\right)) + \operatorname{imag\text{-}part}(\log\left(v\right))) \wedge \\ \quad (\operatorname{imag\text{-}part}(\log\left(u\right)) + \operatorname{imag\text{-}part}(\log\left(v\right))) \\ \quad <\; \pi \Rightarrow \\ \quad \quad (\log\left((u \cdot v)\right) = (\log\left(u\right) + \log\left(v\right))) \end{array}}
\vnbpart{\texttt{cc-log-right-half}}
\vnbstmt{\begin{array}{@{}l@{}} \forall w \in \mathbb{C}. \\ \quad (0 < \operatorname{real\text{-}part}(w)) \Rightarrow  \\ \quad \quad (-(1/2 \cdot \pi) < \operatorname{imag\text{-}part}(\log\left(w\right))) \wedge \\ \quad \quad (\operatorname{imag\text{-}part}(\log\left(w\right)) < (1/2 \cdot \pi)) \end{array}}
\vnbpart{\texttt{cc-pi-gt-three}}
\vnbstmt{(3 < \pi)}
\vnbpart{\texttt{cc-cos-pos-half}}
\vnbstmt{\begin{array}{@{}l@{}} \forall t \in \mathbb{R}. \\ \quad (((-(1/2 \cdot \pi) < t) \wedge (t < (1/2 \cdot \pi))) \Rightarrow (0 < \cos t)) \end{array}}
\vnbpart{\texttt{cc-log-holomorphic-near-one}}
\vnbstmt{\begin{array}{@{}l@{}} \forall r \in \mathbb{R}. \\ \quad ((0 < r) \wedge ((r \cdot (1 + \pi)) \leq \pi)) \Rightarrow  \\ \quad \quad \operatorname{holomorphic\text{-}on}(\operatorname{ball}(\operatorname{cc\text{-}ms}, 1, r), (\lambda z \in \operatorname{ball}(\operatorname{cc\text{-}ms}, 1, r).\; \log\left(z\right))) \end{array}}
\vnbpart{\texttt{cc-log-holomorphic-scale}}
\vnbstmt{\begin{array}{@{}l@{}} \forall r, s \in \mathbb{R}. \\ \quad ((0 < r) \wedge (((r \cdot (1 + \pi)) \leq \pi) \wedge (0 < s))) \Rightarrow  \\ \quad \quad \operatorname{holomorphic\text{-}on}(\operatorname{ball}(\operatorname{cc\text{-}ms}, s, (s \cdot r)), \\ \quad \quad \quad (\lambda z \in \operatorname{ball}(\operatorname{cc\text{-}ms}, s, (s \cdot r)).\; \\ \quad \quad \quad \quad \log\left(z\right))) \end{array}}
%
\end{prop}
%
\begin{proof}
%
The proofs: the file that holds each, the lemmas it cites, its oracles and its bill.
%
\vnbprf{\vnbproofpart{\texttt{cc-log-holomorphic}}}
%
\vnbprf{\textbf{Proof.} \texttt{theorem-\allowbreak{}library/cc-\allowbreak{}log-\allowbreak{}holomorphic.scm} --- 192 recorded steps.}
%
\vnbprf{\textbf{Lemmas.} \texttt{cc-is-metric-space}, \texttt{cc-is-normed-field}, \texttt{cc-zero-in}, \texttt{rr-zero-in}, \texttt{cc-log-holomorphic-at}, \texttt{cc-is-set}, \texttt{cc-log-in-slit}, \texttt{subset-def}, \texttt{cc-ms-open-iff}, \texttt{rr-lt-irrefl}, \texttt{ball-center-in}, \texttt{holomorphic-on-diff}, \texttt{subset-mem-fwd}, \texttt{cc-diff-on-local}.}
%
\vnbprf{\textbf{Oracles.} \texttt{ineq}, \texttt{arith}, \texttt{crs}, \texttt{sos}.  \textbf{Bill:} \texttt{proven modulo 0}.}
%
\vnbprf{\vnbproofpart{\texttt{cc-log-deriv}}}
%
\vnbprf{\textbf{Proof.} \texttt{theorem-\allowbreak{}library/cc-\allowbreak{}log-\allowbreak{}holomorphic.scm} --- 184 recorded steps.}
%
\vnbprf{\textbf{Lemmas.} \texttt{cc-is-metric-space}, \texttt{cc-is-normed-field}, \texttt{cc-zero-in}, \texttt{rr-zero-in}, \texttt{cc-one-in}, \texttt{cc-is-set}, \texttt{cc-log-holomorphic}, \texttt{holomorphic-on-open}, \texttt{holomorphic-on-in-fun}, \texttt{holomorphic-on-diff}, \texttt{cc-log-spec}, \texttt{cc-nf-carr}, \texttt{fun-apply-type-c}, \texttt{cc-exp-converges}, \texttt{cc-exp-deriv}, \texttt{diff-on-chain}, \texttt{diff-on-identity}, \texttt{subset-refl}, \texttt{cc-diff-on-local}, \texttt{cc-nf-small-elements}, \texttt{cc-nf-mul-apply}, \texttt{diff-on-unique}, \texttt{cc-recip-closed}, \texttt{cc-recip-inverse}.}
%
\vnbprf{\textbf{Oracles.} \texttt{arith}, \texttt{crs}, \texttt{ineq}, \texttt{sos}.  \textbf{Bill:} \texttt{proven modulo 0}.}
%
\vnbprf{\vnbproofpart{\texttt{cc-log-holomorphic-at}}}
%
\vnbprf{\textbf{Proof.} \texttt{theorem-\allowbreak{}library/cc-\allowbreak{}log-\allowbreak{}holomorphic.scm} --- 817 recorded steps.}
%
\vnbprf{\textbf{Lemmas.} \texttt{cc-is-metric-space}, \texttt{cc-is-normed-field}, \texttt{cc-zero-in}, \texttt{cc-one-in}, \texttt{rr-one-in}, \texttt{rr-zero-in}, \texttt{pi-pos}, \texttt{rr-subset-cc}, \texttt{cc-log-spec}, \texttt{imag-part-in-rr}, \texttt{rr-abs-closed}, \texttt{rr-sub-in-rr}, \texttt{rr-abs-nonneg}, \texttt{rr-abs-le-parts}, \texttt{rr-abs-lt-of-parts}, \texttt{rr-add-in-rr}, \texttt{rr-mul-in-rr}, \texttt{rr-lt-irrefl}, \texttt{rr-recip-closed}, \texttt{rr-recip-pos}, \texttt{rr-recip-inverse}, \texttt{rr-lt-scale-pos}, \texttt{rr-le-scale-nonneg}, \texttt{cc-log-im-near-one}, \texttt{cc-magnitude-closed}, \texttt{cc-magnitude-nonneg}, \texttt{cc-magnitude-zero-iff}, \texttt{rr-le-ne-lt}, \texttt{cc-recip-closed}, \texttt{cc-recip-inverse}, \texttt{cc-magnitude-mul}, \texttt{cc-mul-closed}, \texttt{ball-membership}, \texttt{cc-sub-in-cc}, \texttt{cc-ms-dist}, \texttt{cc-magnitude-neg}, \texttt{ball-mem-from-le}, \texttt{cc-add-closed}, \texttt{cc-exp-converges}, \texttt{cc-im-add}, \texttt{cc-exp-add}, \texttt{cc-log-char}, \texttt{cc-exp-open-strip-avoids-negatives}, \texttt{subset-def}, \texttt{ball-is-set}, \texttt{ball-is-open}, \texttt{cc-ms-open-iff}, \texttt{holomorphic-const}, \texttt{holomorphic-identity}, \texttt{holomorphic-product}, \texttt{cc-log-holomorphic-near-one}, \texttt{holomorphic-compose}, \texttt{cc-log-in-slit}, \texttt{holomorphic-sum}.}
%
\vnbprf{\textbf{Oracles.} \texttt{ineq}, \texttt{arith}, \texttt{crs}, \texttt{sos}.  \textbf{Bill:} \texttt{proven modulo 0}.}
%
\vnbprf{\vnbproofpart{\texttt{cc-log-im-near-one}}}
%
\vnbprf{\textbf{Proof.} \texttt{theorem-\allowbreak{}library/cc-\allowbreak{}log-\allowbreak{}holomorphic.scm} --- 507 recorded steps.}
%
\vnbprf{\textbf{Lemmas.} \texttt{cc-is-metric-space}, \texttt{cc-is-normed-field}, \texttt{cc-zero-in}, \texttt{cc-one-in}, \texttt{rr-one-in}, \texttt{rr-zero-in}, \texttt{nn-zero-in}, \texttt{pi-pos}, \texttt{rr-subset-cc}, \texttt{rr-lt-irrefl}, \texttt{rr-add-in-rr}, \texttt{rr-mul-in-rr}, \texttt{rr-lt-trichotomy}, \texttt{rr-mul-le-right}, \texttt{nn-in-rr}, \texttt{recip-star-in-rr}, \texttt{power-real-closed}, \texttt{cc-neg-closed}, \texttt{cc-mul-closed}, \texttt{nn-is-set}, \texttt{recip-star-zero}, \texttt{nn-succ-closed}, \texttt{nn-one-le-succ}, \texttt{rr-recip-closed}, \texttt{power-succ}, \texttt{recip-star-value}, \texttt{ball-membership}, \texttt{cc-sub-in-cc}, \texttt{cc-magnitude-closed}, \texttt{cc-ms-dist}, \texttt{cc-magnitude-neg}, \texttt{log0-is-log-on-ball}, \texttt{log0-bound}, \texttt{cc-magnitude-nonneg}, \texttt{real-part-in-rr}, \texttt{cc-abs-re-le-magnitude}, \texttt{rr-abs-le-parts}, \texttt{cc-re-sub}, \texttt{cc-re-im-of}, \texttt{cc-log-in-slit}, \texttt{imag-part-in-rr}, \texttt{cc-abs-im-le-magnitude}, \texttt{rr-sub-in-rr}, \texttt{rr-recip-pos}, \texttt{rr-recip-inverse}, \texttt{rr-abs-closed}, \texttt{rr-lt-scale-pos}.}
%
\vnbprf{\textbf{Oracles.} \texttt{ineq}, \texttt{crs}, \texttt{arith}, \texttt{sos}.  \textbf{Bill:} \texttt{proven modulo 0}.}
%
\vnbprf{\vnbproofpart{\texttt{cc-log-rotate}}}
%
\vnbprf{\textbf{Proof.} \texttt{theorem-\allowbreak{}library/cc-\allowbreak{}log-\allowbreak{}holomorphic.scm} --- 119 recorded steps.}
%
\vnbprf{\textbf{Lemmas.} \texttt{pi-pos}, \texttt{rr-subset-cc}, \texttt{rr-mul-in-rr}, \texttt{real-part-in-rr}, \texttt{rr-zero-in}, \texttt{rr-lt-irrefl}, \texttt{cc-log-spec}, \texttt{cc-log-right-half}, \texttt{cc-i-in}, \texttt{cc-mul-closed}, \texttt{cc-add-closed}, \texttt{cc-exp-converges}, \texttt{cc-cos-converges}, \texttt{cc-sin-converges}, \texttt{cc-exp-add}, \texttt{cc-exp-i}, \texttt{cc-cos-half-pi}, \texttt{cc-sin-half-pi}, \texttt{imag-part-in-rr}, \texttt{cc-im-add}, \texttt{cc-im-real-mul}, \texttt{cc-log-char}.}
%
\vnbprf{\textbf{Oracles.} \texttt{arith}, \texttt{ineq}, \texttt{crs}, \texttt{sos}.  \textbf{Bill:} \texttt{proven modulo 0}.}
%
\vnbprf{\vnbproofpart{\texttt{cc-log-rotate-neg}}}
%
\vnbprf{\textbf{Proof.} \texttt{theorem-\allowbreak{}library/cc-\allowbreak{}log-\allowbreak{}holomorphic.scm} --- 148 recorded steps.}
%
\vnbprf{\textbf{Lemmas.} \texttt{pi-pos}, \texttt{rr-subset-cc}, \texttt{rr-mul-in-rr}, \texttt{real-part-in-rr}, \texttt{rr-zero-in}, \texttt{rr-lt-irrefl}, \texttt{cc-log-spec}, \texttt{cc-log-right-half}, \texttt{rr-neg-closed}, \texttt{cc-i-in}, \texttt{cc-mul-closed}, \texttt{cc-add-closed}, \texttt{cc-sub-in-cc}, \texttt{cc-exp-converges}, \texttt{cc-cos-converges}, \texttt{cc-sin-converges}, \texttt{cc-exp-add}, \texttt{cc-exp-i}, \texttt{cc-cos-even}, \texttt{cc-sin-odd}, \texttt{cc-cos-half-pi}, \texttt{cc-sin-half-pi}, \texttt{imag-part-in-rr}, \texttt{cc-im-add}, \texttt{cc-im-real-mul}, \texttt{cc-log-char}.}
%
\vnbprf{\textbf{Oracles.} \texttt{arith}, \texttt{ineq}, \texttt{crs}, \texttt{sos}.  \textbf{Bill:} \texttt{proven modulo 0}.}
%
\vnbprf{\vnbproofpart{\texttt{cc-log-mul}}}
%
\vnbprf{\textbf{Proof.} \texttt{theorem-\allowbreak{}library/cc-\allowbreak{}log-\allowbreak{}holomorphic.scm} --- 61 recorded steps.}
%
\vnbprf{\textbf{Lemmas.} \texttt{pi-pos}, \texttt{rr-subset-cc}, \texttt{cc-log-spec}, \texttt{cc-add-closed}, \texttt{cc-mul-closed}, \texttt{cc-exp-converges}, \texttt{imag-part-in-rr}, \texttt{cc-im-add}, \texttt{cc-exp-add}, \texttt{cc-log-char}.}
%
\vnbprf{\textbf{Oracles.} \texttt{arith}, \texttt{crs}, \texttt{ineq}, \texttt{sos}.  \textbf{Bill:} \texttt{proven modulo 0}.}
%
\vnbprf{\vnbproofpart{\texttt{cc-log-right-half}}}
%
\vnbprf{\textbf{Proof.} \texttt{theorem-\allowbreak{}library/cc-\allowbreak{}log-\allowbreak{}holomorphic.scm} --- 290 recorded steps.}
%
\vnbprf{\textbf{Lemmas.} \texttt{pi-pos}, \texttt{rr-subset-cc}, \texttt{rr-zero-in}, \texttt{rr-mul-in-rr}, \texttt{real-part-in-rr}, \texttt{rr-lt-irrefl}, \texttt{cc-log-spec}, \texttt{imag-part-in-rr}, \texttt{cc-i-in}, \texttt{cc-mul-closed}, \texttt{cc-add-closed}, \texttt{cc-exp-real-valued}, \texttt{cc-exp-converges}, \texttt{cc-re-im-decompose}, \texttt{cc-exp-add}, \texttt{cc-cos-real-part}, \texttt{cc-cos-sin-real}, \texttt{cc-re-real-mul}, \texttt{cc-exp-real-is-r-exp}, \texttt{r-exp-pos}, \texttt{r-exp-in-rr}, \texttt{rr-neg-closed}, \texttt{cc-cos-even}, \texttt{rr-lt-trichotomy}, \texttt{cc-cos-nonpos-far}, \texttt{rr-le-scale-nonneg}.}
%
\vnbprf{\textbf{Oracles.} \texttt{arith}, \texttt{ineq}, \texttt{crs}, \texttt{sos}.  \textbf{Bill:} \texttt{proven modulo 0}.}
%
\vnbprf{\vnbproofpart{\texttt{cc-pi-gt-three}}}
%
\vnbprf{\textbf{Proof.} \texttt{theorem-\allowbreak{}library/cc-\allowbreak{}log-\allowbreak{}holomorphic.scm} --- 55 recorded steps.}
%
\vnbprf{\textbf{Lemmas.} \texttt{pi-pos}, \texttt{rr-subset-cc}, \texttt{rr-zero-in}, \texttt{cc-pi-gt-two}, \texttt{cc-cos-three-halves-pos}, \texttt{cc-cos-sin-real}, \texttt{rr-mul-in-rr}, \texttt{rr-lt-trichotomy}, \texttt{cc-cos-nonpos-far}, \texttt{rr-lt-irrefl}.}
%
\vnbprf{\textbf{Oracles.} \texttt{arith}, \texttt{ineq}, \texttt{crs}, \texttt{sos}.  \textbf{Bill:} \texttt{proven modulo 0}.}
%
\vnbprf{\vnbproofpart{\texttt{cc-cos-pos-half}}}
%
\vnbprf{\textbf{Proof.} \texttt{theorem-\allowbreak{}library/cc-\allowbreak{}log-\allowbreak{}holomorphic.scm} --- 42 recorded steps.}
%
\vnbprf{\textbf{Lemmas.} \texttt{pi-pos}, \texttt{rr-subset-cc}, \texttt{rr-zero-in}, \texttt{rr-mul-in-rr}, \texttt{cc-cos-sin-real}, \texttt{rr-lt-trichotomy}, \texttt{rr-neg-closed}, \texttt{cc-cos-pos-right}, \texttt{cc-cos-even}, \texttt{cc-cos-zero}.}
%
\vnbprf{\textbf{Oracles.} \texttt{arith}, \texttt{ineq}, \texttt{crs}, \texttt{sos}.  \textbf{Bill:} \texttt{proven modulo 0}.}
%
\vnbprf{\vnbproofpart{\texttt{cc-log-holomorphic-near-one}}}
%
\vnbprf{\textbf{Proof.} \texttt{theorem-\allowbreak{}library/analytic-\allowbreak{}log-\allowbreak{}laws.scm} --- 315 recorded steps.}
%
\vnbprf{\textbf{Lemmas.} \texttt{cc-is-metric-space}, \texttt{cc-is-normed-field}, \texttt{cc-zero-in}, \texttt{cc-one-in}, \texttt{rr-one-in}, \texttt{rr-zero-in}, \texttt{nn-zero-in}, \texttt{pi-pos}, \texttt{rr-subset-cc}, \texttt{rr-add-in-rr}, \texttt{rr-mul-in-rr}, \texttt{rr-lt-trichotomy}, \texttt{rr-mul-le-right}, \texttt{rr-lt-irrefl}, \texttt{ball-is-set}, \texttt{ball-is-open}, \texttt{cc-ms-open-iff}, \texttt{nn-in-rr}, \texttt{recip-star-in-rr}, \texttt{power-real-closed}, \texttt{cc-neg-closed}, \texttt{cc-mul-closed}, \texttt{nn-is-set}, \texttt{recip-star-zero}, \texttt{nn-succ-closed}, \texttt{nn-one-le-succ}, \texttt{rr-recip-closed}, \texttt{power-succ}, \texttt{recip-star-value}, \texttt{cps-log-radius}, \texttt{cps-translate-holomorphic}, \texttt{log0-is-log-on-ball}, \texttt{alg-shift-in-ball}, \texttt{cc-sub-in-cc}, \texttt{cps-holomorphic-ball}, \texttt{holomorphic-on-in-fun}, \texttt{fun-apply-type-c}, \texttt{holomorphic-transfer-ptwise-eq}.}
%
\vnbprf{\textbf{Oracles.} \texttt{crs}, \texttt{ineq}, \texttt{arith}, \texttt{sos}.  \textbf{Bill:} \texttt{proven modulo 0}.}
%
\vnbprf{\vnbproofpart{\texttt{cc-log-holomorphic-scale}}}
%
\vnbprf{\textbf{Proof.} \texttt{theorem-\allowbreak{}library/analytic-\allowbreak{}log-\allowbreak{}laws.scm} --- 767 recorded steps.}
%
\vnbprf{\textbf{Lemmas.} \texttt{cc-is-metric-space}, \texttt{cc-is-normed-field}, \texttt{cc-zero-in}, \texttt{cc-one-in}, \texttt{rr-one-in}, \texttt{rr-zero-in}, \texttt{pi-pos}, \texttt{rr-subset-cc}, \texttt{rr-add-in-rr}, \texttt{rr-mul-in-rr}, \texttt{rr-lt-trichotomy}, \texttt{rr-mul-le-right}, \texttt{rr-lt-irrefl}, \texttt{ball-is-set}, \texttt{ball-is-open}, \texttt{cc-ms-open-iff}, \texttt{rr-recip-closed}, \texttt{rr-recip-pos}, \texttt{rr-recip-inverse}, \texttt{rr-lt-scale-pos}, \texttt{ball-membership}, \texttt{cc-sub-in-cc}, \texttt{cc-mul-closed}, \texttt{cc-magnitude-closed}, \texttt{cc-ms-dist}, \texttt{cc-magnitude-neg}, \texttt{cc-magnitude-mul}, \texttt{rr-magnitude-is-abs}, \texttt{rr-abs-of-nonneg}, \texttt{ball-mem-from-le}, \texttt{real-part-in-rr}, \texttt{cc-abs-re-le-magnitude}, \texttt{rr-abs-le-parts}, \texttt{cc-re-sub}, \texttt{cc-add-closed}, \texttt{cc-re-im-of}, \texttt{holomorphic-const}, \texttt{holomorphic-identity}, \texttt{holomorphic-product}, \texttt{cc-log-holomorphic-near-one}, \texttt{holomorphic-compose}, \texttt{log-in-rr}, \texttt{cc-log-in-slit}, \texttt{cc-log-scale}, \texttt{holomorphic-sum}.}
%
\vnbprf{\textbf{Oracles.} \texttt{crs}, \texttt{ineq}, \texttt{arith}, \texttt{sos}.  \textbf{Bill:} \texttt{proven modulo 0}.}
%
\end{proof}
%
\noindent\textbf{Deviation from the notes.} PROVEN 2026-09-27 (batch 38-B), with the derivative log'(w) = recip(w) (cc-log-deriv). THE ROUTE DIFFERS FROM THE NOTES: no sector cover and no numerical bound on pi is needed -- at every point a of the slit plane, log x = log a + log(x / a) on the ball B(a, |a| r) with r = (pi - |Im log a|) / (2 (1 + pi)), where log(x / a) is the power series near 1 (cc-log-holomorphic-at, cc-log-im-near-one). The notes' rotation laws log(i w) = log w + i pi/2 (cc-log-rotate, from cos $>$ 0 on (-pi/2, pi/2), cc-cos-pos-half) and pi $>$ 3 (cc-pi-gt-three, from cos(3/2) $>$ 0 by six terms of the series) are proven as well. Earlier: PARTIAL 2026-09-27 (batch 38): holomorphic on B(1, pi/(1+pi)) (as the power series log0) and on every scaled ball B(s, s pi/(1+pi)), s $>$ 0, by log(s w) = LOG s + log w. The rotation step log(i w) = log w + i pi/2 and the cover of the slit plane need cos $>$ 0 on (-pi/2, pi/2) and pi $>$ 3, which the tree lacked: batch 38-B.
%
\renewcommand{\thethm}{2.22}
\begin{lem}[{The series of log(1 - z) and log(1 + z)}]\label{thm:cc-log-one-minus}
%
\leavevmode
\vnbpart{\texttt{cc-log-one-minus}}
\vnbstmt{\begin{array}{@{}l@{}} \forall r \in \mathbb{R}, z \in \mathbb{C}. \\ \quad ((0 < r) \wedge (((r \cdot (1 + \pi)) \leq \pi) \wedge (\operatorname{magnitude}(z) < r))) \Rightarrow  \\ \quad \quad \log\left((1 - z)\right) \\ \quad \quad =\; -\operatorname{cc\text{-}series\text{-}limit}((\lambda k \in \mathbb{N}.\; ({k}^{\ast\!-1} \cdot \operatorname{power}(z, k)))) \end{array}}
\vnbpart{\texttt{cc-log-one-plus}}
\vnbstmt{\begin{array}{@{}l@{}} \forall c \in (\mathbb{N} \to \mathbb{C}), r \in \mathbb{R}, z \in \mathbb{C}. \\ \quad (c(0) = 0) \wedge \\ \quad \forall j \in \mathbb{N}. \\ \quad \quad (c((j + 1)) = (\operatorname{power}(-1, j) \cdot \operatorname{recip}((j + 1)))) \wedge \\ \quad (0 < r) \wedge \\ \quad ((r \cdot (1 + \pi)) \leq \pi) \wedge \\ \quad (\operatorname{magnitude}(z) < r) \Rightarrow \\ \quad \quad \log\left((1 + z)\right) \\ \quad \quad =\; \operatorname{cc\text{-}series\text{-}limit}((\lambda k \in \mathbb{N}.\; (c(k) \cdot \operatorname{power}(z, k)))) \end{array}}
\vnbpart{\texttt{log0-is-log-on-ball}}
\vnbstmt{\begin{array}{@{}l@{}} \forall c \in (\mathbb{N} \to \mathbb{C}), r \in \mathbb{R}, z \in \operatorname{ball}(\operatorname{cc\text{-}ms}, 1, r). \\ \quad (c(0) = 0) \wedge \\ \quad \forall j \in \mathbb{N}. \\ \quad \quad (c((j + 1)) = (\operatorname{power}(-1, j) \cdot \operatorname{recip}((j + 1)))) \wedge \\ \quad (0 < r) \wedge \\ \quad ((r \cdot (1 + \pi)) \leq \pi) \Rightarrow \\ \quad \quad \log\left(z\right) \\ \quad \quad =\; \operatorname{cc\text{-}series\text{-}limit}((\lambda k \in \mathbb{N}.\; (c(k) \cdot \operatorname{power}((z - 1), k)))) \end{array}}
%
\end{lem}
%
\begin{proof}
%
The proofs: the file that holds each, the lemmas it cites, its oracles and its bill.
%
\vnbprf{\vnbproofpart{\texttt{cc-log-one-minus}}}
%
\vnbprf{\textbf{Proof.} \texttt{theorem-\allowbreak{}library/cc-\allowbreak{}log-\allowbreak{}holomorphic.scm} --- 452 recorded steps.}
%
\vnbprf{\textbf{Lemmas.} \texttt{cc-is-metric-space}, \texttt{cc-is-normed-field}, \texttt{cc-zero-in}, \texttt{cc-one-in}, \texttt{rr-one-in}, \texttt{rr-zero-in}, \texttt{nn-zero-in}, \texttt{pi-pos}, \texttt{rr-subset-cc}, \texttt{rr-lt-irrefl}, \texttt{rr-add-in-rr}, \texttt{rr-mul-in-rr}, \texttt{rr-lt-trichotomy}, \texttt{rr-mul-le-right}, \texttt{nn-in-rr}, \texttt{recip-star-in-rr}, \texttt{power-real-closed}, \texttt{cc-neg-closed}, \texttt{cc-mul-closed}, \texttt{nn-is-set}, \texttt{recip-star-zero}, \texttt{nn-succ-closed}, \texttt{nn-one-le-succ}, \texttt{rr-recip-closed}, \texttt{power-succ}, \texttt{recip-star-value}, \texttt{cc-magnitude-closed}, \texttt{cc-magnitude-neg}, \texttt{cc-magnitude-nonneg}, \texttt{cc-log-one-plus}, \texttt{cps-log-radius}, \texttt{cc-ms-dist}, \texttt{cc-sub-in-cc}, \texttt{closed-ball-membership}, \texttt{cps-converges-closed-ball}, \texttt{fun-apply-type-c}, \texttt{power-typing-nonneg}, \texttt{cc-series-limit-in-cc}, \texttt{cc-power-mul-base}, \texttt{rr-power-one}, \texttt{cc-series-limit-linear}, \texttt{cc-add-closed}.}
%
\vnbprf{\textbf{Oracles.} \texttt{ineq}, \texttt{crs}, \texttt{arith}, \texttt{sos}.  \textbf{Bill:} \texttt{proven modulo 0}.}
%
\vnbprf{\vnbproofpart{\texttt{cc-log-one-plus}}}
%
\vnbprf{\textbf{Proof.} \texttt{theorem-\allowbreak{}library/analytic-\allowbreak{}log-\allowbreak{}laws.scm} --- 90 recorded steps.}
%
\vnbprf{\textbf{Lemmas.} \texttt{cc-is-metric-space}, \texttt{cc-is-normed-field}, \texttt{cc-zero-in}, \texttt{cc-one-in}, \texttt{rr-one-in}, \texttt{rr-zero-in}, \texttt{cc-add-closed}, \texttt{cc-magnitude-closed}, \texttt{cc-ms-dist}, \texttt{cc-sub-in-cc}, \texttt{cc-magnitude-neg}, \texttt{ball-mem-from-le}, \texttt{log0-is-log-on-ball}.}
%
\vnbprf{\textbf{Oracles.} \texttt{crs}, \texttt{ineq}, \texttt{arith}, \texttt{sos}.  \textbf{Bill:} \texttt{proven modulo 0}.}
%
\vnbprf{\vnbproofpart{\texttt{log0-is-log-on-ball}}}
%
\vnbprf{\textbf{Proof.} \texttt{theorem-\allowbreak{}library/analytic-\allowbreak{}log-\allowbreak{}laws.scm} --- 324 recorded steps.}
%
\vnbprf{\textbf{Lemmas.} \texttt{cc-is-metric-space}, \texttt{cc-is-normed-field}, \texttt{cc-zero-in}, \texttt{cc-one-in}, \texttt{rr-one-in}, \texttt{rr-zero-in}, \texttt{pi-pos}, \texttt{rr-subset-cc}, \texttt{rr-add-in-rr}, \texttt{rr-mul-in-rr}, \texttt{rr-lt-trichotomy}, \texttt{rr-mul-le-right}, \texttt{rr-lt-irrefl}, \texttt{cps-log-radius}, \texttt{ball-membership}, \texttt{cc-sub-in-cc}, \texttt{cc-magnitude-closed}, \texttt{cc-ms-dist}, \texttt{cc-magnitude-neg}, \texttt{alg-shift-in-ball}, \texttt{cps-holomorphic-ball}, \texttt{holomorphic-on-in-fun}, \texttt{fun-apply-type-c}, \texttt{imag-part-in-rr}, \texttt{log0-bound}, \texttt{cc-abs-im-le-magnitude}, \texttt{rr-sub-in-rr}, \texttt{cc-magnitude-nonneg}, \texttt{rr-recip-closed}, \texttt{rr-recip-pos}, \texttt{rr-recip-inverse}, \texttt{rr-abs-closed}, \texttt{rr-abs-le-parts}, \texttt{rr-lt-scale-pos}, \texttt{log0-exp}, \texttt{cc-log-char}.}
%
\vnbprf{\textbf{Oracles.} \texttt{crs}, \texttt{ineq}, \texttt{arith}, \texttt{sos}.  \textbf{Bill:} \texttt{proven modulo 0}.}
%
\end{proof}
%
\noindent\textbf{Deviation from the notes.} PARTIAL 2026-09-27 (batches 38, 38-B): both identities are proven on |z| $<$ r with r (1 + pi) $<$= pi (about 0.758), the ball on which the series log0 is known to be the principal logarithm (log0-is-log-on-ball). The full disc |z| $<$ 1 needs log = log0 on all of B(1, 1), i.e. the identity theorem 3.35 or a lemma that a continuous integer-valued function on a connected set is constant, neither of which the tree has.
%
\renewcommand{\thethm}{3.1}
\begin{prop}[{Path additivity of the line integral}]\label{thm:line-int-concat}
%
\leavevmode
\vnbpart{\texttt{line-int-concat}}
\vnbstmt{\begin{array}{@{}l@{}} \forall u, f, m, g, o, h, a, b, c, d. \\ \quad (u \subseteq \mathbb{C}) \wedge \\ \quad (f \in (u \to \mathbb{C})) \wedge \\ \quad \forall y \in u. \\ \quad \quad \operatorname{is\text{-}continuous\text{-}at}(\operatorname{subspace\text{-}ms}(\operatorname{nf\text{-}metric\text{-}space}(\operatorname{cc\text{-}normed\text{-}field}), u), \\ \quad \quad \quad \operatorname{nf\text{-}metric\text{-}space}(\operatorname{cc\text{-}normed\text{-}field}), f, y) \wedge \\ \quad \operatorname{is\text{-}road}(m, g, a, b) \wedge \\ \quad \operatorname{is\text{-}road}(o, h, c, d) \wedge \\ \quad (m(b) = o(c)) \wedge \\ \quad (\operatorname{trace}(m, a, b) \subseteq u) \wedge \\ \quad (\operatorname{trace}(o, c, d) \subseteq u) \Rightarrow \\ \quad \quad \operatorname{line\text{-}int}(f, \\ \quad \quad \quad \operatorname{juxta}(m, (\lambda s \in \operatorname{ccint}(b, (b + (d - c))).\; o((c + (s - b)))), a, \\ \quad \quad \quad \quad b, (b + (d - c))), \\ \quad \quad \quad \operatorname{juxta}(g, (\lambda s \in \operatorname{ccint}(b, (b + (d - c))).\; h((c + (s - b)))), a, \\ \quad \quad \quad \quad b, (b + (d - c))), \\ \quad \quad \quad a, (b + (d - c))) \\ \quad \quad =\; (\operatorname{line\text{-}int}(f, m, g, a, b) + \operatorname{line\text{-}int}(f, o, h, c, d)) \end{array}}
\vnbpart{\texttt{line-int-juxta}}
\vnbstmt{\begin{array}{@{}l@{}} \forall u, f, m_{1}, d, m_{2}, g, a, b, c. \\ \quad (u \subseteq \mathbb{C}) \wedge \\ \quad (f \in (u \to \mathbb{C})) \wedge \\ \quad \forall y \in u. \\ \quad \quad \operatorname{is\text{-}continuous\text{-}at}(\operatorname{subspace\text{-}ms}(\operatorname{nf\text{-}metric\text{-}space}(\operatorname{cc\text{-}normed\text{-}field}), u), \\ \quad \quad \quad \operatorname{nf\text{-}metric\text{-}space}(\operatorname{cc\text{-}normed\text{-}field}), f, y) \wedge \\ \quad \operatorname{is\text{-}road}(m_{1}, d, a, b) \wedge \\ \quad \operatorname{is\text{-}road}(m_{2}, g, b, c) \wedge \\ \quad (\operatorname{m1}(b) = \operatorname{m2}(b)) \wedge \\ \quad (\operatorname{trace}(m_{1}, a, b) \subseteq u) \wedge \\ \quad (\operatorname{trace}(m_{2}, b, c) \subseteq u) \Rightarrow \\ \quad \quad \operatorname{line\text{-}int}(f, \operatorname{juxta}(m_{1}, m_{2}, a, b, c), \operatorname{juxta}(d, g, a, b, c), a, c) \\ \quad \quad =\; (\operatorname{line\text{-}int}(f, m_{1}, d, a, b) + \operatorname{line\text{-}int}(f, m_{2}, g, b, c)) \end{array}}
\vnbpart{\texttt{juxta-is-road}}
\vnbstmt{\begin{array}{@{}l@{}} \forall m_{1}, d, m_{2}, f, a, b, c. \\ \quad \operatorname{is\text{-}road}(m_{1}, d, a, b) \wedge \\ \quad \operatorname{is\text{-}road}(m_{2}, f, b, c) \wedge \\ \quad (\operatorname{m1}(b) = \operatorname{m2}(b)) \Rightarrow \\ \quad \quad \operatorname{is\text{-}road}(\operatorname{juxta}(m_{1}, m_{2}, a, b, c), \operatorname{juxta}(d, f, a, b, c), a, c) \end{array}}
\vnbpart{\texttt{regulated-on-glue}}
\vnbstmt{\begin{array}{@{}l@{}} \forall a, c, b \in \mathbb{R}, d, g, f. \\ \quad (a < c) \wedge \\ \quad (c < b) \wedge \\ \quad \operatorname{is\text{-}regulated\text{-}on}(d, a, c) \wedge \\ \quad \operatorname{is\text{-}regulated\text{-}on}(g, c, b) \wedge \\ \quad (f \in (\operatorname{ccint}(a, b) \to \mathbb{R})) \wedge \\ \quad \forall t \in \operatorname{ccint}(a, c). \\ \quad \quad ((t < c) \Rightarrow (f(t) \simeq d(t))) \wedge \\ \quad \forall t \in \operatorname{ccint}(c, b). \\ \quad \quad ((c < t) \Rightarrow (f(t) \simeq g(t))) \Rightarrow \\ \quad \quad \operatorname{is\text{-}regulated\text{-}on}(f, a, b) \end{array}}
\vnbpart{\texttt{tri-int-is-line-int}}
\vnbstmt{\begin{array}{@{}l@{}} \forall u, f.; \forall a, b, c \in \mathbb{C}. \\ \quad (u \subseteq \mathbb{C}) \wedge \\ \quad (f \in (u \to \mathbb{C})) \wedge \\ \quad \forall y \in u. \\ \quad \quad \operatorname{is\text{-}continuous\text{-}at}(\operatorname{subspace\text{-}ms}(\operatorname{nf\text{-}metric\text{-}space}(\operatorname{cc\text{-}normed\text{-}field}), u), \\ \quad \quad \quad \operatorname{nf\text{-}metric\text{-}space}(\operatorname{cc\text{-}normed\text{-}field}), f, y) \wedge \\ \quad (\operatorname{trace}(\operatorname{seg\text{-}path}(a, b), 0, 1) \subseteq u) \wedge \\ \quad (\operatorname{trace}(\operatorname{seg\text{-}path}(b, c), 0, 1) \subseteq u) \wedge \\ \quad (\operatorname{trace}(\operatorname{seg\text{-}path}(c, a), 0, 1) \subseteq u) \Rightarrow \\ \quad \quad \operatorname{tri\text{-}int}(f, a, b, c) \\ \quad \quad =\; \operatorname{line\text{-}int}(f, \operatorname{tri\text{-}road}(a, b, c), \operatorname{tri\text{-}road\text{-}deriv}(a, b, c), 0, 3) \end{array}}
\vnbpart{\texttt{line-int-length-bound}}
\vnbstmt{\begin{array}{@{}l@{}} \forall u, f, m, c, a, b.; \forall d \in \mathbb{R}. \\ \quad (u \subseteq \mathbb{C}) \wedge \\ \quad (f \in (u \to \mathbb{C})) \wedge \\ \quad \forall y \in u. \\ \quad \quad \operatorname{is\text{-}continuous\text{-}at}(\operatorname{subspace\text{-}ms}(\operatorname{nf\text{-}metric\text{-}space}(\operatorname{cc\text{-}normed\text{-}field}), u), \\ \quad \quad \quad \operatorname{nf\text{-}metric\text{-}space}(\operatorname{cc\text{-}normed\text{-}field}), f, y) \wedge \\ \quad \operatorname{is\text{-}road}(m, c, a, b) \wedge \\ \quad (\operatorname{trace}(m, a, b) \subseteq u) \wedge \\ \quad \forall \mathit{bny} \in \operatorname{trace}(m, a, b). \\ \quad \quad (\operatorname{magnitude}(f(\mathit{bny})) \leq d) \Rightarrow \\ \quad \quad \operatorname{magnitude}(\operatorname{line\text{-}int}(f, m, c, a, b)) \\ \quad \quad \leq\; (\operatorname{road\text{-}length}(c, a, b) \cdot d) \end{array}}
\vnbpart{\texttt{line-int-adjacent}}
\vnbstmt{\begin{array}{@{}l@{}} \forall u, f, m, c, a, b.; \forall v \in \operatorname{ooint}(a, b). \\ \quad (u \subseteq \mathbb{C}) \wedge \\ \quad (f \in (u \to \mathbb{C})) \wedge \\ \quad \forall y \in u. \\ \quad \quad \operatorname{is\text{-}continuous\text{-}at}(\operatorname{subspace\text{-}ms}(\operatorname{nf\text{-}metric\text{-}space}(\operatorname{cc\text{-}normed\text{-}field}), u), \\ \quad \quad \quad \operatorname{nf\text{-}metric\text{-}space}(\operatorname{cc\text{-}normed\text{-}field}), f, y) \wedge \\ \quad \operatorname{is\text{-}road}(m, c, a, b) \wedge \\ \quad (\operatorname{trace}(m, a, b) \subseteq u) \Rightarrow \\ \quad \quad \operatorname{line\text{-}int}(f, m, c, a, b) \\ \quad \quad =\; (\operatorname{line\text{-}int}(f, \operatorname{restrict}(m, \operatorname{ccint}(a, v)), \operatorname{restrict}(c, \operatorname{ccint}(a, v)), a, v) + \operatorname{line\text{-}int}(f, \operatorname{restrict}(m, \operatorname{ccint}(v, b)), \operatorname{restrict}(c, \operatorname{ccint}(v, b)), v, b)) \end{array}}
\vnbpart{\texttt{road-restrict}}
\vnbstmt{\begin{array}{@{}l@{}} \forall m, c, a, b.; \forall v, \mathit{lidv} \in \mathbb{R}. \\ \quad \operatorname{is\text{-}road}(m, c, a, b) \wedge \\ \quad (a \leq v) \wedge \\ \quad (v < \mathit{lidv}) \wedge \\ \quad (\mathit{lidv} \leq b) \Rightarrow \\ \quad \quad \operatorname{is\text{-}road}(\operatorname{restrict}(m, \operatorname{ccint}(v, \mathit{lidv})), \operatorname{restrict}(c, \operatorname{ccint}(v, \mathit{lidv})), v, \mathit{lidv}) \end{array}}
%
\end{prop}
%
\begin{proof}
%
The proofs: the file that holds each, the lemmas it cites, its oracles and its bill.
%
\vnbprf{\vnbproofpart{\texttt{line-int-concat}}}
%
\vnbprf{\textbf{Proof.} \texttt{theorem-\allowbreak{}library/road-\allowbreak{}juxtaposition-\allowbreak{}laws.scm} --- 240 recorded steps.}
%
\vnbprf{\textbf{Lemmas.} \texttt{is-road-is-path}, \texttt{is-path-in-fun}, \texttt{is-road-dgam-in-fun}, \texttt{is-path-endpoints}, \texttt{rr-sub-in-rr}, \texttt{rr-add-in-rr}, \texttt{rr-one-in}, \texttt{rr-zero-lt-one}, \texttt{rr-is-set}, \texttt{ccint-subset-rr}, \texttt{subclass-of-set-is-set}, \texttt{ccint-parts}, \texttt{ccint-membership}, \texttt{rr-mul-in-rr}, \texttt{fun-apply-type-c}, \texttt{road-affine-reparam}, \texttt{cc-one-mul}, \texttt{trace-affine-reparam}, \texttt{subset-trans}, \texttt{line-int-affine-reparam}, \texttt{rr-leq-reflexive}, \texttt{line-int-juxta}, \texttt{line-int-in-cc}.}
%
\vnbprf{\textbf{Oracles.} \texttt{ineq}, \texttt{crs}, \texttt{arith}.  \textbf{Bill:} \texttt{proven modulo 0}.}
%
\vnbprf{\vnbproofpart{\texttt{line-int-juxta}}}
%
\vnbprf{\textbf{Proof.} \texttt{theorem-\allowbreak{}library/road-\allowbreak{}juxtaposition-\allowbreak{}laws.scm} --- 232 recorded steps.}
%
\vnbprf{\textbf{Lemmas.} \texttt{is-road-is-path}, \texttt{is-path-in-fun}, \texttt{is-road-dgam-in-fun}, \texttt{is-path-endpoints}, \texttt{rr-leq-reflexive}, \texttt{juxta-in-fun}, \texttt{juxta-is-road}, \texttt{juxta-trace-subset}, \texttt{subset-def}, \texttt{union-membership}, \texttt{subset-mem-fwd}, \texttt{subset-trans}, \texttt{ooint-membership}, \texttt{line-int-adjacent}, \texttt{road-restrict}, \texttt{trace-restrict-subset}, \texttt{line-int-interior-congruence}, \texttt{ccint-parts}, \texttt{ccint-membership}, \texttt{restrict-apply}, \texttt{juxta-apply-left}, \texttt{ooint-elt-in-rr}, \texttt{juxta-apply-right-closed}, \texttt{juxta-apply-right}, \texttt{line-int-in-cc}.}
%
\vnbprf{\textbf{Oracles.} \texttt{ineq}, \texttt{crs}, \texttt{arith}.  \textbf{Bill:} \texttt{proven modulo 0}.}
%
\vnbprf{\vnbproofpart{\texttt{juxta-is-road}}}
%
\vnbprf{\textbf{Proof.} \texttt{theorem-\allowbreak{}library/road-\allowbreak{}juxtaposition-\allowbreak{}laws.scm} --- 722 recorded steps.}
%
\vnbprf{\textbf{Lemmas.} \texttt{is-road-is-path}, \texttt{is-path-in-fun}, \texttt{is-road-dgam-in-fun}, \texttt{is-path-endpoints}, \texttt{rr-leq-reflexive}, \texttt{ccint-membership}, \texttt{juxta-in-fun}, \texttt{ccint-subset-rr}, \texttt{rr-is-set}, \texttt{subclass-of-set-is-set}, \texttt{fun-apply-type-c}, \texttt{real-part-in-rr}, \texttt{is-road-re-antiderivative}, \texttt{is-road-re-regulated}, \texttt{primitive-glue}, \texttt{ccint-parts}, \texttt{juxta-apply-left}, \texttt{juxta-apply-right-closed}, \texttt{ooint-elt-in-rr}, \texttt{ooint-membership}, \texttt{juxta-apply-right}, \texttt{regulated-on-glue}, \texttt{primitive-continuous}, \texttt{imag-part-in-rr}, \texttt{is-road-im-antiderivative}, \texttt{is-road-im-regulated}, \texttt{is-path-of-coords}.}
%
\vnbprf{\textbf{Oracles.} \texttt{ineq}, \texttt{crs}, \texttt{arith}.  \textbf{Bill:} \texttt{proven modulo 0}.}
%
\vnbprf{\vnbproofpart{\texttt{regulated-on-glue}}}
%
\vnbprf{\textbf{Proof.} \texttt{theorem-\allowbreak{}library/road-\allowbreak{}juxtaposition-\allowbreak{}laws.scm} --- 391 recorded steps.}
%
\vnbprf{\textbf{Lemmas.} \texttt{ccint-subset-rr}, \texttt{rr-one-in}, \texttt{rr-zero-lt-one}, \texttt{rr-pos-rr-of-lt}, \texttt{ccint-parts}, \texttt{ccint-membership}, \texttt{right-limit-within-local}, \texttt{rr-not-le-lt}, \texttt{rr-sub-in-rr}, \texttt{left-limit-within-local}.}
%
\vnbprf{\textbf{Oracles.} \texttt{ineq}.  \textbf{Bill:} \texttt{proven modulo 0}.}
%
\vnbprf{\vnbproofpart{\texttt{tri-int-is-line-int}}}
%
\vnbprf{\textbf{Proof.} \texttt{theorem-\allowbreak{}library/road-\allowbreak{}juxtaposition-\allowbreak{}laws.scm} --- 277 recorded steps.}
%
\vnbprf{\textbf{Lemmas.} \texttt{rr-zero-lt-one}, \texttt{ccint-membership}, \texttt{seg-path-is-road}, \texttt{seg-shift-is-road}, \texttt{is-road-is-path}, \texttt{is-path-in-fun}, \texttt{is-road-dgam-in-fun}, \texttt{is-path-endpoints}, \texttt{cc-sub-in-cc}, \texttt{rr-subset-cc}, \texttt{seg-path-apply}, \texttt{juxta-is-road}, \texttt{juxta-apply-right}, \texttt{seg-shift-trace}, \texttt{juxta-trace-subset}, \texttt{subset-def}, \texttt{union-membership}, \texttt{subset-mem-fwd}, \texttt{subset-trans}, \texttt{line-int-juxta}, \texttt{seg-shift-line-int}, \texttt{line-int-in-cc}, \texttt{cc-add-closed}.}
%
\vnbprf{\textbf{Oracles.} \texttt{arith}, \texttt{ineq}, \texttt{crs}.  \textbf{Bill:} \texttt{proven modulo 0}.}
%
\vnbprf{\vnbproofpart{\texttt{line-int-length-bound}}}
%
\vnbprf{\textbf{Proof.} \texttt{theorem-\allowbreak{}library/road-\allowbreak{}juxtaposition-\allowbreak{}laws.scm} --- 211 recorded steps.}
%
\vnbprf{\textbf{Lemmas.} \texttt{is-road-is-path}, \texttt{is-path-in-fun}, \texttt{is-road-dgam-in-fun}, \texttt{is-path-endpoints}, \texttt{rr-is-set}, \texttt{ccint-subset-rr}, \texttt{subclass-of-set-is-set}, \texttt{trace-value-in}, \texttt{subset-mem-fwd}, \texttt{fun-apply-type-c}, \texttt{cc-magnitude-closed}, \texttt{rr-mul-in-rr}, \texttt{line-int-abs-bound}, \texttt{continuous-on-compose-regulated}, \texttt{regulated-on-magnitude}, \texttt{real-part-in-rr}, \texttt{imag-part-in-rr}, \texttt{cc-re-im-decompose}, \texttt{rr-subset-cc}, \texttt{cc-i-in}, \texttt{cc-mul-comm}, \texttt{road-length-regulated}, \texttt{regulated-on-mul}, \texttt{regulated-on-has-primitive}, \texttt{road-length-primitive}, \texttt{primitive-in-fun}, \texttt{pw-antiderivative-real-mul}, \texttt{pw-int-monotone}, \texttt{cc-magnitude-nonneg}, \texttt{rr-mul-le-right}, \texttt{pw-int-in-rr}, \texttt{line-int-in-cc}, \texttt{rr-mul-comm}.}
%
\vnbprf{\textbf{Oracles.} \texttt{ineq}, \texttt{crs}, \texttt{arith}.  \textbf{Bill:} \texttt{proven modulo 0}.}
%
\vnbprf{\vnbproofpart{\texttt{line-int-adjacent}}}
%
\vnbprf{\textbf{Proof.} \texttt{theorem-\allowbreak{}library/line-\allowbreak{}int-\allowbreak{}laws.scm} --- 436 recorded steps.}
%
\vnbprf{\textbf{Lemmas.} \texttt{is-road-is-path}, \texttt{is-road-dgam-in-fun}, \texttt{is-path-in-fun}, \texttt{is-path-endpoints}, \texttt{rr-is-set}, \texttt{ccint-subset-rr}, \texttt{subclass-of-set-is-set}, \texttt{trace-value-in}, \texttt{subset-mem-fwd}, \texttt{line-int-exists}, \texttt{line-int-integrand-in-fun}, \texttt{cc-int-adjacent}, \texttt{ooint-elt-in-rr}, \texttt{ooint-membership}, \texttt{rr-leq-reflexive}, \texttt{rr-lt-implies-le}, \texttt{ccint-subset-ccint}, \texttt{pw-antiderivative-restrict}, \texttt{fun-apply-type-c}, \texttt{cc-mul-closed}, \texttt{real-part-in-rr}, \texttt{imag-part-in-rr}, \texttt{restrict-apply}, \texttt{pw-antiderivative-integrand-congruence}, \texttt{cc-int-congruence}.}
%
\vnbprf{\textbf{Oracles.} \texttt{arith}, \texttt{crs}, \texttt{ineq}.  \textbf{Bill:} \texttt{proven modulo 0}.}
%
\vnbprf{\vnbproofpart{\texttt{road-restrict}}}
%
\vnbprf{\textbf{Proof.} \texttt{theorem-\allowbreak{}library/line-\allowbreak{}int-\allowbreak{}laws.scm} --- 353 recorded steps.}
%
\vnbprf{\textbf{Lemmas.} \texttt{is-road-is-path}, \texttt{is-road-dgam-in-fun}, \texttt{is-path-in-fun}, \texttt{is-path-endpoints}, \texttt{rr-is-set}, \texttt{ccint-subset-rr}, \texttt{subclass-of-set-is-set}, \texttt{ccint-subset-ccint}, \texttt{restrict-in-fun}, \texttt{rr-one-in}, \texttt{is-road-re-antiderivative}, \texttt{is-road-im-antiderivative}, \texttt{is-road-re-regulated}, \texttt{is-road-im-regulated}, \texttt{subset-mem-fwd}, \texttt{fun-apply-type-c}, \texttt{real-part-in-rr}, \texttt{imag-part-in-rr}, \texttt{restrict-apply}, \texttt{pw-antiderivative-restrict}, \texttt{pw-antiderivative-real-mul}, \texttt{regulated-on-restrict}, \texttt{primitive-continuous}, \texttt{is-path-of-coords}.}
%
\vnbprf{\textbf{Oracles.} \texttt{crs}, \texttt{ineq}, \texttt{arith}.  \textbf{Bill:} \texttt{proven modulo 0}.}
%
\end{proof}
%
\noindent\textbf{Deviation from the notes.} PROVEN 2026-09-26 (batch 37): the juxtaposition JUXTA(gamma, rho, a, b, c) of roads on adjacent intervals (Dieudonne 9.6; structure-library/road-juxtaposition.scm) is a road (juxta-is-road, with the new regulated-on-glue) and the integral adds (line-int-juxta); the notes' (49) for gamma on [a, b] and rho on [c, d] with gamma(b) = rho(c) is line-int-concat, rho shifted to [b, b + (d - c)]. Also the triangular road TRI-ROAD on [0, 3] with tri-int-is-line-int, and ROAD-LENGTH (the integral of |gamma'|) with the length estimate (53), line-int-length-bound. Earlier: UPDATED 2026-09-23 (batch 24): the SPLIT form is proven -- line-int-adjacent: for a road gamma on [a, b], c in (a, b) and f continuous on a set containing the trace, the integral along gamma is the sum of the integrals along the restrictions to [a, c] and [c, b] (road-restrict: the restriction of a road is a road). The CONCATENATION form of the notes (gamma || rho, two roads glued at a point) is not stated: it needs the glued road, which is primitive-glue (batch 22) applied to each coordinate. Existence of the integral for a road and a continuous f is line-int-exists (batch 23). EARLIER NOTE: Not in the library. The DEFINITIONS exist since 2026-09-21, for functions ON [a, b] (structure-library/path-integral.scm): is-path, is-road(gamma, dgamma, a, b) with the derivative dgamma an explicit argument, trace, cc-int (the notes' equation (44)) and line-int(f, gamma, dgamma, a, b). Proven since: the integral pw-int is well defined (two piecewise primitives have the same endpoint difference), linear, additive over adjacent intervals and monotone; cc-int is linear over CC and additive; the notes' estimate (45) (cc-int-abs-bound); the trace of a path is compact (path-trace-compact); line-int of 1 along a road is gamma(b) - gamma(a); the opposite road and Dieudonne (9.6.1). Existence of a primitive of a piecewise continuous function (Dieudonne 8.7.2) is NOT proven, so the primitives of an integrand's coordinates are hypotheses wherever a value of an integral is asserted. Juxtaposition of two roads (Dieudonne 9.6.2) is not defined yet.
%
\renewcommand{\thethm}{3.2}
\begin{prop}[{Fundamental theorem for line integrals}]\label{thm:line-int-of-derivative}
%
\leavevmode
\vnbpart{\texttt{line-int-of-derivative}}
\vnbstmt{\begin{array}{@{}l@{}} \forall u, f, c, m, d, a, b. \\ \quad \operatorname{holomorphic\text{-}on}(u, f) \wedge \\ \quad (c \in (u \to \mathbb{C})) \wedge \\ \quad \forall z \in u. \\ \quad \quad \operatorname{is\text{-}diff\text{-}on}(\operatorname{cc\text{-}normed\text{-}field}, u, f, z, c(z)) \wedge \\ \quad \operatorname{is\text{-}road}(m, d, a, b) \wedge \\ \quad (\operatorname{trace}(m, a, b) \subseteq u) \Rightarrow \\ \quad \quad (\operatorname{line\text{-}int}(c, m, d, a, b) = (f(m(b)) - f(m(a)))) \end{array}}
\vnbpart{\texttt{holomorphic-chain}}
\vnbstmt{\begin{array}{@{}l@{}} \forall m, g, x, u, f, l, d. \\ \quad (m \subseteq \mathbb{R}) \wedge \\ \quad (g \in (m \to \mathbb{C})) \wedge \\ \quad \operatorname{is\text{-}continuous\text{-}at}(\operatorname{subspace\text{-}ms}(\operatorname{rr\text{-}ms}, m), \operatorname{cc\text{-}ms}, g, x) \wedge \\ \quad \forall y \in m.\; (g(y) \in u) \wedge \\ \quad \operatorname{is\text{-}diff\text{-}on}(\operatorname{cc\text{-}normed\text{-}field}, u, f, g(x), l) \wedge \\ \quad (d \in \mathbb{C}) \wedge \\ \quad \operatorname{has\text{-}deriv\text{-}at}((\lambda t \in m.\; \operatorname{real\text{-}part}(g(t))), x, \operatorname{real\text{-}part}(d)) \wedge \\ \quad \operatorname{has\text{-}deriv\text{-}at}((\lambda t \in m.\; \operatorname{imag\text{-}part}(g(t))), x, \operatorname{imag\text{-}part}(d)) \Rightarrow \\ \quad \quad \operatorname{has\text{-}deriv\text{-}at}((\lambda t \in m.\; \operatorname{real\text{-}part}(f(g(t)))), x, \operatorname{real\text{-}part}((l \cdot d))) \wedge \\ \quad \quad \operatorname{has\text{-}deriv\text{-}at}((\lambda t \in m.\; \operatorname{imag\text{-}part}(f(g(t)))), x, \operatorname{imag\text{-}part}((l \cdot d))) \end{array}}
%
\end{prop}
%
\begin{proof}
%
The proofs: the file that holds each, the lemmas it cites, its oracles and its bill.
%
\vnbprf{\vnbproofpart{\texttt{line-int-of-derivative}}}
%
\vnbprf{\textbf{Proof.} \texttt{theorem-\allowbreak{}library/line-\allowbreak{}int-\allowbreak{}fundamental.scm} --- 686 recorded steps.}
%
\vnbprf{\textbf{Lemmas.} \texttt{cc-is-normed-field}, \texttt{cc-is-metric-space}, \texttt{nf-cc-is-metric-space}, \texttt{nf-cc-agree-pts}, \texttt{nf-cc-agree-dist}, \texttt{ms-agree-pts-sym}, \texttt{ms-agree-dist-sym}, \texttt{cc-nf-carr}, \texttt{rr-is-metric-space}, \texttt{rr-is-set}, \texttt{holomorphic-on-in-fun}, \texttt{holomorphic-on-open}, \texttt{is-road-is-path}, \texttt{is-path-in-fun}, \texttt{is-road-dgam-in-fun}, \texttt{is-path-endpoints}, \texttt{is-path-continuous}, \texttt{is-road-re-antiderivative}, \texttt{is-road-im-antiderivative}, \texttt{ccint-subset-rr}, \texttt{subclass-of-set-is-set}, \texttt{ooint-subset-ccint}, \texttt{subspace-is-metric-space}, \texttt{subspace-pts}, \texttt{rr-leq-reflexive}, \texttt{rr-lt-implies-le}, \texttt{ccint-membership}, \texttt{trace-value-in}, \texttt{subset-mem-fwd}, \texttt{fun-apply-type-c}, \texttt{real-part-in-rr}, \texttt{imag-part-in-rr}, \texttt{cc-mul-closed}, \texttt{diff-on-implies-continuous}, \texttt{cc-compose-continuous-at}, \texttt{ms-pts-is-set}, \texttt{cc-continuous-at-iff-coords}, \texttt{primitive-exceptional-set}, \texttt{countable-union-2}, \texttt{subset-def}, \texttt{union-membership}, \texttt{holomorphic-chain}, \texttt{pw-int-value}, \texttt{cc-i-in}, \texttt{cc-sub-in-cc}, \texttt{cc-re-sub}, \texttt{cc-im-sub}, \texttt{cc-re-im-decompose}, \texttt{rr-subset-cc}, \texttt{cc-mul-comm}, \texttt{eq-sym}.}
%
\vnbprf{\textbf{Oracles.} \texttt{arith}, \texttt{crs}, \texttt{ineq}.  \textbf{Bill:} \texttt{proven modulo 0}.}
%
\vnbprf{\vnbproofpart{\texttt{holomorphic-chain}}}
%
\vnbprf{\textbf{Proof.} \texttt{theorem-\allowbreak{}library/line-\allowbreak{}int-\allowbreak{}fundamental.scm} --- 57 recorded steps.}
%
\vnbprf{\textbf{Lemmas.} \texttt{has-deriv-at-pt-in-rr}, \texttt{rr-is-set}, \texttt{subclass-of-set-is-set}, \texttt{has-deriv-at-iff-within-ooint}, \texttt{has-deriv-at-to-within-ooint}, \texttt{rr-pos-rr-in-rr}, \texttt{ooint-subset-rr}, \texttt{ooint-center-nested}, \texttt{ooint-center}, \texttt{deriv-within-shrink}, \texttt{fun-apply-type-c}, \texttt{real-part-in-rr}, \texttt{deriv-within-restrict-in-fun}, \texttt{restrict-fun-domain-subset}, \texttt{holomorphic-chain-within}, \texttt{has-deriv-at-of-within-ooint}.}
%
\vnbprf{\textbf{Oracles.} \texttt{ineq}, \texttt{arith}, \texttt{crs}.  \textbf{Bill:} \texttt{proven modulo 0}.}
%
\end{proof}
%
\noindent\textbf{Deviation from the notes.} Proven 2026-09-23 (line-int-of-derivative): f holomorphic on the open set U, df in fun(U, CC) its derivative (is-diff-on at every point of U), gamma a road ON [a, b] with derivative dgamma and trace inside U; then line-int(df, gamma, dgamma, a, b) = f(gamma(b)) - f(gamma(a)). The derivative is a PARAMETER df rather than the notes' f', because line-int takes a function as integrand. No primitive of the integrand is assumed: f o gamma is the primitive (holomorphic-chain, the chain rule of a holomorphic f with a road, coordinatewise). Roads are the tree's is-road with a finite exceptional set (restatement over Dieudonne's regulated primitives pending the user's decision).
%
\renewcommand{\thethm}{3.3}
\begin{cor}[{Polynomials integrate to zero over a closed path}]\label{thm:line-int-closed-road-of-derivative}
%
\leavevmode
\vnbpart{\texttt{line-int-closed-road-of-derivative}}
\vnbstmt{\begin{array}{@{}l@{}} \forall u, f, c, m, d, a, b. \\ \quad \operatorname{holomorphic\text{-}on}(u, f) \wedge \\ \quad (c \in (u \to \mathbb{C})) \wedge \\ \quad \forall z \in u. \\ \quad \quad \operatorname{is\text{-}diff\text{-}on}(\operatorname{cc\text{-}normed\text{-}field}, u, f, z, c(z)) \wedge \\ \quad \operatorname{is\text{-}road}(m, d, a, b) \wedge \\ \quad (\operatorname{trace}(m, a, b) \subseteq u) \wedge \\ \quad (m(b) = m(a)) \Rightarrow \\ \quad \quad (\operatorname{line\text{-}int}(c, m, d, a, b) = 0) \end{array}}
\vnbpart{\texttt{affine-closed-road-integral}}
\vnbstmt{\begin{array}{@{}l@{}} \forall p, q, w \in \mathbb{C}, g \in (\mathbb{C} \to \mathbb{C}), m, c, a, b. \\ \quad \forall z \in \mathbb{C}.\; (g(z) \simeq (p + (q \cdot (z - w)))) \wedge \\ \quad \operatorname{is\text{-}road}(m, c, a, b) \wedge \\ \quad (m(b) = m(a)) \Rightarrow \\ \quad \quad (\operatorname{line\text{-}int}(g, m, c, a, b) = 0) \end{array}}
\vnbpart{\texttt{affine-has-primitive-cc}}
\vnbstmt{\begin{array}{@{}l@{}} \forall p, q, w \in \mathbb{C}, f \in (\mathbb{C} \to \mathbb{C}). \\ \quad \forall z \in \mathbb{C}. \\ \quad \quad (f(z) \simeq ((p \cdot z) + ((q \cdot ((z - w) \cdot (z - w))) \cdot \operatorname{recip}(2)))) \Rightarrow \\ \quad \quad \operatorname{holomorphic\text{-}on}(\mathbb{C}, f) \wedge \\ \quad \quad \forall z \in \mathbb{C}. \\ \quad \quad \quad \operatorname{is\text{-}diff\text{-}on}(\operatorname{cc\text{-}normed\text{-}field}, \mathbb{C}, f, z, (p + (q \cdot (z - w)))) \end{array}}
%
\end{cor}
%
\begin{proof}
%
The proofs: the file that holds each, the lemmas it cites, its oracles and its bill.
%
\vnbprf{\vnbproofpart{\texttt{line-int-closed-road-of-derivative}}}
%
\vnbprf{\textbf{Proof.} \texttt{theorem-\allowbreak{}library/goursat-\allowbreak{}bricks.scm} --- 49 recorded steps.}
%
\vnbprf{\textbf{Lemmas.} \texttt{line-int-of-derivative}, \texttt{is-road-is-path}, \texttt{is-path-endpoints}, \texttt{is-path-in-fun}, \texttt{ccint-membership}, \texttt{trace-unfold}, \texttt{fun-apply-type-c}, \texttt{membership-implies-sethood}, \texttt{subset-mem-fwd}, \texttt{holomorphic-on-in-fun}.}
%
\vnbprf{\textbf{Oracles.} \texttt{ineq}, \texttt{crs}, \texttt{arith}.  \textbf{Bill:} \texttt{proven modulo 0}.}
%
\vnbprf{\vnbproofpart{\texttt{affine-closed-road-integral}}}
%
\vnbprf{\textbf{Proof.} \texttt{theorem-\allowbreak{}library/goursat-\allowbreak{}bricks.scm} --- 143 recorded steps.}
%
\vnbprf{\textbf{Lemmas.} \texttt{cc-recip-closed}, \texttt{cc-mul-closed}, \texttt{cc-sub-in-cc}, \texttt{cc-add-closed}, \texttt{cc-is-set}, \texttt{affine-has-primitive-cc}, \texttt{is-road-is-path}, \texttt{is-path-in-fun}, \texttt{trace-subset-cc}, \texttt{line-int-closed-road-of-derivative}.}
%
\vnbprf{\textbf{Oracles.} \texttt{arith}, \texttt{crs}, \texttt{ineq}.  \textbf{Bill:} \texttt{proven modulo 0}.}
%
\vnbprf{\vnbproofpart{\texttt{affine-has-primitive-cc}}}
%
\vnbprf{\textbf{Proof.} \texttt{theorem-\allowbreak{}library/goursat-\allowbreak{}bricks.scm} --- 188 recorded steps.}
%
\vnbprf{\textbf{Lemmas.} \texttt{cc-recip-closed}, \texttt{cc-is-metric-space}, \texttt{carrier-is-open}, \texttt{cc-ms-open-iff}, \texttt{cc-recip-inverse}, \texttt{cc-sub-in-cc}, \texttt{cc-mul-closed}, \texttt{cc-add-closed}, \texttt{cc-magnitude-closed}, \texttt{cc-magnitude-nonneg}, \texttt{cc-magnitude-mul}, \texttt{rr-mul-closed}, \texttt{rr-leq-reflexive}, \texttt{cc-diff-on-quadratic}.}
%
\vnbprf{\textbf{Oracles.} \texttt{arith}, \texttt{crs}, \texttt{ineq}.  \textbf{Bill:} \texttt{proven modulo 0}.}
%
\end{proof}
%
\noindent\textbf{Deviation from the notes.} UPDATED 2026-09-25 (batch 29-B): PROVEN in the form the tree can state: for a closed road (gamma(b) = gamma(a)) the integral of the derivative of any holomorphic f is 0 (line-int-closed-road-of-derivative), and the affine function p + q(z - w) integrates to 0 around a closed road (affine-closed-road-integral, through its primitive affine-has-primitive-cc). The notes' 'polynomial function' has no object in the tree (no complex polynomial function), so the general polynomial case is not stated; the affine case is what Goursat uses. Not in the library. The DEFINITIONS exist since 2026-09-21, for functions ON [a, b] (structure-library/path-integral.scm): is-path, is-road(gamma, dgamma, a, b) with the derivative dgamma an explicit argument, trace, cc-int (the notes' equation (44)) and line-int(f, gamma, dgamma, a, b). Proven since: the integral pw-int is well defined (two piecewise primitives have the same endpoint difference), linear, additive over adjacent intervals and monotone; cc-int is linear over CC and additive; the notes' estimate (45) (cc-int-abs-bound); the trace of a path is compact (path-trace-compact); line-int of 1 along a road is gamma(b) - gamma(a); the opposite road and Dieudonne (9.6.1). Existence of a primitive of a piecewise continuous function (Dieudonne 8.7.2) is NOT proven, so the primitives of an integrand's coordinates are hypotheses wherever a value of an integral is asserted.
%
\renewcommand{\thethm}{3.4}
\begin{lem}[{Splitting a segment integral at an interior point}]\label{thm:segment-int-split}
%
\leavevmode
\vnbpart{\texttt{segment-int-split}}
\vnbstmt{\begin{array}{@{}l@{}} \forall u, f.; \forall z_{0}, w \in \mathbb{C}, t \in \operatorname{ooint}(0, 1). \\ \quad (u \subseteq \mathbb{C}) \wedge \\ \quad (f \in (u \to \mathbb{C})) \wedge \\ \quad \forall y \in u. \\ \quad \quad \operatorname{is\text{-}continuous\text{-}at}(\operatorname{subspace\text{-}ms}(\operatorname{nf\text{-}metric\text{-}space}(\operatorname{cc\text{-}normed\text{-}field}), u), \\ \quad \quad \quad \operatorname{nf\text{-}metric\text{-}space}(\operatorname{cc\text{-}normed\text{-}field}), f, y) \wedge \\ \quad \operatorname{trace}((\lambda x \in \operatorname{ccint}(0, 1).\; (z_{0} + (x \cdot w))), 0, 1) \\ \quad \subseteq\; u \Rightarrow \\ \quad \quad \operatorname{line\text{-}int}(f, (\lambda x \in \operatorname{ccint}(0, 1).\; (z_{0} + (x \cdot w))), \\ \quad \quad \quad (\lambda x \in \operatorname{ccint}(0, 1).\; w), 0, 1) \\ \quad \quad =\; (\operatorname{line\text{-}int}(f, (\lambda x \in \operatorname{ccint}(0, 1).\; (z_{0} + (x \cdot (t \cdot w)))), (\lambda x \in \operatorname{ccint}(0, 1).\; (t \cdot w)), 0, 1) + \operatorname{line\text{-}int}(f, (\lambda x \in \operatorname{ccint}(0, 1).\; ((z_{0} + (t \cdot w)) + (x \cdot ((1 - t) \cdot w)))), (\lambda x \in \operatorname{ccint}(0, 1).\; ((1 - t) \cdot w)), 0, 1)) \end{array}}
\vnbpart{\texttt{segment-int-abs-bound}}
\vnbstmt{\begin{array}{@{}l@{}} \forall u, f.; \forall z_{0}, w \in \mathbb{C}, n \in \mathbb{R}. \\ \quad (u \subseteq \mathbb{C}) \wedge \\ \quad (f \in (u \to \mathbb{C})) \wedge \\ \quad \forall y \in u. \\ \quad \quad \operatorname{is\text{-}continuous\text{-}at}(\operatorname{subspace\text{-}ms}(\operatorname{nf\text{-}metric\text{-}space}(\operatorname{cc\text{-}normed\text{-}field}), u), \\ \quad \quad \quad \operatorname{nf\text{-}metric\text{-}space}(\operatorname{cc\text{-}normed\text{-}field}), f, y) \wedge \\ \quad \operatorname{trace}((\lambda x \in \operatorname{ccint}(0, 1).\; (z_{0} + (x \cdot w))), 0, 1) \\ \quad \subseteq\; u \wedge \\ \quad \forall \mathit{lrpy} \in \operatorname{trace}((\lambda x \in \operatorname{ccint}(0, 1).\; (z_{0} + (x \cdot w))), 0, 1). \\ \quad \quad (\operatorname{magnitude}(f(\mathit{lrpy})) \leq n) \Rightarrow \\ \quad \quad \operatorname{magnitude}(\operatorname{line\text{-}int}(f, (\lambda x \in \operatorname{ccint}(0, 1).\; (z_{0} + (x \cdot w))), \\ \quad \quad \quad \quad (\lambda x \in \operatorname{ccint}(0, 1).\; w), 0, 1)) \\ \quad \quad \leq\; (\operatorname{magnitude}(w) \cdot n) \end{array}}
\vnbpart{\texttt{segment-int-reverse}}
\vnbstmt{\begin{array}{@{}l@{}} \forall u, f.; \forall z_{0}, w \in \mathbb{C}. \\ \quad (u \subseteq \mathbb{C}) \wedge \\ \quad (f \in (u \to \mathbb{C})) \wedge \\ \quad \forall y \in u. \\ \quad \quad \operatorname{is\text{-}continuous\text{-}at}(\operatorname{subspace\text{-}ms}(\operatorname{nf\text{-}metric\text{-}space}(\operatorname{cc\text{-}normed\text{-}field}), u), \\ \quad \quad \quad \operatorname{nf\text{-}metric\text{-}space}(\operatorname{cc\text{-}normed\text{-}field}), f, y) \wedge \\ \quad \operatorname{trace}((\lambda x \in \operatorname{ccint}(0, 1).\; (z_{0} + (x \cdot w))), 0, 1) \\ \quad \subseteq\; u \Rightarrow \\ \quad \quad \operatorname{line\text{-}int}(f, (\lambda x \in \operatorname{ccint}(0, 1).\; ((z_{0} + w) + (x \cdot (-1 \cdot w)))), \\ \quad \quad \quad (\lambda x \in \operatorname{ccint}(0, 1).\; (-1 \cdot w)), 0, 1) \\ \quad \quad =\; (-1 \cdot \operatorname{line\text{-}int}(f, (\lambda x \in \operatorname{ccint}(0, 1).\; (z_{0} + (x \cdot w))), (\lambda x \in \operatorname{ccint}(0, 1).\; w), 0, 1)) \end{array}}
\vnbpart{\texttt{line-int-affine-reparam}}
\vnbstmt{\begin{array}{@{}l@{}} \forall u, f, m, c, a, b.; \forall \mathit{lrpc}, d, l, \mathit{lrpm} \in \mathbb{R}, r, k. \\ \quad (u \subseteq \mathbb{C}) \wedge \\ \quad (f \in (u \to \mathbb{C})) \wedge \\ \quad \forall y \in u. \\ \quad \quad \operatorname{is\text{-}continuous\text{-}at}(\operatorname{subspace\text{-}ms}(\operatorname{nf\text{-}metric\text{-}space}(\operatorname{cc\text{-}normed\text{-}field}), u), \\ \quad \quad \quad \operatorname{nf\text{-}metric\text{-}space}(\operatorname{cc\text{-}normed\text{-}field}), f, y) \wedge \\ \quad \operatorname{is\text{-}road}(m, c, a, b) \wedge \\ \quad (\operatorname{trace}(m, a, b) \subseteq u) \wedge \\ \quad (0 < l) \wedge \\ \quad ((\mathit{lrpm} + (l \cdot \mathit{lrpc})) = a) \wedge \\ \quad ((\mathit{lrpm} + (l \cdot d)) = b) \wedge \\ \quad (\mathit{lrpc} < d) \wedge \\ \quad (r \in (\operatorname{ccint}(\mathit{lrpc}, d) \to \mathbb{C})) \wedge \\ \quad (k \in (\operatorname{ccint}(\mathit{lrpc}, d) \to \mathbb{C})) \wedge \\ \quad \forall z \in \operatorname{ccint}(\mathit{lrpc}, d). \\ \quad \quad (r(z) \simeq m((\mathit{lrpm} + (l \cdot z)))) \wedge \\ \quad \forall z \in \operatorname{ccint}(\mathit{lrpc}, d). \\ \quad \quad (k(z) \simeq (l \cdot c((\mathit{lrpm} + (l \cdot z))))) \Rightarrow \\ \quad \quad \operatorname{line\text{-}int}(f, r, k, \mathit{lrpc}, d) \\ \quad \quad =\; \operatorname{line\text{-}int}(f, m, c, a, b) \end{array}}
\vnbpart{\texttt{line-int-adjacent}}
\vnbstmt{\begin{array}{@{}l@{}} \forall u, f, m, c, a, b.; \forall v \in \operatorname{ooint}(a, b). \\ \quad (u \subseteq \mathbb{C}) \wedge \\ \quad (f \in (u \to \mathbb{C})) \wedge \\ \quad \forall y \in u. \\ \quad \quad \operatorname{is\text{-}continuous\text{-}at}(\operatorname{subspace\text{-}ms}(\operatorname{nf\text{-}metric\text{-}space}(\operatorname{cc\text{-}normed\text{-}field}), u), \\ \quad \quad \quad \operatorname{nf\text{-}metric\text{-}space}(\operatorname{cc\text{-}normed\text{-}field}), f, y) \wedge \\ \quad \operatorname{is\text{-}road}(m, c, a, b) \wedge \\ \quad (\operatorname{trace}(m, a, b) \subseteq u) \Rightarrow \\ \quad \quad \operatorname{line\text{-}int}(f, m, c, a, b) \\ \quad \quad =\; (\operatorname{line\text{-}int}(f, \operatorname{restrict}(m, \operatorname{ccint}(a, v)), \operatorname{restrict}(c, \operatorname{ccint}(a, v)), a, v) + \operatorname{line\text{-}int}(f, \operatorname{restrict}(m, \operatorname{ccint}(v, b)), \operatorname{restrict}(c, \operatorname{ccint}(v, b)), v, b)) \end{array}}
%
\end{lem}
%
\begin{proof}
%
The proofs: the file that holds each, the lemmas it cites, its oracles and its bill.
%
\vnbprf{\vnbproofpart{\texttt{segment-int-split}}}
%
\vnbprf{\textbf{Proof.} \texttt{theorem-\allowbreak{}library/line-\allowbreak{}int-\allowbreak{}reparam.scm} --- 185 recorded steps.}
%
\vnbprf{\textbf{Lemmas.} \texttt{ooint-membership}, \texttt{rr-zero-in}, \texttt{rr-one-in}, \texttt{rr-zero-lt-one}, \texttt{rr-leq-reflexive}, \texttt{rr-lt-implies-le}, \texttt{rr-is-set}, \texttt{rr-subset-cc}, \texttt{rr-sub-in-rr}, \texttt{segment-is-road}, \texttt{is-road-is-path}, \texttt{is-path-in-fun}, \texttt{is-road-dgam-in-fun}, \texttt{cc-mul-closed}, \texttt{cc-add-closed}, \texttt{line-int-adjacent}, \texttt{road-restrict}, \texttt{trace-restrict-subset}, \texttt{subset-trans}, \texttt{ccint-subset-ccint}, \texttt{affine-reparam-ccint-in}, \texttt{ccint-elt-in-rr}, \texttt{subset-mem-fwd}, \texttt{restrict-apply}, \texttt{line-int-affine-reparam}.}
%
\vnbprf{\textbf{Oracles.} \texttt{ineq}, \texttt{crs}, \texttt{arith}.  \textbf{Bill:} \texttt{proven modulo 0}.}
%
\vnbprf{\vnbproofpart{\texttt{segment-int-abs-bound}}}
%
\vnbprf{\textbf{Proof.} \texttt{theorem-\allowbreak{}library/line-\allowbreak{}int-\allowbreak{}reparam.scm} --- 356 recorded steps.}
%
\vnbprf{\textbf{Lemmas.} \texttt{rr-zero-in}, \texttt{rr-one-in}, \texttt{rr-zero-lt-one}, \texttt{rr-is-set}, \texttt{ccint-subset-rr}, \texttt{subclass-of-set-is-set}, \texttt{segment-is-road}, \texttt{is-road-is-path}, \texttt{is-path-in-fun}, \texttt{is-road-dgam-in-fun}, \texttt{cc-magnitude-closed}, \texttt{cc-magnitude-nonneg}, \texttt{rr-mul-closed}, \texttt{ccint-elt-in-rr}, \texttt{rr-subset-cc}, \texttt{cc-mul-closed}, \texttt{cc-add-closed}, \texttt{trace-value-in}, \texttt{subset-mem-fwd}, \texttt{fun-apply-type-c}, \texttt{rr-add-closed}, \texttt{line-int-abs-bound}, \texttt{road-integrand-regulated}, \texttt{regulated-on-magnitude}, \texttt{real-part-in-rr}, \texttt{imag-part-in-rr}, \texttt{cc-re-im-decompose}, \texttt{cc-i-in}, \texttt{cc-mul-comm}, \texttt{cc-magnitude-mul}, \texttt{regulated-on-has-primitive}, \texttt{pw-affine-antiderivative}, \texttt{rr-mul-le-right}, \texttt{pw-int-monotone}, \texttt{pw-int-const}, \texttt{pw-int-in-rr}, \texttt{line-int-in-cc}.}
%
\vnbprf{\textbf{Oracles.} \texttt{crs}, \texttt{ineq}, \texttt{arith}.  \textbf{Bill:} \texttt{proven modulo 0}.}
%
\vnbprf{\vnbproofpart{\texttt{segment-int-reverse}}}
%
\vnbprf{\textbf{Proof.} \texttt{theorem-\allowbreak{}library/line-\allowbreak{}int-\allowbreak{}reparam.scm} --- 66 recorded steps.}
%
\vnbprf{\textbf{Lemmas.} \texttt{rr-zero-in}, \texttt{rr-one-in}, \texttt{rr-neg-closed}, \texttt{rr-subset-cc}, \texttt{cc-mul-closed}, \texttt{cc-add-closed}, \texttt{segment-is-road}, \texttt{is-road-is-path}, \texttt{is-path-in-fun}, \texttt{is-road-dgam-in-fun}, \texttt{ccint-elt-in-rr}, \texttt{ccint-reflect-in}, \texttt{line-int-opposite-road}.}
%
\vnbprf{\textbf{Oracles.} \texttt{crs}, \texttt{ineq}, \texttt{arith}.  \textbf{Bill:} \texttt{proven modulo 0}.}
%
\vnbprf{\vnbproofpart{\texttt{line-int-affine-reparam}}}
%
\vnbprf{\textbf{Proof.} \texttt{theorem-\allowbreak{}library/line-\allowbreak{}int-\allowbreak{}reparam.scm} --- 109 recorded steps.}
%
\vnbprf{\textbf{Lemmas.} \texttt{is-road-is-path}, \texttt{is-path-in-fun}, \texttt{is-road-dgam-in-fun}, \texttt{is-path-endpoints}, \texttt{rr-lt-implies-le}, \texttt{rr-is-set}, \texttt{ccint-subset-rr}, \texttt{subclass-of-set-is-set}, \texttt{line-int-exists}, \texttt{line-int-integrand-in-fun}, \texttt{trace-affine-reparam}, \texttt{subset-trans}, \texttt{rr-subset-cc}, \texttt{affine-reparam-ccint-in}, \texttt{fun-apply-type-c}, \texttt{trace-value-in}, \texttt{subset-mem-fwd}, \texttt{cc-mul-assoc}, \texttt{cc-mul-comm}, \texttt{cc-int-affine-subst}, \texttt{line-int-unfold}.}
%
\vnbprf{\textbf{Oracles.} \texttt{ineq}, \texttt{arith}, \texttt{crs}.  \textbf{Bill:} \texttt{proven modulo 0}.}
%
\vnbprf{\vnbproofpart{\texttt{line-int-adjacent}}}
%
\vnbprf{\textbf{Proof.} \texttt{theorem-\allowbreak{}library/line-\allowbreak{}int-\allowbreak{}laws.scm} --- 436 recorded steps.}
%
\vnbprf{\textbf{Lemmas.} \texttt{is-road-is-path}, \texttt{is-road-dgam-in-fun}, \texttt{is-path-in-fun}, \texttt{is-path-endpoints}, \texttt{rr-is-set}, \texttt{ccint-subset-rr}, \texttt{subclass-of-set-is-set}, \texttt{trace-value-in}, \texttt{subset-mem-fwd}, \texttt{line-int-exists}, \texttt{line-int-integrand-in-fun}, \texttt{cc-int-adjacent}, \texttt{ooint-elt-in-rr}, \texttt{ooint-membership}, \texttt{rr-leq-reflexive}, \texttt{rr-lt-implies-le}, \texttt{ccint-subset-ccint}, \texttt{pw-antiderivative-restrict}, \texttt{fun-apply-type-c}, \texttt{cc-mul-closed}, \texttt{real-part-in-rr}, \texttt{imag-part-in-rr}, \texttt{restrict-apply}, \texttt{pw-antiderivative-integrand-congruence}, \texttt{cc-int-congruence}.}
%
\vnbprf{\textbf{Oracles.} \texttt{arith}, \texttt{crs}, \texttt{ineq}.  \textbf{Bill:} \texttt{proven modulo 0}.}
%
\end{proof}
%
\noindent\textbf{Deviation from the notes.} UPDATED 2026-09-25 (batch 29-A): PROVEN for the segment itself (theorem-library/line-int-reparam.scm): segment-int-split splits the integral along the segment z0 + s w, s in [0,1], at the interior point z0 + t w (0 $<$ t $<$ 1) into the integrals along the two sub-segments, through the affine reparametrisation of the line integral, the notes' (48) (line-int-affine-reparam). The endpoint cases and the restatement with the segment functoid SEG-INT are batch 30. UPDATED 2026-09-23 (batch 24): proven for ANY road, line-int-adjacent (split at an interior point c of [a, b]); the segment of the notes is the special case where gamma is a straight segment (segment-is-road), and the restriction of a segment is not restated as a segment. EARLIER NOTE: Not in the library. The DEFINITIONS exist since 2026-09-21, for functions ON [a, b] (structure-library/path-integral.scm): is-path, is-road(gamma, dgamma, a, b) with the derivative dgamma an explicit argument, trace, cc-int (the notes' equation (44)) and line-int(f, gamma, dgamma, a, b). Proven since: the integral pw-int is well defined (two piecewise primitives have the same endpoint difference), linear, additive over adjacent intervals and monotone; cc-int is linear over CC and additive; the notes' estimate (45) (cc-int-abs-bound); the trace of a path is compact (path-trace-compact); line-int of 1 along a road is gamma(b) - gamma(a); the opposite road and Dieudonne (9.6.1). Existence of a primitive of a piecewise continuous function (Dieudonne 8.7.2) is NOT proven, so the primitives of an integrand's coordinates are hypotheses wherever a value of an integral is asserted.
%
\renewcommand{\thethm}{3.5}
\begin{lem}[{The barycentric-subdivision sum}]\label{thm:tri-int-subdivision}
%
\leavevmode
\vnbpart{\texttt{tri-int-subdivision}}
\vnbstmt{\begin{array}{@{}l@{}} \forall u, f.; \forall a, b, c \in \mathbb{C}. \\ \quad (u \subseteq \mathbb{C}) \wedge \\ \quad (f \in (u \to \mathbb{C})) \wedge \\ \quad \forall y \in u. \\ \quad \quad \operatorname{is\text{-}continuous\text{-}at}(\operatorname{subspace\text{-}ms}(\operatorname{nf\text{-}metric\text{-}space}(\operatorname{cc\text{-}normed\text{-}field}), u), \\ \quad \quad \quad \operatorname{nf\text{-}metric\text{-}space}(\operatorname{cc\text{-}normed\text{-}field}), f, y) \wedge \\ \quad (\operatorname{conv3}(a, b, c) \subseteq u) \Rightarrow \\ \quad \quad \operatorname{tri\text{-}int}(f, a, b, c) \\ \quad \quad =\; (((\operatorname{tri\text{-}int}(f, a, \operatorname{cc\text{-}mid}(a, b), \operatorname{cc\text{-}mid}(a, c)) + \operatorname{tri\text{-}int}(f, \operatorname{cc\text{-}mid}(a, b), b, \operatorname{cc\text{-}mid}(b, c))) + \operatorname{tri\text{-}int}(f, \operatorname{cc\text{-}mid}(a, c), \operatorname{cc\text{-}mid}(b, c), c)) + \operatorname{tri\text{-}int}(f, \operatorname{cc\text{-}mid}(a, b), \operatorname{cc\text{-}mid}(b, c), \operatorname{cc\text{-}mid}(a, c))) \end{array}}
\vnbpart{\texttt{seg-int-split-at}}
\vnbstmt{\begin{array}{@{}l@{}} \forall u, f.; \forall a, b \in \mathbb{C}, c \in \operatorname{trace}(\operatorname{seg\text{-}path}(a, b), 0, 1). \\ \quad (u \subseteq \mathbb{C}) \wedge \\ \quad (f \in (u \to \mathbb{C})) \wedge \\ \quad \forall y \in u. \\ \quad \quad \operatorname{is\text{-}continuous\text{-}at}(\operatorname{subspace\text{-}ms}(\operatorname{nf\text{-}metric\text{-}space}(\operatorname{cc\text{-}normed\text{-}field}), u), \\ \quad \quad \quad \operatorname{nf\text{-}metric\text{-}space}(\operatorname{cc\text{-}normed\text{-}field}), f, y) \wedge \\ \quad (\operatorname{trace}(\operatorname{seg\text{-}path}(a, b), 0, 1) \subseteq u) \Rightarrow \\ \quad \quad \operatorname{seg\text{-}int}(f, a, b) \\ \quad \quad =\; (\operatorname{seg\text{-}int}(f, a, c) + \operatorname{seg\text{-}int}(f, c, b)) \end{array}}
\vnbpart{\texttt{segment-int-reverse}}
\vnbstmt{\begin{array}{@{}l@{}} \forall u, f.; \forall z_{0}, w \in \mathbb{C}. \\ \quad (u \subseteq \mathbb{C}) \wedge \\ \quad (f \in (u \to \mathbb{C})) \wedge \\ \quad \forall y \in u. \\ \quad \quad \operatorname{is\text{-}continuous\text{-}at}(\operatorname{subspace\text{-}ms}(\operatorname{nf\text{-}metric\text{-}space}(\operatorname{cc\text{-}normed\text{-}field}), u), \\ \quad \quad \quad \operatorname{nf\text{-}metric\text{-}space}(\operatorname{cc\text{-}normed\text{-}field}), f, y) \wedge \\ \quad \operatorname{trace}((\lambda x \in \operatorname{ccint}(0, 1).\; (z_{0} + (x \cdot w))), 0, 1) \\ \quad \subseteq\; u \Rightarrow \\ \quad \quad \operatorname{line\text{-}int}(f, (\lambda x \in \operatorname{ccint}(0, 1).\; ((z_{0} + w) + (x \cdot (-1 \cdot w)))), \\ \quad \quad \quad (\lambda x \in \operatorname{ccint}(0, 1).\; (-1 \cdot w)), 0, 1) \\ \quad \quad =\; (-1 \cdot \operatorname{line\text{-}int}(f, (\lambda x \in \operatorname{ccint}(0, 1).\; (z_{0} + (x \cdot w))), (\lambda x \in \operatorname{ccint}(0, 1).\; w), 0, 1)) \end{array}}
%
\end{lem}
%
\begin{proof}
%
The proofs: the file that holds each, the lemmas it cites, its oracles and its bill.
%
\vnbprf{\vnbproofpart{\texttt{tri-int-subdivision}}}
%
\vnbprf{\textbf{Proof.} \texttt{theorem-\allowbreak{}library/goursat.scm} --- 176 recorded steps.}
%
\vnbprf{\textbf{Lemmas.} \texttt{conv3-vertices-in}, \texttt{conv3-mid-in}, \texttt{cc-add-closed}, \texttt{rr-subset-cc}, \texttt{cc-recip-closed}, \texttt{cc-mul-closed}, \texttt{cc-mid-unfold}, \texttt{conv3-seg-trace}, \texttt{subset-trans}, \texttt{seg-int-in-cc}, \texttt{cc-mid-in-seg-trace}, \texttt{tri-int-unfold}, \texttt{seg-int-split-at}, \texttt{seg-int-reverse-at}.}
%
\vnbprf{\textbf{Oracles.} \texttt{arith}, \texttt{crs}, \texttt{ineq}.  \textbf{Bill:} \texttt{proven modulo 0}.}
%
\vnbprf{\vnbproofpart{\texttt{seg-int-split-at}}}
%
\vnbprf{\textbf{Proof.} \texttt{theorem-\allowbreak{}library/goursat.scm} --- 269 recorded steps.}
%
\vnbprf{\textbf{Lemmas.} \texttt{seg-trace-elim}, \texttt{ccint-elt-in-rr}, \texttt{ccint-membership}, \texttt{rr-subset-cc}, \texttt{cc-sub-in-cc}, \texttt{cc-mul-closed}, \texttt{cc-add-closed}, \texttt{subset-mem-fwd}, \texttt{seg-int-in-cc}, \texttt{rr-zero-in}, \texttt{rr-one-in}, \texttt{cc-zero-in}, \texttt{seg-trace-intro}, \texttt{in-transport-eq}, \texttt{seg-int-degenerate}, \texttt{neq-sym}, \texttt{rr-le-ne-lt}, \texttt{ooint-membership}, \texttt{segment-is-road}, \texttt{is-road-is-path}, \texttt{is-path-in-fun}, \texttt{is-road-dgam-in-fun}, \texttt{seg-trace-bridge}, \texttt{seg-int-bridge}, \texttt{segment-int-split}.}
%
\vnbprf{\textbf{Oracles.} \texttt{ineq}, \texttt{crs}, \texttt{arith}.  \textbf{Bill:} \texttt{proven modulo 0}.}
%
\vnbprf{\vnbproofpart{\texttt{segment-int-reverse}}}
%
\vnbprf{\textbf{Proof.} \texttt{theorem-\allowbreak{}library/line-\allowbreak{}int-\allowbreak{}reparam.scm} --- 66 recorded steps.}
%
\vnbprf{\textbf{Lemmas.} \texttt{rr-zero-in}, \texttt{rr-one-in}, \texttt{rr-neg-closed}, \texttt{rr-subset-cc}, \texttt{cc-mul-closed}, \texttt{cc-add-closed}, \texttt{segment-is-road}, \texttt{is-road-is-path}, \texttt{is-path-in-fun}, \texttt{is-road-dgam-in-fun}, \texttt{ccint-elt-in-rr}, \texttt{ccint-reflect-in}, \texttt{line-int-opposite-road}.}
%
\vnbprf{\textbf{Oracles.} \texttt{crs}, \texttt{ineq}, \texttt{arith}.  \textbf{Bill:} \texttt{proven modulo 0}.}
%
\end{proof}
%
\noindent\textbf{Deviation from the notes.} PROVEN 2026-09-25 (batch 30, theorem-library/goursat.scm). The triangle integral TRI-INT(f, a, b, c) is DEFINED as the sum of the three segment integrals (the user's decision, docs/goursat-definitions-2026-09-25.md), so the triangular path is not a road object; the subdivision S(triangle) is written out as the four vertex triples with CC-MID. Hypothesis: f continuous on a set containing CONV3(a, b, c).
%
\renewcommand{\thethm}{3.6}
\begin{prop}[{Cauchy's theorem for triangles (Goursat)}]\label{thm:goursat-triangle-except}
%
\leavevmode
\vnbpart{\texttt{goursat-triangle-except}}
\vnbstmt{\begin{array}{@{}l@{}} \forall u, f, x.; \forall a, b, c \in u. \\ \quad \operatorname{is\text{-}open}(\operatorname{nf\text{-}metric\text{-}space}(\operatorname{cc\text{-}normed\text{-}field}), u) \wedge \\ \quad (f \in (u \to \mathbb{C})) \wedge \\ \quad \forall y \in u. \\ \quad \quad \operatorname{is\text{-}continuous\text{-}at}(\operatorname{subspace\text{-}ms}(\operatorname{nf\text{-}metric\text{-}space}(\operatorname{cc\text{-}normed\text{-}field}), u), \\ \quad \quad \quad \operatorname{nf\text{-}metric\text{-}space}(\operatorname{cc\text{-}normed\text{-}field}), f, y) \wedge \\ \quad (x \in u) \wedge \\ \quad \operatorname{holomorphic\text{-}on}(\operatorname{difference}(u, \operatorname{singleton}(x)), \\ \quad \quad \operatorname{restrict}(f, \operatorname{difference}(u, \operatorname{singleton}(x)))) \wedge \\ \quad (\operatorname{conv3}(a, b, c) \subseteq u) \Rightarrow \\ \quad \quad (\operatorname{tri\text{-}int}(f, a, b, c) = 0) \end{array}}
\vnbpart{\texttt{goursat-vertex}}
\vnbstmt{\begin{array}{@{}l@{}} \forall u, f, x.; \forall b, c \in u. \\ \quad \operatorname{is\text{-}open}(\operatorname{nf\text{-}metric\text{-}space}(\operatorname{cc\text{-}normed\text{-}field}), u) \wedge \\ \quad (f \in (u \to \mathbb{C})) \wedge \\ \quad \forall y \in u. \\ \quad \quad \operatorname{is\text{-}continuous\text{-}at}(\operatorname{subspace\text{-}ms}(\operatorname{nf\text{-}metric\text{-}space}(\operatorname{cc\text{-}normed\text{-}field}), u), \\ \quad \quad \quad \operatorname{nf\text{-}metric\text{-}space}(\operatorname{cc\text{-}normed\text{-}field}), f, y) \wedge \\ \quad (x \in u) \wedge \\ \quad \operatorname{holomorphic\text{-}on}(\operatorname{difference}(u, \operatorname{singleton}(x)), \\ \quad \quad \operatorname{restrict}(f, \operatorname{difference}(u, \operatorname{singleton}(x)))) \wedge \\ \quad (\operatorname{conv3}(x, b, c) \subseteq u) \Rightarrow \\ \quad \quad (\operatorname{tri\text{-}int}(f, x, b, c) = 0) \end{array}}
\vnbpart{\texttt{tri-int-fan}}
\vnbstmt{\begin{array}{@{}l@{}} \forall u, f.; \forall a, b, c \in \mathbb{C}, x \in \operatorname{conv3}(a, b, c). \\ \quad (u \subseteq \mathbb{C}) \wedge \\ \quad (f \in (u \to \mathbb{C})) \wedge \\ \quad \forall y \in u. \\ \quad \quad \operatorname{is\text{-}continuous\text{-}at}(\operatorname{subspace\text{-}ms}(\operatorname{nf\text{-}metric\text{-}space}(\operatorname{cc\text{-}normed\text{-}field}), u), \\ \quad \quad \quad \operatorname{nf\text{-}metric\text{-}space}(\operatorname{cc\text{-}normed\text{-}field}), f, y) \wedge \\ \quad (\operatorname{conv3}(a, b, c) \subseteq u) \Rightarrow \\ \quad \quad \operatorname{tri\text{-}int}(f, a, b, c) \\ \quad \quad =\; ((\operatorname{tri\text{-}int}(f, a, b, x) + \operatorname{tri\text{-}int}(f, b, c, x)) + \operatorname{tri\text{-}int}(f, c, a, x)) \end{array}}
\vnbpart{\texttt{tri-int-vertex-split}}
\vnbstmt{\begin{array}{@{}l@{}} \forall u, f.; \forall a, b, c \in \mathbb{C}, t \in \operatorname{ccint}(0, 1). \\ \quad (u \subseteq \mathbb{C}) \wedge \\ \quad (f \in (u \to \mathbb{C})) \wedge \\ \quad \forall y \in u. \\ \quad \quad \operatorname{is\text{-}continuous\text{-}at}(\operatorname{subspace\text{-}ms}(\operatorname{nf\text{-}metric\text{-}space}(\operatorname{cc\text{-}normed\text{-}field}), u), \\ \quad \quad \quad \operatorname{nf\text{-}metric\text{-}space}(\operatorname{cc\text{-}normed\text{-}field}), f, y) \wedge \\ \quad (\operatorname{conv3}(a, b, c) \subseteq u) \Rightarrow \\ \quad \quad \operatorname{tri\text{-}int}(f, a, b, c) \\ \quad \quad =\; ((\operatorname{tri\text{-}int}(f, a, (a + (t \cdot (b - a))), (a + (t \cdot (c - a)))) + \operatorname{tri\text{-}int}(f, (a + (t \cdot (b - a))), b, c)) + \operatorname{tri\text{-}int}(f, (a + (t \cdot (b - a))), c, (a + (t \cdot (c - a))))) \end{array}}
\vnbpart{\texttt{conv3-avoids-vertex}}
\vnbstmt{\begin{array}{@{}l@{}} \forall a, b, c \in \mathbb{C}, t \in \mathbb{R}. \\ \quad \neg (\operatorname{imag\text{-}part}(((b - a) \cdot \operatorname{conjugate}((c - a)))) = 0) \wedge \\ \quad (0 < t) \wedge \\ \quad (t \leq 1) \Rightarrow \\ \quad \quad \neg (a \in \operatorname{conv3}((a + (t \cdot (b - a))), b, c)) \wedge \\ \quad \quad \neg (a \in \operatorname{conv3}((a + (t \cdot (b - a))), c, (a + (t \cdot (c - a))))) \end{array}}
\vnbpart{\texttt{collinear-cases}}
\vnbstmt{\begin{array}{@{}l@{}} \forall a, b, c \in \mathbb{C}. \\ \quad (\operatorname{imag\text{-}part}(((b - a) \cdot \operatorname{conjugate}((c - a)))) = 0) \Rightarrow  \\ \quad \quad (c \in \operatorname{trace}(\operatorname{seg\text{-}path}(a, b), 0, 1)) \vee \\ \quad \quad (a \in \operatorname{trace}(\operatorname{seg\text{-}path}(b, c), 0, 1)) \vee \\ \quad \quad (b \in \operatorname{trace}(\operatorname{seg\text{-}path}(c, a), 0, 1)) \end{array}}
\vnbpart{\texttt{tri-int-collinear}}
\vnbstmt{\begin{array}{@{}l@{}} \forall u, f.; \forall a, b, c \in \mathbb{C}. \\ \quad (u \subseteq \mathbb{C}) \wedge \\ \quad (f \in (u \to \mathbb{C})) \wedge \\ \quad \forall y \in u. \\ \quad \quad \operatorname{is\text{-}continuous\text{-}at}(\operatorname{subspace\text{-}ms}(\operatorname{nf\text{-}metric\text{-}space}(\operatorname{cc\text{-}normed\text{-}field}), u), \\ \quad \quad \quad \operatorname{nf\text{-}metric\text{-}space}(\operatorname{cc\text{-}normed\text{-}field}), f, y) \wedge \\ \quad (\operatorname{conv3}(a, b, c) \subseteq u) \wedge \\ \quad (c \in \operatorname{trace}(\operatorname{seg\text{-}path}(a, b), 0, 1)) \vee \\ \quad (a \in \operatorname{trace}(\operatorname{seg\text{-}path}(b, c), 0, 1)) \vee \\ \quad (b \in \operatorname{trace}(\operatorname{seg\text{-}path}(c, a), 0, 1)) \Rightarrow \\ \quad \quad (\operatorname{tri\text{-}int}(f, a, b, c) = 0) \end{array}}
\vnbpart{\texttt{tri-int-restrict}}
\vnbstmt{\begin{array}{@{}l@{}} \forall u, f, w.; \forall a, b, c \in \mathbb{C}. \\ \quad (u \subseteq \mathbb{C}) \wedge \\ \quad (f \in (u \to \mathbb{C})) \wedge \\ \quad \forall y \in u. \\ \quad \quad \operatorname{is\text{-}continuous\text{-}at}(\operatorname{subspace\text{-}ms}(\operatorname{nf\text{-}metric\text{-}space}(\operatorname{cc\text{-}normed\text{-}field}), u), \\ \quad \quad \quad \operatorname{nf\text{-}metric\text{-}space}(\operatorname{cc\text{-}normed\text{-}field}), f, y) \wedge \\ \quad (w \subseteq u) \wedge \\ \quad (\operatorname{conv3}(a, b, c) \subseteq w) \Rightarrow \\ \quad \quad \operatorname{tri\text{-}int}(\operatorname{restrict}(f, w), a, b, c) \\ \quad \quad =\; \operatorname{tri\text{-}int}(f, a, b, c) \end{array}}
\vnbpart{\texttt{goursat-triangle}}
\vnbstmt{\begin{array}{@{}l@{}} \forall u, f.; \forall a, b, c \in u. \\ \quad (\operatorname{holomorphic\text{-}on}(u, f) \wedge (\operatorname{conv3}(a, b, c) \subseteq u)) \Rightarrow  \\ \quad \quad (\operatorname{tri\text{-}int}(f, a, b, c) = 0) \end{array}}
%
\end{prop}
%
\begin{proof}
%
The proofs: the file that holds each, the lemmas it cites, its oracles and its bill.
%
\vnbprf{\vnbproofpart{\texttt{goursat-triangle-except}}}
%
\vnbprf{\textbf{Proof.} \texttt{theorem-\allowbreak{}library/goursat-\allowbreak{}exceptional.scm} --- 161 recorded steps.}
%
\vnbprf{\textbf{Lemmas.} \texttt{nf-cc-is-metric-space}, \texttt{cc-open-subset-cc}, \texttt{holomorphic-on-open}, \texttt{holomorphic-on-in-fun}, \texttt{holomorphic-on-continuous}, \texttt{subset-def}, \texttt{difference-membership}, \texttt{subset-mem-fwd}, \texttt{conv3-vertices-in}, \texttt{conv3-seg-trace}, \texttt{subset-trans}, \texttt{seg-int-in-cc}, \texttt{conv3-sub-hull}, \texttt{goursat-vertex}, \texttt{tri-int-unfold}, \texttt{tri-int-fan}, \texttt{singleton-membership}, \texttt{in-transport-eq}, \texttt{goursat-triangle}, \texttt{tri-int-restrict}.}
%
\vnbprf{\textbf{Oracles.} \texttt{crs}, \texttt{arith}, \texttt{ineq}.  \textbf{Bill:} \texttt{proven modulo 0}.}
%
\vnbprf{\vnbproofpart{\texttt{goursat-vertex}}}
%
\vnbprf{\textbf{Proof.} \texttt{theorem-\allowbreak{}library/goursat-\allowbreak{}exceptional.scm} --- 615 recorded steps.}
%
\vnbprf{\textbf{Lemmas.} \texttt{nf-cc-is-metric-space}, \texttt{cc-open-subset-cc}, \texttt{holomorphic-on-open}, \texttt{holomorphic-on-in-fun}, \texttt{holomorphic-on-continuous}, \texttt{subset-def}, \texttt{difference-membership}, \texttt{cc-is-set}, \texttt{subclass-of-set-is-set}, \texttt{subset-mem-fwd}, \texttt{tri-int-in-cc}, \texttt{collinear-cases}, \texttt{tri-int-collinear}, \texttt{cc-is-metric-space}, \texttt{fun-apply-type-c}, \texttt{cc-sub-in-cc}, \texttt{cc-magnitude-closed}, \texttt{cc-magnitude-nonneg}, \texttt{rr-one-in}, \texttt{rr-add-in-rr}, \texttt{rr-mul-in-rr}, \texttt{rr-leq-mul-nonneg}, \texttt{rr-pos-rr-in-rr}, \texttt{rr-lt-of-pos-rr}, \texttt{subspace-pts}, \texttt{subspace-dist}, \texttt{nf-cc-agree-dist}, \texttt{cc-ms-dist}, \texttt{cc-magnitude-sub-sym}, \texttt{cc-magnitude-le-add}, \texttt{rr-pos-ne-zero}, \texttt{rr-recip-closed}, \texttt{rr-recip-inverse}, \texttt{rr-recip-pos}, \texttt{rr-lt-implies-le}, \texttt{rr-mul-pos}, \texttt{rr-min-pos}, \texttt{rr-mul-le-right}, \texttt{ccint-membership}, \texttt{rr-subset-cc}, \texttt{cc-mul-closed}, \texttt{cc-add-closed}, \texttt{conv3-vertex-shrink}, \texttt{subset-trans}, \texttt{tri-int-vertex-split}, \texttt{conv3-avoids-vertex}, \texttt{singleton-membership}, \texttt{in-transport-eq}, \texttt{conv3-vertices-in}, \texttt{goursat-triangle}, \texttt{tri-int-restrict}, \texttt{cc-magnitude-mul}, \texttt{rr-magnitude-is-abs}, \texttt{rr-abs-of-nonneg}, \texttt{conv3-near-vertex}, \texttt{tri-int-abs-bound}, \texttt{rr-le-all-pos-nonpos}, \texttt{cc-magnitude-zero-iff}.}
%
\vnbprf{\textbf{Oracles.} \texttt{ineq}, \texttt{arith}, \texttt{crs}.  \textbf{Bill:} \texttt{proven modulo 0}.}
%
\vnbprf{\vnbproofpart{\texttt{tri-int-fan}}}
%
\vnbprf{\textbf{Proof.} \texttt{theorem-\allowbreak{}library/goursat-\allowbreak{}exceptional.scm} --- 54 recorded steps.}
%
\vnbprf{\textbf{Lemmas.} \texttt{conv3-vertices-in}, \texttt{subset-mem-fwd}, \texttt{conv3-seg-trace}, \texttt{subset-trans}, \texttt{seg-int-in-cc}, \texttt{tri-int-unfold}, \texttt{seg-int-reverse-at}.}
%
\vnbprf{\textbf{Oracles.} \texttt{crs}, \texttt{ineq}, \texttt{arith}.  \textbf{Bill:} \texttt{proven modulo 0}.}
%
\vnbprf{\vnbproofpart{\texttt{tri-int-vertex-split}}}
%
\vnbprf{\textbf{Proof.} \texttt{theorem-\allowbreak{}library/goursat-\allowbreak{}exceptional.scm} --- 162 recorded steps.}
%
\vnbprf{\textbf{Lemmas.} \texttt{ccint-elt-in-rr}, \texttt{ccint-membership}, \texttt{rr-subset-cc}, \texttt{cc-sub-in-cc}, \texttt{cc-mul-closed}, \texttt{cc-add-closed}, \texttt{conv3-vertex-shrink}, \texttt{conv3-vertices-in}, \texttt{subset-mem-fwd}, \texttt{conv3-seg-trace}, \texttt{subset-trans}, \texttt{seg-int-in-cc}, \texttt{seg-trace-intro}, \texttt{rr-one-in}, \texttt{rr-sub-in-rr}, \texttt{in-transport-eq}, \texttt{tri-int-unfold}, \texttt{seg-int-split-at}, \texttt{seg-int-reverse-at}.}
%
\vnbprf{\textbf{Oracles.} \texttt{ineq}, \texttt{crs}, \texttt{arith}.  \textbf{Bill:} \texttt{proven modulo 0}.}
%
\vnbprf{\vnbproofpart{\texttt{conv3-avoids-vertex}}}
%
\vnbprf{\textbf{Proof.} \texttt{theorem-\allowbreak{}library/goursat-\allowbreak{}exceptional.scm} --- 422 recorded steps.}
%
\vnbprf{\textbf{Lemmas.} \texttt{rr-lt-implies-le}, \texttt{conv3-member-iff}, \texttt{ccint-elt-in-rr}, \texttt{ccint-membership}, \texttt{rr-mul-in-rr}, \texttt{rr-add-in-rr}, \texttt{rr-subset-cc}, \texttt{cc-sub-in-cc}, \texttt{cc-mul-closed}, \texttt{cc-add-closed}, \texttt{cc-affinely-independent}, \texttt{rr-leq-mul-nonneg}, \texttt{rr-lt-irrefl}.}
%
\vnbprf{\textbf{Oracles.} \texttt{crs}, \texttt{ineq}, \texttt{arith}.  \textbf{Bill:} \texttt{proven modulo 0}.}
%
\vnbprf{\vnbproofpart{\texttt{collinear-cases}}}
%
\vnbprf{\textbf{Proof.} \texttt{theorem-\allowbreak{}library/goursat-\allowbreak{}exceptional.scm} --- 757 recorded steps.}
%
\vnbprf{\textbf{Lemmas.} \texttt{cc-i-in}, \texttt{cc-sub-in-cc}, \texttt{cc-cross-coords}, \texttt{cc-generated-by-rr}, \texttt{cc-re-im-of}, \texttt{rr-mul-in-rr}, \texttt{rr-sub-in-rr}, \texttt{rr-add-in-rr}, \texttt{rr-subset-cc}, \texttt{cc-mul-closed}, \texttt{rr-sq-nonneg}, \texttt{rr-no-zero-divisors}, \texttt{rr-one-in}, \texttt{ccint-membership}, \texttt{seg-trace-intro}, \texttt{cc-add-closed}, \texttt{in-transport-eq}, \texttt{neq-sym}, \texttt{rr-le-ne-lt}, \texttt{rr-recip-closed}, \texttt{rr-recip-inverse}, \texttt{rr-recip-pos}, \texttt{rr-lt-implies-le}, \texttt{rr-leq-mul-nonneg}, \texttt{rr-mul-le-right}, \texttt{rr-not-le-lt}, \texttt{rr-pos-ne-zero}.}
%
\vnbprf{\textbf{Oracles.} \texttt{ineq}, \texttt{crs}, \texttt{arith}.  \textbf{Bill:} \texttt{proven modulo 0}.}
%
\vnbprf{\vnbproofpart{\texttt{tri-int-collinear}}}
%
\vnbprf{\textbf{Proof.} \texttt{theorem-\allowbreak{}library/goursat-\allowbreak{}exceptional.scm} --- 57 recorded steps.}
%
\vnbprf{\textbf{Lemmas.} \texttt{conv3-vertices-in}, \texttt{conv3-seg-trace}, \texttt{subset-trans}, \texttt{seg-int-in-cc}, \texttt{tri-int-unfold}, \texttt{seg-int-split-at}, \texttt{seg-int-reverse-at}.}
%
\vnbprf{\textbf{Oracles.} \texttt{crs}, \texttt{ineq}, \texttt{arith}.  \textbf{Bill:} \texttt{proven modulo 0}.}
%
\vnbprf{\vnbproofpart{\texttt{tri-int-restrict}}}
%
\vnbprf{\textbf{Proof.} \texttt{theorem-\allowbreak{}library/goursat-\allowbreak{}exceptional.scm} --- 71 recorded steps.}
%
\vnbprf{\textbf{Lemmas.} \texttt{restrict-in-fun}, \texttt{seg-trace-in-conv3}, \texttt{tri-int-unfold}, \texttt{subset-trans}, \texttt{subset-mem-fwd}, \texttt{restrict-apply}, \texttt{seg-int-in-cc}, \texttt{seg-int-congruence}, \texttt{cc-add-closed}.}
%
\vnbprf{\textbf{Oracles.} \texttt{ineq}, \texttt{arith}, \texttt{crs}.  \textbf{Bill:} \texttt{proven modulo 0}.}
%
\vnbprf{\vnbproofpart{\texttt{goursat-triangle}}}
%
\vnbprf{\textbf{Proof.} \texttt{theorem-\allowbreak{}library/goursat.scm} --- 1162 recorded steps.}
%
\vnbprf{\textbf{Lemmas.} \texttt{holomorphic-on-open}, \texttt{cc-open-subset-cc}, \texttt{holomorphic-on-in-fun}, \texttt{holomorphic-on-continuous}, \texttt{subset-mem-fwd}, \texttt{goursat-nest}, \texttt{bary-in-conv3}, \texttt{real-part-in-rr}, \texttt{rr-subset-cc}, \texttt{cc-mul-closed}, \texttt{imag-part-in-rr}, \texttt{cc-add-closed}, \texttt{rr-one-in}, \texttt{cc-sub-in-cc}, \texttt{holomorphic-on-diff}, \texttt{cc-nf-carr}, \texttt{fun-apply-type-c}, \texttt{cc-is-normed-field}, \texttt{nf-open-is-set}, \texttt{cc-neg-closed}, \texttt{affine-holomorphic}, \texttt{holomorphic-sum}, \texttt{tri-int-in-cc}, \texttt{cc-magnitude-closed}, \texttt{cc-magnitude-nonneg}, \texttt{rr-add-in-rr}, \texttt{rr-mul-in-rr}, \texttt{rr-leq-mul-nonneg}, \texttt{rr-pos-rr-in-rr}, \texttt{rr-lt-of-pos-rr}, \texttt{rr-pos-ne-zero}, \texttt{rr-recip-closed}, \texttt{rr-recip-pos}, \texttt{rr-recip-inverse}, \texttt{rr-mul-pos}, \texttt{rr-lt-implies-le}, \texttt{neq-sym}, \texttt{caratheodory-remainder-small}, \texttt{small-dyadic}, \texttt{power-two-pos}, \texttt{power-real-closed}, \texttt{recip-power-two-halves}, \texttt{conv3-sub-hull}, \texttt{subset-trans}, \texttt{seg-trace-in-conv3}, \texttt{cc-magnitude-sub-sym}, \texttt{cc-magnitude-zero}, \texttt{bary-side-near}, \texttt{rr-le-scale-nonneg}, \texttt{tri-int-abs-bound-sides}, \texttt{tri-int-sum}, \texttt{tri-int-affine-zero}, \texttt{bary-lip}, \texttt{rr-mul-le-right}, \texttt{rr-power-pos}, \texttt{four-power-halving}, \texttt{rr-zero-in}, \texttt{rr-le-of-le-add-all-pos}, \texttt{cc-magnitude-zero-iff}.}
%
\vnbprf{\textbf{Oracles.} \texttt{crs}, \texttt{ineq}, \texttt{arith}.  \textbf{Bill:} \texttt{proven modulo 0}.}
%
\end{proof}
%
\noindent\textbf{Deviation from the notes.} PROVEN 2026-09-26 (batch 35) WITH THE EXCEPTIONAL POINT: the exceptional-point hypothesis is the notes' (f continuous on the open set u, x in u, f holomorphic on u minus \{x\}), spelled with restrict(f, difference(u, singleton(x))) because HOLOMORPHIC-ON types its function on exactly its set. Proof: x off the hull, Goursat on u minus \{x\}; x in the hull, the fan through x and the vertex case three times; the vertex case is Rudin 10.13 (shrink the triangle towards the vertex; the two remaining triangles avoid x by affine independence; the collinear triangle through the three placements). Earlier: PROVEN 2026-09-25 (batch 30) for f HOLOMORPHIC ON ALL OF U: holomorphic-on(U, f), a, b, c in U, CONV3(a, b, c) subset U implies TRI-INT(f, a, b, c) = 0 (goursat-triangle). The notes' hypothesis 'continuous on U and holomorphic on U minus a point x' is the strengthening of a later batch, needed by the Cauchy integral formula. The proof departs from the notes at the common point of the nested triangles: no Cantor intersection and no compactness; the chosen sub-triangles are driven on the standard triangle and the vertex sequence telescopes (cc-telescoping-limit-bound), its limit lying in CONV3 by definition. No diameter and no length are defined: the estimates are per segment with the perimeter as a term.
%
\renewcommand{\thethm}{3.7}
\begin{cor}[{Segment integrals in a convex set}]\label{thm:cauchy-convex-segments-except}
%
\leavevmode
\vnbpart{\texttt{cauchy-convex-segments-except}}
\vnbstmt{\begin{array}{@{}l@{}} \forall u, f, x.; \forall a, b, c \in u. \\ \quad \operatorname{is\text{-}convex}(u) \wedge \\ \quad \operatorname{is\text{-}open}(\operatorname{nf\text{-}metric\text{-}space}(\operatorname{cc\text{-}normed\text{-}field}), u) \wedge \\ \quad (f \in (u \to \mathbb{C})) \wedge \\ \quad \forall y \in u. \\ \quad \quad \operatorname{is\text{-}continuous\text{-}at}(\operatorname{subspace\text{-}ms}(\operatorname{nf\text{-}metric\text{-}space}(\operatorname{cc\text{-}normed\text{-}field}), u), \\ \quad \quad \quad \operatorname{nf\text{-}metric\text{-}space}(\operatorname{cc\text{-}normed\text{-}field}), f, y) \wedge \\ \quad (x \in u) \wedge \\ \quad \operatorname{holomorphic\text{-}on}(\operatorname{difference}(u, \operatorname{singleton}(x)), \\ \quad \quad \operatorname{restrict}(f, \operatorname{difference}(u, \operatorname{singleton}(x)))) \Rightarrow \\ \quad \quad \operatorname{seg\text{-}int}(f, a, c) \\ \quad \quad =\; (\operatorname{seg\text{-}int}(f, a, b) + \operatorname{seg\text{-}int}(f, b, c)) \end{array}}
\vnbpart{\texttt{cauchy-convex-segments}}
\vnbstmt{\begin{array}{@{}l@{}} \forall u, f.; \forall a, b, c \in u. \\ \quad (\operatorname{is\text{-}convex}(u) \wedge \operatorname{holomorphic\text{-}on}(u, f)) \Rightarrow  \\ \quad \quad \operatorname{seg\text{-}int}(f, a, c) \\ \quad \quad =\; (\operatorname{seg\text{-}int}(f, a, b) + \operatorname{seg\text{-}int}(f, b, c)) \end{array}}
%
\end{cor}
%
\begin{proof}
%
The proofs: the file that holds each, the lemmas it cites, its oracles and its bill.
%
\vnbprf{\vnbproofpart{\texttt{cauchy-convex-segments-except}}}
%
\vnbprf{\textbf{Proof.} \texttt{theorem-\allowbreak{}library/goursat-\allowbreak{}exceptional.scm} --- 67 recorded steps.}
%
\vnbprf{\textbf{Lemmas.} \texttt{nf-cc-is-metric-space}, \texttt{cc-open-subset-cc}, \texttt{holomorphic-on-open}, \texttt{holomorphic-on-in-fun}, \texttt{holomorphic-on-continuous}, \texttt{subset-def}, \texttt{difference-membership}, \texttt{subset-mem-fwd}, \texttt{conv3-in-convex}, \texttt{goursat-triangle-except}, \texttt{seg-trace-in-convex}, \texttt{seg-int-in-cc}, \texttt{seg-int-reverse-at}, \texttt{tri-int-unfold}, \texttt{cc-add-closed}.}
%
\vnbprf{\textbf{Oracles.} \texttt{crs}, \texttt{arith}, \texttt{ineq}.  \textbf{Bill:} \texttt{proven modulo 0}.}
%
\vnbprf{\vnbproofpart{\texttt{cauchy-convex-segments}}}
%
\vnbprf{\textbf{Proof.} \texttt{theorem-\allowbreak{}library/goursat.scm} --- 54 recorded steps.}
%
\vnbprf{\textbf{Lemmas.} \texttt{holomorphic-on-open}, \texttt{cc-open-subset-cc}, \texttt{holomorphic-on-in-fun}, \texttt{holomorphic-on-continuous}, \texttt{subset-mem-fwd}, \texttt{conv3-in-convex}, \texttt{goursat-triangle}, \texttt{seg-trace-in-convex}, \texttt{seg-int-in-cc}, \texttt{seg-int-reverse-at}, \texttt{tri-int-unfold}, \texttt{cc-add-closed}.}
%
\vnbprf{\textbf{Oracles.} \texttt{crs}, \texttt{ineq}, \texttt{arith}.  \textbf{Bill:} \texttt{proven modulo 0}.}
%
\end{proof}
%
\noindent\textbf{Deviation from the notes.} PROVEN 2026-09-26 (batch 35) with the exceptional point; the exceptional-point hypothesis is the notes' (f continuous on the open set u, x in u, f holomorphic on u minus \{x\}), spelled with restrict(f, difference(u, singleton(x))) because HOLOMORPHIC-ON types its function on exactly its set. Earlier: PROVEN 2026-09-25 (batch 30) for f holomorphic on the whole convex open set U (is-convex(U), holomorphic-on(U, f), a, b, c in U implies SEG-INT(f, a, c) = SEG-INT(f, a, b) + SEG-INT(f, b, c)); the exceptional point x as for 3.6.
%
\renewcommand{\thethm}{3.8}
\begin{prop}[{A primitive on a convex set}]\label{thm:cauchy-convex-primitive-except}
%
\leavevmode
\vnbpart{\texttt{cauchy-convex-primitive-except}}
\vnbstmt{\begin{array}{@{}l@{}} \forall u, f, x.; \forall o \in u. \\ \quad \operatorname{is\text{-}convex}(u) \wedge \\ \quad \operatorname{is\text{-}open}(\operatorname{nf\text{-}metric\text{-}space}(\operatorname{cc\text{-}normed\text{-}field}), u) \wedge \\ \quad (f \in (u \to \mathbb{C})) \wedge \\ \quad \forall y \in u. \\ \quad \quad \operatorname{is\text{-}continuous\text{-}at}(\operatorname{subspace\text{-}ms}(\operatorname{nf\text{-}metric\text{-}space}(\operatorname{cc\text{-}normed\text{-}field}), u), \\ \quad \quad \quad \operatorname{nf\text{-}metric\text{-}space}(\operatorname{cc\text{-}normed\text{-}field}), f, y) \wedge \\ \quad (x \in u) \wedge \\ \quad \operatorname{holomorphic\text{-}on}(\operatorname{difference}(u, \operatorname{singleton}(x)), \\ \quad \quad \operatorname{restrict}(f, \operatorname{difference}(u, \operatorname{singleton}(x)))) \Rightarrow \\ \quad \quad \exists g \in (u \to \mathbb{C}). \\ \quad \quad \quad \operatorname{holomorphic\text{-}on}(u, g) \wedge \\ \quad \quad \quad \forall z \in u. \\ \quad \quad \quad \quad \operatorname{is\text{-}diff\text{-}on}(\operatorname{cc\text{-}normed\text{-}field}, u, g, z, f(z)) \end{array}}
\vnbpart{\texttt{cc-diff-on-of-eps-bound}}
\vnbstmt{\begin{array}{@{}l@{}} \forall s.; \forall g \in (s \to \mathbb{C}), a \in s, l \in \mathbb{C}. \\ \quad \operatorname{is\text{-}open}(\operatorname{nf\text{-}metric\text{-}space}(\operatorname{cc\text{-}normed\text{-}field}), s) \wedge \\ \quad \forall e. \\ \quad \quad \operatorname{pos\text{-}rr}(e) \Rightarrow  \\ \quad \quad \quad \exists d. \\ \quad \quad \quad \quad \operatorname{pos\text{-}rr}(d) \wedge \\ \quad \quad \quad \quad \forall y \in s. \\ \quad \quad \quad \quad \quad (\operatorname{magnitude}((y - a)) < d) \Rightarrow  \\ \quad \quad \quad \quad \quad \quad \operatorname{magnitude}(((g(y) - g(a)) - ((y - a) \cdot l))) \\ \quad \quad \quad \quad \quad \quad \leq\; (e \cdot \operatorname{magnitude}((y - a))) \Rightarrow \\ \quad \quad \operatorname{is\text{-}diff\text{-}on}(\operatorname{cc\text{-}normed\text{-}field}, s, g, a, l) \end{array}}
\vnbpart{\texttt{cc-continuous-at-add-const}}
\vnbstmt{\begin{array}{@{}l@{}} \forall u.; \forall f, g \in (u \to \mathbb{C}), k \in \mathbb{C}, y \in u. \\ \quad (u \subseteq \mathbb{C}) \wedge \\ \quad \forall z \in u.\; (g(z) \simeq (f(z) + k)) \wedge \\ \quad \forall y \in u. \\ \quad \quad \operatorname{is\text{-}continuous\text{-}at}(\operatorname{subspace\text{-}ms}(\operatorname{nf\text{-}metric\text{-}space}(\operatorname{cc\text{-}normed\text{-}field}), u), \\ \quad \quad \quad \operatorname{nf\text{-}metric\text{-}space}(\operatorname{cc\text{-}normed\text{-}field}), f, y) \Rightarrow \\ \quad \quad \operatorname{is\text{-}continuous\text{-}at}(\operatorname{subspace\text{-}ms}(\operatorname{nf\text{-}metric\text{-}space}(\operatorname{cc\text{-}normed\text{-}field}), u), \\ \quad \quad \quad \operatorname{nf\text{-}metric\text{-}space}(\operatorname{cc\text{-}normed\text{-}field}), g, y) \end{array}}
\vnbpart{\texttt{cauchy-convex-primitive}}
\vnbstmt{\begin{array}{@{}l@{}} \forall u, f.; \forall o \in u. \\ \quad (\operatorname{is\text{-}convex}(u) \wedge \operatorname{holomorphic\text{-}on}(u, f)) \Rightarrow  \\ \quad \quad \exists g \in (u \to \mathbb{C}). \\ \quad \quad \quad \operatorname{holomorphic\text{-}on}(u, g) \wedge \\ \quad \quad \quad \forall z \in u. \\ \quad \quad \quad \quad \operatorname{is\text{-}diff\text{-}on}(\operatorname{cc\text{-}normed\text{-}field}, u, g, z, f(z)) \end{array}}
%
\end{prop}
%
\begin{proof}
%
The proofs: the file that holds each, the lemmas it cites, its oracles and its bill.
%
\vnbprf{\vnbproofpart{\texttt{cauchy-convex-primitive-except}}}
%
\vnbprf{\textbf{Proof.} \texttt{theorem-\allowbreak{}library/goursat-\allowbreak{}exceptional.scm} --- 480 recorded steps.}
%
\vnbprf{\textbf{Lemmas.} \texttt{nf-cc-is-metric-space}, \texttt{cc-open-subset-cc}, \texttt{holomorphic-on-open}, \texttt{holomorphic-on-in-fun}, \texttt{holomorphic-on-continuous}, \texttt{subset-def}, \texttt{difference-membership}, \texttt{cc-is-normed-field}, \texttt{nf-open-is-set}, \texttt{subset-mem-fwd}, \texttt{seg-trace-in-convex}, \texttt{seg-int-in-cc}, \texttt{fun-apply-type-c}, \texttt{cc-neg-closed}, \texttt{cc-zero-in}, \texttt{cc-sub-in-cc}, \texttt{cc-mul-closed}, \texttt{cc-add-closed}, \texttt{affine-holomorphic}, \texttt{cc-continuous-at-add-const}, \texttt{rr-pos-rr-in-rr}, \texttt{rr-lt-of-pos-rr}, \texttt{subspace-pts}, \texttt{cc-is-metric-space}, \texttt{subspace-dist}, \texttt{nf-cc-agree-dist}, \texttt{cc-ms-dist}, \texttt{cc-magnitude-closed}, \texttt{cc-magnitude-sub-sym}, \texttt{cauchy-convex-segments-except}, \texttt{seg-int-sum}, \texttt{seg-int-const}, \texttt{cc-magnitude-nonneg}, \texttt{seg-trace-elim}, \texttt{ccint-elt-in-rr}, \texttt{ccint-membership}, \texttt{rr-subset-cc}, \texttt{rr-mul-in-rr}, \texttt{cc-magnitude-mul}, \texttt{rr-magnitude-is-abs}, \texttt{rr-abs-of-nonneg}, \texttt{rr-mul-le-right}, \texttt{seg-int-abs-bound}, \texttt{cc-diff-on-of-eps-bound}.}
%
\vnbprf{\textbf{Oracles.} \texttt{crs}, \texttt{ineq}, \texttt{arith}.  \textbf{Bill:} \texttt{proven modulo 0}.}
%
\vnbprf{\vnbproofpart{\texttt{cc-diff-on-of-eps-bound}}}
%
\vnbprf{\textbf{Proof.} \texttt{theorem-\allowbreak{}library/goursat-\allowbreak{}exceptional.scm} --- 550 recorded steps.}
%
\vnbprf{\textbf{Lemmas.} \texttt{cc-is-normed-field}, \texttt{nf-cc-is-metric-space}, \texttt{cc-nf-carr}, \texttt{nf-open-is-set}, \texttt{subspace-is-metric-space}, \texttt{subspace-pts}, \texttt{subset-mem-fwd}, \texttt{fun-apply-type-c}, \texttt{cc-sub-in-cc}, \texttt{cc-recip-closed}, \texttt{cc-recip-inverse}, \texttt{cc-mul-closed}, \texttt{cc-magnitude-closed}, \texttt{rr-pos-rr-in-rr}, \texttt{rr-lt-of-pos-rr}, \texttt{rr-mul-in-rr}, \texttt{rr-lt-implies-le}, \texttt{rr-pos-ne-zero}, \texttt{neq-sym}, \texttt{subspace-dist}, \texttt{nf-cc-agree-dist}, \texttt{cc-ms-dist}, \texttt{cc-magnitude-sub-sym}, \texttt{cc-magnitude-zero}, \texttt{cc-magnitude-mul}, \texttt{cc-magnitude-neg}, \texttt{cc-magnitude-nonneg}, \texttt{cc-magnitude-zero-iff}, \texttt{rr-le-ne-lt}, \texttt{rr-mul-le-cancel-pos}, \texttt{cc-nf-sub}, \texttt{cc-nf-mul-apply}.}
%
\vnbprf{\textbf{Oracles.} \texttt{crs}, \texttt{arith}, \texttt{ineq}.  \textbf{Bill:} \texttt{proven modulo 0}.}
%
\vnbprf{\vnbproofpart{\texttt{cc-continuous-at-add-const}}}
%
\vnbprf{\textbf{Proof.} \texttt{theorem-\allowbreak{}library/goursat-\allowbreak{}exceptional.scm} --- 137 recorded steps.}
%
\vnbprf{\textbf{Lemmas.} \texttt{nf-cc-is-metric-space}, \texttt{cc-is-set}, \texttt{subclass-of-set-is-set}, \texttt{cc-is-metric-space}, \texttt{cc-nf-carr}, \texttt{subspace-pts}, \texttt{fun-apply-type-c}, \texttt{nf-cc-agree-dist}, \texttt{cc-ms-dist}, \texttt{cc-sub-in-cc}, \texttt{cc-add-closed}.}
%
\vnbprf{\textbf{Oracles.} \texttt{crs}, \texttt{arith}, \texttt{ineq}.  \textbf{Bill:} \texttt{proven modulo 0}.}
%
\vnbprf{\vnbproofpart{\texttt{cauchy-convex-primitive}}}
%
\vnbprf{\textbf{Proof.} \texttt{theorem-\allowbreak{}library/goursat.scm} --- 646 recorded steps.}
%
\vnbprf{\textbf{Lemmas.} \texttt{holomorphic-on-open}, \texttt{cc-open-subset-cc}, \texttt{holomorphic-on-in-fun}, \texttt{holomorphic-on-continuous}, \texttt{holomorphic-on-open-cc-ms}, \texttt{cc-is-normed-field}, \texttt{nf-open-is-set}, \texttt{subset-mem-fwd}, \texttt{seg-trace-in-convex}, \texttt{seg-int-in-cc}, \texttt{holomorphic-on-diff}, \texttt{fun-apply-type-c}, \texttt{cc-nf-carr}, \texttt{rr-one-in}, \texttt{caratheodory-remainder-small}, \texttt{rr-pos-rr-in-rr}, \texttt{cc-magnitude-closed}, \texttt{rr-add-in-rr}, \texttt{cc-magnitude-nonneg}, \texttt{cc-neg-closed}, \texttt{cc-zero-in}, \texttt{cc-sub-in-cc}, \texttt{cc-mul-closed}, \texttt{cc-add-closed}, \texttt{affine-holomorphic}, \texttt{holomorphic-sum}, \texttt{cauchy-convex-segments}, \texttt{seg-int-sum}, \texttt{seg-int-const}, \texttt{rr-mul-in-rr}, \texttt{seg-trace-elim}, \texttt{ccint-elt-in-rr}, \texttt{ccint-membership}, \texttt{rr-subset-cc}, \texttt{cc-magnitude-mul}, \texttt{rr-magnitude-is-abs}, \texttt{rr-abs-of-nonneg}, \texttt{rr-mul-le-right}, \texttt{cc-magnitude-triangle}, \texttt{rr-le-scale-nonneg}, \texttt{seg-int-abs-bound}, \texttt{cc-diff-on-local-quadratic}.}
%
\vnbprf{\textbf{Oracles.} \texttt{arith}, \texttt{ineq}, \texttt{crs}.  \textbf{Bill:} \texttt{proven modulo 0}.}
%
\end{proof}
%
\noindent\textbf{Deviation from the notes.} PROVEN 2026-09-26 (batch 35) with the exceptional point, for a given base point o in u; the derivative of the primitive comes from the continuity of f alone through the new brick cc-diff-on-of-eps-bound (the converse of caratheodory-remainder-small). the exceptional-point hypothesis is the notes' (f continuous on the open set u, x in u, f holomorphic on u minus \{x\}), spelled with restrict(f, difference(u, singleton(x))) because HOLOMORPHIC-ON types its function on exactly its set. Earlier: PROVEN 2026-09-25 (batch 30) for f holomorphic on the whole convex open set U, for a given a0 in U (the notes' 'let a in U'): there is g in fun(U, CC), holomorphic on U, with is-diff-on(cc-normed-field, U, g, z, f(z)) at every z in U; the witness is the notes' (68), b |-$>$ SEG-INT(f, a0, b), and its differentiability comes from a local quadratic bound (cc-diff-on-local-quadratic) obtained from the Caratheodory remainder at eps = 1. The exceptional point x as for 3.6.
%
\renewcommand{\thethm}{3.9}
\begin{cor}[{Cauchy's theorem on a convex set}]\label{thm:cauchy-convex-closed-except}
%
\leavevmode
\vnbpart{\texttt{cauchy-convex-closed-except}}
\vnbstmt{\begin{array}{@{}l@{}} \forall u, f, x, m, c, a, b. \\ \quad \operatorname{is\text{-}convex}(u) \wedge \\ \quad \operatorname{is\text{-}open}(\operatorname{nf\text{-}metric\text{-}space}(\operatorname{cc\text{-}normed\text{-}field}), u) \wedge \\ \quad (f \in (u \to \mathbb{C})) \wedge \\ \quad \forall y \in u. \\ \quad \quad \operatorname{is\text{-}continuous\text{-}at}(\operatorname{subspace\text{-}ms}(\operatorname{nf\text{-}metric\text{-}space}(\operatorname{cc\text{-}normed\text{-}field}), u), \\ \quad \quad \quad \operatorname{nf\text{-}metric\text{-}space}(\operatorname{cc\text{-}normed\text{-}field}), f, y) \wedge \\ \quad (x \in u) \wedge \\ \quad \operatorname{holomorphic\text{-}on}(\operatorname{difference}(u, \operatorname{singleton}(x)), \\ \quad \quad \operatorname{restrict}(f, \operatorname{difference}(u, \operatorname{singleton}(x)))) \wedge \\ \quad \operatorname{is\text{-}road}(m, c, a, b) \wedge \\ \quad (\operatorname{trace}(m, a, b) \subseteq u) \wedge \\ \quad (m(b) = m(a)) \Rightarrow \\ \quad \quad (\operatorname{line\text{-}int}(f, m, c, a, b) = 0) \end{array}}
\vnbpart{\texttt{cauchy-convex-closed}}
\vnbstmt{\begin{array}{@{}l@{}} \forall u, f, m, c, a, b. \\ \quad \operatorname{is\text{-}convex}(u) \wedge \\ \quad \operatorname{holomorphic\text{-}on}(u, f) \wedge \\ \quad \operatorname{is\text{-}road}(m, c, a, b) \wedge \\ \quad (\operatorname{trace}(m, a, b) \subseteq u) \wedge \\ \quad (m(b) = m(a)) \Rightarrow \\ \quad \quad (\operatorname{line\text{-}int}(f, m, c, a, b) = 0) \end{array}}
\vnbpart{\texttt{line-int-closed-road-of-derivative}}
\vnbstmt{\begin{array}{@{}l@{}} \forall u, f, c, m, d, a, b. \\ \quad \operatorname{holomorphic\text{-}on}(u, f) \wedge \\ \quad (c \in (u \to \mathbb{C})) \wedge \\ \quad \forall z \in u. \\ \quad \quad \operatorname{is\text{-}diff\text{-}on}(\operatorname{cc\text{-}normed\text{-}field}, u, f, z, c(z)) \wedge \\ \quad \operatorname{is\text{-}road}(m, d, a, b) \wedge \\ \quad (\operatorname{trace}(m, a, b) \subseteq u) \wedge \\ \quad (m(b) = m(a)) \Rightarrow \\ \quad \quad (\operatorname{line\text{-}int}(c, m, d, a, b) = 0) \end{array}}
%
\end{cor}
%
\begin{proof}
%
The proofs: the file that holds each, the lemmas it cites, its oracles and its bill.
%
\vnbprf{\vnbproofpart{\texttt{cauchy-convex-closed-except}}}
%
\vnbprf{\textbf{Proof.} \texttt{theorem-\allowbreak{}library/goursat-\allowbreak{}exceptional.scm} --- 50 recorded steps.}
%
\vnbprf{\textbf{Lemmas.} \texttt{nf-cc-is-metric-space}, \texttt{cc-open-subset-cc}, \texttt{holomorphic-on-open}, \texttt{holomorphic-on-in-fun}, \texttt{holomorphic-on-continuous}, \texttt{subset-def}, \texttt{difference-membership}, \texttt{is-road-is-path}, \texttt{is-path-in-fun}, \texttt{is-path-endpoints}, \texttt{ccint-membership}, \texttt{trace-value-in}, \texttt{subset-mem-fwd}, \texttt{cauchy-convex-primitive-except}, \texttt{line-int-closed-road-of-derivative}.}
%
\vnbprf{\textbf{Oracles.} \texttt{ineq}, \texttt{arith}, \texttt{crs}.  \textbf{Bill:} \texttt{proven modulo 0}.}
%
\vnbprf{\vnbproofpart{\texttt{cauchy-convex-closed}}}
%
\vnbprf{\textbf{Proof.} \texttt{theorem-\allowbreak{}library/goursat.scm} --- 37 recorded steps.}
%
\vnbprf{\textbf{Lemmas.} \texttt{holomorphic-on-open}, \texttt{cc-open-subset-cc}, \texttt{holomorphic-on-in-fun}, \texttt{holomorphic-on-continuous}, \texttt{is-road-is-path}, \texttt{is-path-in-fun}, \texttt{is-path-endpoints}, \texttt{ccint-membership}, \texttt{trace-value-in}, \texttt{subset-mem-fwd}, \texttt{cauchy-convex-primitive}, \texttt{line-int-closed-road-of-derivative}.}
%
\vnbprf{\textbf{Oracles.} \texttt{ineq}, \texttt{crs}, \texttt{arith}.  \textbf{Bill:} \texttt{proven modulo 0}.}
%
\vnbprf{\vnbproofpart{\texttt{line-int-closed-road-of-derivative}}}
%
\vnbprf{\textbf{Proof.} \texttt{theorem-\allowbreak{}library/goursat-\allowbreak{}bricks.scm} --- 49 recorded steps.}
%
\vnbprf{\textbf{Lemmas.} \texttt{line-int-of-derivative}, \texttt{is-road-is-path}, \texttt{is-path-endpoints}, \texttt{is-path-in-fun}, \texttt{ccint-membership}, \texttt{trace-unfold}, \texttt{fun-apply-type-c}, \texttt{membership-implies-sethood}, \texttt{subset-mem-fwd}, \texttt{holomorphic-on-in-fun}.}
%
\vnbprf{\textbf{Oracles.} \texttt{ineq}, \texttt{crs}, \texttt{arith}.  \textbf{Bill:} \texttt{proven modulo 0}.}
%
\end{proof}
%
\noindent\textbf{Deviation from the notes.} PROVEN 2026-09-26 (batch 35) with the exceptional point; a closed path is a road on [a, b] with pgam(b) = pgam(a). the exceptional-point hypothesis is the notes' (f continuous on the open set u, x in u, f holomorphic on u minus \{x\}), spelled with restrict(f, difference(u, singleton(x))) because HOLOMORPHIC-ON types its function on exactly its set. Earlier: PROVEN 2026-09-25 (batch 30) for f holomorphic on the whole convex open set U: for a road on [a, b] with trace in U and gamma(b) = gamma(a), the integral of f along it is 0 (from 3.8 and Cor 3.3 for a closed road). 'Closed path' is the hypothesis gamma(b) = gamma(a); there is no closed-road predicate. The exceptional point x as for 3.6.
%
\section{Deviations}
%
Every entry whose library statement deviates from the notes, or whose
status is not \texttt{proven}.
%
\begin{longtable}{@{}ll>{\raggedright\arraybackslash}p{4.2cm}>{\raggedright\arraybackslash}p{4.2cm}l@{}}
\hline
document & ref & title & library name(s) & status \\
\hline
\endfirsthead
\hline
document & ref & title & library name(s) & status \\
\hline
\endhead
Calculus & 1.1 & Order compatible with addition and multiplication on Z & \texttt{rr-\allowbreak{}lt-\allowbreak{}add}, \texttt{rr-\allowbreak{}mul-\allowbreak{}le-\allowbreak{}right}, \texttt{rr-\allowbreak{}mul-\allowbreak{}pos} & partial \\
Calculus & 1.2 & Trichotomy of the integer ordering & \texttt{rr-\allowbreak{}lt-\allowbreak{}trichotomy}, \texttt{zz-\allowbreak{}trichotomy} & proven \\
Calculus & 1.3 & 1 is the least natural number & \texttt{nn-\allowbreak{}one-\allowbreak{}le-\allowbreak{}succ}, \texttt{zz-\allowbreak{}discrete}, \texttt{nn-\allowbreak{}pos-\allowbreak{}is-\allowbreak{}succ} & partial \\
Calculus & 1.4 & No integer strictly between m and m + 1 & \texttt{zz-\allowbreak{}lt-\allowbreak{}succ-\allowbreak{}le}, \texttt{zz-\allowbreak{}discrete} & proven \\
Calculus & 1.9 & Uniqueness of the least upper bound & --- & absent \\
Calculus & 1.11 & The nested interval property & \texttt{rr-\allowbreak{}nested-\allowbreak{}intervals} & proven \\
Calculus & 1.13 & Order isomorphism of (-1, 1) with the extended reals & --- & absent \\
Calculus & 1.19 & Bolzano-Weierstrass on a closed interval & \texttt{rr-\allowbreak{}interval-\allowbreak{}seq-\allowbreak{}has-\allowbreak{}convergent-\allowbreak{}subseq}, \texttt{rr-\allowbreak{}bolzano-\allowbreak{}weierstrass} & proven \\
Calculus & 2.1 & Derivative of the monomial x\textasciicircum{}n & \texttt{deriv-\allowbreak{}power}, \texttt{deriv-\allowbreak{}const} & proven \\
Calculus & 2.3 & Higher derivatives of x\textasciicircum{}n & \texttt{nth-\allowbreak{}deriv-\allowbreak{}monomial}, \texttt{nth-\allowbreak{}deriv-\allowbreak{}monomial-\allowbreak{}vanish} & proven \\
Calculus & 2.4 & Differentiability implies continuity & \texttt{diff-\allowbreak{}implies-\allowbreak{}continuous}, \texttt{diff-\allowbreak{}on-\allowbreak{}implies-\allowbreak{}continuous} & proven \\
Calculus & 2.5 & Derivative of a sum & \texttt{deriv-\allowbreak{}sum}, \texttt{diff-\allowbreak{}on-\allowbreak{}sum} & proven \\
Calculus & 2.6 & Derivative of a product & \texttt{deriv-\allowbreak{}product}, \texttt{diff-\allowbreak{}on-\allowbreak{}product} & proven \\
Calculus & 2.8 & Chain Rule & \texttt{has-\allowbreak{}deriv-\allowbreak{}at-\allowbreak{}chain}, \texttt{deriv-\allowbreak{}chain}, \texttt{diff-\allowbreak{}on-\allowbreak{}chain} & proven \\
Calculus & 2.10 & Vanishing of the derivative at an interior extremum & \texttt{interval-\allowbreak{}local-\allowbreak{}max-\allowbreak{}right-\allowbreak{}deriv-\allowbreak{}nonpos}, \texttt{interval-\allowbreak{}local-\allowbreak{}max-\allowbreak{}left-\allowbreak{}deriv-\allowbreak{}nonneg}, \texttt{interval-\allowbreak{}local-\allowbreak{}max-\allowbreak{}deriv-\allowbreak{}zero}, \texttt{interior-\allowbreak{}max-\allowbreak{}deriv-\allowbreak{}zero} & proven \\
Calculus & 2.11 & Generalized Mean Value Theorem & \texttt{generalized-\allowbreak{}mvt-\allowbreak{}on-\allowbreak{}interval}, \texttt{generalized-\allowbreak{}mvt} & proven \\
Calculus & 2.12 & Rolle's Lemma & \texttt{rolle-\allowbreak{}on-\allowbreak{}interval}, \texttt{rolle} & proven \\
Calculus & 2.13 & Mean Value Theorem & \texttt{mvt-\allowbreak{}on-\allowbreak{}interval}, \texttt{mvt} & proven \\
Calculus & 2.14 & Increment bounds from derivative bounds & \texttt{mvt-\allowbreak{}upper-\allowbreak{}bound-\allowbreak{}on-\allowbreak{}interval}, \texttt{mvt-\allowbreak{}lower-\allowbreak{}bound-\allowbreak{}on-\allowbreak{}interval}, \texttt{mvt-\allowbreak{}abs-\allowbreak{}bound-\allowbreak{}on-\allowbreak{}interval}, \texttt{mvt-\allowbreak{}upper-\allowbreak{}bound}, \texttt{mvt-\allowbreak{}lower-\allowbreak{}bound}, \texttt{mvt-\allowbreak{}abs-\allowbreak{}bound} & proven \\
Calculus & 2.15 & Zero derivative implies constant & \texttt{deriv-\allowbreak{}zero-\allowbreak{}constant-\allowbreak{}on-\allowbreak{}interval}, \texttt{deriv-\allowbreak{}zero-\allowbreak{}implies-\allowbreak{}constant}, \texttt{deriv-\allowbreak{}zero-\allowbreak{}const-\allowbreak{}up} & proven \\
Calculus & 2.16 & Uniqueness of the interpolating Taylor polynomial at 0 & \texttt{taylor-\allowbreak{}interpolant-\allowbreak{}at-\allowbreak{}zero} & proven \\
Calculus & 2.18 & Taylor's formula with a general auxiliary function G & \texttt{taylor-\allowbreak{}formula-\allowbreak{}on-\allowbreak{}interval}, \texttt{taylor-\allowbreak{}lagrange} & proven \\
Calculus & 2.19 & Taylor's formula with Lagrange remainder & \texttt{taylor-\allowbreak{}lagrange-\allowbreak{}on-\allowbreak{}interval}, \texttt{taylor-\allowbreak{}lagrange} & proven \\
Calculus & 2.20 & Supremum bound for the Taylor remainder & \texttt{taylor-\allowbreak{}remainder-\allowbreak{}sup-\allowbreak{}bound} & proven \\
Calculus & 2.21 & Continuous remainder factor S\_n & \texttt{taylor-\allowbreak{}remainder-\allowbreak{}factor}, \texttt{taylor-\allowbreak{}lagrange-\allowbreak{}subinterval} & proven \\
Calculus & 2.22 & Taylor's formula is exact for a polynomial & \texttt{taylor-\allowbreak{}exact-\allowbreak{}polynomial} & proven \\
Calculus & 2.24 & Mean Value Theorem for a directional derivative & \texttt{dir-\allowbreak{}deriv-\allowbreak{}mvt}, \texttt{dir-\allowbreak{}deriv-\allowbreak{}increment} & proven \\
Calculus & 2.25 & Homogeneity of the directional derivative & \texttt{dir-\allowbreak{}deriv-\allowbreak{}homogeneous} & proven \\
Calculus & 2.26 & Additivity of the directional derivative under continuity & \texttt{dir-\allowbreak{}deriv-\allowbreak{}additive} & proven \\
Calculus & 2.27 & Linearity of the directional derivative on the span & \texttt{dir-\allowbreak{}deriv-\allowbreak{}linear} & partial \\
Calculus & 2.28 & Continuous partials give a directional derivative in every direction & --- & absent \\
Calculus & 2.29 & Continuous partials imply differentiability & --- & absent \\
Calculus & 3.3 & Unions and finite intersections of open sets & \texttt{union-\allowbreak{}of-\allowbreak{}opens-\allowbreak{}open}, \texttt{inter-\allowbreak{}of-\allowbreak{}opens-\allowbreak{}open} & partial \\
Calculus & 3.7 & A monotone bounded sequence converges to its sup (resp. inf) & \texttt{monotone-\allowbreak{}convergence-\allowbreak{}rr} & partial \\
Calculus & 3.9 & Closed sets are exactly those stable under limits of sequences & \texttt{closed-\allowbreak{}limit-\allowbreak{}in}, \texttt{closed-\allowbreak{}set-\allowbreak{}limit-\allowbreak{}in}, \texttt{closed-\allowbreak{}set-\allowbreak{}limit-\allowbreak{}in-\allowbreak{}tail} & partial \\
Calculus & 3.10 & A set is closed if and only if it contains its accumulation points & --- & absent \\
Calculus & 3.12 & Four equivalent characterizations of a compact metric space & \texttt{compact-\allowbreak{}iff-\allowbreak{}fip}, \texttt{compact-\allowbreak{}iff-\allowbreak{}cluster-\allowbreak{}point}, \texttt{compact-\allowbreak{}iff-\allowbreak{}seq-\allowbreak{}compact}, \texttt{compact-\allowbreak{}iff-\allowbreak{}tb-\allowbreak{}complete} & proven \\
Calculus & 3.14 & Continuity at a point is sequential continuity & \texttt{continuous-\allowbreak{}at-\allowbreak{}iff-\allowbreak{}sequential}, \texttt{continuous-\allowbreak{}at-\allowbreak{}implies-\allowbreak{}sequential} & proven \\
Calculus & 3.15 & Continuity equals closedness (openness) of preimages & \texttt{continuous-\allowbreak{}implies-\allowbreak{}closed-\allowbreak{}preimage}, \texttt{closed-\allowbreak{}preimage-\allowbreak{}implies-\allowbreak{}continuous}, \texttt{continuous-\allowbreak{}implies-\allowbreak{}open-\allowbreak{}preimage}, \texttt{open-\allowbreak{}preimage-\allowbreak{}implies-\allowbreak{}continuous} & proven \\
Calculus & 3.16 & Subsequence principle & \texttt{subsequence-\allowbreak{}principle} & proven \\
Calculus & 3.19 & The continuous image of a compact set is compact & \texttt{continuous-\allowbreak{}image-\allowbreak{}of-\allowbreak{}compact}, \texttt{compact-\allowbreak{}image} & proven \\
Calculus & 3.20 & A continuous bijection from a compact space is a homeomorphism & --- & absent \\
Calculus & 3.21 & A continuous function on a compact space is uniformly continuous & \texttt{continuous-\allowbreak{}uniformly-\allowbreak{}continuous-\allowbreak{}on-\allowbreak{}ccint} & partial \\
Calculus & 3.22 & Continuity plus uniform continuity off a compact set gives uniform continuity & --- & absent \\
Calculus & 3.23 & A continuous function constant outside a compact set is uniformly continuous & --- & absent \\
Calculus & 3.26 & Unique uniformly continuous extension from a dense subspace & --- & absent \\
Calculus & 3.27 & min(d,1) is a metric, with the same convergent sequences & \texttt{trunc-\allowbreak{}metric-\allowbreak{}is-\allowbreak{}metric-\allowbreak{}space}, \texttt{trunc-\allowbreak{}metric-\allowbreak{}dist}, \texttt{trunc-\allowbreak{}metric-\allowbreak{}id-\allowbreak{}uniformly-\allowbreak{}continuous}, \texttt{trunc-\allowbreak{}metric-\allowbreak{}id-\allowbreak{}uniformly-\allowbreak{}continuous-\allowbreak{}back}, \texttt{bdd-\allowbreak{}metric-\allowbreak{}converges-\allowbreak{}iff} & partial \\
Calculus & 3.28 & Tychonov Theorem for Metric Spaces & \texttt{compact-\allowbreak{}countable-\allowbreak{}product}, \texttt{seq-\allowbreak{}compact-\allowbreak{}countable-\allowbreak{}product} & proven \\
Calculus & 3.33 & On a compact space pointwise and uniform convergence agree on an equicontinuous set & \texttt{ptwise-\allowbreak{}cauchy-\allowbreak{}compact-\allowbreak{}equicont-\allowbreak{}unif}, \texttt{unif-\allowbreak{}cauchy-\allowbreak{}cont-\allowbreak{}implies-\allowbreak{}uniform-\allowbreak{}limit}, \texttt{uniform-\allowbreak{}limit-\allowbreak{}continuous} & partial \\
Calculus & 3.34 & Ascoli-Arzela & \texttt{ascoli-\allowbreak{}arzela-\allowbreak{}sequential}, \texttt{ascoli-\allowbreak{}sequential-\allowbreak{}from-\allowbreak{}diagonal}, \texttt{ascoli-\allowbreak{}pointwise-\allowbreak{}diagonal} & proven \\
Calculus & 3.36 & Equicontinuous sequence converging on a dense set converges everywhere, with continuous limit & \texttt{ascoli-\allowbreak{}dense-\allowbreak{}bridge}, \texttt{equicont-\allowbreak{}dense-\allowbreak{}conv-\allowbreak{}implies-\allowbreak{}unif-\allowbreak{}cauchy}, \texttt{equicont-\allowbreak{}dense-\allowbreak{}conv-\allowbreak{}ptwise-\allowbreak{}cauchy} & proven \\
Calculus & 3.37 & The isometry group of a compact metric space is compact & --- & absent \\
Calculus & 3.38 & In the isometry group uniform and pointwise convergence agree & --- & absent \\
Calculus & 4.5 & Regulated functions are exactly the uniform limits of step functions & \texttt{regulated-\allowbreak{}on-\allowbreak{}is-\allowbreak{}step-\allowbreak{}limit}, \texttt{regulated-\allowbreak{}on-\allowbreak{}step-\allowbreak{}approx}, \texttt{uniform-\allowbreak{}limit-\allowbreak{}of-\allowbreak{}regulated}, \texttt{continuous-\allowbreak{}is-\allowbreak{}regulated} & partial \\
Calculus & 4.8 & The antiderivable functions on [a,b] form a vector space & \texttt{antiderivable-\allowbreak{}add}, \texttt{antiderivable-\allowbreak{}scale}, \texttt{antiderivable-\allowbreak{}sub} & partial \\
Calculus & 4.10 & Two antiderivatives of the same function differ by a constant & \texttt{pw-\allowbreak{}antiderivative-\allowbreak{}differ-\allowbreak{}by-\allowbreak{}some-\allowbreak{}constant}, \texttt{pw-\allowbreak{}antiderivative-\allowbreak{}differ-\allowbreak{}by-\allowbreak{}constant}, \texttt{antiderivative-\allowbreak{}differ-\allowbreak{}by-\allowbreak{}constant} & proven \\
Calculus & 4.11 & The endpoint difference is the same for all antiderivatives & \texttt{pw-\allowbreak{}antiderivative-\allowbreak{}endpoint-\allowbreak{}difference}, \texttt{antiderivative-\allowbreak{}endpoint-\allowbreak{}difference} & proven \\
Calculus & 4.12 & The definite integral is a linear functional & \texttt{pw-\allowbreak{}antiderivative-\allowbreak{}linear}, \texttt{pw-\allowbreak{}antiderivative-\allowbreak{}sum}, \texttt{pw-\allowbreak{}antiderivative-\allowbreak{}real-\allowbreak{}mul}, \texttt{c-\allowbreak{}int-\allowbreak{}add}, \texttt{c-\allowbreak{}int-\allowbreak{}scale} & proven \\
Calculus & 4.13 & The integral of a sum is the sum of the integrals & \texttt{pw-\allowbreak{}antiderivative-\allowbreak{}sum}, \texttt{c-\allowbreak{}int-\allowbreak{}add}, \texttt{antiderivative-\allowbreak{}add} & proven \\
Calculus & 4.16 & Differentiating a sequence whose derivatives converge uniformly & \texttt{primitive-\allowbreak{}uniform-\allowbreak{}limit}, \texttt{antiderivative-\allowbreak{}family-\allowbreak{}uniform-\allowbreak{}limit}, \texttt{antiderivable-\allowbreak{}uniform-\allowbreak{}limit}, \texttt{diff-\allowbreak{}at-\allowbreak{}ptwise-\allowbreak{}limit} & partial \\
Calculus & 4.17 & The antiderivable functions are closed under uniform limits & \texttt{antiderivable-\allowbreak{}uniform-\allowbreak{}limit} & proven \\
Calculus & 5.2 & Bernstein approximation of a continuous function on [0,1] & \texttt{bernstein-\allowbreak{}uniform-\allowbreak{}approximation} & proven \\
Calculus & 5.4 & The two Bernstein estimates & \texttt{bernstein-\allowbreak{}partition}, \texttt{bernstein-\allowbreak{}variance-\allowbreak{}bound} & proven \\
Calculus & 5.5 & The four Bernstein identities & \texttt{bernstein-\allowbreak{}partition}, \texttt{bernstein-\allowbreak{}moment-\allowbreak{}1}, \texttt{bernstein-\allowbreak{}moment-\allowbreak{}2}, \texttt{bernstein-\allowbreak{}variance} & proven \\
Calculus & 5.7 & Weierstrass approximation on an arbitrary interval & \texttt{bernstein-\allowbreak{}uniform-\allowbreak{}approximation-\allowbreak{}ccint} & partial \\
Linear algebra & 3.3 & Action of the matrix unit E\_n[k,l] on the right & \texttt{matunit-\allowbreak{}col-\allowbreak{}shift} & proven \\
Linear algebra & 3.5 & The elementary column operations & \texttt{elem-\allowbreak{}f-\allowbreak{}action}, \texttt{elem-\allowbreak{}g-\allowbreak{}action}, \texttt{elem-\allowbreak{}h-\allowbreak{}action} & proven \\
Linear algebra & 3.6 & The elementary column matrices are invertible, with explicit inverses & \texttt{elem-\allowbreak{}f-\allowbreak{}inverse}, \texttt{elem-\allowbreak{}g-\allowbreak{}inverse}, \texttt{elem-\allowbreak{}h-\allowbreak{}inverse}, \texttt{elem-\allowbreak{}f-\allowbreak{}invertible}, \texttt{elem-\allowbreak{}g-\allowbreak{}invertible} & proven \\
Linear algebra & 3.8 & Column reduction of the first row over a field & \texttt{clear-\allowbreak{}first-\allowbreak{}row}, \texttt{class-\allowbreak{}min-\allowbreak{}pivot} & partial \\
Linear algebra & 3.10 & Column reduction to a step matrix with zeros above the diagonal & \texttt{smith-\allowbreak{}diagonalization} & partial \\
Linear algebra & 3.12 & Over a field, invertible iff column equivalent to the identity & --- & absent \\
Linear algebra & 3.13 & The elementary column matrices generate the group of invertible matrices & --- & absent \\
Linear algebra & 3.15 & A permutation of the indeterminates is a ring automorphism of R[x\_1,...,x\_n] & --- & absent \\
Linear algebra & 3.18 & det is multilinear and alternating, and det of the transpose equals det & \texttt{det-\allowbreak{}row-\allowbreak{}linear}, \texttt{det-\allowbreak{}col-\allowbreak{}linear}, \texttt{det-\allowbreak{}swap-\allowbreak{}rows}, \texttt{det-\allowbreak{}swap-\allowbreak{}cols}, \texttt{det-\allowbreak{}alternating-\allowbreak{}rows}, \texttt{det-\allowbreak{}equal-\allowbreak{}cols-\allowbreak{}zero}, \texttt{det-\allowbreak{}transpose} & proven \\
Linear algebra & 3.19 & Determinant under an elementary column operation & \texttt{det-\allowbreak{}elem-\allowbreak{}f-\allowbreak{}col}, \texttt{det-\allowbreak{}elem-\allowbreak{}g-\allowbreak{}col}, \texttt{det-\allowbreak{}elem-\allowbreak{}h-\allowbreak{}col}, \texttt{det-\allowbreak{}mul-\allowbreak{}elem-\allowbreak{}f-\allowbreak{}col}, \texttt{det-\allowbreak{}mul-\allowbreak{}elem-\allowbreak{}g-\allowbreak{}col}, \texttt{det-\allowbreak{}mul-\allowbreak{}elem-\allowbreak{}h-\allowbreak{}col} & proven \\
Linear algebra & 3.20 & det(AV) = det A det V for invertible V over a field & \texttt{det-\allowbreak{}mul-\allowbreak{}invertible}, \texttt{det-\allowbreak{}multiplicative} & proven \\
Linear algebra & 3.22 & Expansion by minors along a row and along a column & \texttt{det-\allowbreak{}expand-\allowbreak{}row}, \texttt{det-\allowbreak{}expand-\allowbreak{}col}, \texttt{det-\allowbreak{}cofactor} & proven \\
Linear algebra & 3.23 & The determinant of a triangular matrix is the product of its diagonal entries & --- & absent \\
Linear algebra & 3.25 & The formal adjoint satisfies adj(A) A = det(A) I & --- & absent \\
Linear algebra & 3.26 & Over a field, invertible iff the determinant is nonzero & --- & absent \\
Linear algebra & 3.27 & det is multiplicative & \texttt{det-\allowbreak{}multiplicative}, \texttt{det-\allowbreak{}alternating-\allowbreak{}form} & proven \\
Linear algebra & 3.29 & Behaviour of the sigma-th order subdeterminants under F and G operations & --- & absent \\
Linear algebra & 3.31 & The gcd of the sigma-th subdeterminants is an invariant over Z & --- & absent \\
Linear algebra & 3.36 & Integer normal form: existence and uniqueness of the diagonal form & \texttt{smith-\allowbreak{}normal-\allowbreak{}form}, \texttt{smith-\allowbreak{}diagonalization} & partial \\
Linear algebra & 3.37 & The reducing matrices U and V are products of elementary matrices & \texttt{smith-\allowbreak{}normal-\allowbreak{}form} & partial \\
Linear algebra & 3.38 & An invertible integer matrix is a product of elementary matrices & --- & absent \\
Linear algebra & 3.40 & An invertible change of sequence preserves generation and freeness & \texttt{generates-\allowbreak{}transport}, \texttt{free-\allowbreak{}transport}, \texttt{spans-\allowbreak{}transport} & proven \\
Linear algebra & 3.41 & A relation-free sequence is no longer than a generating sequence & \texttt{free-\allowbreak{}length-\allowbreak{}le-\allowbreak{}generators} & partial \\
Linear algebra & 3.43 & GL\_n(Z) acts freely and transitively on the free generating sequences & --- & absent \\
Linear algebra & 3.44 & Adapted basis for a subgroup, with the divisibility chain & \texttt{submodule-\allowbreak{}of-\allowbreak{}free-\allowbreak{}is-\allowbreak{}free} & partial \\
Linear algebra & 3.45 & The norm N\_F(E): the absolute determinant is basis-independent & --- & absent \\
Linear algebra & 3.46 & A subgroup of a free abelian group of rank n is free of rank at most n & \texttt{submodule-\allowbreak{}of-\allowbreak{}free-\allowbreak{}is-\allowbreak{}free} & asserted \\
Linear algebra & 3.47 & A subgroup of an n-generated group is generated by at most n elements & \texttt{submodule-\allowbreak{}fg}, \texttt{spans-\allowbreak{}submodule-\allowbreak{}fg} & proven \\
Linear algebra & 3.48 & Uniqueness of the principal sequence of a subgroup & --- & absent \\
Linear algebra & 3.49 & The quotient E/F for a diagonal subgroup & --- & absent \\
Linear algebra & 3.50 & Structure of a finitely generated abelian group & --- & absent \\
Linear algebra & 3.51 & Every finitely generated abelian group is a direct sum of cyclic groups & --- & absent \\
Linear algebra & 3.52 & The cardinality of E/F is the norm N\_E(F) & --- & absent \\
Complex analysis & 1.2 & Uniqueness of the least upper bound & \texttt{esup-\allowbreak{}unique} & partial \\
Complex analysis & 1.3 & C as 2x2 real matrices & --- & absent \\
Complex analysis & 1.5 & Unions and finite intersections of open sets & \texttt{union-\allowbreak{}of-\allowbreak{}opens-\allowbreak{}open}, \texttt{inter-\allowbreak{}of-\allowbreak{}opens-\allowbreak{}open} & partial \\
Complex analysis & 1.7 & Monotone convergence for real sequences & \texttt{monotone-\allowbreak{}convergence-\allowbreak{}rr} & partial \\
Complex analysis & 1.9 & Closed sets and limits of sequences & \texttt{closed-\allowbreak{}limit-\allowbreak{}in}, \texttt{closed-\allowbreak{}set-\allowbreak{}limit-\allowbreak{}in} & partial \\
Complex analysis & 1.11 & Characterisations of compactness & \texttt{compact-\allowbreak{}iff-\allowbreak{}fip}, \texttt{compact-\allowbreak{}iff-\allowbreak{}cluster-\allowbreak{}point}, \texttt{compact-\allowbreak{}iff-\allowbreak{}tb-\allowbreak{}complete} & proven \\
Complex analysis & 1.15 & C is an algebraically closed field & --- & absent \\
Complex analysis & 1.16 & Convergence in the extended plane & --- & absent \\
Complex analysis & 1.18 & Linear maps on C\textasciicircum{}2 and linear fractional transformations & --- & absent \\
Complex analysis & 1.19 & The linear fractional transformations form a group & --- & absent \\
Complex analysis & 1.20 & SL(2,C) onto the linear fractional transformations & --- & absent \\
Complex analysis & 1.21 & The Cayley transform & --- & absent \\
Complex analysis & 1.22 & The Hopf map is continuous and onto S\textasciicircum{}2 & --- & absent \\
Complex analysis & 1.24 & The Hopf fibration factors through a bijection S\textasciicircum{}2 -$>$ P & --- & absent \\
Complex analysis & 2.2 & Differentiable implies continuous & \texttt{diff-\allowbreak{}on-\allowbreak{}implies-\allowbreak{}continuous} & proven \\
Complex analysis & 2.3 & Sum and product rules & \texttt{diff-\allowbreak{}on-\allowbreak{}sum}, \texttt{diff-\allowbreak{}on-\allowbreak{}product}, \texttt{holomorphic-\allowbreak{}sum}, \texttt{holomorphic-\allowbreak{}product} & proven \\
Complex analysis & 2.4 & Chain Rule & \texttt{diff-\allowbreak{}on-\allowbreak{}chain} & proven \\
Complex analysis & 2.5 & Root Test & \texttt{root-\allowbreak{}test-\allowbreak{}converges}, \texttt{root-\allowbreak{}test-\allowbreak{}diverges} & proven \\
Complex analysis & 2.7 & Radius of convergence by the root test & \texttt{cps-\allowbreak{}radius-\allowbreak{}formula} & proven \\
Complex analysis & 2.8 & Normal convergence inside the radius & \texttt{cps-\allowbreak{}normally-\allowbreak{}convergent-\allowbreak{}inside}, \texttt{cps-\allowbreak{}converges-\allowbreak{}inside} & proven \\
Complex analysis & 2.9 & Polynomially bounded factors do not shrink the radius & \texttt{cps-\allowbreak{}radius-\allowbreak{}poly-\allowbreak{}factor} & proven \\
Complex analysis & 2.10 & A power series is holomorphic and differentiates termwise & \texttt{cps-\allowbreak{}holomorphic-\allowbreak{}ball}, \texttt{cps-\allowbreak{}radius-\allowbreak{}derived}, \texttt{cps-\allowbreak{}diff-\allowbreak{}on-\allowbreak{}ball} & proven \\
Complex analysis & 2.12 & Principle of Analytic Continuation & --- & absent \\
Complex analysis & 2.13 & Finitely many zeros in a closed ball & --- & absent \\
Complex analysis & 2.14 & Taylor coefficients of an analytic function & \texttt{taylor-\allowbreak{}coefficients-\allowbreak{}at}, \texttt{taylor-\allowbreak{}coefficients}, \texttt{taylor-\allowbreak{}coefficient-\allowbreak{}formula}, \texttt{cps-\allowbreak{}translate-\allowbreak{}nth-\allowbreak{}deriv}, \texttt{analytic-\allowbreak{}holomorphic-\allowbreak{}ball}, \texttt{cps-\allowbreak{}taylor-\allowbreak{}coefficients} & proven \\
Complex analysis & 2.15 & Ratio Test & \texttt{ratio-\allowbreak{}test-\allowbreak{}converges-\allowbreak{}limsup}, \texttt{ratio-\allowbreak{}test-\allowbreak{}diverges-\allowbreak{}liminf}, \texttt{ratio-\allowbreak{}test-\allowbreak{}converges} & proven \\
Complex analysis & 2.16 & The exponential series has infinite radius & \texttt{cc-\allowbreak{}exp-\allowbreak{}radius}, \texttt{cps-\allowbreak{}radius-\allowbreak{}pos-\allowbreak{}inf}, \texttt{cc-\allowbreak{}cos-\allowbreak{}radius}, \texttt{cc-\allowbreak{}sin-\allowbreak{}radius} & proven \\
Complex analysis & 2.17 & exp is a continuous homomorphism into C* & \texttt{cc-\allowbreak{}exp-\allowbreak{}add}, \texttt{cc-\allowbreak{}exp-\allowbreak{}zero}, \texttt{cc-\allowbreak{}exp-\allowbreak{}neg}, \texttt{cc-\allowbreak{}exp-\allowbreak{}nonzero}, \texttt{cc-\allowbreak{}exp-\allowbreak{}maps-\allowbreak{}to-\allowbreak{}nonzero}, \texttt{cc-\allowbreak{}exp-\allowbreak{}continuous}, \texttt{cc-\allowbreak{}exp-\allowbreak{}deriv}, \texttt{cc-\allowbreak{}exp-\allowbreak{}holomorphic} & proven \\
Complex analysis & 2.18 & The kernel of exp is 2 pi i Z & \texttt{cc-\allowbreak{}exp-\allowbreak{}one-\allowbreak{}iff}, \texttt{pi-\allowbreak{}spec}, \texttt{cc-\allowbreak{}exp-\allowbreak{}kernel-\allowbreak{}generator}, \texttt{cc-\allowbreak{}exp-\allowbreak{}kernel-\allowbreak{}generator-\allowbreak{}unique}, \texttt{cc-\allowbreak{}exp-\allowbreak{}two-\allowbreak{}pi-\allowbreak{}i}, \texttt{cc-\allowbreak{}exp-\allowbreak{}i-\allowbreak{}pi} & proven \\
Complex analysis & 2.19 & exp is a bijection on a horizontal strip & \texttt{cc-\allowbreak{}exp-\allowbreak{}bijection-\allowbreak{}strip}, \texttt{cc-\allowbreak{}exp-\allowbreak{}bijection-\allowbreak{}open-\allowbreak{}strip}, \texttt{cc-\allowbreak{}exp-\allowbreak{}surjective-\allowbreak{}strip}, \texttt{cc-\allowbreak{}exp-\allowbreak{}injective-\allowbreak{}strip}, \texttt{cc-\allowbreak{}exp-\allowbreak{}open-\allowbreak{}strip-\allowbreak{}onto}, \texttt{cc-\allowbreak{}exp-\allowbreak{}open-\allowbreak{}strip-\allowbreak{}avoids-\allowbreak{}negatives}, \texttt{cc-\allowbreak{}polar-\allowbreak{}form}, \texttt{cc-\allowbreak{}exp-\allowbreak{}real-\allowbreak{}is-\allowbreak{}r-\allowbreak{}exp} & proven \\
Complex analysis & 2.20 & The logarithm series has radius 1 & \texttt{cps-\allowbreak{}log-\allowbreak{}radius-\allowbreak{}one}, \texttt{cps-\allowbreak{}log-\allowbreak{}series}, \texttt{cps-\allowbreak{}log-\allowbreak{}radius}, \texttt{cps-\allowbreak{}log-\allowbreak{}radius-\allowbreak{}le-\allowbreak{}one}, \texttt{cc-\allowbreak{}series-\allowbreak{}terms-\allowbreak{}small}, \texttt{cc-\allowbreak{}geometric-\allowbreak{}alternating} & proven \\
Complex analysis & 2.21 & The logarithm series agrees with ln on (0,2) & \texttt{log0-\allowbreak{}is-\allowbreak{}ln}, \texttt{log0-\allowbreak{}exp}, \texttt{log0-\allowbreak{}is-\allowbreak{}log-\allowbreak{}on-\allowbreak{}ball}, \texttt{cc-\allowbreak{}log-\allowbreak{}holomorphic-\allowbreak{}near-\allowbreak{}one} & proven \\
Complex analysis & 2.3.3.1 & The principal branch of the logarithm & \texttt{cc-\allowbreak{}log-\allowbreak{}spec}, \texttt{cc-\allowbreak{}exp-\allowbreak{}log}, \texttt{cc-\allowbreak{}log-\allowbreak{}exp}, \texttt{cc-\allowbreak{}log-\allowbreak{}char}, \texttt{cc-\allowbreak{}log-\allowbreak{}unique}, \texttt{cc-\allowbreak{}log-\allowbreak{}real}, \texttt{cc-\allowbreak{}log-\allowbreak{}one}, \texttt{cc-\allowbreak{}log-\allowbreak{}polar}, \texttt{cc-\allowbreak{}log-\allowbreak{}scale}, \texttt{cc-\allowbreak{}log-\allowbreak{}in-\allowbreak{}slit} & proven \\
Complex analysis & 2.3.3.1-hol & log is holomorphic on C minus (-inf, 0] & \texttt{cc-\allowbreak{}log-\allowbreak{}holomorphic}, \texttt{cc-\allowbreak{}log-\allowbreak{}deriv}, \texttt{cc-\allowbreak{}log-\allowbreak{}holomorphic-\allowbreak{}at}, \texttt{cc-\allowbreak{}log-\allowbreak{}im-\allowbreak{}near-\allowbreak{}one}, \texttt{cc-\allowbreak{}log-\allowbreak{}rotate}, \texttt{cc-\allowbreak{}log-\allowbreak{}rotate-\allowbreak{}neg}, \texttt{cc-\allowbreak{}log-\allowbreak{}mul}, \texttt{cc-\allowbreak{}log-\allowbreak{}right-\allowbreak{}half}, \texttt{cc-\allowbreak{}pi-\allowbreak{}gt-\allowbreak{}three}, \texttt{cc-\allowbreak{}cos-\allowbreak{}pos-\allowbreak{}half}, \texttt{cc-\allowbreak{}log-\allowbreak{}holomorphic-\allowbreak{}near-\allowbreak{}one}, \texttt{cc-\allowbreak{}log-\allowbreak{}holomorphic-\allowbreak{}scale} & proven \\
Complex analysis & 2.22 & The series of log(1 - z) and log(1 + z) & \texttt{cc-\allowbreak{}log-\allowbreak{}one-\allowbreak{}minus}, \texttt{cc-\allowbreak{}log-\allowbreak{}one-\allowbreak{}plus}, \texttt{log0-\allowbreak{}is-\allowbreak{}log-\allowbreak{}on-\allowbreak{}ball} & partial \\
Complex analysis & 3.1 & Path additivity of the line integral & \texttt{line-\allowbreak{}int-\allowbreak{}concat}, \texttt{line-\allowbreak{}int-\allowbreak{}juxta}, \texttt{juxta-\allowbreak{}is-\allowbreak{}road}, \texttt{regulated-\allowbreak{}on-\allowbreak{}glue}, \texttt{tri-\allowbreak{}int-\allowbreak{}is-\allowbreak{}line-\allowbreak{}int}, \texttt{line-\allowbreak{}int-\allowbreak{}length-\allowbreak{}bound}, \texttt{line-\allowbreak{}int-\allowbreak{}adjacent}, \texttt{road-\allowbreak{}restrict} & proven \\
Complex analysis & 3.2 & Fundamental theorem for line integrals & \texttt{line-\allowbreak{}int-\allowbreak{}of-\allowbreak{}derivative}, \texttt{holomorphic-\allowbreak{}chain} & proven \\
Complex analysis & 3.3 & Polynomials integrate to zero over a closed path & \texttt{line-\allowbreak{}int-\allowbreak{}closed-\allowbreak{}road-\allowbreak{}of-\allowbreak{}derivative}, \texttt{affine-\allowbreak{}closed-\allowbreak{}road-\allowbreak{}integral}, \texttt{affine-\allowbreak{}has-\allowbreak{}primitive-\allowbreak{}cc} & partial \\
Complex analysis & 3.4 & Splitting a segment integral at an interior point & \texttt{segment-\allowbreak{}int-\allowbreak{}split}, \texttt{segment-\allowbreak{}int-\allowbreak{}abs-\allowbreak{}bound}, \texttt{segment-\allowbreak{}int-\allowbreak{}reverse}, \texttt{line-\allowbreak{}int-\allowbreak{}affine-\allowbreak{}reparam}, \texttt{line-\allowbreak{}int-\allowbreak{}adjacent} & proven \\
Complex analysis & 3.5 & The barycentric-subdivision sum & \texttt{tri-\allowbreak{}int-\allowbreak{}subdivision}, \texttt{seg-\allowbreak{}int-\allowbreak{}split-\allowbreak{}at}, \texttt{segment-\allowbreak{}int-\allowbreak{}reverse} & proven \\
Complex analysis & 3.6 & Cauchy's theorem for triangles (Goursat) & \texttt{goursat-\allowbreak{}triangle-\allowbreak{}except}, \texttt{goursat-\allowbreak{}vertex}, \texttt{tri-\allowbreak{}int-\allowbreak{}fan}, \texttt{tri-\allowbreak{}int-\allowbreak{}vertex-\allowbreak{}split}, \texttt{conv3-\allowbreak{}avoids-\allowbreak{}vertex}, \texttt{collinear-\allowbreak{}cases}, \texttt{tri-\allowbreak{}int-\allowbreak{}collinear}, \texttt{tri-\allowbreak{}int-\allowbreak{}restrict}, \texttt{goursat-\allowbreak{}triangle} & proven \\
Complex analysis & 3.7 & Segment integrals in a convex set & \texttt{cauchy-\allowbreak{}convex-\allowbreak{}segments-\allowbreak{}except}, \texttt{cauchy-\allowbreak{}convex-\allowbreak{}segments} & proven \\
Complex analysis & 3.8 & A primitive on a convex set & \texttt{cauchy-\allowbreak{}convex-\allowbreak{}primitive-\allowbreak{}except}, \texttt{cc-\allowbreak{}diff-\allowbreak{}on-\allowbreak{}of-\allowbreak{}eps-\allowbreak{}bound}, \texttt{cc-\allowbreak{}continuous-\allowbreak{}at-\allowbreak{}add-\allowbreak{}const}, \texttt{cauchy-\allowbreak{}convex-\allowbreak{}primitive} & proven \\
Complex analysis & 3.9 & Cauchy's theorem on a convex set & \texttt{cauchy-\allowbreak{}convex-\allowbreak{}closed-\allowbreak{}except}, \texttt{cauchy-\allowbreak{}convex-\allowbreak{}closed}, \texttt{line-\allowbreak{}int-\allowbreak{}closed-\allowbreak{}road-\allowbreak{}of-\allowbreak{}derivative} & proven \\
\hline
\end{longtable}
%
\subsection*{Totals}
%
\begin{tabular}{@{}lrrrrr@{}}
\hline
document & entries & proven & asserted & partial & absent \\
\hline
Calculus & 66 & 42 & 0 & 13 & 11 \\
Linear algebra & 33 & 10 & 1 & 6 & 16 \\
Complex analysis & 44 & 27 & 0 & 6 & 11 \\
\hline
\end{tabular}
%
